%%%%%%%%%%%%%%%%%%%%%%%%%%%%%%%%%%%%%%%%%%%%%%%%%%%%%%%%%%%%%%%%%%%%%%
% Lớp 11 | Các số đặc trưng đo xu thế trung tâm cho mẫu số liệu | Trung vị và tứ phân vị của mẫu số liệu ghép nhóm
% Dạng 1D5?2-3: Tứ phân vị 
% Số câu: 205  (sinh tự động bởi tach_cau.py; từ nay thầy sửa trực tiếp file này)
%%%%%%%%%%%%%%%%%%%%%%%%%%%%%%%%%%%%%%%%%%%%%%%%%%%%%%%%%%%%%%%%%%%%%%

%%%%%%%%%%%%%%%% TRẮC NGHIỆM 4 PHƯƠNG ÁN %%%%%%%%%%%%%%%%

% nguồn: Ok_Toan11_TracNghiem_Lan2
\begin{ex}%[1D5N2-3]%[ma:648b633f]
	Chiều  cao của học sinh lớp $5$ được mô tả như bảng sau
	\begin{center}
		\begin{tabular}{|c|c|}
			\hline 
			Lớp  chiều cao & Tần số \\
			 \hline
			$[98; 103)$ &  $6$ \\
			\hline 
			$[103; 108)$  & $7$ \\
			\hline 
			$[108; 113)$  & $9$ \\
			\hline 
			$[113; 118)$ &  $5$ \\
			\hline 
			$[118; 123)$ &  $6$ \\
			\hline 
			$[123;  128)$ & $4$ \\
			\hline 
			$[128; 133)$ &  $2$ \\
			\hline 
			$[133; 138)$  & $2$ \\
			 \hline
			$[138;  143)$ & $3$ \\
			 \hline
			$[143;  148]$ & $1$ \\
			 \hline
		\end{tabular}
	\end{center}
	Tứ phân vị thứ nhất của bảng phân bố ghép nhóm trên thuộc khoảng nào sau đây? 
	 \choice
		{$[95; 100]$} 
		{\True $[101;  107]$}
		{$[108;  112]$}
		{$[113;115]$} 
	 \loigiai{
		Ta có $Q_1=103+\dfrac{\dfrac{45}{4}-6}{7} \cdot 5=106{,}75$. 
	} 
\end{ex}

% nguồn: Ok_Toan11_TracNghiem_Lan2
\begin{ex}%[1D5N2-3]%[ma:d562f671]
	Điểm của $32$ học sinh trong kì thi Tiếng Anh (thang $100$  điểm) được phân bố như sau
	\begin{center}
		\begin{tabular}{|c|c|}
			 \hline
			 Lớp điểm & Tần số $n_{i}$ \\
			\hline 
			$[40;  50)$ & $4$ \\
			 \hline
			$[50;  60)$ & $6$ \\
			\hline 
			$[60; 70)$  & $10$ \\
			\hline 
			$[70; 80)$ & $6$  \\
			\hline 
			$[80; 90)$ & $4$  \\
			\hline 
			$[90;  100]$ & $2$ \\
			 \hline
		\end{tabular}
	\end{center}
	Tứ phân vị  thứ ba của bảng phân bố ghép nhóm bằng
	 \choice
		{$\dfrac{310}{4}$} 
		{$\dfrac{315}{3}$} 
		{$\dfrac{220}{4}$} 
		{\True $\dfrac{230}{3}$} 
	\loigiai{ 
		Ta  có $Q_3=70+\dfrac{24-20}{6} \cdot 10=\dfrac{230}{3}$.
	} 
\end{ex}

% nguồn: Ok_Toan11_TracNghiem_Lan2
\begin{ex}%[1D5N2-3]%[ma:cf1f4401]
	Thời  gian luyện tập trong một ngày (tính theo giờ) của một số vận động viên được ghi lại ở bảng sau
	\begin{center}
		\begin{tabular}{|c|c|c|c|c|c|}
			\hline 
			Thời gian luyện tập  (giờ) & $[0; 2)$ & $[2; 4)$ & $[4; 6)$ & $[6; 8)$ & $[8; 10)$ \\
			\hline 
			Số vận động viên & $3$ & $8$ & $12$  & $12$ & $4$ \\
			 \hline
		\end{tabular}
	\end{center}
	Hãy xác định tứ phân vị thứ ba của mẫu  số liệu trên.
	\choice 
		 {$3{,}6875$}
		{$5{,}417$} 
		{\True $7{,}042$} 
		 {$7,68$}
	\loigiai{ 
		Tứ phân vị thứ ba của dãy số liệu $x_{1}$, $x_{2}$,$\ldots$, $x_{39}$ là $x_{30}  \in[6; 8)$.\\
		Do đó tứ phân vị thứ ba của  mẫu số liệu ghép nhóm là\\
		\centerline{$Q_{3}=6+\dfrac{\dfrac{3\cdot39}{4}-(3+8+12)}{12} \cdot(8-6)=\dfrac{169}{24}  \approx 7{,}042$.}
	 }
\end{ex}

% nguồn: Ok_Toan11_TracNghiem_Lan2
\begin{ex}%[1D5N2-3]%[ma:2ab74ec2]
	Bộ ba tứ phân vị của  mẫu số liệu ghép nhóm
	 \choice
		{\True Là giá trị xấp xỉ cho tứ phân vị của mẫu số liệu gốc và được sử dụng làm giá trị đo xu thế trung tâm của  mẫu số liệu}
		{Chia mẫu số liệu đã sắp xếp theo thứ tự không giảm thành $3$ phần  đều nhau}
		{Là giá trị xấp xỉ cho mẫu số  liệu gốc và có thể lấy làm giá trị đại diện cho mẫu số liệu}
		{Xác định chính xác bộ ba tứ phân vị của  mẫu số liệu gốc}
	 \loigiai{
		Bộ ba tứ phân vị của mẫu số liệu ghép nhóm là giá  trị xấp xỉ cho tứ phân vị của mẫu số liệu gốc và được sử dụng làm giá trị đo xu thế trung tâm của mẫu số liệu.
	} 
\end{ex}

% nguồn: Ok_Toan11_TracNghiem_Lan2
\begin{ex}%[1D5N2-3]%[ma:9bd0a113]
	Cho mẫu số liệu ghép nhóm thống kê thời gian sử dụng điện thoại trước khi ngủ của một  người trong $120$ ngày như ở bảng sau
	\begin{center}
		\begin{tabular}{|c|c|}
			 \hline
			Nhóm & Tần số \\ 
			\hline 
			$[0; 4)$  & $13$ \\
			$[4; 8)$ & $29$  \\
			$[8; 12)$  & $48$ \\
			\hline 
			$[12;  16)$ & $22$ \\
			 \hline
			 $[16; 20)$ & $8$ \\
			\hline 
			& $n=120$ \\ 
			 \hline
		\end{tabular}
	\end{center}
	Giá trị các tứ phân vị thứ nhất,  thứ hai và thứ ba lần lượt là
	 \choice
		{$9{,}5$; $12$; $6{,}3$} 
		{\True $6{,}3$; $9{,}5$; $12$} 
		 {$9{,}5$; $6{,}3$; $12$}
		{$12$;  $6{,}3$; $9{,}5$}
	 \loigiai{
		Bảng tần số ghép nhóm bao gồm  cả tần số tích luỹ được cho như ở bảng sau đây
		\begin{center}
			\begin{tabular}{|c|c|c|}
				\hline 
				Nhóm & Tần  số & Tần số tích luỹ \\
				 \hline
				$[0; 4)$ & $13$ & $13$ \\ 
				 \hline
				$[4; 8)$ &  $29$ & $42$ \\
				 \hline
				$[8; 12)$  & $48$ & $90$ \\
				 \hline
				$[12;  16)$ & $22$ & $112$ \\
				\hline 
				$[16; 20)$ &  $8$ & $120$ \\
				 \hline
				& $n=120$ &  \\
				 \hline
			\end{tabular}
		\end{center}
		 Ta có $\dfrac{n}{2}=60$, $\dfrac{n}{4}=30$, $\dfrac{3n}{4}=90$.\\
		Vì  $42< 60< 90$ nên nhóm $3$ là nhóm đầu tiên có tần số tích luỹ lớn hơn hoặc bằng $60$.\\
		Suy ra trung vị là $M_e=8+\left(\dfrac{60-42}{48}\right) \cdot 4=9{,}5$ hay tứ phân vị thứ hai là $Q_2=M_e=9{,}5$.\\ 
		Do $13< 30< 42$ nên nhóm $2$ là nhóm đầu  tiên có tần số tích luỹ lớn hơn hoặc bằng $30$.\\
		Suy ra tứ phân vị thứ nhất  là $Q_1=4+\left(\dfrac{30-13}{29}\right) \cdot 4\approx 6{,}3$.\\
		Do $42< 90\leq 90$  nên nhóm $3$ là nhóm đầu tiên có tần số tích luỹ lớn hơn hoặc bằng $90$.\\
		Suy ra tứ phân vị thứ ba là  $Q_3=8+\left(\dfrac{90-42}{48}\right) \cdot 4=12$.
	 }
\end{ex}

% nguồn: Ok_Toan11_TracNghiem_Lan2
\begin{ex}%[1D5N2-3]%[ma:584272f3]
	[Dự án 11  HVA 2024-2025]
	Doanh thu bán hàng trong $20$ ngày được lựa chọn ngẫu nhiên của một của hàng  được ghi lại ở bảng sau
	\begin{center}
		\begin{tabular}{|c|c|c|c|c|c|}
			\hline 
			Doanh thu & $[5; 7)$ & $[7; 9)$ &  $[9; 11)$ & $[11; 13)$ & $[13; 15)$ \\
			 \hline
			Số ngày & $2$ & $7$ & $7$ &  $3$ & $1$ \\
			\hline 
		\end{tabular}
	\end{center}
	Tứ phân vị thứ nhất của mẫu số liệu gần nhất với giá trị nào trong  các giá trị dưới đây?
	 \choice
		 {$7$}
		{$7{,}6$} 
		 {\True $8$}
		{$8{,}6$} 
	 \loigiai{
		Gọi $x_1$,  $x_2$, $\ldots$, $x_{20}$ là doanh thu bán hàng trong $20$ ngày xếp theo thứ tự không giảm.\\
		Khi đó $x_1$, $x_2 \in[5;  7)$, $x_3$, $\ldots$, $x_9 \in[7; 9)$, $x_9$, $\ldots$, $x_{16} \in[9; 11)$, $x_{17}$, $\ldots$, $x_{19} \in[11; 13)$, $x_{20} \in[13; 15)$. \\
		Nên tứ phân vị thứ nhất của mẫu số liệu thuộc nhóm $[7; 9)$.  \\
		Do đó $n=20$, $n_m=7$, $C=2$, $u_m=7$,  $u_{m+1}=9$. \\
		Ta có  $Q_1=7+\dfrac{\dfrac{1 \cdot 20}{4}-2}{7}(9-7) \approx 7{,}86$.\\
		Trong  các phương án thì $7{,}86$ có kết quả gần với $8$ nhất.
	 }
\end{ex}

% nguồn: Ok_Toan11_TracNghiem_Lan2
\begin{ex}%[1D5N2-3]%[ma:c446810d]
	 Đo chiều cao các em học sinh khối $10$ ta thu được kết quả
	\begin{center}
		\begin{tabular}{|c|c|}
			\hline 
			Chiều cao (cm) & Số học sinh \\ 
			\hline 
			 $[150; 152)$ & $5$ \\
			\hline 
			$[152; 154)$ & $18$ \\ 
			 \hline
			$[154;  156)$ & $40$ \\
			\hline 
			$[156; 158)$ & $26$  \\
			 \hline
			$[158; 160)$ & $8$  \\
			 \hline
			$[160; 162)$ & $3$ \\ 
			\hline 
		\end{tabular}
	\end{center}
	Tứ phân vị thứ nhất của mẫu số liệu ghép  nhóm trên là
	 \choice
		{$152{,}2$} 
		{$153{,}3$} 
		{\True  $154{,}1$}
		 {$151{,}5$}
	 \loigiai{
		Ta  có $Q_1=154+\dfrac{25-23}{40} \cdot 2 \approx 154{,}1$.
	} 
\end{ex}

% nguồn: Ok_Toan11_TracNghiem_Lan2
\begin{ex}%[1D5N2-3]%[ma:fe8ad540]
	Khảo sát thời gian  tập thể dục của một số học sinh khối $11$ thu được mẫu số liệu ghép nhóm sau
	\begin{center}
		\begin{tabular}{|c|c|c|c|c|c|}
			\hline 
			Thời gian (phút) & $[0; 20)$ & $[20; 40)$ &  $[40; 60)$ & $[60; 80)$ & $[80; 100)$ \\
			 \hline
			Số học sinh & $5$ & $9$ & $12$ &  $10$ & $6$ \\
			 \hline
		\end{tabular}
	\end{center}
	Nhóm chứa tứ phân vị thứ ba của mẫu  số liệu trên là
	 \choice
		 {$[40; 60)$}
		{$[20; 40)$} 
		 {\True $[60; 80)$}
		{$[80;  100)$}
	 \loigiai{
		 Ta có $n=42$.\\
		$\dfrac{3\cdot n}{4}=\dfrac{3\cdot 42}{4}=31{,}5$. 
		Nên  tứ phân vị thứ ba của mẫu số liệu trên là $x_{32}$.\\
		Mà $x_{32} \in[60; 80)$.\\ 
		Vậy nhóm chứa tứ phân vị thứ ba của  mẫu số liệu trên là nhóm $[60;80)$.
	} 
\end{ex}

% nguồn: Ok_Toan11_TracNghiem_Lan2
\begin{ex}%[1D5N2-3]%[ma:ecb1ee39]
	Khảo sát thời gian chạy bộ trong một ngày của một số học sinh khối $11$ thu được mẫu số liệu ghép nhóm sau 
	\begin{center}
		\begin{tabular}{|c|c|c|c|c|c|}
			 \hline
			Thời gian (phút) & $[0; 20)$ & $[20; 40)$ & $[40;  60)$ & $[60; 80)$ & $[80; 100)$ \\
			 \hline
			Số học sinh & $5$ & $9$ & $12$ & $10$ & $6$ \\ 
			 \hline
		\end{tabular}
	\end{center}
	Nhóm chứa tứ phân vị  thứ nhất là
	 \choice
		{$[0;  20)$}
		{\True  $[20; 40)$}
		 {$[40; 60)$}
		{$[60;  80)$}
	\loigiai{ 
		$\dfrac{n}{4}=\dfrac{42}{4}=10{,}5$;  $5+9=14$.\\
		Nhóm chứa trung vị thứ  nhất $[20; 40)$.
	} 
\end{ex}

% nguồn: Ok_Toan11_TracNghiem_Lan2
\begin{ex}%[1D5N2-3]%[ma:fe0a68e2]
	Chiều dài của  $60$ lá dương xỉ trường thành
	\begin{center}
		\begin{tabular}{|c|c|}
			\hline 
			Lớp của chiều dài $(\mathrm{cm})$ &  Tần số \\
			\hline 
			$[10; 20)$ & $8$  \\
			 \hline
			$[20;  30)$ & $18$ \\
			\hline 
			 $[30; 40)$ & $24$ \\
			 \hline
			$[40;  50)$ & $10$ \\
			\hline 
		\end{tabular}
	\end{center}
	Tứ phân vị thứ hai của  mẫu số liệu ghép nhóm trên là
	\choice 
		 {\True $\dfrac{95}{3}$}
		 {$\dfrac{91}{3}$}
		 {$\dfrac{89}{3}$}
		 {$\dfrac{93}{3}$}
	 \loigiai{
		Ta có  $Q_2=30+\dfrac{30-26}{24} \cdot 10=\dfrac{95}{3} \approx 31{,}67$.
	 }
\end{ex}

% nguồn: Ok_Toan11_TracNghiem_Lan2
\begin{ex}%[1D5N2-3]%[ma:7a824213]
	Tìm hiểu thời gian hoàn thành một bài tập của một số học sinh thu được  kết quả sau
	\begin{center}
		\begin{tabular}{|c|c|c|c|c|c|}
			 \hline
			Thời gian (phút) & $[0; 4)$  & $[4; 8)$ & $[8; 12)$ & $[12; 16)$ & $[16; 20)$ \\
			\hline 
			Số học sinh & $2$ & $4$ & $7$ &  $4$ & $3$ \\
			 \hline
		\end{tabular}
	\end{center}
	Tứ phân vị thứ ba của mẫu số liệu ghép nhóm  này là
	 \choice
		 {$Q_3=13$}
		{\True  $Q_3=14$}
		{$Q_3=15$} 
		 {$Q_3=12$}
	 \loigiai{
		Cỡ mẫu $n=2+4+7+4+3=20$.\\ 
		Tứ phân vị thứ ba $Q_3$ là $\dfrac{x_{15}+x_{16}}{2}$. Do $x_{15}$, $x_{16}$ đều thuộc nhóm $[12; 16)$ nên  nhóm này chứa $Q_3$.\\
		Do đó $p=4$, $a_4=12$, $m_4=4$, $m_1+m_2+m_3=2+4+7=13$,  $a_5-a_4=4$. \\
		Ta có $Q_3=12+\dfrac{\dfrac{3\cdot 20}{4}-13}{4}\cdot4=14$. 
	 }
\end{ex}

% nguồn: Ok_Toan11_TracNghiem_Lan2
\begin{ex}%[1D5N2-3]%[ma:32015a31]
	Doanh thu bán hàng trong $20$ ngày được lựa chọn ngẫu nhiên của một cửa hàng  được ghi lại ở bảng sau (đơn vị: triệu đồng).
	\begin{center}
		\begin{tabular}{|c|c|c|c|c|c|}
			 \hline
			Doanh thu & $[5; 7)$ & $[7; 9)$ &  $[9; 11)$ & $[11; 13)$ & $[13; 15)$ \\
			 \hline
			Số ngày & $2$ & $7$  & $7$ & $3$ & $1$ \\
			 \hline
		\end{tabular}
	\end{center}
	Tứ phân vị thứ nhất của mẫu số liệu gần nhất với giá trị nào trong các  giá trị dưới đây?
	 \choice
		{$7$} 
		 {$7{,}6$}
		{\True  $8$}
		 {$8{,}6$}
	\loigiai{ 
		Gọi $x_{1}$, $x_{2}$,$\ldots$, $x_{20}$ là doanh thu bán hàng trong $20$ ngày xếp theo thứ tự không  giảm.\\
		Khi đó: $x_{1}$, $x_{2} \in[5; 7)$; $x_{3}$,$\ldots$, $x_{9} \in[7; 9)$; $x_{10}$,$\ldots$, $x_{16} \in[9; 11)$; $x_{17}$,$\ldots$, $x_{19} \in[11; 13)$; $x_{20} \in[13;  15)$.\\
		Do đó, tứ phân vị thứ  nhất của mẫu số liệu thuộc nhóm $[7; 9)$.\\
		Ta có $n=20$, $n_{m}=7$, $C=2$,  $u_{m}=7$, $u_{m+1}=9$.\\
		$Q_{1}=7+\dfrac{\dfrac{1\cdot20}{4}-2}{7}(9-7) \approx  7{,}86$.
	} 
\end{ex}

% nguồn: Ok_TracNghiemLop11_Lan1
\begin{ex}%[1D5N2-3]%[ma:8aabf09a]
	Điều tra $42$ học sinh của một lớp $11$ về số giờ tự  học ở nhà, người ta có bảng sau đây
	\begin{center}
		\begin{tabular}{|c|c|c|}
			\hline 
			Lớp & Tần số & Tần số tích lũy \\ 
			\hline 
			$\left[1; 2\right)$ &  $8$ & $8$ \\
			 \hline
			$\left[2; 3\right)$ & $10$ &  $18$ \\
			\hline 
			$\left[3; 4\right)$ & $12$ & $30$  \\
			 \hline
			$\left[4; 5\right)$  & $9$ & $39$ \\
			 \hline
			$\left[5; 6\right)$  & $3$ & $42$ \\
			 \hline
		\end{tabular}
	\end{center}
	Tứ phân vị thứ nhất của mẫu số liệu trên là 
	\choice 
		{$2{,}25$} 
		{$3{,}25$} 
		{\True $2{,}5$} 
		{$2{,}75$} 
	\loigiai{ 
		Ta có số phần tử của  mẫu là\\
		$n=42\Rightarrow  \dfrac{n}{4}=10{,}5$.\\
		 Suy ra nhóm $2$ là nhóm đầu tiên có tần số tích lũy lớn hơn hoặc bằng $10{,}5$.\\
		Xét nhóm $2$ là nhóm $\left[2; 3\right)$ có $s=2$; $h=1$; $n_2=10$ và nhóm $1$ là nhóm  $\left[1; 2\right)$ có $cf_1=8$.\\
		Áp dụng công thức ta có trung vị  của mẫu số liệu là
		$$Q_1=2+\left(\dfrac{10{,}5-8}{10}  \right)\cdot 1=2{,}25.$$
	} 
\end{ex}

% nguồn: Ok_TracNghiemLop11_Lan1
\begin{ex}%[1D5N2-3]%[ma:6de0f9a3]
	Khảo sát thời gian tập  thể dục của một số học sinh khối 11 thu được mẫu số liệu ghép nhóm sau:
	\begin{center}
		\begin{tabular}{|c|c|c|c|c|c|}
			 \hline
			Thời gian  (Phút) & $\left[0;20\right)$ & $\left[20;40\right)$ &$\left[40;60\right)$ &$\left[60;80\right)$ &$\left[80;100\right)$ \\
			 \hline
			Số học sinh &  $5$ &$9$ &$12$ &$10$ &$6$ \\
			\hline 
		\end{tabular}
	\end{center}
	Nhóm chứa  tứ phân vị thứ ba của mẫu số liệu trên là
	\choice 
		{$\left[40;60\right)$} 
		{$\left[20;40\right)$} 
		{\True $\left[60;80\right)$} 
		{$\left[80;100\right)$} 
	 \loigiai{
		Ta có $n=42$ nên tứ phân vị thứ ba của mẫu số liệu trên là  $Q_3=x_{33}$.\\
		Mà  $x_{33}\in \left[60;80\right)$.\\
		Vậy nhóm chứa tứ phân  vị thứ ba của mẫu số liệu trên là nhóm $\left[60;80\right)$.
	 }
\end{ex}

% nguồn: Ok_TracNghiemLop11_Lan1
\begin{ex}%[1D5N2-3]%[ma:e3c7e487]
	Khảo sát thời gian chạy bộ trong một ngày của một số học sinh khối $11$ thu được mẫu số liệu ghép nhóm sau 
	\begin{center}
		{\tabcolsep =  2mm
		\begin{tabular}{|c|c|c|c|c|c|}
			 \hline
			Thời gian (phút) & $[0; 20)$ & $[20; 40)$ & $[40; 60)$ & $[60; 80)$ & $[80;  100)$ \\
			\hline 
			Số học sinh & $5$ & $9$ & $12$ & $10$ &  $6$\\
			 \hline
		\end{tabular}
		 }
	\end{center}
	Nhóm chứa tứ phân vị thứ nhất của mẫu  số liệu trên là
	\choice 
		{$[40;  60)$}
		 {\True $[20; 40)$}
		{$[60;  80)$}
		{$[80;  100)$}
	 \loigiai{
		Ta có $n=42$ nên tứ phân vị thứ nhất  của mẫu số liệu trên là $Q_1 = x_{11}$.\\
		Mà $x_{11}  \in [20; 40)$.\\
		Vậy  nhóm chứa trung vị của mẫu số liệu trên là nhóm $[20; 40)$.
	} 
\end{ex}

% nguồn: Ok_TracNghiemLop11_Lan1
\begin{ex}%[1D5N2-3]%[ma:0d33c016]
	Khảo sát thời gian  (phút) tập thể dục trong ngày của một số học sinh khối $11$ thu được mẫu số liệu ghép nhóm sau:
	\begin{center}
		\begin{tabular}{|c|c|c|c|c|c|}
			\hline \begin{tabular}{c}
			 Thời gian \\ (phút)
		\end{tabular} & {$[0; 20)$} & {$[20; 40)$} & {$[40; 60)$} & {$[60; 80)$} & {$[80; 100)$} \\
		\hline Số học sinh & $5$ & $9$ &  $12$ & $10$ & $8$ \\
		 \hline
	\end{tabular}
\end{center}
Tứ phân  vị thứ nhất $Q_1$ của mẫu số liệu ghép nhóm này là
\choice 
	 {\True $\dfrac{290}{9}$}
	{$\dfrac{500}{9}$} 
	 {$\dfrac{710}{9}$}
	{$\dfrac{175}{6}$} 
 \loigiai{Cỡ mẫu: $n=5+9+12+10+6=42$.\\
	Nhóm chứa tứ phân vị thứ nhất: $[20; 40)$. Suy ra: $u_m=20$  và $u_{m+1}=40$.\\
	 Ta có $n_m=9, C=5$.\\
	Vậy tứ phân vị  thứ nhất $Q_1$ của mẫu số liệu ghép nhóm là $$Q_1=20+\dfrac{\dfrac{42}{4}-5}{9} \cdot(40-20)=\dfrac{290}{9}.$$
 }
\end{ex}

% nguồn: Ok_TracNghiemLop11_Lan1
\begin{ex}%[1D5N2-3]%[ma:d5dc004d]
	Tổng hợp tiền lương tháng (đơn vị: triệu đồng) của một số nhân văn phòng được ghi lại như  sau:
	\begin{longtable}{|c|c|c|c|c|}
		 \hline
		Lương tháng  (triệu đồng) &$[6;8)$&$[8;10)$&$[10;12)$&$[12;14)$\\
		\hline 
		Số nhân viên&$3$&$6$&$8$&$7$\\ 
		 \hline
	\end{longtable}
	Giá  trị của tứ phân vị thứ nhất bằng
	 \choice
		 {$11$}
		{\True $9$} 
		 {$8$}
		{$10$} 
	 \loigiai{
		Cỡ  mẫu $n=24$.\\
		Gọi $x_1,x_2,\ldots,x_{24}$  lần lượt là số tiền lương của $24$ nhân viên xếp theo thứ tự không giảm.\\
		Tứ phân vị thứ  nhất của mẫu số liệu $x_1,x_2,\ldots,x_{24}$ là $Q_1=\dfrac{1}{2}(x_6+x_7)\in [8;10)$.\\
		Do đó giá trị của tứ phân vị thứ  nhất là $Q_1=8+\dfrac{\dfrac{24}{4}-3}{6}\cdot (10-8)=9$.
	} 
\end{ex}

% nguồn: Ok_TracNghiemLop11_Lan1
\begin{ex}%[1D5N2-3]%[ma:51a69ccc]
	Doanh thu bán hàng trong $20$ ngày được lựa  chọn ngẫu nhiên của một của hàng được ghi lại ở bảng sau (đơn vị: triệu đồng)
	\begin{center}
		{\tabcolsep  = 2mm
		\begin{tabular}{|c|c|c|c|c|c|}
			 \hline
			\textbf{Doanh thu} & $[5; 7)$ & $[7; 9)$ & $[9; 11)$ &  $[11; 13)$ & $[13; 15)$ \\
			\hline \textbf{Số  ngày} & $2$ & $7$ & $7$ & $3$ & $1$\\
			 \hline
		\end{tabular}
		} 
	\end{center}
	Tứ phân vị thứ ba của mẫu số liệu  gần nhất với giá trị nào trong các giá trị dưới đây?
	 \choice
		 {$10$}
		{\True  $11$}
		{$12$} 
		{$13$} 
	\loigiai{ 
	Gọi $x_1, x_2, \cdots , x_{20}$ là doanh thu bán hàng trong $20$  ngày xếp theo thứ tự không giảm.
	Khi đó, $x_1, x_2 \in [5;7)$; $x_3, \cdots  , x_9 \in [7; 9)$; $x_9, \cdots , x_{16} \in [9; 11)$; $x_{17}, x_{18} , x_{19} \in [11; 13)$; $x_{20} \in [13; 15)$.\\
	Do đó, tứ phân vị thứ  ba của mẫu số liệu thuộc nhóm $[9; 11)$.\\
	Ta có $n=20$, $n_m=7$,  $C=2+7=9$, $u_m=9$, $u_{m+1}=11$.
	 $$Q_1 = 9+\dfrac{\dfrac{3 \cdot 20}{4}-9}{7}(11-9) \approx 10{,}71 \approx 11.$$
	} 
\end{ex}

% nguồn: Ok_Toan11_TracNghiem_Lan2
\begin{ex}%[1D5H2-3]%[ma:e6f07b18]
	[Dự án  11 HVA 2024-2025]
	Để kiểm tra thời gian sử dụng pin của chiếc laptop mới, bạn An thống kê thời gian sử dụng laptop của mình từ lúc sạc đầy pin cho tới khi hết pin ở  bảng sau
	\begin{center}
		\begin{tabular}{|c|c|c|c|c|c|}
			 \hline
			Thời gian sử dụng & $[6; 8)$ & $[8; 10)$ & $[10; 12)$ & $[12; 14)$ & $[14;  16)$ \\
			\hline 
			Số lần sạc &  $2$ & $5$ & $8$ & $5$ & $1$ \\
			 \hline
		\end{tabular}
	\end{center}
	Tứ phân vị thứ nhất của mẫu  số là giá trị nào trong các giá trị sau?
	 \choice
		{\True $9{,}3$} 
		{$10{,}3$} 
		 {$10{,}2$}
		 {$9{,}2$}
	 \loigiai{
		Gọi $x_1$, $x_2$, $\ldots$, $x_{21}$ là thời gian sử dụng được xếp theo thứ tự không  giảm.\\
		Tứ phân vị thứ nhất của dãy số liệu $x_1$, $x_2$, $\ldots$, $x_{21}$ là $\dfrac{1}{2}\left(x_5+x_6\right) \in[8; 10)$. Do đó tứ phân vị thứ nhất của mẫu số liệu  ghép nhóm trên là
		 $$Q_1=8+\dfrac{\dfrac{1\cdot 21}{4}-2}{5} \cdot(10-8)=9{,}3.$$
	} 
\end{ex}

% nguồn: Ok_Toan11_TracNghiem_Lan2
\begin{ex}%[1D5H2-3]%[ma:a0277e98]
	Doanh thu bán  hàng trong 20 ngày được lựa chọn ngẫu nhiên của một của hàng được ghi lại ở bảng sau (đơn vị: triệu đồng):
	\begin{longtable}{|*{6}{c|}}
		\hline 
		Doanh thu & [$5$; $7$) & [$7$; $9$) & [$9$; $11$)  & [$11$; $13$) & [$13$; $15$) \\
		 \hline
		Số ngày &  $2$& $7$ & $7$ & $3$ & $1$ \\
		 \hline
	\end{longtable}
	Tứ phân vị thứ ba của mẫu số liệu gần nhất với giá  trị nào trong các giá trị dưới đây?
	\choice 
		{$10$} 
		 {\True $11$}
		{$12$} 
		{$13$} 
	\loigiai{ 
		Gọi $x_1, x_2, \ldots, x_{20}$ là  doanh thu bán hàng trong 20 ngày xếp theo thứ tự không giảm.\\
		Khi đó: $x_1, x_2 \in \left[5; 7\right)$, $x_3, \ldots, x_9 \in \left[7; 9\right)$, $ x_9, \ldots, x_{16} \in \left[9; 11\right)$, $x_{17}, \ldots, x_{19} \in [11; 13)$, $x_{20} \in  \left[13; 15\right)$.\\
		Do đó tứ phân vị thứ ba của mẫu số liệu thuộc nhóm $\left[9; 11\right)$ 
		\begin{eqnarray*}
			&n=20,n_m=7,C=9,u_m=9,u_{m+1}=11;  \\
			&Q_3=9+\dfrac{\dfrac{3\cdot20}{4}-9}{7} \left(11-9\right)\approx 10{,}71\approx  11.
		\end{eqnarray*}
	 }
\end{ex}

% nguồn: Ok_Toan11_TracNghiem_Lan2
\begin{ex}%[1D5H2-3]%[ma:f3f0aa91]
	Người ta tiến hành phỏng vấn $40$ người về một mẫu áo khoác. Người điều tra yêu cầu cho điểm mẫu áo đó theo  thang điểm là $100$. Kết quả được trình bày trong bảng ghép nhóm sau
	\begin{center}
		\begin{tabular}{|c|c|c|c|c|c|c|}
			 \hline
			Nhóm & $[50; 60)$ & $[60; 70)$ & $[70; 80)$ & $[80; 90)$  & $[90; 100)$ & \\
			 \hline
			Tần số &  $4$ & $5$ & $23$ & $6$ & $2$ & $n=40$ \\
			\hline 
		\end{tabular}
	\end{center}
	Tứ phân  vị của mẫu số liệu ghép nhóm trên (làm tròn đến hàng đơn vị) là
	 \choice
		{$Q_{1} \approx 71$; $Q_{2} \approx 76$; $Q_{3}  \approx 78$}
		{$Q_{1}  \approx 71$; $Q_{2} \approx 75$; $Q_{3} \approx 78$}
		{$Q_{1} \approx 70$; $Q_{2} \approx 76$; $Q_{3} \approx 79$} 
		 {\True $Q_{1} \approx 70$; $Q_{2} \approx 75$; $Q_{3} \approx 79$}
	\loigiai{ 
		$Q_{1}=70+\dfrac{10-9}{23}  \cdot 10=70$; $Q_{2}=70+\dfrac{20-9}{23} \cdot 10=75$; $Q_{3}=70+\dfrac{30-9}{23} \cdot 10=79$.
	} 
\end{ex}

% nguồn: Ok_Toan11_TracNghiem_Lan2
\begin{ex}%[1D5H2-3]%[ma:db031d0b]
	Đo chiều cao của $500$ học sinh  trong một trường THPT ta thu được kết quả như sau:
	\begin{center}
		\begin{tabular}{|c|c|c|c|c|c|}
			\hline 
			Chiều cao & $\left[150; 154\right)$ & $\left[154; 158\right)$ & $\left[158; 162\right)$ & $\left[162; 166\right)$ & $\left[166; 170\right)$ \\ 
			\hline 
			Tần số & 25 & 50 & 200 & 175  & 50 \\
			\hline 
		\end{tabular}
	\end{center}
	Tứ phân vị thứ  nhất của mẫu số liệu trên thuộc khoảng nào trong các khoảng dưới đây
	 \choice
		 {$\left[150; 154\right)$}
		{$\left[154; 158\right)$} 
		{\True $\left[158; 162\right)$} 
		{$\left[162;  166\right)$}
	 \loigiai{
		Gọi $x_1$, $x_2$, $\ldots$, $x_{500}$  là chiều cao của $500$ học sinh xếp theo thứ tự không giảm. Tứ phân vị thứ nhất của mẫu số liệu là $\dfrac{1}{2} \left(x_{125}+x_{126} \right)\in \left[158; 162\right)$.
	 }
\end{ex}

% nguồn: Ok_Toan11_TracNghiem_Lan2
\begin{ex}%[1D5H2-3]%[ma:41c98e70]
	[Dự án 11  HVA 2024-2025]
	Cho mẫu số liệu ghép nhóm về thống kê điểm số của $100$ học sinh tham dự kỳ thi  học sinh giỏi toán, ta có bảng số liệu sau
	\begin{center}
		\begin{tabular}{|c|c|c|c|c|c|c|}
			 \hline
			Điểm & $[8; 10)$ & $[10; 12)$ &  $[12; 14)$ & $[14; 16)$ & $[16; 18)$ & $[18; 20)$ \\
			 \hline
			Số học sinh  & $6$ & $21$ & $30$ & $25$ & $14$ & $4$ \\
			\hline 
		\end{tabular}
	\end{center}
	Tìm tứ phân vị thứ ba của mẫu số liệu ghép nhóm  trên.
	 \choice
		 {\True $15{,}44$}
		 {$13{,}29$}
		 {$13{,}92$}
		 {$14{,}54$}
	 \loigiai{
		 Cỡ mẫu $n=6+21+30+25+14+4=100$.\\
		Tứ phân vị thứ ba  $Q_3$ là $\dfrac{x_{75}+x_{76}}{2}$.\\
		Do $x_{75}$, $x_{76}$ đều thuộc nhóm  $[14; 16)$ nên nhóm này chứa $Q_3$.\\
		 Do đó $p=4$, $a_4=14$, $m_4=25$, $m_1+m_2+m_3=2+4+7=57$, $a_5-a_4=2$.\\
		Ta có $Q_3=14+\dfrac{\dfrac{3  \cdot 100}{4}-27}{25} \cdot 2=15{,}44$.
	 }
\end{ex}

% nguồn: Ok_Toan11_TracNghiem_Lan2
\begin{ex}%[1D5H2-3]%[ma:0e597521]
	Khảo sát thời gian tập thể dục của một số học sinh khối 11 thu được mẫu số liệu ghép nhóm  sau
	\begin{longtable}{|c|c|c|c|c|c|}
		\hline 
		Thời gian (phút) &  $[0; 20)$ & $[20; 40)$ & $[40; 60)$ & $[60; 80)$ & $[80; 100)$ \\
		 \hline
		 Số học sinh & $5$ & $9$ & $12$ & $10$ & $6$ \\
		\hline 
	\end{longtable}
	Nhóm chứa  tứ phân vị thứ ba của mẫu số liệu trên là
	 \choice
		 {$\left[40; 60\right)$}
		 {$\left[20; 40\right)$}
		{\True  $\left[60; 80\right)$}
		 {$\left[80; 100\right)$}
	\loigiai{ 
		Ta có $n=42$ nên tứ phân  vị thứ ba của mẫu số liệu trên là $Q_3=x_{33}$.\\
		Mà $x_{33}  \in \left[60; 80\right)$.\\
		Vậy nhóm chứa tứ phân vị thứ  ba của mẫu số liệu trên là nhóm $[60; 80)$.
	} 
\end{ex}

% nguồn: Ok_Toan11_TracNghiem_Lan2
\begin{ex}%[1D5H2-3]%[ma:d70a3873]
	Một  phòng khám thống kê số bệnh nhân đến khám mỗi ngày trong tháng $6$ năm $2020$ ở bảng sau:
	\begin{center}
		\begin{tabular}{|c|c|c|c|c|c|}
			\hline 
			Số bệnh nhân  &$ \left[1;10\right)$ & $\left[10;20\right)$ & $\left[20;30\right)$ & $\left[30;40\right)$ & $\left[40;50\right)$ \\
			 \hline
			Số người & $7$ & $5$ &  $10$ & $6$ & $2$ \\
			\hline 
		\end{tabular}
	\end{center}
	Số bệnh nhân đến khám nhiều hơn $23$ người trên ngày chiếm tỉ lệ  bao nhiêu phần trăm?
	 \choice
		{$25\%$} 
		{\True $50\%$} 
		 {$75\%$}
		{$10\%$} 
	 \loigiai{
		\begin{itemize}
			\item Gọi $x_1$, $x_2$, $\cdots$, $x_{30}$ là số bệnh nhân đến khám mỗi ngày và giả sử dãy này đã được sắp xếp theo thứ tự không giảm. Khi đó  tứ phân vị thứ hai của dãy số liệu là $\dfrac{x_{15}+x_{16}}{2}\in\left[20;30\right)$.
			\item $Q_2=20+\dfrac{\dfrac{30}{2}-\left(7+5\right)}{10}\cdot10=23$. 
		\end{itemize}
		Như vậy có  khoảng $50\%$ số ngày có nhiều hơn $23$ bệnh nhân đến khám.
	 }
\end{ex}

% nguồn: Ok_Toan11_TracNghiem_Lan2
\begin{ex}%[1D5H2-3]%[ma:fd9f77ea]
	Cân nặng của một người trưởng thành được lựa chọn ngẫu nhiên trong $30$ người được  ghi lại ở bảng sau:
	\begin{center}
		\begin{tabular}{|c|c|c|c|c|c|}
			\hline 
			Cân nặng &$ \left[50;60\right)$ & $\left[60;70\right)$ & $\left[70;80\right)$ &  $\left[80;90\right)$ & $\left[90;100\right)$ \\
			\hline 
			Số người & $7$ & $16$ &  $4$ & $2$ & $1$ \\
			\hline 
		\end{tabular}
	\end{center}
	Tứ phân  vị thứ nhất của mẫu số liệu trên gần nhất với giá trị nào trong các giá trị sau?
	 \choice
		 {$63$}
		{\True  $60$}
		{$64$} 
		{$64$} 
	\loigiai{ 
		Ta có tứ phân vị thứ  nhất thuộc khoảng $\left[60;70\right)$
		Nên $Q_1=60+\dfrac{\dfrac{30}{4}-7}{16}\cdot10=\dfrac{965}{16} \approx 60{,}31$. 
	 }
\end{ex}

% nguồn: Ok_Toan11_TracNghiem_Lan2
\begin{ex}%[1D5H2-3]%[ma:31e641ab]
	Tổng lượng mưa trong tháng $8$ đo được tại một trạm quan trắc đặt tại  Vũng Tàu từ năm $2002$ đến năm $2020$ được ghi lại như dưới đây (đơn vị: mm)
	\begin{center}
		\begin{tabular}{|c|c|c|c|c|c|c|c|c|c|}
			\hline 
			$121{,}8$ & $158{,}3$ & $334{,}9$ & $200{,}9$ & $165{,}6$ & $161{,}5$ &  $194{,}3$ & $220{,}7$ & $189{,}8$ & $234{,}2$ \\
			\hline 
			$165{,}9$ & $165{,}9$ & $134$  & $173$ & $169$ & $189$ & $254$ & $168$ & $255$ & \\
			\hline 
		\end{tabular}
		\begin{flushright}
			(Nguồn: Tổng cục Thống  kê)
		\end{flushright}
	\end{center}
	Hoàn thiện bảng tần số ghép nhóm theo mẫu sau và tìm tứ phân vị thứ  hai của mẫu số liệu góp nhóm đó.
	\begin{center}
		\begin{tabular}{|c|c|c|c|c|}
			 \hline
			Tổng lượng mưa trong tháng $8$ (mm) & $[120; 175)$  & $[175; 230)$ & $[230; 285)$ & $[285; 340)$ \\
			\hline 
			Số năm & $x$ & $y$  & $z$ & $t$ \\
			 \hline
		\end{tabular}
	\end{center}
	 \choice
		{\True $x=10$; $y=5$; $z=3$; $t=1$;  $Q_{2}=172{,}5$}
		{$x=9$;  $y=6$; $z=3$; $t=1$; $Q_{2}=172{,}5$}
		{$x=10$; $y=5$;  $z=2$; $t=2$; $Q_{2}=182{,}5$}
		{$x=10$; $y=4$;  $z=4$; $t=1$; $Q_{2}=162{,}5$}
	 \loigiai{
		\begin{center}
			\begin{tabular}{|c|c|c|c|c|}
				\hline 
				Tổng lượng mưa trong tháng $8$ (mm) &  $[120; 175)$ & $[175; 230)$ & $[230; 285)$ & $[285; 340)$ \\
				\hline 
				Số  năm & $10$ & $5$ & $3$ & $1$ \\
				 \hline
			\end{tabular}
		\end{center}
		Gọi $x_{1}$, $x_{2}$, $x_{3}$,$\ldots$, $x_{19}$ lần lượt là số năm theo thứ tự  không giảm.\\
		Tứ phân vị thứ hai của dãy số liệu là $x_{10}$ thuộc nhóm $[120; 175)$ nên tứ phân  vị thứ hai của mẫu số liệu là $Q_{2}=120+\dfrac{\dfrac{19}{2}-0}{10}(175-120)=172{,}5$.
	} 
\end{ex}

% nguồn: Ok_Toan11_TracNghiem_Lan2
\begin{ex}%[1D5H2-3]%[ma:57044ab0]
	Thời gian truy cập Internet mỗi buổi  tối của một số học sinh được cho trong bảng sau
	\begin{longtable}{|c|c|c|c|c|c|}
		 \hline
		Thời gian (phút) & $[9{,}5; 12{,}5)$ & $[12{,}5; 15{,}5)$ & $[15{,}5; 18{,}5)$ & $[18{,}5; 21{,}5)$  & $[21{,}5; 24{,}5)$ \\
		 \hline
		Số  học sinh & $3$ & $12$ & $15$ & $24$ & $2$ \\
		\hline 
	\end{longtable}
	Tính trung vị của mẫu số liệu  ghép nhóm này
	 \choice
		 {$17$}
		{\True $18{,}1$} 
		{$19{,}3$} 
		{$18{,}5$} 
	\loigiai{ 
		Cỡ mẫu  là $n=3+12+15+24+2=56$.\\
		Gọi $x_1, \ldots, x_{56}$ là thời gian vào internet của 56 học sinh và giả sử dãy này được sắp xếp  theo thứ tự tăng dần. \\
		Khi  đó: trung vị là $\dfrac{x_{28}+x_{29}}{2}$. \\
		Do 2 giá trị $x_{28}, x_{29}$  thuộc nhóm $\left[15{,}5; 18{,}5\right)$ nên nhóm này chứa trung vị. \\
		Do đó  $p=3$; $ a_3=15{,}5$; $ m_3=15$; $ m_1+m_2=3+12=15$; $ a_4-a_3=3$ và ta có
		\begin{align*}
			M_e=15{,}5+\dfrac{\dfrac{56}{2}-15}{15}  \cdot 3=18{,}1.
		\end{align*}
	 }
\end{ex}

% nguồn: Ok_Toan11_TracNghiem_Lan2
\begin{ex}%[1D5H2-3]%[ma:2c4b50ac]
	Kiểm tra điện lượng của một số viên pintiểu do một hãng  sản xuất thu được kết quả sau
	\begin{longtable}{|p{2.6cm}|*{5}{c|}}
		 \hline
		Điện lượng  \newline (nghìn mAh)
		& [$0{,}9$; $0{,}95$) & [$0{,}95$; $1{,}0$) & [$1{,}0$; $1{,}05$) & [$1{,}05$; $1{,}1$) & [$1{,}1$;  $1{,}15$) \\
		\hline 
		Số viên pin & $10$ & $20$ &  $35$ & $15$ & $5$ \\
		 \hline
	\end{longtable}
	Số trung vị  của mẫu số liệu ghép nhóm trên là
	 \choice
		 {\True $1{,}02$}
		 {$1{,}07$}
		 {$1{,}15$}
		 {$1{,}2$}
	\loigiai{ 
		Số viên  pin: $n=10+20+35+15+5=85$.\\
		Gọi $x_1$, $ x_2$, $ x_3$, $ \ldots$, $ x_{85}$ lần lượt là điện lượng  của các viên pintheo thứ tự không giảm.\\
		Do $x_1,  \ldots, x_{10} \in \left[0{,}9; 0{,}95\right)$;
		 $x_{11}, \ldots, x_{30} \in \left[0{,}95; 1{,}0\right)$;
		$x_{31}, \ldots, x_{65}  \in \left[1{,}0; 1{,}05\right)$;
		$x_{66}, \ldots,  x_{80} \in \left[1{,}05; 1{,}1\right)$;
		$x_{81}, \ldots, x_{85} \in  \left[1{,}1; 1{,}15\right)$ nên
		trung vị của dãy số  liệu là $x_{43}$ thuộc nhóm $\left[1{,}0; 1{,}05\right)$.\\
		Ta xác định được $n=85$; $ p=3$; $ m_3=35$; $  m_1+m_2=10+20=30$; $a_3=1$; $ a_4=1{,}05$.\\
		 Vậy trung vị của mẫu số liệu ghép nhóm là $M_e=1{,}0+\dfrac{\dfrac{85}{2}-30}{35} \left(1{,}05-1{,}0\right)=1{,}02$.
	} 
\end{ex}

% nguồn: Ok_Toan11_TracNghiem_Lan2
\begin{ex}%[1D5H2-3]%[ma:dcbed4ee]
	Ghi lại tốc độ bóng trong 200 lần giao bóng của  một vận động viên môn quần vợt cho kết quả như bảng sau
	\begin{longtable}{|c|c|}
		 \hline
		Tốc  độ $v$ (km/h) & Số lần \\
		\hline 
		$150 \le v < 155$  & $18$ \\
		\hline 
		$155 \le v < 160$ & $28$ \\ 
		 \hline
		$160 \le v <  165$ & $35$ \\
		 \hline
		$165 \le v < 170$  & $43$ \\
		 \hline
		 $170 \le v < 175$ & $41$ \\
		\hline 
		$175 \le v < 180$  & $35$ \\
		\hline 
	\end{longtable}
	Trung vị của mẫu số liệu  ghép nhóm này là
	 \choice
		{\True  $167{,}21$}
		 {$170{,}21$}
		 {$175$}
		{$165$} 
	\loigiai{ 
		Cỡ mẫu là $n=18+28+35+43+41+35=200$.\\ 
		Gọi $x_1, \ldots,  x_{200}$ là tốc độ giao bóng của $200$ lần và giả sử dãy này được sắp xếp theo thứ tự tăng dần.\\
		Khi đó: trung vị là $\dfrac{x_{100}+x_{101}}{2}$. Do 2 giá trị $x_{100},  x_{101}$ thuộc nhóm $\left[165; 170\right)$ nên nhóm này chứa trung vị.\\
		Suy ra $p=4$; $ a_4=165$; $ m_4=43$; $ m_1+m_2+m_3=18+28+35=81$; $ a_5-a_4=170-165=5$.\\ 
		 Ta có trung vị $M_e=165+\dfrac{\dfrac{200}{2}-81}{43} \cdot 5=167{,}21$.
	} 
\end{ex}

% nguồn: Ok_Toan11_TracNghiem_Lan2
\begin{ex}%[1D5H2-3]%[ma:1a26c01a]
	Số điểm một cầu thủ bóng rổ  ghi được trong 20 trận đấu được cho ở bảng sau:
	\begin{center}
		\begin{tabular}{|c|c|c|c|c|c|c|c|c|c|}
			 \hline
			 25&23&21&13&8&14&15&18&22&11\\
			\hline 
			24&12&14&14&18&6&8&25&10&11\\ 
			 \hline
		\end{tabular}
	\end{center}
	Tứ phân vị của dãy số liệu trên lần  lượt là
	\choice 
		{\True $Q_1=11;Q_2=14;Q_3=21{,}5$} 
		 {$Q_1=10;Q_2=14;Q_3=21$}
		{$Q_1=12;Q_2=15;Q_3=22$} 
		 {$Q_1=11;Q_2=14;Q_3=23$}
	 \loigiai{
		Ta có dãy điểm số  xếp theo thứ tự không giảm:
		$$6;8;8;10;11;11;12;13;14;14;14;15;18;18;21;22;23;24;25;25.$$ 
		Tứ phân vị thứ nhất là  $11$.\\
		Tứ phân vị thứ hai là  $14$.\\
	Tứ phân vị thứ ba là  $21{,}5$.}
\end{ex}

% nguồn: Ok_Toan11_TracNghiem_Lan2
\begin{ex}%[1D5H2-3]%[ma:c24ad983]
	Tìm tứ phân vị thứ nhất cho  mẫu số liệu ghép nhóm sau:
	\begin{center}
		\begin{tabular}{|c|c|}
			\hline 
			 Tốc độ $v$ (km/h)& Số lần \\
			\hline 
			$150\le v<155$& $18$  \\
			 \hline
			$155\le v<160$& $28$ \\ 
			\hline 
			$160\le v<165$&  $35$\\
			 \hline
			$165\le  v<170$& $43$\\
			 \hline
			$170\le  v<175$& $41$ \\
			\hline 
			$175\le  v<180$& $35$ \\
			\hline 
		\end{tabular}
	\end{center}
	 \choice
		 {$150{,}57$}
		 {\True $160{,}57$}
		{$170{,}57$} 
		{$180{,}57$} 
	 \loigiai{
		\begin{itemize}
			\item Cỡ mẫu là $n=200$. 
			\item Gọi  $x_1$, $x_2$, $\cdots$, $x_{200}$ lần lượt là tốc độ $v$ theo thứ tự không giảm.
			\item Ta có $x_1$, $x_2$, $\cdot$, $x_{18}\in\left[150;155\right)$; $x_{19}$, $x_{20}$, $\cdots$, $ x_{46}\in\left[155;160\right)$; $x_{47}$, $x_{48}$, $\cdots$, $x_{81}\in\left[160;165\right)$; $x_{82}$, $x_{83}$, $\cdots$,  $x_{124}\in\left[165;170\right)$; $x_{125}$, $x_{126}$, $\cdot$, $x_{165}\in\left[170;175\right)$; $x_{166}$, $x_{167}$, $\cdot$, $x_{200}\in \left[175;180\right)$.
			\item Tứ phân vị thứ nhất  của mẫu số liệu $Q_1=\dfrac{1}{2}\left(x_{50}+x_{51}\right)\in\left[160;165\right)$ nên
			\[Q_1=160+\dfrac{\dfrac{1\cdot 200}{4}-\left(18+28\right)}{35}\cdot  \left(165-160\right)=160{,}57.\]
		\end{itemize}
	} 
\end{ex}

% nguồn: Ok_Toan11_TracNghiem_Lan2
\begin{ex}%[1D5H2-3]%[ma:d57829b4]
	Đo chiều cao của $500$ học sinh trong một trường THPT ta thu  được kết quả như sau:
	\begin{center}
		\begin{tabular}{|c|c|c|c|c|c|}
			\hline 
			Chiều cao & $\left[150; 154\right)$ & $\left[154; 158\right)$ & $\left[158; 162\right)$ & $\left[162; 166\right)$  & $\left[166; 170\right)$ \\
			\hline 
			Tần số & $25$ & $50$ & $200$ & $175$ &  $50$ \\
			 \hline
		\end{tabular}
	\end{center}
	Tứ phân vị thứ hai của mẫu số  liệu trên gần với giá trị nào trong các giá trị sau?
	 \choice
		 {$161{,}4$}
		 {\True $161{,}5$}
		 {$160{,}8$}
		{$165{,}5$} 
	 \loigiai{
		Gọi $x_1$, $x_2$, $\ldots$, $x_{500}$ là chiều cao của $500$ học sinh xếp theo thứ tự không giảm. Tứ phân vị thứ hai của mẫu số liệu  là $\dfrac{1}{2} \left(x_{250}+x_{251} \right)\in \left[158; 162\right)$. Do đó tứ phân vị thứ hai của mẫu số liệu ghép nhóm là
		\[Q_2=158+\dfrac{\dfrac{2\cdot500}{4}-75}{200} \left(162-158\right)=161{,}5.\] 
	 }
\end{ex}

% nguồn: Ok_Toan11_TracNghiem_Lan2
\begin{ex}%[1D5H2-3]%[ma:1dba141d]
	Trong một đề tài nghiên cứu về bệnh A, người ta ghi lại lại tuổi của  bệnh nhân mắc bệnh này, số liệu thông kê được trình bày trong bảng sau
	\begin{longtable}{|*{6}{c|}}
		\hline 
		Độ tuổi & [$15$; $25$) & [$25$; $35$) & [$35$;  $45$) & [$45$; $55$) & [$55$; $65$) \\
		\hline 
		Số bệnh nhân & $10$ & $12$ & $14$ & $9$  & $5$ \\
		 \hline
	\end{longtable}
	Trung vị của  mẫu số liệu ghép nhóm này bằng
	 \choice
		{$27{,}14$} 
		{$47{,}14$} 
		{\True $37{,}14$} 
		 {$57{,}14$}
	\loigiai{ 
		 Cỡ mẫu là $n=10+12+14+9+5=50$.\\
		Gọi $x_1, \ldots, x_{50}$ là độ tuổi của $50$ bệnh nhân và giả sử dãy này được sắp xếp  theo thứ tự tăng dần.\\
		Khi đó: trung vị là $\dfrac{x_{25}+x_{26}}{2}$. Do $2$ giá trị $x_{25}, x_{26}$ thuộc nhóm $\left[35; 45\right)$ nên nhóm  này chứa trung vị. Suy ra $p=3$; $ a_3=35$; $ m_3=14$; $ m_1+m_2=10+12=22$; $a_4-a_3=45-35=10$.\\
		Ta  có trung vị $M_e=35+\dfrac{\dfrac{50}{2}-22}{14} \cdot 10\approx 37{,}14$.
	} 
\end{ex}

% nguồn: Ok_Toan11_TracNghiem_Lan2
\begin{ex}%[1D5H2-3]%[ma:cb25d4a0]
	Lương tháng của một số nhân viên một văn  phòng được ghi lại như sau (đơn vị: triệu đồng).
	\begin{center}
		\begin{tabular}{|c|c|c|c|c|}
			 \hline
			Lương tháng (triệu đồng) & $[6; 8)$ & $[8; 10)$ & $[10; 12)$ & $[12;  14)$ \\
			 \hline
			Số nhân viên & $3$ & $6$ & $8$ & $7$ \\ 
			 \hline
		\end{tabular}
	\end{center}
	Tìm tứ phân  vị của dãy số liệu trên.
	 \choice
		{\True $Q_{1}=9$;  $Q_{2}=10{,}75$; $Q_{3}=12{,}3$}
		{$Q_{1}=9$; $Q_{2}=10{,}75$;  $Q_{3}=14{,}3$}
		 {$Q_{1}=9$; $Q_{2}=11{,}75$; $Q_{3}=12{,}3$}
		{$Q_{1}=10$;  $Q_{2}=10{,}75$; $Q_{3}=12{,}3$}
	\loigiai{ 
		Gọi $x_{1}$, $x_{2}$, $x_{3}$,$\ldots$, $x_{24}$ lần lượt là số nhân viên theo thứ tự không giảm.\\ 
		Do $x_{1}$,$\ldots$, $x_{3} \in[6; 8)$; $x_{4}$,$\ldots$, $x_{9} \in[8; 10)$; $x_{10}$,$\ldots$, $x_{17} \in[10; 12)$; $x_{18}$,$\ldots$,  $x_{24} \in[12; 14)$.\\
		Tứ phân vị thứ hai của dãy số liệu là $\dfrac{1}{2}(x_{12}+x_{13})$ thuộc nhóm $[10;  12)$ nên tứ phân vị thứ hai của mẫu số liệu là $Q_{2}=10+\dfrac{\dfrac{24}{2}-9}{8}(12-10)=10{,}75$.\\
		Tứ phân vị thứ nhất của dãy số liệu là $\dfrac{1}{2}(x_{6}+x_{7})$ thuộc nhóm $[8; 10)$ nên tứ phân vị thứ nhất của mẫu  số liệu là $9$.\\
		Tứ phân vị thứ ba của dãy số liệu là $\dfrac{1}{2}(x_{18}+x_{19})$ thuộc nhóm $[12; 14)$ nên tứ  phân vị thứ ba của mẫu số liệu là $Q_{3}=12+\dfrac{\dfrac{3\cdot24}{4}-17}{7}(14-12)=12{,}3$.
	} 
\end{ex}

% nguồn: Ok_Toan11_TracNghiem_Lan2
\begin{ex}%[1D5H2-3]%[ma:baaa9251]
	Doanh thu bán hàng trong $20$  ngày được lựa chọn ngẫu nhiên của một cửa hàng được ghi lại ở bảng sau:
	\begin{center}
		\begin{tabular}{|c|c|c|c|c|c|}
			 \hline
			Doanh thu& $\left[5; 7\right)$ & $\left[7; 9\right)$ & $\left[9; 11\right)$ & $\left[11; 13\right)$ & $\left[13; 15\right)$ \\ 
			 \hline
			Số ngày &  $2$ & $7$ & $7$ & $3$ & $1$ \\
			\hline 
		\end{tabular}
	\end{center}
	Tứ phân vị thứ ba của mẫu số liệu trên gần nhất với giá trị  nào dưới đây?
	 \choice
		{$10$} 
		{\True  $11$}
		 {$12$}
		 {$13$}
	\loigiai{ 
		\begin{itemize}
			\item Sắp xếp các giá trị (số ngày)  theo thứ tự không giảm $x_1$, $x_2$, $\cdots$, $x_{20}$. Vì $n=20$ là số chẵn nên $Q_2=\dfrac{1}{2}\left(x_{10}+x_{11}\right)$.
			\item Tứ phân vị thứ ba là trung vị của nửa số liệu bên phải của mẫu số liệu ban đầu  gồm $x_{11}$, $x_{12}$, $\cdots$, $x_{20}$. Vì $n=10$ là số chẵn nên $Q_3=\dfrac{1}{2}\left(x_{15}+x_{16}\right)$.
			\item Ta có $x_{15}$, $x_{16}\in\left[9;11\right)$ nên $Q_1=9+\dfrac{\dfrac{3\cdot20}{4}-9}{7}\left(11-9\right)\approx 10{,}71$. 
		\end{itemize}
	 }
\end{ex}

% nguồn: Ok_Toan11_TracNghiem_Lan2
\begin{ex}%[1D5H2-3]%[ma:60e33fcb]
	[Dự án 11 HVA 2024-2025] 
	 Khảo sát về cân nặng của các học sinh lớp $11D3$ người ta được một mẫu dữ liệu ghép nhóm như sau
	\begin{center}
		\begin{tabular}{|c|c|c|c|c|c|c|}
			 \hline
			Cân nặng & $[30; 40)$ & $[40; 50)$ & $[50;  60)$ & $[60; 70)$ & $[70; 80)$ & $[80; 90)$ \\
			 \hline
			Số học sinh & $2$ & $10$ & $16$ & $8$  & $2$ & $2$ \\
			 \hline
		\end{tabular}
	\end{center}
	Tứ phân vị thứ  nhất của mẫu số liệu trên là
	\choice 
		{$50$} 
		{\True $48$} 
		{$40$} 
		 {$45$}
	 \loigiai{
		Ta có Số  phần tử của mẫu là $n=40\Rightarrow \dfrac{n}{4}=10$.\\
		 Suy ra nhóm $[40; 50)$ chứa tứ phân vị thứ nhất.\\
		Do đó tứ phân vị thứ  nhất của mẫu số liệu trên là $Q_1=40+\dfrac{10-2}{10} \cdot 10=48$.
	} 
\end{ex}

% nguồn: Ok_Toan11_TracNghiem_Lan2
\begin{ex}%[1D5H2-3]%[ma:327885d5]
	Chiều cao của một số cây được cho trong  bảng sau:
	\begin{center}
		\begin{tabular}{|c|c|c|c|c|c|}
			 \hline
			Chiều cao (cm) & $\left[0;10\right)$ & $\left[10;20\right)$  & $\left[20;30\right)$ & $\left[30;40\right)$ & $\left[40;50\right)$ \\
			 \hline
			Số cây & $4$ & $6$ & $7$ & $5$ &  $3$ \\
			 \hline
		\end{tabular}
	\end{center}
	Tứ phân vị thứ nhất của mẫu  số liệu ghép nhóm này là
	\choice 
		{$Q_1=13{,}5$} 
		{$Q_1=13{,}9$} 
		 {$Q_1=15{,}75$}
		 {\True $Q_1=13{,}75$}
	 \loigiai{
		\begin{itemize}
			 \item Cỡ mẫu $n=4+6+7+5+3=25$.
			\item Giả sử $x_1$,  $x_2$, $\cdots$, $x_25$ lần lượt là chiều cao của $25$ cây được sắp xếp theo thứ tự không giảm.
			\item Ta có  $Q_1=x_7\in\left[10;20\right)$. Suy ra
			\[Q_1=10+\dfrac{1\cdot\dfrac{25}{4}-4}{6}\cdot\left(20-10\right)=13{,}75.\] 
		\end{itemize}
	} 
\end{ex}

% nguồn: Ok_Toan11_TracNghiem_Lan2
\begin{ex}%[1D5H2-3]%[ma:9f19e54e]
	Lương tháng của một số nhân viên văn phòng được  ghi lại như sau
	\begin{center}
		\begin{tabular}{|c|c|c|c|c|}
			 \hline
			Lương tháng (triệu đồng) &$ \left[6;8\right)$ & $\left[8;10\right)$ & $\left[10;12\right)$ & $\left[12;14\right)$ \\ 
			\hline 
			Số nhân viên & $3$ &  $6$ & $8$ & $7$ \\
			\hline 
		\end{tabular}
	\end{center}
	 Tứ phân vị của dãy số liệu trên lần lượt là
	\choice 
		{$Q_1=9$; $Q_2=10$,  $Q_3\approx 12$}
		 {\True $Q_1=9$; $Q_2=10{,}75$, $Q_3\approx 12{,}3$}
		{$Q_1=9$; $Q_2=11$, $Q_3\approx  13$}
		{$Q_1=10{,}73$; $Q_2=11$,  $Q_3\approx 13$}
	 \loigiai{
		\begin{itemize}
			\item Gọi  $x_1$, $x_2$, $\cdots$, $x_{24}$ lần lượt là số nhân viên có lương tháng theo thứ tự không giảm.
			\item Ta có  $x_1$, $x_2$, $x_3\in\left[6;8\right)$; $x_4$, $x_5$, $\cdots$, $x_9\left[8;10\right)$; $x_{10}$, $x_{11}$, $\cdots$, $x_{17}\left[10;12\right)$; $x_{18}$, $x_{19}$, $\cdots$, $x_{24}\left[12;14\right)$.
			\item Tứ phân vị thứ hai của mẫu số liệu trên là $\dfrac{1}{2}\left(x_{12}+x_{13}\right)\in\left[10;12\right)$  nên ta có
			\[Q_2=10+\dfrac{\dfrac{24}{2}-9}{8}\cdot\left(12-10\right)=10{,}75.\] 
			\item Tứ phân vị thứ nhất của dãy số liệu trên là  là $\dfrac{1}{2}\left(x_{6}+x_{7}\right)\in\left[8;10\right)$ nên ta có
			\[Q_1=8+\dfrac{\dfrac{24}{4}-3}{6}\cdot\left(10-8\right)=9.\] 
			 \item Tứ phân vị thứ ba của dãy số liệu trên là là $\dfrac{1}{2}\left(x_{18}+x_{19}\right)\in\left[12;14\right)$ nên ta có
			\[Q_3=12+\dfrac{\dfrac{3\cdot24}{4}-\left(3+6+8\right)}{7}\cdot\left(14-12\right)\approx  12{,}3.\]
		\end{itemize}
	 }
\end{ex}

% nguồn: Ok_Toan11_TracNghiem_Lan2
\begin{ex}%[1D5H2-3]%[ma:f8f58aad]
	Tìm hiểu thời gian  hoàn thành một bài tập của một số học sinh thu được kết quả sau:
	\begin{center}
		\begin{tabular}{|c|c|c|c|c|c|}
			 \hline
			Thời gian (phút)  &$ \left[0;4\right)$ & $\left[4;8\right)$ & $\left[8;12\right)$ & $\left[12;16\right)$ & $\left[16;20\right)$ \\
			\hline 
			Số học sinh & $2$ & $4$ & $7$ & $4$  & $3$ \\
			\hline 
		\end{tabular}
	\end{center}
	Tứ phân vị thứ ba của mẫu  số liệu ghép nhóm này là
	\choice 
		 {$Q_3=13$}
		 {\True $Q_3=14$}
		{$Q_3=15$} 
		{$Q_3=12$} 
	\loigiai{ 
		\begin{itemize}
			\item Cỡ  mẫu là $n=20$.
			\item Gọi $x_1$, $x_2$, $\cdots$, $x_{20}$ lần lượt là thời gian hoàn thành  bài tập của học sinh theo thứ tự không giảm.
			\item Ta có $x_1$, $x_2\in\left[0;4\right)$; $x_3$, $x_{4}$, $x_5$, $ x_{6}\in\left[4;8\right)$;  $x_{7}$, $x_{8}$, $\cdots$, $x_{13}\in\left[8;12\right)$; $x_{14}$, $x_{15}$, $x_{16}$, $x_{17}\in\left[12;16\right)$; $x_{18}$, $x_{19}$, $x_{20}\in\left[16;20\right)$.
			\item Tứ phân vị thứ  ba của mẫu số liệu $Q_3=\dfrac{1}{2}\left(x_{15}+x_{16}\right)\in\left[12;16\right)$ nên
			\[Q_3=12+\dfrac{\dfrac{3\cdot 20}{4}-\left(2+4+7\right)}{4}\cdot  \left(16-12\right)=14.\]
		\end{itemize}
	} 
\end{ex}

% nguồn: Ok_Toan11_TracNghiem_Lan2
\begin{ex}%[1D5H2-3]%[ma:dd7d26fc]
	Kết quả khảo sát  chiều cao của học sinh nam lớp 11A cho kết quả như bảng sau
	\begin{longtable}{|*{6}{c|}}
		\hline 
		Chiều cao học sinh (cm) & [$150$; $155$) & [$155$; $160$) & [$160$; $165$) & [$165$; $170$)  & [$170$; $175$) \\
		 \hline
		Số học sinh &  $1$ & $7$ & $12$ & $3$ & $2$ \\
		\hline 
	\end{longtable}
	Trung  vị của mẫu ghép nhóm trên là
	 \choice
		 {$165$}
		 {$160$}
		{$166{,}875$} 
		 {\True $161{,}875$}
	\loigiai{ 
		Cỡ mẫu  là $n=1+7+12+3+2=25$.\\
		Gọi $x_1, \ldots,  x_{25}$ là chiều cao của 25 học sinh nam được sắp xếp theo thứ tự tăng dần.\\
		Khi  đó: trung vị là $x_{13}$. Do giá trị $x_{13}$ thuộc nhóm $x_{13} \in \left[160; 165\right)$ nên nhóm này chứa trung vị.\\
		Suy ra $p=3$; $  a_3=160$; $ m_3=12$; $ m_1+m_2=1+7=8$; $ a_4-a_3=165-160=5$.\\
		Ta có trung vị $M_e=160+\dfrac{\dfrac{25}{2}-8}{12} \cdot  5=161{,}875$.
	} 
\end{ex}

% nguồn: Ok_Toan11_TracNghiem_Lan2
\begin{ex}%[1D5H2-3]%[ma:2906eef8]
	 Cân nặng của lợn con giống $A$ và giống $B$ được thống kê như bảng sau
	\begin{center}
		\begin{tabular}{|c|c|c|c|c|}
			\hline 
			 Cân nặng $(\mathrm{kg})$ & $[1{,}0; 1{,}1)$ & $[1{,}1; 1{,}2)$ & $[1{,}2; 1{,}3)$ & $[1{,}3; 1{,}4)$ \\
			\hline 
			Số con giống $A$ & $8$ & $28$  & $32$ & $17$ \\
			 \hline
			Số con  giống $B$ & $13$ & $14$ & $24$ & $14$ \\
			\hline 
		\end{tabular}
	\end{center}
	Hãy ước lượng trung vị và tứ phân vị thứ nhất của cân nặng lợn con mới sinh giống $A$ và của cân nặng lợn con mới sinh giống $B$. 
	\choice 
		{\True $M_{A}=1{,}22$; $Q_{1 A}=1{,}15$; $M_{B}=1{,}223$; $Q_{1  B}=1{,}12$}
		{$M_{A}=1{,}22$; $Q_{1 A}=1{,}45$; $M_{B}=1{,}223$; $Q_{1  B}=1{,}12$}
		{$M_{A}=1{,}22$; $Q_{1 A}=1{,}15$;  $M_{B}=1{,}43$; $Q_{1 B}=1{,}12$}
		{$M_{A}=1{,}02$; $Q_{1 A}=1{,}15$; $M_{B}=1{,}223$;  $Q_{1 B}=1{,}12$}
	\loigiai{ 
		Gọi $x_{1}$, $x_{2}$, $x_{3}$,$\ldots$, $x_{85}$ lần lượt là số lợn con giống  $A$ theo thứ tự không giảm.\\
		Do $x_{1}$,$\ldots$, $x_{8} \in[1{,}0; 1{,}1)$; $x_{9}$,$\ldots$, $x_{36} \in[1{,}1; 1{,}2)$; $x_{37}$,$\ldots$, $x_{68} \in[1{,}2; 1{,}3)$;  $x_{69}$,{\ldots}, $x_{85} \in[1{,}3; 1{,}4)$.\\
		 Trung vị của mẫu số liệu lợn con giống $A$ thuộc nhóm $[1{,}2; 1{,}3)$\\
		\centerline{$M_{A}=1{,}2+\dfrac{\dfrac{85}{2}-36}{32}  \cdot(1{,}3-1{,}2)=1{,}22.$}
		Gọi $y_{1}$, $y_{2}$, $y_{3}$,$\ldots$, $y_{65}$ lần lượt là số lợn con giống  $B$ theo thứ tự không giảm.\\
		Do $y_{1}$,$\ldots$, $y_{13}  \in[1{,}0; 1{,}1)$; $y_{14}$,$\ldots$, $y_{27} \in[1{,}1; 1{,}2)$; $y_{28}$,$\ldots$, $y_{51} \in[1{,}2; 1{,}3)$; $y_{52}$,$\ldots$, $y_{65} \in[1{,}3; 1{,}4)$.\\
		Trung vị của  mẫu số liệu lợn con giống $B$ thuộc nhóm $[1{,}2; 1{,}3)$\\
		 \centerline{$M_{B}=1{,}2+\dfrac{\dfrac{65}{2}-27}{24} \cdot(1{,}3-1{,}2)=1{,}223.$}
		Vậy cân nặng trung bình của lợn con giống $A$ nhỏ hơn giống $B$.\\ 
		Tứ phân vị thứ  nhất của dãy số liệu giống $A$ là $\dfrac{1}{2}(x_{21}+x_{22})$ thuộc nhóm $[1{,}1; 1{,}2)$ nên tứ phân vị thứ nhất của mẫu số liệu là $Q_{1 A}=1{,}1+\dfrac{\dfrac{85}{4}-8}{28}(1{,}2-1{,}1)=1{,}15$.\\
		Tứ phân vị thứ nhất của dãy số liệu giống $B$ là $\dfrac{1}{2}(y_{16}+y_{17})$ thuộc nhóm $[1{,}1; 1{,}2)$ nên tứ phân vị thứ  nhất của mẫu số liệu là $Q_{1 B}=1{,}1+\dfrac{\dfrac{65}{4}-13}{14}(1{,}2-1{,}1)=1{,}12$.
	} 
\end{ex}

% nguồn: Ok_Toan11_TracNghiem_Lan2
\begin{ex}%[1D5H2-3]%[ma:41c67736]
	[Dự án 11  HVA 2024-2025]
	Cho mẫu số liệu ghép nhóm về chiều cao của $25$ cây  dừa như ở bảng sau
	\begin{center}
		\begin{tabular}{|c|c|c|c|c|c|}
			\hline 
			Chiều cao (cm) & $[0; 10)$  & $[10; 20)$ & $[20; 30)$ & $[30; 40)$ & $[40; 50)$ \\
			\hline 
			Số cây &  $4$ & $6$ & $7$ & $5$ & $3$ \\
			 \hline
		\end{tabular}
	\end{center}
	Tứ phân vị thứ nhất của mẫu số liệu ghép nhóm này  là
	\choice 
		 {$Q_1=13{,}5$}
		 {$Q_1=13{,}9$}
		 {$Q_1=15{,}75$}
		{\True  $Q_1=13{,}75$}
	\loigiai{ 
		Cỡ mẫu  $n=4+6+7+5+3=25$.\\
		Tứ phân vị thứ nhất $Q_1$ là  $\dfrac{x_6+x_7}{2}$.\\
		Do $x_6$, $x_7$ đều  thuộc nhóm $[10; 20)$ nên nhóm này chứa $Q_1$.\\
		Khi đó  $p=2$, $a_2=10$, $m_2=6$, $m_1=4$, $a_3-a_2=10$. \\
		Ta có  $Q_1=10+\dfrac{\dfrac{25}{4}-4}{6} \cdot 10=13{,}75$.\\
	} 
\end{ex}

% nguồn: Ok_Toan11_TracNghiem_Lan2
\begin{ex}%[1D5H2-3]%[ma:74a07501]
	[Dự án 11 HVA 2024-2025] 
	Cho mẫu  số liệu cho bởi tần số ghép nhóm dưới đây
	\begin{center}
		\begin{tabular}{|c|c|c|c|c|c|c|}
			\hline 
			Nhóm & $[4; 6)$ & $[6; 8)$ & $[8; 10)$ & $[10; 12)$ & $[12; 14)$ &  $[14; 16)$ \\
			\hline 
			Tần số & $3$ & $8$ &  $18$ & $11$ & $12$ & $8$ \\
			\hline 
		\end{tabular}
	\end{center}
	Tứ phân vị thứ nhất của mẫu số  gần nhất giá trị nào trong các giá trị sau?
	 \choice
		 {\True $8{,}4$}
		{$8{,}9$} 
		{$8{,}0$} 
		 {$8{,}7$}
	 \loigiai{
		Gọi $x_1$, $x_2$, $\ldots$, $x_{60}$ là  số liệu được xếp theo thứ tự không giảm.\\
		Tứ phân vị thứ nhất của dãy số liệu $x_1$, $x_2$, $\ldots$, $x_{60}$ là $\dfrac{1}{2}\left(x_{15}+x_{16}\right) \in[8; 10)$. Do đó tứ phân vị thứ nhất của mẫu số  liệu ghép nhóm trên là $$Q_1=8+\dfrac{\dfrac{1 \cdot 60}{4}-(3+8)}{18} \cdot (10-8) \approx 8{,}4.$$
	 }
\end{ex}

% nguồn: Ok_Toan11_TracNghiem_Lan2
\begin{ex}%[1D5H2-3]%[ma:e526f586]
	Số điểm một  cầu thủ bóng rổ ghi được trong $20$ trận đấu được cho ở bảng sau:
	\begin{center}
		\begin{tabular}{|c|c|c|c|c|c|c|c|c|c|}
			\hline 
			$25$& $23$ & $21$ & $13$ & $8$ & $14$ & $15$ & $18$ & $22$ &  $11$ \\
			\hline 
			$24$& $12$ & $14$  & $14$ & $18$ & $6$ & $8$ & $25$ & $10$ & $11$ \\
			\hline 
		\end{tabular}
	\end{center}
	Tứ phân vị của dãy số liệu  trên lần lượt là
	\choice 
		{\True  $Q_1=11$, $Q_2=14$, $Q_3=21{,}5$}
		{$Q_1=10$, $Q_2=14$,  $Q_3=21$}
		 {$Q_1=12$, $Q_2=15$, $Q_3=22$}
		{$Q_1=11$, $Q_2=14$,  $Q_3=23$}
	 \loigiai{
		Ta có dãy điểm số  xếp theo thứ tự không giảm:
		$6$; $8$; $8$; $10$; $11$; $11$; $12$; $13$; $14$; $14$; $14$; $15$; $18$; $18$; $21$; $22$;  $23$; $24$; $25$; $25$.
		\begin{itemize}
			\item Tứ phân vị thứ nhất là:  $11$.
			\item Tứ phân vị thứ  hai là: $14$.
			\item Tứ phân vị thứ ba là:  $21{,}5$.
		\end{itemize}
	} 
\end{ex}

% nguồn: Ok_Toan11_TracNghiem_Lan2
\begin{ex}%[1D5H2-3]%[ma:8122e1b8]
	Trung tâm ngoại ngữ thống kê bảng điểm môn  Tiếng Anh của một khóa học trong bảng bên dưới.
	\begin{center}
		\begin{tabular}{|c|c|c|c|c|c|}
			 \hline
			Điểm & $\left[0; 2\right)$ & $\left[2; 4\right)$ & $\left[4;  6\right)$ & $\left[6; 8\right)$ & $\left[8; 10\right)$ \\
			\hline 
			Tầ Số học viên & $10$ & $30$ & $55$  & $42$ & $9$ \\
			\hline 
		\end{tabular}
	\end{center}
	Tứ phân vị thứ ba của mẫu số liệu trên gần nhất với giá trị nào trong  các giá trị sau?
	 \choice
		{$6{,}1$} 
		{\True  $6{,}7$}
		 {$6{,}9$}
		 {$6{,}5$}
	 \loigiai{
		Gọi $x_1$, $x_2$, $\ldots$, $x_{146}$ là số liệu được xếp theo  thứ tự không giảm.
		Tứ phân vị thứ ba của dãy số liệu  $x_1$, $x_2$, $\ldots$, $x_{146}$ là $x_{110} \in \left[6; 8\right)$. Do đó tứ phân vị thứ ba của mẫu số liệu ghép nhóm trên là
		\[Q_3=6+\dfrac{\dfrac{3\cdot146}{4}-\left(10+30+55\right)}{42}\cdot\left(8-6\right)\approx  6{,}7.\]
	} 
\end{ex}

% nguồn: Ok_Toan11_TracNghiem_Lan2
\begin{ex}%[1D5H2-3]%[ma:88c75d87]
	Khảo sát thời gian tập thể dục của một số học sinh khối  11 thu được mẫu số liệu ghép nhóm sau
	\begin{longtable}{|*{6}{c|}}
		\hline 
		 Thời gian (giờ)
		& [$0$;  $20$) & [$20$; $40$) & [$40$; $60$) & [$60$; $80$) & [$80$; $100$) \\
		\hline 
		Số học sinh & $5$ & $9$ & $13$ & $10$  & $6$ \\
		 \hline
	\end{longtable}
	Trung vị mẫu số liệu trên gần số nào nhất? 
	 \choice
		{$41{,}23$} 
		{$50{,}44$} 
		{\True $51{,}54$} 
		{$40{,}55$} 
	 \loigiai{
		Ta  có $M_e=u_m+\dfrac{\dfrac{n}{2}-C}{n_m} \cdot \left(u_{m+1}-u_m \right)$.
		\[\left. \begin{array}{l}{n=43} \\ {n_m=13} \\ {C=5+9=14} \\ {u_{m+1}=60} \\ {u_m=40}
\end{array}\right\}\Rightarrow M_e=40+\dfrac{\dfrac{43}{2}-14}{13} \cdot \left(60-40\right)\approx 51{,}54.\]
 }
\end{ex}

% nguồn: Ok_Toan11_TracNghiem_Lan2
\begin{ex}%[1D5H2-3]%[ma:b91af049]
	[Dự án 11 HVA 2024-2025] 
	Doanh thu bán hàng trong $20$ ngày được lựa chọn  ngẫu nhiên trong tháng $8$ tại một cửa hàng được ghi lại ở bảng sau
	\begin{center}
		\begin{tabular}{|c|c|c|c|c|c|}
			\hline 
			Doanh thu & $[7; 10)$ & $[10; 13)$ & $[13; 16)$ &  $[16; 19)$ & $[19; 22)$ \\
			\hline 
			Số ngày & $3$ & $4$ & $6$ & $5$ &  $2$ \\
			 \hline
		\end{tabular}
	\end{center}
	Tứ phân  vị thứ ba của mẫu số liệu là giá trị nào trong các giá trị sau?
	 \choice
		 {$16{,}2$}
		 {\True $17{,}2$}
		 {$18{,}2$}
		 {$19{,}2$}
	 \loigiai{
		Gọi $x_1$, $x_2$, $\ldots$, $x_{20}$ là doanh thu trong $20$ ngày được xếp theo thứ  tự không giảm.\\
		Tứ phân vị thứ ba của dãy số liệu $x_1$, $x_2, $\ldots, $x_{20}$ là $\dfrac{1}{2}\left(x_{15}+x_{16}\right)  \in[16; 19)$. Do đó tứ phân vị thứ ba của mẫu số liệu ghép nhóm trên là
		$$Q_3=16+\dfrac{\dfrac{3\cdot  20}{4}-(3+4+6)}{5}\cdot(19-16)=17{,}2.$$
	 }
\end{ex}

% nguồn: Ok_Toan11_TracNghiem_Lan2
\begin{ex}%[1D5H2-3]%[ma:0528a776]
	Lương  tháng của một số nhân viên văn phòng được ghi lại như sau (đơn vị: triệu đồng)
	\begin{center}
		\begin{tabular}{|c|c|c|c|c|}
			 \hline
			Lương tháng  (triệu đồng)&$[6;8)$&$[8;10)$&$[10;12)$&$[12;14)$\\
			 \hline
			Số nhân  viên&3&6&8&7\\
			\hline 
		\end{tabular}
	\end{center}
	Tứ phân vị của dãy số liệu trên lần lượt là 
	 \choice
		{$Q_1=9;Q_2=10;Q_3=12$} 
		{\True $Q_1=9;Q_2=10{,}75;Q_3=12{,}3$} 
		 {$Q_1=9;Q_2=11;Q_3=13$}
		 {$Q_1=10{,}73;Q_2=11;Q_3=13$}
	\loigiai{ 
		Gọi  $x_1;x_2;x_3;\ldots;x_{24}$ lần lượt là số nhân viên theo thứ tự không gian.\\
		Do $x_1,\ldots,x_3\in \left[6;8\right);x_4,\ldots,x_9\in \left[8;10\right);x_{10},\ldots,x_{17}\in \left[10;12\right);x_{18},\ldots,x_{24}\in \left[12;14\right)$.\\ 
		Tứ phân vị  thứ hai của dãy số liệu là $\dfrac{1}{2}\left(x_{12}+x_{13}\right)$ thuộc nhóm $\left[10;12\right)$ nên tứ phân vị thứ hai
		của mẫu số liệu  là $Q_2=10+\dfrac{\dfrac{24}{2}-9}{8}\left(12-10\right)=10{,}75$.\\
		Tứ phân vị thứ nhất của dãy số liệu là $\dfrac{1}{2}\left(x_6+x_7\right)$ thuộc nhóm $\left[8;10\right)$  nên tứ phân vị thứ nhất
		của mẫu số  liệu là $Q_1=8+\dfrac{\dfrac{24}{4}-3}{6}\left(10-8\right)=9$.\\
		Tứ phân vị thứ ba của dãy số liệu là $\dfrac{1}{2}\left(x_{18}+x_{19}\right)$ thuộc  nhóm $\left[12;14\right)$ nên tứ phân vị thứ ba của
	mẫu số  liệu là $Q_3=12+\dfrac{\dfrac{3\cdot24}{4}-17}{7}\left(14-12\right)=12{,}3$.}
\end{ex}

% nguồn: Ok_Toan11_TracNghiem_Lan2
\begin{ex}%[1D5H2-3]%[ma:d59e0f6b]
	 [Dự án 11 HVA 2024-2025]
	Khảo sát thời gian tập thể dục của  một số học sinh khối $11$ thu được mẫu số liệu ghép nhóm sau
	\begin{center}
		\begin{tabular}{|c|c|c|c|c|c|}
			\hline 
			 Thời gian (phút) & $[0; 20)$ & $[20; 40)$ & $[40; 60)$ & $[60; 80)$ & $[80; 100)$ \\
			\hline 
			Số học sinh & $5$ & $9$ & $12$ & $10$ & $6$ \\ 
			\hline 
		\end{tabular}
	\end{center}
	Tứ phân vị thứ ba của mẫu  số liệu trên là
	\choice 
		 {$115$}
		 {$117$}
		 {\True $71$}
		{$112$} 
	\loigiai{ 
		Số học sinh  được khảo sát là $n=5+9+12+10+6=42$\\
		Gọi $x_1$, $x_2$, $\ldots$,  $x_{42}$ là thời gian tập thể dục của $42$ học sinh lớp $11$ được xếm theo thứ tự không giảm.\\
		Tứ phân vị thứ ba của dãy số liệu $x_1$, $x_2$, $\ldots$, $x_{42}$  là $Q_3=x_{31} \in[60; 80)$. Do đó tứ phân vị của mẫu số liệu ghép nhóm trên là
		$$Q_3=60+\dfrac{\dfrac{3 \cdot 42}{4}-26}{10}  \cdot (80-60)=71.$$
	} 
\end{ex}

% nguồn: Ok_Toan11_TracNghiem_Lan2
\begin{ex}%[1D5H2-3]%[ma:b7607800]
	Cho bảng tần số ghép nhóm số liệu thống kê cân nặng của $40$ học sinh  lớp $11A$ trong một trường trung học phổ thông (đơn vị: kilôgam).
	\begin{center}
		\begin{tabular}{|c|c|}
			 \hline
			Nhóm  & Tần số \\
			 \hline
			$[30; 40)$ &  $2$ \\
			\hline 
			$[40; 50)$  & $10$ \\
			\hline 
			$[50; 60)$ & $16$  \\
			\hline 
			$[60; 70)$ & $8$  \\
			 \hline
			$[70;  80)$ & $2$ \\
			\hline 
			$[80; 90)$ & $2$ \\ 
			 \hline
		\end{tabular}
	\end{center}
	Hãy ước lượng các tứ  phân vị của mẫu số liệu ghép số trên.
	\choice 
		{$Q_{1}=49$~(kg);  $Q_{2}=50$~(kg); $Q_{3}=52{,}5$~(kg)}
		{\True  $Q_{1}=48$~(kg); $Q_{2}=55$~(kg); $Q_{3}=62{,}5$~(kg)}
		{$Q_{1}=47$~(kg);  $Q_{2}=54$~(kg); $Q_{3}=63{,}5$~(kg)}
		{$Q_{1}=46$~(kg); $Q_{2}=53$~(kg);  $Q_{3}=64{,}5$~(kg)}
	\loigiai{ 
		Số phần tử  của mẫu là $n=40$.\\
		Áp dụng  công thức, ta có tứ phân vị thứ nhất là $Q_{1}=40+\left(\dfrac{10-2}{10}\right) \cdot 10=48$~(kg).\\
		Áp  dụng công thức, ta có tứ phân vị thứ hai là $Q_{2}=M_{e}=50+\left(\dfrac{20-12}{16}\right) \cdot 10=55$~(kg).\\
		 Áp dụng công thức, ta có tứ phân vị thứ ba là $Q_{3}=60+\left(\dfrac{30-28}{8}\right) \cdot 10=62{,}5$~(kg).\\
		Vậy tứ phân vị  của mẫu số liệu trên là $Q_{1}=48$~(kg); $Q_{2}=55$~(kg); $Q_{3}=62{,}5$~(kg).
	 }
\end{ex}

% nguồn: Ok_Toan11_TracNghiem_Lan2
\begin{ex}%[1D5H2-3]%[ma:6efeeca3]
	Lương  tháng của một số giáo viên của một trường THPT được ghi lại như sau (đơn vị: triệu đồng):
	\begin{longtable}{|p{2.5cm}|*{5}{c|}}
		\hline 
		Lương  tháng \newline
		(triệu đồng) 
		& [$3{,}5$; $5{,}5$) & [$5{,}5$; $7{,}5$) & [$7{,}5$; $9{,}5$) & [$9{,}5$; $11{,}5$) & [$11{,}5$; $13{,}5$] \\ 
		\hline 
		Số giáo viên & $8$ & $11$ &  $13$ & $4$ & $9$ \\
		\hline 
	\end{longtable}
	Tứ phân  vị thứ hai của mẫu số liệu ghép nhóm trên là
	 \choice
		 {$Q_2=8$}
		 {$Q_2=8{,}5$}
		 {\True $Q_2=8{,}04$}
		 {$Q_2=8{,}4$}
	 \loigiai{
		Cỡ mẫu là  $n=8+11+13+4+9=45$.\\
		 Gọi $x_1, x_2, \ldots, x_{45}$ là số tiền lương của $45$ giáo viên và giả sử dãy này được sắp xếp theo thứ tự không giảm.\\
		Tứ phân  vị thứ hai $x_{23}$ thuộc nhóm $\left[7{,}5; 9{,}5\right)$.\\
		Khi  đó $Q_2=7{,}5+\dfrac{\dfrac{2\cdot 45}{4}-19}{13} \left(9{,}5-7{,}5\right)\approx 8{,}04$.
	 }
\end{ex}

% nguồn: Ok_Toan11_TracNghiem_Lan2
\begin{ex}%[1D5H2-3]%[ma:9666d50b]
	[Dự án  11 HVA 2024-2025]
	Cho mẫu số  liệu cho bởi tần số ghép nhóm dưới đây
	\begin{center}
		\begin{tabular}{|c|c|c|c|c|c|c|}
			\hline 
			Nhóm & $[4; 6)$ & $[6; 8)$ & $[8; 10)$ & $[10; 12)$ & $[12;  14)$ & $[14; 16)$ \\
			 \hline
			Tần số & $3$ & $8$  & $18$ & $11$ & $12$ & $8$ \\
			\hline 
		\end{tabular}
	\end{center}
	Tứ phân vị thứ ba của mẫu số liệu gần nhất  với giá trị nào trong các giá trị sau?
	 \choice
		 {$12{,}2$}
		{\True $12{,}8$} 
		{$12{,}0$} 
		{$13{,}2$} 
	\loigiai{ 
		Gọi $x_1$, $x_2$, $\ldots$, $x_{60}$ là số liệu được xếp  theo thứ tự không giảm.\\
		Tứ phân vị thứ ba của dãy số liệu $x_1$, $x_2$, $\ldots$, $x_{60}$là $\dfrac{1}{2}\left(x_{45}+x_{46}\right) \in[12; 14)$. Do đó tứ phân vị  thứ ba của mẫu số liệu ghép nhóm trên là $$Q_3=12+\dfrac{\dfrac{3\cdot 60}{4}-(3+8+18+11)}{12} \cdot(14-12) \approx 12{,}8.$$
	} 
\end{ex}

% nguồn: Ok_Toan11_TracNghiem_Lan2
\begin{ex}%[1D5H2-3]%[ma:9ec7b10a]
	Kiểm tra điện lượng của một số viên pintiểu do một hãng sản  xuất thu được kết quả sau
	\begin{center}
		\begin{tabular}{|c|c|c|c|c|c|}
			\hline 
			Điện lượng (nghìn mAh) & $[0{,}9; 0{,}95)$ & $[0{,}95; 1{,}0)$  & $[1{,}0; 1{,}05)$ & $[1{,}05; 1{,}1)$ & $[1{,}1; 1{,}15)$ \\
			\hline 
			Số  viên pin & 10 & 20 & 35 & 15 & 5 \\
			 \hline
		\end{tabular}
	\end{center}
	Hãy ước lượng tứ  phân vị của mẫu số liệu ghép nhóm trên.
	\choice 
		{$Q_{1}=0{,}58$;  $Q_{2}=1{,}02$; $Q_{3}=1{,}048$}
		 {$Q_{1}=0{,}98$; $Q_{2}=1{,}02$; $Q_{3}=1{,}248$}
		 {$Q_{1}=0{,}98$; $Q_{2}=1{,}22$; $Q_{3}=1{,}048$}
		{\True $Q_{1}=0{,}98$; $Q_{2}=1{,}02$;  $Q_{3}=1{,}048$}
	 \loigiai{
		Gọi $x_{1}$,  $x_{2}$, $x_{3}$,$\ldots$, $x_{85}$ lần lượt là số viên pintheo thứ tự không giảm.\\
		Do $x_{1}$,$\ldots$,  $x_{10}\in[0{,}9; 0{,}95)$; $x_{11}$,$\ldots$, $x_{30} \in[0{,}95; 1{,}0)$; $x_{31}$,$\ldots$, $x_{65} \in[1{,}0; 1{,}05)$; $x_{66}$,$\ldots$, $x_{80} \in[1{,}05; 1{,}1)$; $x_{81}$,$\ldots$, $x_{85} \in[1{,}1; 1{,}15)$.\\
		Tứ phân vị thứ hai của dãy số liệu là $\dfrac{1}{2}\left(x_{42}+x_{43}\right)$ thuộc nhóm $[1{,}0; 1{,}05)$ nên tứ phân vị thứ hai của mẫu số liệu là $Q_{2}=1{,}0+\dfrac{\dfrac{85}{2}-30}{35}(1{,}05-1{,}0)=1{,}02$.\\ 
		Tứ phân vị thứ  nhất của dãy số liệu là $\dfrac{1}{2}\left(x_{21}+x_{22}\right)$ thuộc nhóm $[0{,}95; 1{,}0)$ nên tứ phân vị thứ nhất của mẫu số liệu là $Q_{1}=0{,}95+\dfrac{\dfrac{85}{4}-10}{20}(1{,}0-0{,}95)=0{,}98$.\\
		Tứ phân vị thứ ba của dãy số liệu là $\dfrac{1}{2}\left(x_{63}+x_{64}\right)$ thuộc nhóm $[1{,}0;  1{,}05)$ nên tứ phân vị thứ ba của mẫu số liệu là $Q_{3}=1{,}0+\dfrac{\dfrac{3.85}{4}-30}{35}(1{,}05-1{,}0)=1{,}048$.
	} 
\end{ex}

% nguồn: Ok_Toan11_TracNghiem_Lan2
\begin{ex}%[1D5H2-3]%[ma:5bfaed8c]
	Doanh  thu bán hàng trong $20$ ngày được lựa chọn ngẫu nhiên của một của hàng được ghi lại ở bảng sau (đơn vị: triệu đồng):
	\begin{center}
		\begin{tabular}{|c|c|c|c|c|c|}
			\hline 
			Doanh thu & $\left[5;7\right)$ & $\left[7;9\right) $ & $\left[9;11\right) $  & $\left[11;13\right) $ & $\left[13;15\right) $ \\
			\hline 
			Số ngày & $2$ & $7$ & $7$ & $3$  & $1$ \\
			 \hline
		\end{tabular}
	\end{center}
	 Tứ phân vị thứ nhất của mẫu số liệu gần nhất với giá trị nào trong các giá trị dưới đây?
	 \choice
		 {$7$}
		 {$7{,}6$}
		{\True  $8$}
		{$8{,}6$} 
	\loigiai{ 
		\begin{itemize}
			\item  Cỡ mẫu là $n=20$.
			\item Gọi $x_1$, $x_2$, $\cdots$, $x_{20}$ lần lượt là doanh thu (triệu đồng) $20$ ngày được sắp xếp theo thứ tự không  giảm.
			\item Ta có $x_1$, $x_2\in\left[5;7\right)$; $x_3$, $x_4$, $\cdots$, $ x_{9}\in\left[7;9\right)$; $x_{10}$, $x_{11}$, $\cdots$,  $x_{16}\in\left[9;11\right)$; $x_{17}$, $x_{18}$, $x_{19}\in\left[11;13\right)$; $x_{20}\in\left[13;15\right)$.
			\item Tứ phân vị  thứ nhất của mẫu số liệu $Q_1=\dfrac{1}{2}\left(x_{5}+x_{6}\right)\in\left[7;9\right)$ nên
			\[Q_1=7+\dfrac{\dfrac{1\cdot 20}{4}-2}{7}\cdot  \left(9-7\right)\approx 7{,}86.\]
		\end{itemize}
	 }
\end{ex}

% nguồn: Ok_Toan11_TracNghiem_Lan2
\begin{ex}%[1D5H2-3]%[ma:93d01c72]
	 [Dự án 11 HVA 2024-2025]
	Để kiểm tra thời gian  sử dụng pin của chiếc laptop mới, bạn An thống kê thời gian sử dụng laptop của mình từ lúc sạc đầy pin cho tới khi hết pin ở bảng sau
	\begin{center}
		\begin{tabular}{|c|c|c|c|c|c|}
			 \hline
			Thời gian sử dụng &  $[6; 8)$ & $[8; 10)$ & $[10; 12)$ & $[12; 14)$ & $[14; 16)$ \\
			\hline 
			Số lần sạc & $2$ & $5$ & $8$  & $5$ & $1$ \\
			 \hline
		\end{tabular}
	\end{center}
	Tứ phân vị thứ ba của mẫu số liệu là giá trị nào trong các giá  trị sau?
	 \choice
		 {$12{,}2$}
		{\True $12{,}3$} 
		 {$12{,}4$}
		{$12{,}5$} 
	 \loigiai{
		Gọi $x_1$, $x_2$, $\ldots$, $x_{21}$ là thời gian sử dụng  được xếp theo thứ tự không giảm.\\
		Tứ phân vị thứ ba của dãy số liệu $x_1$, $x_2$, $\ldots$, $x_{21}$ là $\dfrac{1}{2}\left(x_{16}+x_{17}\right) \in[12; 14)$. Do đó tứ phân vị thứ ba của mẫu  số liệu ghép nhóm trên là $$Q_3=12+\dfrac{\dfrac{3\cdot 21}{4}-(2+5+8)}{5} \cdot(14-12)=12{,}3.$$
	} 
\end{ex}

% nguồn: Ok_Toan11_TracNghiem_Lan2
\begin{ex}%[1D5H2-3]%[ma:1f1cfbf2]
	Thời gian hoàn thành bài tập về nhà môn Toán của một số bạn học sinh đưa cho trong bảng  sau:
	\begin{center}
		\begin{tabular}{|c|c|c|c|c|c|}
			\hline 
			Thời gian (phút)& $\left[9{,}5;12{,}5\right)$ & $\left[12{,}5;15{,}5\right)$ & $\left[15{,}5;18{,}5\right)$ &  $\left[18{,}5;21{,}5\right)$ & $\left[21{,}5;24{,}5\right)$ \\
			 \hline
			Số học sinh & $3$ & $12$ & $15$ & $24$ & $2$  \\
			\hline 
		\end{tabular}
	\end{center}
	 Tìm tứ phân vị thứ ba của mẫu số liệu ghép nhóm trên là
	 \choice
		{$19$} 
		{\True  $20$}
		{$21$} 
		{$22$} 
	\loigiai{ 
		\begin{itemize}
			\item Cỡ mẫu là  $n=56$.
			\item Gọi  $x_1$, $x_2$, $\cdots$, $x_{56}$ lần lượt là số nhọc sinh có thời gian làm bài tập theo thứ tự không giảm.
			\item Ta có $x_1$, $x_2$, $x_3\in\left[9{,}5;12{,}5\right)$; $x_4$, $x_5$, $\cdots$, $ x_{15}\in\left[12{,}5;15{,}5\right)$; $x_{16}$, $x_{17}$, $\cdots$, $x_{30}\in\left[15{,}5;18{,}5\right)$; $x_{31}$,  $x_{32}$, $\cdots$, $x_{54}\in\left[18{,}5;21{,}5\right)$; $x_{55}$, $x_{56}\in\left[21{,}5;24{,}5\right)$.
			\item Tứ phân  vị thứ ba của mẫu số liệu $Q_3=\dfrac{1}{2}\left(x_{42}+x_{43}\right)\in\left[18{,}5;21{,}5\right)$ nên
			\[Q_3=18{,}5+\dfrac{\dfrac{3\cdot56}{4}-\left(3+12+15\right)}{24}\cdot  3=20.\]
		\end{itemize}
	 }
\end{ex}

% nguồn: Ok_Toan11_TracNghiem_Lan2
\begin{ex}%[1D5H2-3]%[ma:439a176e]
	Khảo sát thời gian (giờ) tự học mỗi ngày của một số học sinh khối 11 được  cho trong bảng sau
	\begin{longtable}{|*{6}{c|}}
		\hline 
		 Thời gian (giờ)
		& [$0{,}5$; $1{,}5$) & [$1{,}5$; $2{,}5$) & [$2{,}5$; $3{,}5$) & [$3{,}5$; $4{,}5$) & [$4{,}5$;  $5{,}5$) \\
		 \hline
		Số học  sinh & $6$ & $18$ & $27$ & $20$ & $8$ \\
		 \hline
	\end{longtable}
	Tứ phân vị của mẫu số  liệu ghép nhóm trên thứ tự là
	 \choice
		{$2{,}3$; $ 3{,}3$; $ 4{,}0$} 
		{$1{,}5$; $  2{,}5$; $ 3{,}5$}
		{$2{,}1$; $ 3{,}2$; $  4{,}1$}
		{\True $2{,}3$; $ 3{,}1$; $ 3{,}9$} 
	 \loigiai{
		 Cỡ mẫu là $n=6+18+27+20+8=79$.\\
		Gọi $x_1, x_2, \ldots, x_{79}$ là số giờ của $79$ học sinh và giả sử  dãy này được sắp xếp theo thứ tự không giảm.\\
		Ta có $x_1, x_2,  \ldots, x_6$ thuộc nhóm $\left[0{,}5; 1{,}5\right)$,
		 $x_7, x_8, \ldots, x_{24}$ thuộc nhóm $\left[1{,}5; 2{,}5\right)$,
		$x_{25}, x_{26}, \ldots, x_{51}$ thuộc nhóm $\left[2{,}5;  3{,}5\right)$,
		$x_{52}, x_{53}, \ldots, x_{71}$  thuộc nhóm $\left[3{,}5; 4{,}5\right)$,
		$x_{72},  x_{73}, \ldots, x_{79}$ thuộc nhóm $\left[4{,}5; 6\right]$.\\
		Tứ phân vị thứ nhất $x_{20}$ thuộc nhóm $\left[1{,}5; 2{,}5\right)$. Khi  đó $Q_1=1{,}5+\dfrac{\dfrac{79}{4}-6}{18} \left(2{,}5-1{,}5\right)\approx 2{,}3$.\\
		Tứ phân vị thứ hai $x_{40}$ thuộc nhóm $\left[2{,}5; 3{,}5\right)$. Khi đó  $Q_2=2{,}5+\dfrac{\dfrac{2\cdot 79}{4}-24}{27} \left(3{,}5-2{,}5\right)\approx 3{,}1$.\\
		Tứ phân vị thứ ba $x_{60}$ thuộc nhóm $\left[3{,}5;  4{,}5\right)$. Khi đó $Q_3=3{,}5+\dfrac{\dfrac{3\cdot 79}{4}-51}{20} \left(4{,}5-3{,}5\right)\approx 3{,}9$.
	 }
\end{ex}

% nguồn: Ok_Toan11_TracNghiem_Lan2
\begin{ex}%[1D5H2-3]%[ma:2025fff1]
	Đo  cân nặng của $1$ lớp gồm $40$ học sinh lớp $12$B. Kết quả được ghi lại trong bảng sau:
	\begin{center}
		\begin{tabular}{|c|c|c|c|c|c|c|c|c|}
			 \hline
			Khối lượng (kg) & $\left[40;45\right)$ & $\left[45;50\right)$ & $\left[50;55\right)$ & $\left[55;60\right)$ &  $\left[60;65\right)$ & $\left[65;70\right)$ & $\left[70;75\right)$ & $\left[75;80\right)$ \\
			\hline 
			Số học sinh & $4$ & $13$ & $7$ & $5$  & $6$ & $2$ & $1$ & $2$ \\
			 \hline
		\end{tabular}
	\end{center}
	Tứ phân vị thứ hai của mẫu số liệu ghép nhóm thuộc khoảng  nào sau đây?
	\choice 
		 {$\left[40;45\right]$}
		{$\left[45;50\right]$} 
		{\True $\left[50;55\right]$} 
		 {$\left[55;60\right]$}
	\loigiai{ 
		\begin{itemize}
			\item Cỡ  mẫu $n=40$.
			\item Gọi $x_1$, $x_2$, $\cdots$, $x_{40}$  lần lượt là số huy chương nhận được của các quốc gia được sắp xếp theo thứ tự không giảm.
			\item $Q_2=\dfrac{1}{2}\left(x_{20}+x_{21}\right)\in\left[50;  55\right)$. Suy ra
			\[Q_2=50+\dfrac{2\cdot\dfrac{40}{4}-\left(4+13\right)  }{7}\cdot\left(55-50\right)=52{,}14.\]
		\end{itemize}
	} 
\end{ex}

% nguồn: Ok_Toan11_TracNghiem_Lan2
\begin{ex}%[1D5H2-3]%[ma:e973db86]
	Một hãng ô  tô thống kê lại số lần gặp sự cố về động cơ của $100$ chiếc xe cùng loại sau $2$ năm sử dụng đầu tiên ở bảng sau
	\begin{center}
		\begin{tabular}{|c|c|c|c|c|c|}
			 \hline
			Số lần gặp sự cố & $[1; 2]$ & $[3; 4]$ & $[5; 6]$ &  $[7; 8]$ & $[9; 10]$ \\
			\hline 
			Số xe & $17$ & $33$ &  $25$ & $20$ & $5$ \\
			 \hline
		\end{tabular}
	\end{center}
	Hãy ước lượng tứ phân vị thứ  nhất của mẫu số liệu ghép số trên.
	\choice 
		{$2{,}64$} 
		 {$2{,}89$}
		{$2{,}73$} 
		 {\True $2{,}98$}
	\loigiai{ 
		Do số lần gặp sự cố là số nguyên nên ta hiệu chỉnh lại như  sau
		\begin{center}
			\begin{tabular}{|c|c|c|c|c|c|}
				\hline 
				Số lần gặp sự cố & $[0{,}5; 2{,}5)$ & $[2{,}5; 4{,}5)$ & $[4{,}5; 6{,}5)$ & $[6{,}5; 8{,}5)$ & $[8{,}5; 10{,}5)$ \\ 
				 \hline
				Số xe & $17$ & $33$ & $25$ & $20$ & $5$ \\ 
				\hline 
			\end{tabular}
		\end{center}
		Gọi $x_{1}$, $x_{2}$, $\ldots$, $x_{100}$ là mẫu số liệu được xếp theo thứ tự không giảm.\\ 
		Ta có $x_{1}$, $\ldots$, $x_{17} \in [0{,}5; 2{,}5)$; $x_{18}$, $\ldots$,  $x_{50} \in [2{,}5; 4{,}5)$; $x_{51}$, $\ldots$, $x_{75} \in [4{,}5; 6{,}5)$;
		$x_{76}$, $\ldots$, $x_{95} \in [6{,}5; 8{,}5)$; $x_{96}$, $\ldots$, $x_{100} \in [8{,}5; 10{,}5)$.\\ 
		Tứ phân vị thứ nhất của dãy số liệu $x_{1}$, $x_{2}$, $\ldots$, $x_{100}$ là $\dfrac{1}{2}\left(x_{25}+x_{26}\right)$. Do $x_{25}$ và $x_{26}$ thuộc nhóm $[2{,}5; 4{,}5)$ nên tứ phân vị thứ nhất của mẫu số liệu  ghép nhóm là\\
		\centerline{$Q_{1}=2{,}5+\dfrac{\dfrac{1 \cdot 100}{4}-17}{33} \cdot  (4{,}5-2{,}5)=\dfrac{197}{66} \approx 2{,}98.$}
	 }
\end{ex}

% nguồn: Ok_Toan11_TracNghiem_Lan2
\begin{ex}%[1D5H2-3]%[ma:05df9d42]
	Tìm hiểu thời gian hoàn thành một bài tập (đơn vị: phút) của một số học sinh thu được kết quả sau 
	\begin{longtable}{|*{6}{c|}}
		\hline 
		Thời gian  (phút)
		& [$0$;  $4$) & [$4$; $8$) & [$8$; $12$) & [$12$; $16$) & [$16$; $20$) \\
		\hline 
		Số học sinh & $2$ & $4$ & $7$ & $4$ & $3$  \\
		\hline 
	\end{longtable}
	Tứ phân vị thứ  ba của mẫu số liệu ghép nhóm này là
	\choice 
		 {$Q_3=13$}
		{\True $Q_3=14$} 
		 {$Q_3=15$}
		 {$Q_3=12$}
	\loigiai{ 
		Cỡ mẫu:  $n=2+4+7+4+3=20$.\\
		 Tứ phân vị thứ ba $Q_3$ là $\dfrac{x_{15}+x_{16}}{2}$.\\
		Do $x_{15}, x_{16}$ đều thuộc  nhóm $\left[12; 16\right)$ nên nhóm này chứa $Q_3$.\\
		 Do đó $p=4$, $a_4=12$, $m_4=4$, $m_1+m_2+m_3=2+4+7=13$, $a_5-a_4=4$.\\
		 Ta có $Q_3=12+\dfrac{\dfrac{3\cdot 20}{4}-13}{4} \cdot 4=14$.
	} 
\end{ex}

% nguồn: Ok_Toan11_TracNghiem_Lan2
\begin{ex}%[1D5H2-3]%[ma:11f635d0]
	Doanh thu bán  hàng trong $20$ ngày được lựa chọn ngẫu nhiên của một cửa hàng được ghi lại ở bảng sau:
	\begin{center}
		\begin{tabular}{|c|c|c|c|c|c|}
			\hline 
			Doanh thu& $\left[5; 7\right)$ & $\left[7; 9\right)$ & $\left[9; 11\right)$ & $\left[11; 13\right)$  & $\left[13; 15\right)$ \\
			 \hline
			Số ngày & $2$ & $7$ & $7$ & $3$ & $1$  \\
			\hline 
		\end{tabular}
	\end{center}
	Tứ phân vị thứ nhất của mẫu số liệu trên gần nhất với  giá trị nào dưới đây?
	\choice 
		{$7$} 
		 {$7{,}6$}
		{\True  $8$}
		 {$8{,}6$}
	 \loigiai{
		\begin{itemize}
			\item Sắp xếp các giá trị (số ngày) theo thứ tự không giảm $x_1$, $x_2$, $\cdots$, $x_{20}$.  Vì $n=20$ là số chẵn nên $Q_2=\dfrac{1}{2}\left(x_{10}+x_{11}\right)$.
			\item Tứ phân  vị thứ nhất là trung vị của nửa số liệu bên trái của mẫu số liệu ban đầu gồm $x_1$, $x_2$, $\cdots$, $x_{10}$. Vì $n=10$ là số chẵn nên $Q_1=\dfrac{1}{2}\left(x_5+x_6\right)$.
			 \item Ta có $x_5$, $x_6\in\left[7;9\right)$ nên $Q_1=7+\dfrac{\dfrac{1\cdot20}{4}-2}{7}\left(9-7\right)\approx 7{,}86$.
		\end{itemize}
	} 
\end{ex}

% nguồn: Ok_TracNghiemLop11_Lan1
\begin{ex}%[1D5H2-3]%[ma:11479cc0]
	Doanh thu bán hàng trong $20$ ngày được lựa chọn ngẫu nhiên của một cửa hàng được ghi lại ở  bảng sau:
	\begin{center}
		\begin{tabular}{|c|c|c|c|c|c|}
			 \hline
			Doanh thu& $[5; 7)$ & $[7; 9)$ & $[9; 11)$ & $[11; 13)$ & $[13;  15)$ \\
			\hline 
			 Số ngày & $2$ & $7$ & $7$ & $3$ & $1$ \\
			\hline 
		\end{tabular}
	\end{center}
	Tứ phân vị thứ ba của mẫu  số liệu trên gần nhất với giá trị nào dưới đây?
	 \choice
		{$13$} 
		{$12$} 
		{\True $11$} 
		 {$10$}
	\loigiai{ 
		\begin{itemize}
			\item Sắp xếp các giá trị (số ngày) theo thứ tự không giảm  $x_1$, $x_2$, $\cdots$, $x_{20}$.\\
			Vì $n=20$ là số chẵn  nên $Q_2=\dfrac{1}{2}\left(x_{10}+x_{11}\right)$.
			\item Tứ phân vị thứ ba là trung vị của nửa số liệu bên phải của mẫu  số liệu ban đầu gồm $Q_2$, $x_{11}$, $x_{12}$, $\cdots$, $x_{20}$.\\
			Vì $n=10$ là số chẵn nên  $Q_3=\dfrac{1}{2}\left(x_{15}+x_{16}\right)$.
			\item Ta có $x_{15}$, $x_{16}\in\left[9;11\right)$  nên $Q_3=9+\dfrac{\dfrac{3\cdot20}{4}-9}{7}\cdot \left(11-9\right)\approx 10{,}71$.
		\end{itemize}
	} 
\end{ex}

% nguồn: Ok_TracNghiemLop11_Lan1
\begin{ex}%[1D5H2-3]%[ma:54847f48]
	Kiểm tra điện lượng của một số viên pintiểu do một hãng sản xuất thu được  kết quả sau
	\begin{center}
		\begin{tabular}{|l|c|c|c|c|c|}
			\hline Điện  lượng (nghìn mAh) & {$ [0{,}9; 0{,}95) $} & {$ [0{,}95; 1{,}0) $} & {$ [1{,}0; 1{,}05) $} & {$ [1{,}05; 1{,}1) $} & {$ [1{,}1; 1{,}15) $} \\
			\hline Số viên pin & $10$ & $20$ & $35$ &  $15$ & $5$ \\
			\hline 
		\end{tabular}
	\end{center}
	Hãy ước lượng tứ phân  vị của mẫu số liệu ghép nhóm trên.
	\choice 
		{$Q_{1}=0{,}58; Q_{2}=1{,}02;  Q_{3}=1{,}048$}
		{$Q_{1}=0{,}98;  Q_{2}=1{,}02; Q_{3}=1{,}248$}
		{$Q_{1}=0{,}98; Q_{2}=1{,}22; Q_{3}=1{,}048$} 
		{\True  $ Q_{1}=0{,}98; Q_{2}=1{,}02; Q_{3}=1{,}048$}
	\loigiai{ 
		Gọi $ x_{1}; x_{2}; x_{3}; \ldots; x_{85} $ lần lượt là số viên pintheo  thứ tự không giảm.\\
		Do $x_{1}, \ldots, x_{10} \in[0{,}9; 0{,}95); x_{11}, \ldots, x_{30} \in[0{,}95; 1{,}0);  x_{31}, \ldots, x_{65} \in[1{,}0; 1{,}05)$.\\
		$x_{66}, \ldots, x_{80} \in[1{,}05; 1{,}1); x_{81}, \ldots, x_{85}  \in[1{,}1; 1{,}15)$.\\
		Tứ phân vị thứ hai của dãy số liệu là  $\dfrac{1}{2}\left(x_{42}+x_{43}\right)$ thuộc nhóm $[1{,}0; 1{,}05)$ nên tứ phân vị thứ hai của mẫu số liệu là $ Q_{2}=1{,}0+\dfrac{\dfrac{85}{2}-30}{35}\cdot (1{,}05-1{,}0)=1{,}02$.\\
		Tứ phân vị thứ nhất của dãy  số liệu là $\dfrac{1}{2}\left(x_{21}+x_{22}\right)$ thuộc nhóm $[0{,}95; 1{,}0)$ nên tứ phân vị thứ nhất của mẫu số liệu là $ Q_{1}=0{,}95+\dfrac{\dfrac{85}{4}-10}{20}\cdot (1{,}0-0{,}95)=0{,}98$.\\
		 Tứ phân vị thứ ba của dãy số liệu là $\dfrac{1}{2}\left(x_{63}+x_{64}\right)$ thuộc nhóm $ [1{,}0; 1{,}05) $ nên tứ phân vị thứ ba của mẫu số liệu là $ Q_{3}=1{,}0+\dfrac{\dfrac{3\cdot 85}{4}-30}{35}\cdot (1{,}05-1{,}0)=1{,}048$.
	} 
\end{ex}

% nguồn: Ok_TracNghiemLop11_Lan1
\begin{ex}%[1D5H2-3]%[ma:4a4355aa]
	Điều tra  $42$ học sinh của một lớp $11$ về số giờ tự học ở nhà, người ta có bảng sau đây
	\begin{center}
		\begin{tabular}{|c|c|c|}
			\hline 
			Lớp & Tần số & Tần số tích lũy  \\
			 \hline
			$\left[1; 2\right)$ & $8$ &  $8$ \\
			 \hline
			$\left[2;  3\right)$ & $10$ & $18$ \\
			 \hline
			$\left[3;  4\right)$ & $12$ & $30$ \\
			 \hline
			 $\left[4; 5\right)$ & $9$ & $39$ \\
			\hline 
			$\left[5; 6\right)$ &  $3$ & $42$ \\
			 \hline
		\end{tabular}
	\end{center}
	Nhận xét  nào đúng về tứ phân vị của mẫu số liệu trên
	 \choice
		 {Tứ phân vị của mẫu số liệu trên luôn giảm}
		{\True Tứ  phân vị của mẫu số liệu trên luôn tăng}
		{Tứ phân vị của mẫu số liệu trên  luôn cách đều nhau}
		{Tứ phân vị của mẫu số liệu trên không  tăng}
	\loigiai{ 
		\begin{itemize}
			\item Ta có số phần tử của mẫu  là\\
			$n=42\Rightarrow  \dfrac{n}{4}=10{,}5$.\\
			Suy ra nhóm $2$ là nhóm đầu tiên  có tần số tích lũy lớn hơn hoặc bằng $10{,}5$.\\
			Xét nhóm  $2$ là nhóm $\left[2; 3\right)$ có $s=2$; $h=1$; $n_2=10$ và nhóm $1$ là nhóm $\left[1; 2\right)$ có $cf_1=8$.\\
			Áp dụng công thức tứ phân vị thứ nhất của mẫu số  liệu có\\
			$$Q_1=2+\left(\dfrac{10{,}5-8}{10}  \right)\cdot 1=2{,}25.$$
			\item Ta có số phần  tử của mẫu là\\
			$n=42\Rightarrow \dfrac{n}{2}=21$.\\ 
			Mà $cf_2=18< 21< cf_3=30$ suy ra nhóm $3$ là nhóm đầu tiên có tần số tích lũy lớn hơn hoặc bằng  $21$.\\
			Xét nhóm $3$ là nhóm $\left[3; 4\right)$ có $r=3$; $d=1$;  $n_3=12$ và nhóm $2$ là nhóm $\left[2; 3\right)$ có $cf_2=18$.\\
			Áp dụng công thức  ta có trung vị của mẫu số liệu là
			$$M_e=3+\left(\dfrac{21-18}{12}  \right)\cdot 1=3{,}25.$$
			Vậy tứ phân vị thứ $2$ là $Q_2=M_e=3{,}25$. 
			 \item Ta có số phần tử của mẫu là\\
			 $\dfrac{3n}{4}=31{,}5$.
			Suy ra nhóm $4$ là nhóm đầu tiên có tần số tích lũy lớn  hơn hoặc bằng $31{,}5$.\\
			Xét nhóm $4$ là nhóm $\left[4; 5\right)$ có $t=4$; $l=1$; $n_4=9$ và nhóm $3$  là nhóm $\left[3; 4\right)$ có $cf_3=30$.\\
			Áp dụng công thức tứ phân vị thứ  ba của mẫu số liệu có
			$$Q_3=4+\left(\dfrac{31{,}5-30}{9} \right)\cdot  1=\dfrac{25}{6}.$$
		\end{itemize}
		Vậy $Q_1=2{,}25$; $Q_2=3{,}1$; $Q_3=\dfrac{25}{6} \approx 4{,}166$ nên tứ phân vị  luôn tăng.
	 }
\end{ex}

% nguồn: Ok_TracNghiemLop11_Lan1
\begin{ex}%[1D5H2-3]%[ma:44a05ce6]
	Thời gian luyện tập trong một ngày (tính theo giờ) của một số vận  động viên được ghi lại ở bảng sau
	\begin{center}
		\begin{tabular}{|l|c|c|c|c|c|}
			\hline Thời gian luyện tập (giờ) & {$ [0; 2) $} & {$ [2; 4) $} & {$ [4; 6)$} & {$ [6; 8)  $} & {$ [8; 10)$} \\
			\hline Số  vận động viên & $3$ & $8$ & $12$ & $12$ & $4$ \\
			\hline 
		\end{tabular}
	\end{center}
	Hãy xác định tứ phân vị thứ $3$ của mẫu số liệu trên. 
	 \choice
		{$3{,}6875$} 
		 {$5{,}417$}
		{\True $7{,}042$} 
		{$7{,}68$} 
	\loigiai{ 
		Tứ phân vị thứ ba của dãy số liệu $ x_{1}; x_{2}; \ldots;  x_{39} $ là $ x_{30} \in[6; 8) $.
		Do đó tứ phân  vị thứ ba của mẫu số liệu ghép nhóm là
		 $
		Q_{3}=6+\dfrac{\dfrac{3\cdot 39}{4}-(3+8+12)}{12} \cdot (8-6)=\dfrac{169}{24} \approx  7{,}042$.
	 }
\end{ex}

% nguồn: Ok_TracNghiemLop11_Lan1
\begin{ex}%[1D5H2-3]%[ma:3a457bae]
	Điểm thi của $32$ học sinh trong  kì thi Tiếng Anh (thang $100$ điểm) được phân bố như sau
	\begin{center}
		\begin{tabular}{|c|c|}
			\hline  Lớp điểm & Tần số $n_i$ \\
			 \hline$[40; 50)$ & $4$ \\
			\hline$[50; 60)$ & $6$ \\ 
			\hline$[60; 70)$ & $10$ \\ 
			\hline$[70; 80)$  & $6$ \\
			\hline$[80; 90)$ & $4$ \\ 
			\hline$[90;  100]$ & $2$ \\
			\hline 
		\end{tabular}
	\end{center}
	Tứ phân vị thứ ba của bảng phân bố ghép nhóm bằng 
	 \choice
		{$\dfrac{310}{4}$} 
		 {$\dfrac{315}{3}$}
		{$\dfrac{220}{4}$} 
		 {\True $\dfrac{230}{3}$}
	\loigiai{ 
	Ta có  $Q_3=70+\dfrac{24-20}{6}\cdot 10=\dfrac{230}{3}$.}
\end{ex}

% nguồn: Ok_TracNghiemLop11_Lan1
\begin{ex}%[1D5H2-3]%[ma:98d65dc2]
	Người ta tiến hành phỏng vấn $40$ người về  một mẫu áo khoác. Người điều tra yêu cầu cho điểm mẫu áo đó theo thang điểm là $100$. Kết quả được trình bày trong bảng ghép nhóm sau
	\begin{center}
		\begin{tabular}{|c|c|c|c|c|c|c|}
			 \hline
			Nhóm & $\left[50;60\right)$ & $\left[60;70\right)$ & $\left[70;80\right)$ & $\left[80;90\right)$ &  $\left[90;100\right)$ & \\
			\hline 
			Tần  số & $4$ & $5$ & $23$ & $6$ & $2$ & $N=40$ \\
			\hline 
		\end{tabular}
	\end{center}
	Tứ phân vị  của mẫu số liệu ghép nhóm trên (làm tròn đến hàng đơn vị) là
	 \choice
		{$Q_1  \approx 71$, $Q_2 \approx 76$, $Q_3 \approx 78$}
		{$Q_1 \approx  71$, $Q_2 \approx 75$, $Q_3 \approx 78$}
		 {$Q_1 \approx 70$, $Q_2 \approx 76$, $Q_3 \approx 79$}
		{\True $Q_1 \approx 70$, $Q_2 \approx 75$, $Q_3 \approx  79$}
	 \loigiai{
		Xét nửa mẫu  số liệu bên trái $x_1, x_2, x_3, \ldots, x_{20}$ có trung vị $\dfrac{x_{10}+x_{11}}{2} \in [70;80)$.\\
		Ta có  $n_i=23$; $C_1=9$; $u_i=70$; $u_{i+1}=80$.\\
		 Tứ phân vị thứ nhất của mẫu số liệu ghép nhóm là
		\[Q_1=70+\dfrac{10-9}{23}\cdot 10 \approx 70.\] 
		Tứ  phân vị thứ hai của mẫu số liệu ghép nhóm là
		\[Q_2=70+\dfrac{2\cdot10-9}{23}\cdot  10 \approx 75.\]
		Xét nửa mẫu số liệu bên phải $x_{21}, x_{22}, x_{23}, \ldots, x_{40}$ có trung vị $\dfrac{x_{30}+x_{31}}{2}  \in [70;80)$.\\
		 Ta có $n_j=23$; $C_3=9$; $u_j=70$; $u_{j+1}\approx 80$.\\
		Tứ phân vị thứ ba  của mẫu số liệu ghép nhóm là
		\[Q_3=70+\dfrac{30-9}{23}\cdot  10 \approx 79.\]
	} 
\end{ex}

% nguồn: Ok_TracNghiemLop11_Lan1
\begin{ex}%[1D5H2-3]%[ma:977cb747]
	Khảo sát thời gian tập thể dục của một số học sinh khối 11 thu được mẫu số liệu ghép nhóm sau: 
	\begin{center}
		\begin{tabular}{|l|c|c|c|c|c|}
			\hline Thời  gian (phút) & {$[0; 20)$} & {$[20; 40)$} & {$[40; 60)$} & {$[60; 80)$} & {$[80; 100)$} \\
			\hline Số học sinh & $5$ & $9$ &  $12$ & $10$ & $6$ \\
			 \hline
		\end{tabular}
	\end{center}
	Nhóm chứa tứ phân vị  thứ ba của mẫu số liệu trên là
	 \choice
		 {$[40; 60)$}
		{$[20; 40)$} 
		{\True $[60;  80)$}
		 {$[80; 100)$}
	 \loigiai{
		Ta có  $n=42$.\\
		Nên tứ  phân vị thứ nhất của mẫu số liệu trên là $Q_3=x_{33}$.\\
		Mà $x_{33} \in[60; 80)$ nên nhóm chứa tứ  phân vị thứ nhất của mẫu số liệu trên là nhóm $[60; 80)$.
	 }
\end{ex}

% nguồn: Ok_TracNghiemLop11_Lan1
\begin{ex}%[1D5H2-3]%[ma:ee351835]
	Mức lương hàng tháng  ở một công ty được Công đoàn thu thập theo bảng sau (đơn vị triệu đồng).
	\begin{center}
		\begin{tabular}{|c|c|c|c|c|c|}
			\hline Mức lương & {$[6; 10)$}  & {$[10; 15)$} & {$[15; 20)$} & {$[20; 25)$} & {$[25; 30)$} \\
			\hline Nhân viên  & $17$ & $38$ & $27$ & $21$ & $7$ \\
			\hline 
		\end{tabular}
	\end{center}
	Tứ phân vị thứ ba  của mẫu số liệu ghép nhóm thuộc nửa khoảng nào sau đây?
	 \choice
		{\True $[20; 25)$} 
		{$[25;  30)$}
		 {$[25; 30)$}
		{$[15; 20)$} 
	 \loigiai{
		Tứ phân vị thứ ba thuộc $[20; 25)$. 
	 }
\end{ex}

% nguồn: Ok_TracNghiemLop11_Lan1
\begin{ex}%[1D5H2-3]%[ma:3730f16d]
	Người ta tiến hành phỏng vấn $40$ người về một mẫu áo sơ mi mới. Người điều tra yêu cầu cho điểm mẫu áo đó theo thang điểm là $100$. Kết quả được trình bày trong \textit{Bảng  dưới}.
	\begin{center}
		 {\renewcommand\arraystretch{1.25}\tabcolsep=5mm
		\begin{tabular}{|c|c|c|c|c|c|}
			\hline 
			Nhóm& 
			 $[50;60)$&
			$[60;70)$& 
			$[70;80)$& 
			 $[80;90)$&
			 $[90;100)$\\
			\hline 
			 Tần số&
			$4$ & 
			 $5$&
			$23$& 
			 $6$&
			 $2$\\
			 \hline
		\end{tabular}}
	\end{center}
	Tứ phân vị  của mẫu số liệu ghép nhóm trên (làm tròn đến hàng đơn vị là)
	 \choice
		{$Q_1\approx  71; Q_2\approx 76; Q_3\approx 78$}
		 {$Q_1\approx 71; Q_2\approx 75; Q_3\approx 78$}
		{$Q_1\approx 70; Q_2\approx 76; Q_3\approx 79$} 
		{\True $Q_1\approx 70; Q_2\approx 75; Q_3\approx 79$} 
	 \loigiai{
		Số phần tử  của mẫu là $n=40$.
		\begin{itemize}
			\item Ta có: $\dfrac{n}{4}=10$ mà $9<10<32$. \\Suy ra nhóm $3$ là nhóm đầu tiên có tần số tích lũy lớn hơn hoặc bằng $10$. \\Xét  nhóm $3$ là nhóm $[70;80)$ có $s=70;h=10;n_3=23$ và nhóm $2$ là nhóm $[60;70)$ có $cf_2=9$.\\
			Áp dụng công thức, ta  có tứ phân vị thứ nhất là
			$$Q_1=70+\left(\dfrac{10-9}{23}\right) \cdot 10 \approx  70.$$
			\item Ta có tứ phân vị thứ hai là 
			$$Q_2=M_e \approx 75.$$ 
			\item Ta có $\dfrac{3n}{4}=30$ mà $9<30<32$. \\ 
			Suy ra nhóm $3$ là nhóm đầu tiên có  tần số tích lũy lớn hơn hoặc bằng $30$.\\
			Xét nhóm $3$ là  nhóm $[70;80)$ có $t=70;l=10;n_3=23$ và nhóm $2$ là nhóm $[60;70)$ có $cf_2=9$.\\
			Áp dụng công thức, ta có tứ phân vị  thứ ba là
			 $$Q_3=70+\left(\dfrac{30-9}{23}\right) \cdot 10 \approx 79.$$
		\end{itemize}
	 }
\end{ex}

% nguồn: Ok_TracNghiemLop11_Lan1
\begin{ex}%[1D5H2-3]%[ma:b9641065]
	Thống kê điểm học kì môn toán của các học sinh  lớp $11$A của một trường THPT, người ta thu được số liệu sau
	\begin{center}
		\begin{tabular}{|c|c|c|c|c|c|c|c|c|c|c|c|c|c|c|}
			\hline 
			$3$ & $5{,}5$ & $5$ & $4$  & $4{,}5$ & $4{,}5$ & $3$ & $5$ & $4$ & $4{,}5$ & $4{,}5$ & $6{,}5$ & $6{,}5$ & $7{,}5$ & $3{,}5$ \\
			\hline 
			$5$ & $6$ & $7$  & $8$ & $8$ & $7$ & $4{,}5$ & $6$ & $5$ & $7$ & $4$ & $5{,}5$ & $7{,}5$ & $8{,}5$ & $9{,}5$ \\
			 \hline
			$4$ & $3{,}5$ & $5$ & $8{,}5$ & $6{,}5$ & $4{,}5$ & $7{,}5$  & $7$ & $4{,}5$ & $3$ & $7$ & $5{,}5$ & $5{,}5$ & $6{,}5$ & $9$ \\
			 \hline
		\end{tabular}
	\end{center}
	Xác định tứ phân vị của mẫu số  liệu khi ta ghép lớp thành các nhóm có độ dài là $1$ như sau
	 \[\left[3; 4\right), \left[4; 5\right),\ldots, \left[9; 10\right).\]
	 \choice
		{\True $4{,}6; 5{,}7;  7{,}3$}
		{$4{,}6; 5{,}7;  7{,}4$}
		{$5{,}6; 6{,}7;  8{,}3$}
		{$4{,}7;  5{,}7; 7{,}4$}
	 \loigiai{
		Ta có bảng tần số ghép lớp, tần số tích  lũy sau
		\begin{center}
			\begin{tabular}{|c|c|c|}
				 \hline
				Lớp & Tần  số & Tần số tích lũy \\
				 \hline
				$\left[3; 4\right)$ & $5$ & $5$  \\
				 \hline
				$\left[4; 5\right)$  & $11$ & $16$ \\
				 \hline
				$\left[5; 6\right)$ & $9$ &  $25$ \\
				 \hline
				$\left[6; 7\right)$ & $6$  & $31$ \\
				 \hline
				$\left[7; 8\right)$ & $8$ &  $39$ \\
				 \hline
				$\left[8; 9\right)$ & $4$ & $43$  \\
				 \hline
				$\left[9; 10\right)$ & $2$ & $45$  \\
				\hline 
				 & $n=45$ & \\
				\hline 
			\end{tabular}
		\end{center}
		\begin{itemize}
			\item Ta có số  phần tử của mẫu là\\
			$n=45\Rightarrow \dfrac{n}{4}=11{,}25$.\\ 
			Suy ra nhóm $2$ là nhóm đầu tiên có tần số tích lũy lớn  hơn hoặc bằng $11{,}25$.\\
			Xét nhóm $2$ là nhóm $\left[4; 5\right)$ có $s=4$; $h=1$; $n_2=11$ và nhóm $1$ là nhóm $\left[3; 4\right)$ có  $cf_1=5$.\\
			Áp dụng  công thức tứ phân vị thứ nhất của mẫu số liệu có
			 \[Q_1=4+\left(\dfrac{11,25-5}{11} \right)\cdot 1\approx 4{,}6.\]
			\item Ta có số phần  tử của mẫu là\\
			 $n=45\Rightarrow \dfrac{n}{2}=22{,}5$.
			Mà $cf_2=16< 22{,}5< cf_3=25$ suy ra nhóm $3$ là nhóm đầu tiên có tần  số tích lũy lớn hơn hoặc bằng $22{,}5$.\\
			Xét nhóm $3$ là nhóm $\left[5; 6\right)$ có $r=5$; $d=1$; $n_3=9$ và nhóm $2$ là nhóm $\left[4;  5\right)$ có $cf_2=16$.\\
			Áp dụng công thức ta có trung  vị của mẫu số liệu là
			\[M_e=5+\left(\dfrac{22{,}5-16}{9} \right)\cdot  1\approx 5{,}7.\]
			Suy ra $Q_2=M_e  \approx 5{,}7$.
			\item Ta có số phần tử của mẫu  là\\
			 $n=45\Rightarrow \dfrac{3n}{4}=33{,}75$.
			Suy ra nhóm $5$ là nhóm đầu tiên có tần số tích lũy lớn hơn hoặc bằng $33{,}75$.\\ 
			Xét nhóm $5$ là nhóm $\left[7; 8\right)$ có $t=7$; $l=1$; $n_4=8$ và nhóm $4$ là  nhóm $\left[6; 7\right)$ có $cf_4=31$.\\
			Áp dụng công thức tứ phân vị  thứ nhất của mẫu số liệu có
			 \[Q_3=7+\left(\dfrac{33{,}75-31}{8} \right)\cdot 1\approx 7{,}3.\]
		\end{itemize}
	} 
\end{ex}

% nguồn: Ok_TracNghiemLop11_Lan1
\begin{ex}%[1D5H2-3]%[ma:257177d4]
	Đo cân nặng của $1$ lớp gồm $40$ học sinh lớp $12$B.  Kết quả được ghi lại trong bảng sau:
	\begin{center}
		\begin{tabular}{|c|c|c|c|c|c|c|c|c|}
			 \hline
			Khối lượng (kg) & $[40;45)$ & $[45;50)$ & $[50;55)$ & $[55;60)$ & $[60;65)$ & $[65;70)$ & $[70;75)$ & $[75;80)$  \\
			 \hline
			Số học sinh & $4$ & $13$ & $7$ & $5$ & $6$ & $2$ & $1$  & $2$ \\
			 \hline
		\end{tabular}
	\end{center}
	 Tứ phân vị thứ hai của mẫu số liệu ghép nhóm thuộc khoảng nào sau đây?
	 \choice
		 {$[40;45]$}
		{$[45;50]$} 
		 {\True $[50;55]$}
		{$[55;60]$} 
	\loigiai{ 
		\begin{itemize}
			 \item Cỡ mẫu $n=40$.
			\item Gọi $x_1$, $x_2$, $\ldots$, $x_{40}$ lần lượt là số huy chương nhận được của các quốc gia được sắp xếp theo thứ tự  không giảm.
			\item $Q_2=\dfrac{1}{2}\left(x_{20}+x_{21}\right)\in\left[50; 55\right)$.\\ 
			Suy ra 
			 \[Q_2=50+\dfrac{2\cdot\dfrac{40}{4}-\left(4+13\right) }{7}\cdot\left(55-50\right)=52{,}14.\]
		\end{itemize}
	 }
\end{ex}

% nguồn: Ok_TracNghiemLop11_Lan1
\begin{ex}%[1D5H2-3]%[ma:43415102]
	Trong mẫu số liệu ghép nhóm, tứ phân vị thứ ba $Q_3$  thể hiện:
	 \choice
		{Có $25\%$ số giá  trị nhỏ hơn $Q_3$}
		{\True Có $75\%$ số  giá trị nhỏ hơn $Q_3$}
		{Có $25\%$ số giá trị lớn hơn $Q_3$} 
		{Có $75\%$ số giá  trị lớn hơn $Q_3$}
	 \loigiai{
		Trong mẫu số liệu ghép nhóm, tứ phân vị thứ ba $Q_3$ thể hiện có $75\%$ số giá trị nhỏ  hơn $Q_3$.
	 }
\end{ex}

% nguồn: Ok_TracNghiemLop11_Lan1
\begin{ex}%[1D5H2-3]%[ma:b0f56af2]
	Cho bảng tần  số ghép nhóm số liệu thống kê cân nặng của $40$ học sinh lớp $11A$ trong một trường trung học phổ thông (đơn vị: kilôgam)
	\begin{center}
		\begin{tabular}{|l|c|c|c|c|c|c|}
			\hline Nhóm & {$[30;  40)$} & {$[40; 50)$} & {$[50; 60)$} & {$[60; 70)$} & {$ [70; 80)$}& {$[80; 90)$} \\
			\hline Tần số & $2$ & $10$  & $16$ & $8$ & $2$ & $2$ \\
			\hline 
		\end{tabular}
	\end{center}
	Hãy ước lượng các tứ phân  vị của mẫu số liệu ghép nhóm trên.
	 \choice
		{$Q_{1}=49$ (kg); $Q_{2}=50$ (kg); $Q_{3}=52{,}5$  (kg)}
		{\True $ Q_{1}=48$  (kg); $Q_{2}=55$ (kg); $Q_{3}=62{,}5$ (kg)}
		{$Q_{1}=47$ (kg); $Q_{2}=54$  (kg); $Q_{3}=63{,}5$ (kg)}
		{$Q_{1}=46$ (kg); $Q_{2}=53$ (kg);  $Q_{3}=64{,}5$ (kg)}
	\loigiai{ 
		Số phần tử của mẫu là $n=40$.\\ 
		Áp dụng công thức, ta có tứ phân vị thứ  nhất là\\
		$ Q_{1}=40+\left(\dfrac{10-2}{10}\right) \cdot 10=48$  (kg).\\
		Áp dụng công thức, ta  có tứ phân vị thứ hai là\\
		$ Q_{2}=\mathrm{M}_{e}=50+\left(\dfrac{20-12}{16}\right) \cdot 10=55$  (kg).\\
		Áp dụng công  thức, ta có tứ phân vị thứ ba là\\
		$Q_{3}=60+\left(\dfrac{30-28}{8}\right) \cdot  10=62{,}5$ (kg).\\
		Vậy tứ phân vị của mẫu số liệu trên  là\\
		$Q_{1}=48$ (kg);  $Q_{2}=55$ (kg); $Q_{3}=62{,}5$ (kg).
	} 
\end{ex}

% nguồn: Ok_TracNghiemLop11_Lan1
\begin{ex}%[1D5H2-3]%[ma:6a7a2938]
	Lương tháng của một số nhân viên một văn phòng được ghi  lại như sau (đơn vị: triệu đồng)
	\begin{center}
		\begin{tabular}{|l|c|c|c|c|}
			\hline Lương tháng (triệu  đồng) & {$ [6; 8) $} & {$ [8; 10) $} & {$ [10; 12) $} & {$ [12; 14) $} \\
			\hline Số nhân viên & $3$ & $6$ & $8$ &  $7$ \\
			\hline 
		\end{tabular}
	\end{center}
	Tìm tứ phân vị của mẫu  số liệu trên.
	\choice 
		{\True $Q_{1}=9; Q_{2}=10{,}75; Q_{3}=12{,}3$} 
		 {$Q_{1}=9; Q_{2}=10{,}75; Q_{3}=14{,}3$}
		{$Q_{1}=9;  Q_{2}=11{,}75; Q_{3}=12{,}3$}
		 {$Q_{1}=10; Q_{2}=10{,}75; Q_{3}=12{,}3$}
	\loigiai{ 
		Gọi $ x_{1}; x_{2}; x_{3}; \ldots; x_{24} $  lần lượt là lương tháng của các nhân viên theo thứ tự không giảm.\\
		Do $ x_{1}, \ldots, x_{3} \in[6; 8); x_{4}, \ldots,  x_{9} \in[8; 10); x_{10}, \ldots, x_{17} \in[10; 12); x_{18}, \ldots, x_{24} \in[12; 14)$.\\
		Tứ phân vị thứ hai của dãy số liệu là $ \dfrac{1}{2}\left(x_{12}+x_{13}\right)$ thuộc nhóm $[10;12)$ nên tứ phân vị thứ hai của mẫu số liệu là $ Q_{2}=10+\dfrac{\dfrac{24}{2}-9}{8}\cdot (12-10)=10{,}75  $.\\
		Tứ phân vị thứ nhất của dãy số liệu là $ \dfrac{1}{2}\left(x_{6}+x_{7}\right)$ thuộc nhóm  $ [8; 10) $ nên tứ phân vị thứ nhất của mẫu số liệu là $ Q_{1}=8+\dfrac{\dfrac{24}{4}-3}{6}\cdot (10-8)=9$.\\
		Tứ phân vị thứ  ba của dãy số liệu là $ \dfrac{1}{2}\left(x_{18}+x_{19}\right)$ thuộc nhóm $[12;14)$ nên tứ phân vị thứ ba của mẫu số liệu là $Q_{3}=12+\dfrac{\dfrac{3\cdot 24}{4}-17}{7}\cdot (14-12)=12{,}3$.
	 }
\end{ex}

% nguồn: Ok_TracNghiemLop11_Lan1
\begin{ex}%[1D5H2-3]%[ma:d6a596d6]
	Cân nặng của  lợn con giống $A$ và giống $B$ được thống kê như bảng sau
	\begin{center}
		\begin{tabular}{|c|c|c|c|c|}
			 \hline
			Cân nặng ({kg}) & $[1{,}0; 1{,}1)$ & $[1{,}1; 1{,}2)$ & $[1{,}2; 1{,}3)$ & $[1{,}3; 1{,}4)$ \\ 
			\hline 
			Số con giống  $A$ & $8$ & $28$ & $32$ & $17$ \\
			 \hline
			Số con giống $B$ & $13$ &  $14$ & $24$ & $14$ \\
			\hline 
		\end{tabular}
	\end{center}
	Hãy ước lượng trung vị và tứ phân vị thứ nhất của cân nặng lợn con mới sinh giống $A$ và  của cân nặng lợn con mới sinh giống $B$.
	 \choice
		{\True $M_{A}=1{,}22$; $Q_{1  A}=1{,}15$; $M_{B}=1{,}223$; $Q_{1 B}=1{,}12$}
		{$M_{A}=1{,}22$;  $Q_{1 A}=1{,}45$; $M_{B}=1{,}223$; $Q_{1 B}=1{,}12$}
		 {$M_{A}=1{,}22$; $Q_{1 A}=1{,}15$; $M_{B}=1{,}43$; $Q_{1 B}=1{,}12$}
		{$M_{A}=1{,}02$;  $Q_{1 A}=1{,}15$; $M_{B}=1{,}223$; $Q_{1 B}=1{,}12$}
	 \loigiai{
		Gọi $x_{1}$, $x_{2}$,  $x_{3}$,$\ldots$, $x_{85}$ lần lượt là số lợn con giống $A$ theo thứ tự không giảm.\\
		 Do $x_{1}$,$\ldots$, $x_{8} \in[1{,}0; 1{,}1)$; $x_{9}$,$\ldots$, $x_{36} \in[1{,}1; 1{,}2)$; $x_{37}$,$\ldots$, $x_{68} \in[1{,}2; 1{,}3)$; $x_{69}$,{\ldots}, $x_{85} \in[1{,}3; 1{,}4)$.\\
		Trung vị của mẫu số liệu lợn con giống $A$ thuộc nhóm $[1{,}2; 1{,}3)$\\ 
		 \centerline{$M_{A}=1{,}2+\dfrac{\dfrac{85}{2}-36}{32} \cdot(1{,}3-1{,}2)=1{,}22.$}
		Gọi $y_{1}$, $y_{2}$, $y_{3}$,$\ldots$,  $y_{65}$ lần lượt là số lợn con giống $B$ theo thứ tự không giảm.\\
		Do $y_{1}$,$\ldots$, $y_{13} \in[1{,}0; 1{,}1)$; $y_{14}$,$\ldots$, $y_{27} \in[1{,}1; 1{,}2)$; $y_{28}$,$\ldots$, $y_{51}  \in[1{,}2; 1{,}3)$; $y_{52}$,$\ldots$, $y_{65} \in[1{,}3; 1{,}4)$.\\
		Trung vị của mẫu số liệu lợn con giống $B$ thuộc  nhóm $[1{,}2; 1{,}3)$\\
		 \centerline{$M_{B}=1{,}2+\dfrac{\dfrac{65}{2}-27}{24} \cdot(1{,}3-1{,}2)=1{,}223.$}
		Vậy cân nặng trung bình của lợn con  giống $A$ nhỏ hơn giống $B$.\\
		Tứ phân vị thứ nhất của dãy số liệu giống $A$ là $\dfrac{1}{2}(x_{21}+x_{22})$ thuộc nhóm $[1{,}1; 1{,}2)$ nên  tứ phân vị thứ nhất của mẫu số liệu là $$Q_{1 A}=1{,}1+\dfrac{\dfrac{85}{4}-8}{28}\cdot (1{,}2-1{,}1)=1{,}15.$$
		Tứ phân vị thứ nhất của dãy số liệu giống $B$ là $\dfrac{1}{2}(y_{16}+y_{17})$ thuộc nhóm $[1{,}1; 1{,}2)$ nên tứ phân vị thứ nhất của mẫu số liệu  là $$Q_{1B}=1{,}1+\dfrac{\dfrac{65}{4}-13}{14}\cdot (1{,}2-1{,}1)=1{,}12.$$
	} 
\end{ex}

% nguồn: Ok_TracNghiemLop11_Lan1
\begin{ex}%[1D5H2-3]%[ma:d3cd06d2]
	Doanh thu bán  hàng trong $20$ ngày được lựa chọn ngẫu nhiên của một của hàng được ghi lại ở bảng sau (đơn vị: triệu đồng)
	\begin{center}
		{\tabcolsep  = 2mm
		\begin{tabular}{|c|c|c|c|c|c|}
			 \hline
			Doanh thu  & $[5; 7)$ & $[7; 9)$ & $[9; 11)$ & $[11; 13)$ & $[13; 15)$ \\
			 \hline
			 Số ngày & $2$ & $7$ & $7$ & $3$ & $1$\\
			 \hline
		\end{tabular}
		 }
	\end{center}
	Mốt của mẫu số liệu trên thuộc khoảng  nào trong các khoảng dưới đây?
	 \choice
		{$[7;  9)$}
		{\True  $[9; 11)$}
		{$[11;  13)$}
		{$[13;  15)$}
	 \loigiai{
	Có $2$ nhóm chứa mốt của mẫu số liệu trên đó là $[7; 9)$ và $[9; 11)$. \\ 
	Do  đó, ta xét
	\begin{itemize}
		\item Xét nhóm $[7;  9)$ ta có
		$\mathrm{M}_0 = 7  + \dfrac{7-2}{(7-2)+(7-7)} (9-7) =9$.
		\item Xét nhóm  $[9; 11)$ ta có
		$\mathrm{M}_0 = 9 + \dfrac{7-7}{(7-7)+(7-3)} (11-9)  =9$.
	\end{itemize}
	Vậy mốt của mẫu số  liệu trên là $9$ thuộc nhóm $[9; 11)$.
	 }
\end{ex}

% nguồn: Ok_Toan11_TracNghiem_Lan2
\begin{ex}%[1D5V2-3]%[ma:73693e45]
	Kết quả bài  kiểm tra toán Giữa HKI của lớp 11 A được ghi lại ở bảng sau
	\begin{longtable}{|*{6}{c|}}
		\hline 
		 Điểm số & [$0$; $2$) & [$2$; $4$) & [$4$; $6$) & [$6$; $8$) & [$8$; $10{,}5$) \\
		 \hline
		Số học sinh &  $1$ & $5$ & $22$ & $10$ & $7$ \\
		 \hline
	\end{longtable}
	Dựa vào bảng số liệu trên, giáo viên toán có thể nhận định $50\%$ học sinh trong lớp có điểm từ bao nhiêu  trở lên?
	 \choice
		{$4{,}5$} 
		{$5{,}0$} 
		{\True  $5{,}5$}
		 {$6{,}0$}
	\loigiai{ 
		Số học sinh trong lớp 10A1 là:  $n=1+5+22+10+7=45$.\\
		Gọi $x_1$; $ x_2$; $ x_3$; $  \ldots$; $ x_{45}$ lần lượt là điểm số của các học sinh theo thứ tự không giảm.\\
		Do $x_1 \in \left[0; 2\right)$; $ x_2, \ldots, x_6 \in \left[2; 4\right)$; $ x_7,  \ldots, x_{28} \in \left[4; 6\right)$; $ x_{29}, \ldots, x_{38} \in \left[6; 8\right)$; $ x_{39}$; $ \ldots$; $ x_{45} \in \left[8; 10{,}5\right)$.\\
		Trung vị của dãy số liệu là $x_{23} \in \left[4;  6\right)$.\\
		Ta có  $n=45$; $ p=3$; $ m_3=22$; $ m_1+m_2=1+5=6$; $ a_3=4$; $ a_4=6$.\\
		Trung vị của mẫu số liệu ghép  nhóm là: $M_e=4+\dfrac{\dfrac{45}{2}-6}{22} \cdot \left(6-4\right)=5{,}5$.\\
		Vậy có $50\%$ học  sinh trong lớp 10A1 từ $5{,}5$ điểm trở lên.
	} 
\end{ex}

% nguồn: Ok_Toan11_TracNghiem_Lan2
\begin{ex}%[1D5V2-3]%[ma:1d545f4d]
	Phỏng vấn một số học sinh khối 11 về thời gian (giờ) ngủ của  một buổi tối, thu được bảng số liệu sau:
	\begin{center}\begin{tabular}{|c|c|c|c|c|c|}
		\hline 
		Thời  gian (giờ)&$[4;5)$&$[5;6)$&$[6;7)$&$[7;8)$&$[8;9)$\\
		 \hline
		Số  lượng&6&12&13&10&3\\
		 \hline
	\end{tabular}
\end{center}
Hãy cho biết $75 \% $ học sinh khối 11 ngủ nhiều nhất  bao nhiêu giờ?
\choice 
	{$6$  giờ}
	 {$6{,}5$ giờ}
	{$7$ giờ} 
	 {\True $7{,}2$ giờ}
 \loigiai{
	 Cỡ mẫu là $n=6+12+13+10+3=44$.\\
	 Tứ phân vị thứ hai $Q_2$ là giá trị của $\dfrac{x_{22}+x_{23}}{2}$.\\
	Tứ  phân vị thứ nhất $Q_1$ là giá trị của $\dfrac{x_{11}+x_{12}}{2}$.\\
	Tứ phân vị  thứ ba $Q_3$ là giá trị của $\dfrac{x_{33}+x_{34}}{2}$.\\
	Do  $x_{33};x_{34}$ thuộc nhóm $\left[7;8\right)$ nên nhóm này chứa $Q_3$.\\
	Do  đó $p=4;a_4=7;m_4=10;m_1+m_2+m_3=6+12+13=31;a_5-a_4=1$ và ta có
	\begin{eqnarray*}
		Q_3&=&a_p+\dfrac{\dfrac{3n}{4}-\left(m_1+\ldots +m_{p-1}\right)}{m_p}\cdot \left(a_{p+1}-a_p\right)\\&=&a_4+\dfrac{\dfrac{132}{4}-\left(m_1+m_2+m_3\right)}{m_4}\cdot \left(a_5-a_4\right)\\&=&7+\dfrac{33-31}{10}\cdot 1=7{,}2. 
	\end{eqnarray*}
Vậy tứ phân vị thứ ba $Q_3=7{,}2$ nên  $75 \% $ học sinh khối 11 ngủ nhiều nhất là $7{,}2$ giờ.}
\end{ex}

% nguồn: Ok_Toan11_TracNghiem_Lan2
\begin{ex}%[1D5V2-3]%[ma:a88d1c01]
	Một phòng khám thống kê số  bệnh nhân đến khám mỗi ngày trong tháng 6 năm 2020 ở bảng sau:
	\begin{center}
		\begin{tabular}{|c|c|c|c|c|c|}
			\hline 
			Số bệnh  nhân&$[1;10)$&$[10;20)$&$[20;30)$&$[30;40)$&$[40;50)$\\
			\hline 
			Số ngày&7&5&10&6&2\\ 
			 \hline
		\end{tabular}
	\end{center}
	Số bệnh nhân đến khám nhiều hơn 23 người trên ngày chiếm tỉ lệ  bao nhiêu phần trăm?
	\choice 
		{$25  \% $}
		 {\True $50 \% $}
		 {$75 \% $}
		{$10  \% $}
	 \loigiai{
		Gọi $x_1,x_2,\ldots,x_{30}$ là số bênh nhân đến khám mỗi ngày và giả sử dãy này đã được sắp xếp theo thứ tự  tăng dần. Khi đó tứ phân vị thứ hai của dãy số liệu là $\dfrac{x_{15}+x_{16}}{2}$ thuộc $\left[20;30\right)$.\\
		$Q_2=20+\dfrac{\dfrac{30}{2}-\left(7+5\right)}{10}\cdot  10=23$.\\
	Như vậy có khoảng $50 \% $ số  ngày có nhiều hơn 23 bệnh nhân đến khám.}
\end{ex}

% nguồn: Ok_TracNghiemLop11_Lan1
\begin{ex}%[1D5V2-3]%[ma:48f4e1d7]
	Trong một cuộc đua Marathon được tổ chức ở thành  phố $A$ người ta thống kê lại được như sau
	\begin{center}
		\begin{tabular}{|l|c|c|c|c|c|}
			 \hline
			Thời gian  & $\left[120;140\right)$ & $\left[140;160\right)$ & $\left[160;180\right)$ &
			$\left[180;200\right)$ &  $\left[200;220\right)$ \\
			 \hline
			Số người &  $4$ & $6$ & $10$ & $15$ & $25$ \\
			\hline 
		\end{tabular}
	\end{center}
	Tứ phân vị thứ nhất của mẫu số liệu ghép nhóm trên  là
	\choice 
		 {$150$}
		 {$160$}
		 {\True $170$}
		 {$180$}
	 \loigiai{
		Ta có $n=4+6+10+15+25=60$. \\ 
		Trung vị  là $M_e=180+\dfrac{\dfrac{60}{2}-20}{15}\cdot \left(200-180\right)=\dfrac{580}{3}\approx 193{,}33$.\\
		Tứ phân vị thứ nhất là  $Q_1=160+\dfrac{\dfrac{60}{4}-10}{10}\cdot \left(180-160\right)=170$.
	 }
\end{ex}

% nguồn: Ok_TracNghiemLop11_Lan1
\begin{ex}%[1D5V2-3]%[ma:b33d53b5]
	Tổng  lượng mưa trong tháng $8$ đo được tại một trạm quan trắc đặt tại Vũng Tàu từ năm $2002$ đến năm $2020$ được ghi lại như dưới đây (đơn vị: mm)
	\begin{center}
		$ 
		\begin{array}{llllllllll}
			121{,}8 & 158{,}3 & 334{,}9 &  200{,}9 & 165{,}6 & 161{,}5 & 194{,}3 & 220{,}7 & 189{,}8 & 234{,}2 \\
			165{,}9 & 165{,}9 &  134 & 173 & 169 & 189 & 254 & 168 & 255 &
		\end{array}
		$\\ 
		\textit{(Nguồn: Tổng cục Thống  kê)}.
	\end{center}
	Hoàn thiện bảng  tần số ghép nhóm theo mẫu sau và tìm tứ phân vị thứ hai của mẫu số liệu ghép nhóm đó.
	\begin{center}
		\begin{tabular}{|c|c|c|c|c|}
			\hline Tổng lượng mưa trong tháng $8$ (mm)  & [120;175) & [175;230) & [230;285) & [285;340) \\
			\hline Số năm & $ x $ & $ y $  & $z$ & $t$ \\
			\hline 
		\end{tabular}
	\end{center}
	 \choice
		{\True $  x=10; y=5; z=3; t=1; Q_{2}=172{,}5 $}
		{ $ x=9; y=6; z=3;  t=1; Q_{2}=172{,}5 $}
		{ $ x=10; y=5; z=2; t=2; Q_{2}=182{,}5  $}
		{ $ x=10; y=4;  z=4; t=1; Q_{2}=162{,}5 $}
	 \loigiai{
		Gọi $ x_{1}; x_{2}; x_{3}; \ldots; x_{19} $ lần lượt là  số năm theo thứ tự không giảm.\\
		Tứ phân vị thứ hai của dãy số liệu là $ x_{10} $ thuộc nhóm $[120;175)$ nên tứ phân vị thứ hai của mẫu  số liệu là $ Q_{2}=120+\dfrac{\dfrac{19}{2}-0}{10}\cdot (175-120)=172{,}5 $.
	 }
\end{ex}

% nguồn: Ok_TracNghiemLop11_Lan1
\begin{ex}%[1D5V2-3]%[ma:d1f3d9ba]
	Một hãng ô  tô thống kê lại số lần gặp sự cố về động cơ của $100$ chiếc xe cùng loại sau $2$ năm sử dụng đầu tiên ở bảng sau
	\begin{center}
		\begin{tabular}{|l|c|c|c|c|c|}
			\hline Số lần gặp sự cố  & {$[1; 2]$} & {$[3; 4]$} & {$[5; 6]$} & {$[7; 8]$} & {$[9; 10]$} \\
			\hline Số xe & $17$ & $33$ & $25$ & $20$  & $5$ \\
			\hline 
		\end{tabular}
	\end{center}
	Hãy ước lượng  tứ phân vị thứ nhất của mẫu số liệu ghép số trên.
	 \choice
		 {$2{,}64$}
		{$2,89 $} 
		{$2,73$} 
		{\True $2,98$} 
	 \loigiai{
		 Do số lần gặp sự cố là số nguyên nên ta hiệu chỉnh lại như sau:
		\begin{center}
			\begin{tabular}{|l|c|c|c|c|c|}
				\hline Số lần gặp sự cố & {$[0{,}5; 2{,}5)$}  & {$[2{,}5; 4{,}5)$} & {$ [4{,}5; 6{,}5)$} & {$[6{,}5; 8{,}5)$} & {$[8{,}5; 10{,}5)$} \\
				\hline Số  xe & $17$ & $33$ & $25$ & $20$ & $5$ \\
				\hline 
			\end{tabular}
		\end{center}
		Gọi $x_{1}; x_{2}; \ldots; x_{100}$ là mẫu số liệu được xếp  theo thứ tự không giảm.\\
		Ta có $x_{1}, \ldots, x_{17} \in[0{,}5;  2{,}5); x_{18}, \ldots, x_{50} \in[2{,}5; 4{,}5); x_{51}, \ldots, x_{75} \in[4{,}5; 6{,}5)$;\\
		$x_{76}, \ldots, x_{95} \in[6{,}5; 8{,}5); x_{96}, \ldots, x_{100} \in[8{,}5;  10{,}5)$.\\
		Tứ phân vị thứ nhất của dãy số  liệu $x_{1}; x_{2}; \ldots; x_{100}$ là $ \dfrac{1}{2}\left(x_{25}+x_{26}\right)$.\\
		 Do $ x_{25}$ và $x_{26}$ thuộc nhóm $[2{,}5; 4{,}5)$ nên tứ phân vị thứ nhất của mẫu số liệu ghép nhóm là
		$ 
		Q_{1}=2{,}5+\dfrac{\dfrac{1\cdot 100}{4}-17}{33}\cdot  (4{,}5-2{,}5)=\dfrac{197}{66} \approx 2{,}98$.
	 }
\end{ex}

% nguồn: Ok_TracNghiemLop11_Lan1
\begin{ex}%[1D5B2-3]%[ma:de9650be]
	Một người thống kê lại  thời gian thực hiện các cuộc gọi điện thoại của người đó trong một tuần cho trong bảng sau
	\begin{center}
		\begin{tabular}{|l|c|c|c|c|c|}
			\hline Số  bệnh nhân &{$[0;60)$}&{$[60;120)$}&{$[120;180)$}&{$[180;240)$}&{$[240;300)$}\\
			 \hline Số ngày & $8$ & $10$ & $7$ & $5$ & $2$ \\
			\hline 
		\end{tabular}
	\end{center}
	Tìm tứ phân vị thứ  hai của mẫu số liệu.
	\choice 
		{$Q_2\approx  80{,}25$}
		{$Q_2\approx  100{,}32$}
		 {$Q_2=115$}
		 {\True $Q_2=108$}
	\loigiai{ Cỡ  mẫu là $n=32$.\\
		Gọi $x_1, \ldots, x_{32}$ là mẫu số liệu được sắp xếp theo thứ tự không giảm. Khi đó, trung  vị là $\dfrac{x_{16}+x_{17}}{2}$.\\
		Do đó, tứ phân vị thứ hai $Q_2$ là $\dfrac{x_{16}+x_{17}}{2}$. Do $x_{16} $, $x_{17} \in [60;120)$ nên nhóm này  chứa $Q_2$. \\Do đó, $p= 2; \;a_2=60;\; m_2=10; \;m_1=8; \;a_3-a_2=60$ và ta có
		 $$
		Q_3=60+\dfrac{\dfrac{ 32}{2}-8}{10}\cdot 60=108. 
		 $$
	 }
\end{ex}


%%%%%%%%%%%%%%%% ĐÚNG/SAI %%%%%%%%%%%%%%%%

% nguồn: Ok_dungsai_lop11_Lan1
\begin{ex}%[1D5H2-3]%[ma:1c92f7d4]
 Một hãng xe ô tô thống kê lại số lần gặp sự cố về động cơ của 100 chiếc xe  cùng loại sau $2$ năm sử dụng đầu tiên ở bảng sau
 \begin{center}
  \begin{tabular}{|c|c|c|c|c|c|}
   \hline 
   Số lần gặp sự cố & $[1;2]$ & $[3;4]$ & $[5;6]$  & $[7;8]$ & $[9;10]$ \\
    \hline
   Số xe & $17$ & $33$ & $25$ & $20$ & $5$ \\ 
    \hline
  \end{tabular}
 \end{center}
  Khi đó
 \choiceTF 
  {\True Cỡ mẫu của mẫu số liệu  là $n=100$}
  {\True Tứ phân vị thứ hai của mẫu số  liệu ghép nhóm là $Q_2=4{,}5$}
  {Tứ  phân vị thứ nhất của mẫu số liệu ghép nhóm là $Q_1 \approx 1{,}98$}
  {\True  Tứ phân vị thứ ba của mẫu số liệu ghép nhóm là $Q_3=6{,}5$}
  \loigiai{
  Số lần gặp sự cố là số nguyên nên ta có thể sử dụng bảng tần  số ghép nhóm sau
  \begin{center}
   \begin{tabular}{|c|c|c|c|c|c|}
    \hline 
    Số  lần gặp sự cố & $[0{,}5;2{,}5)$ & $[2{,}5;4{,}5)$ & $[4{,}5;6{,}5)$ & $[6{,}5;8{,}5)$ & $[8{,}5;10{,}5)$ \\
     \hline
    Số xe & $17$ & $33$ & $25$ & $20$ & $5$ \\ 
     \hline
   \end{tabular}
  \end{center}
  \begin{itemchoice}
   \itemch \textbf{Đúng.} \\ 
   Cỡ mẫu của mẫu số  liệu là $n=100$.
   \itemch \textbf{Đúng.} \\ 
   Gọi $x_1$, $x_2$, $x_3$, $\ldots$, $x_{100}$ là mẫu số liệu được xếp theo thứ tự không  giảm.\\
   Trung vị của của  mẫu số liệu là $\dfrac{x_{50}+x_{51}}{2}$ với $x_{50} \in [2{,}5;4{,}5),x_{51} \in [4{,}5;6{,}5)$.\\
   Suy  ra tứ phân vị thứ hai cũng là trung vị của mẫu số liệu ghép nhóm là $Q_2=4{,}5$.
   \itemch \textbf{Sai.}  \\
   Xét nửa mẫu số liệu bên trái $x_1$, $x_2$, $x_3$, $\ldots$, $x_{50}$ có trung vị $\dfrac{x_{25}+x_{26}}{2} \in  [2{,}5;4{,}5)$.
   Ta có $n_m=33$; $C=17$; $u_{m+1}=4{,}5$;  $u_m=2{,}5$.
   Vì vậy, tứ phân vị thứ nhất  của mẫu số liệu ghép nhóm là
   \[Q_1=2{,}5+\dfrac{\dfrac{100}{4}-17}{33}  (4{,}5-2{,}5)=\dfrac{197}{66} \approx 2{,}98.
   \] 
   \itemch  \textbf{Đúng.} \\
   Xét nửa mẫu số liệu  bên phải $x_{51}$, $x_{52}$, $\ldots$, $x_{100}$ có trung vị
   $\dfrac{x_{75}+x_{76}}{2}$ và \break $  x_{75} \in [4{,}5;6{,}5)$, $x_{76} \in [6{,}5;8{,}5)$ nên tứ phân vị thứ ba của mẫu số liệu ghép nhóm $Q_3=6{,}5$.
   Vậy  các tứ phân vị của mẫu số liệu ghép nhóm là
   \[Q_1  \approx 2{,}98; \,Q_2\approx 4{,}5;\, Q_3\approx6{,}5.\]
  \end{itemchoice}
  }
\end{ex}

% nguồn: Ok_dungsai_lop11_Lan1
\begin{ex}%[1D5H2-3]%[ma:97a1af34]
 Hãy tìm các tứ phân vị của mẫu số liệu được cho dưới dạng bảng tần số ghép  nhóm sau
 \begin{center}
  \begin{tabular}{|c|c|c|c|c|c|}
    \hline
   Nhóm  & $[0;2)$ & $[2;4)$ & $[4;6)$ & $[6;8)$ & $[8;10)$ \\
   \hline 
   Tần số & $3$ & $8$ & $12$ & $12$ &  $4$ \\
   \hline 
  \end{tabular}
 \end{center}
  \choiceTF
  {Cỡ mẫu của mẫu số liệu  là $n=38$}
  {\True  Tứ phân vị thứ hai của mẫu số liệu ghép nhóm là $Q_2\approx 5{,}42$}
  {Tứ phân vị thứ nhất của mẫu số liệu ghép nhóm là $Q_1\approx  2{,}69$}
  {\True Tứ phân vị thứ ba của  mẫu số liệu ghép nhóm là $Q_3=7{,}04$}
 \loigiai{ 
  \begin{itemchoice}
    \itemch Sai. Cỡ mẫu của mẫu số liệu là $n=3+8+12+12+4=39$.
   \itemch Đúng. Gọi $x_1{,} x_2{,} \ldots{,} x_{39}$ là mẫu số liệu được sắp xếp theo thứ tự không  giảm.
   Trung vị của mẫu số liệu này  là $x_{20} \in [4;6)$.
    Ta có $n_m=12$; $C_2=3+8=11$; $u_m=4$; $u_{m+1}=6$.
   Tứ phân vị  thứ hai chính là trung vị của mẫu số liệu ghép nhóm là
   \[Q_2=M_e=4+\dfrac{\dfrac{39}{2}-11}{12} (6-4)=\dfrac{65}{12} \approx 5{,}42. 
   \] 
    \itemch Sai. Xét nửa mẫu số liệu bên trái $x_1{,} x_2{,} \ldots{,} x_{19}$ có trung vị $x_{10} \in [2;4)$.
   Ta có  $n_i=8$; $C_1=3$; $x_i=2$; $x_{i+1}=4$.\\
   Suy ra tứ phân vị  thứ nhất của mẫu số liệu là $Q_1=2+\dfrac{\dfrac{39}{4}-3}{8} (4-2)=\dfrac{59}{16} \approx 3{,}69$.
    \itemch Đúng. Xét nửa mẫu số liệu bên phải $x_{21}{,} x_{22}{,} \ldots{,} x_{39}$ có trung vị $x_{30} \in [6;8)$.
   Ta có $n_j=12$;  $C_3=3+8+12=23$; $x_j=6$; $x_j=8$.\\
   Suy ra tứ phân vị thứ ba của mẫu số liệu là $Q_3=6+\dfrac{\dfrac{3\cdot39}{4}-23}{12} (8-6)=\dfrac{169}{24} \approx 7{,}04$. 
  \end{itemchoice}
  Vậy các tứ phân vị của mẫu  số liệu ghép nhóm là $Q_1\approx 3{,}69$; $Q_2=5{,}42$; $Q_3=7{,}04$.
  }
\end{ex}

% nguồn: Ok_dungsai_lop11_Lan1
\begin{ex}%[1D5H2-3]%[ma:0138e56f]
 Thống kê điểm trung bình môn Toán của một số học sinh lớp 11 được cho ở  bảng sau:
 \begin{center}
  \begin{tabular}{|c|c|c|c|c|c|c|c|}
    \hline
   Khoảng điểm & $[6{,}5;7)$ & $[7;7{,}5)$ & $[7{,}5;8)$ & $[8;8{,}5)$ & $[8{,}5;9)$ & $[9;9{,}5)$  & $[9{,}5;10)$ \\
    \hline
   Số học sinh & 8 & 10 & 16 & 24 & 13  & 7 & 4 \\
   \hline 
  \end{tabular}
 \end{center}
 Khi đó 
  \choiceTF
  {Cỡ  mẫu của mẫu số liệu là $n = 80$}
   {\True Tứ phân vị thứ hai của mẫu số liệu ghép nhóm là $Q_2 \approx 8{,}15$}
  {\True Tứ phân vị thứ  nhất của mẫu số liệu ghép nhóm là $Q_1 \approx 7{,}58$}
  {\True Tứ phân vị  thứ ba của mẫu số liệu ghép nhóm là $Q_3 \approx 8{,}63$}
  \loigiai{
  \begin{itemchoice}
   \itemch \textbf{Sai}. \\ 
   Vì cỡ mẫu số  liệu là $n =80$.
   \itemch \textbf{Đúng}.\\ 
   Gọi $x_1, x_2, \ldots, x_{82}$ lần  lượt là điểm trung bình môn Toán của các học sinh sắp xếp theo thứ tự không giảm. Ta có $x_1, \ldots, x_8 \in [6{,}5; 7)$; $x_9, \ldots, x_{18} \in [7; 7{,}5)$; $x_{19}, \ldots, x_{34} \in [8; 8{,}5)$; $x_{35}, \ldots, x_{58} \in [8{,}5; 9)$; $x_{59}, \ldots, x_{71} \in [7; 7{,}5)$; $x_{72}, \ldots, x_{78} \in [9; 9{,}5)$; $x_{79}, \ldots, x_{82} \in [9{,}5; 10)$.
   Tứ phân vị thứ hai của mẫu số liệu là $\dfrac{1}{2}(x_{41} +  x_{42})$ mà $x_{41}, x_{42} \in [8; 8{,}5)$ nên tứ phân vị thứ hai của mẫu số liệu là $Q_2 = 8 + \dfrac{\dfrac{82}{2} - 34}{24} (8{,}5 - 8) \approx 8{,}15$.
   \itemch  \textbf{Đúng}.\\
   Tứ phân vị thứ  nhất của mẫu số liệu là $x_{21} \in [7{,}5; 8)$. Do đó tứ phân vị thứ nhất của mẫu số liệu là $$Q_{1}= 7{,}5 + \dfrac{\dfrac{82}{4}-18}{16}(8-7{,}5) \approx 7{,}58.$$
   \itemch  \textbf{Đúng}. \\
   Tứ phân vị thứ ba của mẫu số liệu là $x_{62} \in [8; 5{,}9)$. Do đó tứ phân vị thứ ba của mẫu số liệu là $$Q_{3}= 8{,}5  + \dfrac{\dfrac{3\cdot 82}{4}-58}{13}(9-8{,}5) \approx 8{,}63.$$
  \end{itemchoice}
  }
\end{ex}

% nguồn: Ok_dungsai_lop11_Lan1
\begin{ex}%[1D5H2-3]%[ma:10e3d6dd]
 Cho bảng số liệu ghép  nhóm sau
 \begin{center}
  \begin{tabular}{|c|c|c|c|c|c|}
   \hline 
   Nhóm &  $[0;2)$ & $[2;4)$ & $[4;6)$ & $[6;8)$ & $[8;10)$ \\
   \hline 
   Tần số & $3$ & $8$ & $12$ & $12$ &  $4$ \\
   \hline 
  \end{tabular}
 \end{center}
 Khi  đó
  \choiceTF
  {Cỡ mẫu của mẫu số liệu là $n=38$} 
  {\True Tứ phân vị thứ hai của mẫu số liệu ghép  nhóm là $Q_2 \approx 5{,}42$}
  {\True Tứ phân vị thứ  nhất của mẫu số liệu ghép nhóm là $Q_1 \approx 2{,}69$}
  {\True Tứ phân vị thứ ba của mẫu số liệu ghép  nhóm là $Q_3\approx7{,}04$}
  \loigiai{
  \begin{itemchoice}
   \itemch \textbf{Sai.} \\ 
   Cỡ mẫu của mẫu số liệu là $n=3+8+12+12+4=39$. 
   \itemch \textbf{Đúng.} \\ 
    Gọi $x_1$, $x_2$, $\ldots$, $x_{39}$ là mẫu số liệu được sắp xếp theo thứ tự không giảm.\\
   Trung vị của mẫu số liệu này là  $x_{20} \in [4;6)$.\\
   Ta có $n_m=12$; $C_2=3+8=11$; $u_m=4$;  $u_{m+1}=6$.\\
   Tứ phân vị thứ hai chính là  trung vị của mẫu số liệu ghép nhóm là
   \[Q_2=M_e=4+\dfrac{\dfrac{39}{2}-11}{12}  (6-4)=\dfrac{65}{12} \approx 5{,}42.
    \]
   \itemch \textbf{Đúng.}  \\
   Xét nửa mẫu  số liệu bên trái $x_1$, $x_2$, $\ldots$, $x_{19}$ có trung vị $x_{10} \in [2;4)$.\\
   Ta có $n_i=8$; $C_1=3$; $x_i=2$; $x_{i+1}=4$. Suy ra  tứ phân vị thứ nhất của mẫu số liệu là
    \[Q_1=2+\dfrac{\dfrac{39}{4}-3}{8} (4-2)=\dfrac{59}{16} \approx 3,69.\]
    \itemch \textbf{Đúng.} \
   Xét nửa mẫu số liệu bên  phải $x_{21}$, $x_{22}$, $\ldots$, $x_{39}$ có trung vị $x_{30} \in [6;8)$.\\
   Ta  có $n_j=12$; $C_3=3+8+12=23$; $x_j=6;x_j=8$.\\
   Suy  ra tứ phân vị thứ ba của mẫu số liệu là $Q_3=6+\dfrac{\dfrac{3.39}{4}-23}{12} (8-6)=\dfrac{169}{24} \approx 7{,}04$.\\
   Vậy các tứ  phân vị của mẫu số liệu ghép nhóm là $Q_1 \approx 3{,}69$; $Q_2\approx5{,}42$; $Q_3\approx7{,}04$.
  \end{itemchoice}
  }
\end{ex}

% nguồn: Ok_dungsai_lop11_Lan1
\begin{ex}%[1D5H2-3]%[ma:b1ab760e]
 Tiền lương tháng của một số nhân viên ở 01 công ty được  ghi lại như sau (đơn vị: triệu đồng)
 \begin{center}
  \begin{tabular}{llllllllllllllll}
   13{,}8 & 6{,}7 & 8{,}4 & 11{,}9 & 11{,}1 & 8{,}3 &  13{,}2 & 11{,}2 &
   8{,}9 & 10{,}7 & 6{,}5 &  13{,}1 & \\
   12{,}5 & 9{,}6 & 11{,}7 & 12{,}7 & 10{,}0 & 10{,}0  & 12{,}2 & 9{,}8 &
   10{,}9 & 6{,}7 & 13{,}6  & 9{,}2 & \\
  \end{tabular}
 \end{center}
  Khi đó
  \choiceTF
  {\True Tứ phân  vị thứ hai của dãy số liệu là $Q_2=10{,}8$}
  {Tứ phân vị thứ nhất của dãy số liệu là  $Q_1=6{,}05$}
  {\True Tổng hợp lại  dãy số liệu trên vào bảng tần số ghép nhóm theo mẫu sau
  \begin{tabular}{|c|c|c|c|c|}
   \hline 
   Lương tháng (triệu đồng) & $[6; 8)$ & $[8;10)$ & $[10; 12)$  & $[12;14)$\\
    \hline
   Số nhân viên & 3 & 6 & 8 & 7 \\ 
    \hline
 \end{tabular}}
  {Ước lượng tứ phân vị của số liệu ở bảng tần số ghép nhóm trên ta được $Q_2=6{,}75$}
 \loigiai{ 
  Sắp xếp mẫu số  liệu theo thứ tự không giảm ta được
  \begin{center}
   \begin{tabular}{llllllllllll}
    6{,}5 & 6{,}7  & 6{,}7 & 8{,}3 & 8{,}4 & 8{,}9 & 9{,}2 & 9{,}6 & 9{,}8 & 10{,}0 & 10{,}0 & 10{,}7 \\
    10{,}9 & 11{,}1 & 11{,}2 & 11{,}7 & 11{,}9 & 12{,}2 & 12{,}5 &  12{,}7 & 13{,}1 & 13{,}2 & 13{,}6 & 13{,}8 \\
   \end{tabular}
  \end{center}
  \begin{itemchoice}
   \itemch \textbf{Đúng}.  \\
   Vì cỡ của mẫu số liệu là $24$ nên tứ phân vị thứ  hai của mẫu số liệu là $$Q_2= \dfrac{x_{12}+x_{13}}{2}= \dfrac{
    10{,}7+10{,}9}{2} =10{,}8.$$
   \itemch  \textbf{Sai}. \\
   Tứ phân vị thứ nhất là  trung vị của mẫu số liệu
   \begin{center}
    \begin{tabular}{llllllllllll}
     6{,}5 & 6{,}7  & 6{,}7 & 8{,}3 & 8{,}4 & 8{,}9 & 9{,}2 & 9{,}6 & 9{,}8 & 10{,}0 & 10{,}0 & 10{,}7 \\
    \end{tabular}
   \end{center}
   Vì cỡ mẫu là $12$  nên trung vị của mẫu số liệu trên là $$Q_1= \dfrac{x_6+x_7}{2}= \dfrac{8{,}9+9{,}2}{2}= 9{,}05.$$
   Vậy tứ phân  vị thứ nhất của mẫu số liệu là $Q_1= 9{,}05$.
   \itemch \textbf{Đúng}.  \\
   Vì  ta có bảng tần số ghép nhóm sau
   \begin{center}
    \begin{tabular}{|c|c|c|c|c|}
     \hline 
     Lương tháng (triệu đồng) & $[6; 8)$ & $[8;10)$ & $[10; 12)$ &  $[12;14)$\\
      \hline
     Số nhân viên & 3 & 6  & 8 & 7 \\
     \hline 
    \end{tabular}
   \end{center}
   \itemch \textbf{Đúng}.  Gọi $x_1, x_2, \ldots , x_{24}$ lần lượt là lương tháng của mỗi nhân viên được xếp theo thứ tự không giảm.\\
   Do $x_1, x_2, x_3 \in [6;8); x_4, \ldots , x_9 \in [8;10); x_{10}, \ldots ,x_{17} \in [10;12); x_{18}, \ldots , x_{24} \in [12;14)$ nên trung vị của mẫu là $$\dfrac{1}{2}(x_{12}+x_{13}) \in  [10;12).$$
   Ta xác định được $n=24, n_m =8 , C=3+6=9, u_m=  10, u_{m+1}=12$. Suy ra tứ phân vị thứ hai của mẫu số liệu ghép nhóm trên là
   $$Q_2= 10 + \dfrac{\dfrac{24}{2}-9}{8}(12-10)=10{,}75.$$ 
  \end{itemchoice}
 } 
\end{ex}

% nguồn: Ok_dungsai_lop11_Lan1
\begin{ex}%[1D5H2-3]%[ma:1e526498]
 Kiểm tra điện lượng của một số viên pintiểu do một hãng sản xuất thu được kết quả sau \begin{center}
 \begin{tabular}{|c|c|c|c|c|c|}
  \hline  Điện lượng (Nghìn mAh) & {$[0,9; 0,95)$} & {$[0,95; 1,0)$} & {$[1,0; 1,05)$} & {$[1,05; 1,1)$} & {$[1,1; 1,15)$} \\
  \hline Số pin &  10 & 20 & 35 & 15 & 5 \\
   \hline
 \end{tabular}
\end{center}
\choiceTF 
 {\True Số trung bình của  dãy số liệu là $1{,}016$ }
 {Nhóm chứa mốt của dãy số liệu là $[1{,}05; 1{,}1)$} 
 {\True Tứ phân  vị thứ nhất của mẫu số liệu nhóm là $Q_1=0{,}98$}
 {Tứ  phân vị thứ ba của mẫu số liệu nhóm là $Q_3=1{,}248$}
\loigiai{ 
 \begin{center}
  \begin{tabular}{|c|c|c|c|c|c|}
   \hline Điện lượng (Nghìn mAh) & {$[0,9; 0,95)$} & {$[0,95; 1,0)$} & {$[1,0; 1,05)$} &  {$[1,05; 1,1)$} & {$[1,1; 1,15)$} \\
   \hline Giá trị đại diện &  0,925 & 0,975 & 1,025 & 1,075 & 1,125 \\
   \hline Sồ trận & 10 & 20  & 35 & 15 & 5 \\
    \hline
  \end{tabular}
 \end{center}
 \begin{itemchoice}
  \itemch \textbf{Đúng.} Số trung bình của dãy số liệu là 
   $$
  \dfrac{0{,}925\cdot 10+0{,}975\cdot  20+1{,}025\cdot 35+1{,}075\cdot 15+1{,}125\cdot 5}{85}=1{,}016 .
   $$
  \itemch \textbf{Sai.} Nhóm chứa mốt của dãy số liệu là $[1; 1{,}05)$ 
   $$
  M_{o}=1+\dfrac{35-20}{(35-20)+(35-15)}(1,05-1)=1{,}02. 
   $$
  \itemch \textbf{Đúng.} Gọi $x_1; x_2; \ldots; x_{85}$ lần lượt là điện lượng mỗi viên pin xếp theo thứ tự không giảm. \\Do $x_1; x_2; \ldots; x_{10} \in[0,9; 0,95); x_{11}; \ldots; x_{30} \in[0,95; 1,0);  x_{31}; \ldots; x_{65} \in[1,0; 1,05)$;\\
  $x_{66}; \ldots; x_{80} \in[1,05; 1,1); x_{81}; \ldots; x_{85} \in[1,1; 1,15)$ nên trung vị của mẫu  số liệu $x_1; x_2; \ldots; x_{85}$ là $x_{43} \in[1; 1{,}05)$.\\
  Ta xác định  được $n=85, n_m=35, C=30, u_m=1, u_{m+1}=1,05$.\\
  Suy ra tứ  phân vị thứ hai của mẫu số liệu là $Q_2=1+\dfrac{\tfrac{85}{2}-30}{35}(1,05-1)=1,02$.\\
   Tứ phân vị thứ nhất của mẫu số liệu là $\dfrac{1}{2}\left(x_{21}+x_{22}\right)$.\\
   Do $x_{21}, x_{22} \in[0,95; 1)$ nên tứ phân vị thứ nhất của mẫu số liệu nhóm là
  $$ 
   Q_1=0{,}95+\dfrac{\tfrac{85}{4}-10}{20}(1-0,95)=0{,}98.
  $$ 
  \itemch \textbf{Sai.} Tứ phân vị thứ  ba của mẫu số liệu là $\dfrac{1}{2}\left(x_{64}+x_{65}\right)$.\\
  Do $x_{64}, x_{65} \in[1; 1,05)$ nên tứ phân vị thứ ba của mẫu số liệu nhóm là $$Q_3=1+\dfrac{\tfrac{3\cdot  85}{4}-30}{35}(1{,}05-1)=1{,}048.$$
\end{itemchoice}}
\end{ex}

% nguồn: Ok_dungsai_lop11_Lan1
\begin{ex}%[1D5H2-3]%[ma:7131cadb]
 Cân nặng của một số lợn con mới sinh thuộc hai giống $A$ và $B$ được  cho ở bảng đây (đơn vị: kg)
 \begin{center}
  \begin{tabular}{|c|c|c|c|c|}
   \hline 
   Cân  nặng kg & $[1,0;1,1)$ & $[1,1;1,2)$ & $[1,2;1,3)$ & $[1,3;1,4)$ \\
    \hline
   Số con  giống $A$ & $8$ & $28$ & $32$ & $17$ \\
    \hline
   Số con giống $B$ & $13$ & $14$ &  $24$ & $14$ \\
    \hline
  \end{tabular}
 \end{center}
  \choiceTF
  {\True Cân nặng trung bình của giống  $A$ là $1{,}22$}
  {\True Cân nặng trung bình của giống $B$  là $1{,}21$}
  {\True Tứ phân vị thứ nhất của mẫu số liệu lợn  con giống $A$ là $Q_{1A}=1{,}15$}
  {Tứ phân vị  thứ nhất của mẫu số liệu lợn con giống $B$ là $Q_{1B}=1{,}62$}
  \loigiai{
  \begin{itemchoice}
   \itemch Đúng. Cân nặng trung bình của giống $A$  là
   \[\dfrac{1{,}05\cdot 8+1{,}15\cdot 28+1{,}25\cdot 32+1{,}35\cdot  17}{85}=1{,}22.
   \] 
   \itemch Đúng. Cân nặng trung bình của giống  $B$ là
    \[\dfrac{1{,}05\cdot 13+1{,}15\cdot 14+1{,}25\cdot 24+1{,}35\cdot 14}{65}=1{,}21.
   \] 
   \itemch Đúng. Gọi $x_1; x_2; \ldots; x_{85}$ lần lượt là cân nặng của mỗi con  lợn giống $A$ theo thứ tự không giảm.\\
   Tứ phân vị thứ nhất của dãy số liệu lợn con giống $A$ là $\dfrac{1}{2}  \left(x_{21}+x_{22} \right)\in [1{,}1;1{,}2)$ nên tứ phân vị thứ nhất của mẫu số liệu lợn con giống $A$ là
   \[Q_{1A}=1{,}1+\dfrac{\dfrac{85}{4}-8}{28}  (1{,}2-1{,}1)=1{,}15.
    \]
   \itemch Sai. Gọi $y_1;y_2;\ldots;y_{65}$ lần lượt là cân nặng của  mỗi con lợn giống $B$ theo thứ tự không giảm.\\
   Tứ phân vị thứ nhất của dãy số liệu lợn  con giống $B$ là $\dfrac{1}{2} \left(y_{16}+y_{17} \right)\in [1{,}1;1{,}2)$ nên tứ phân vị thứ nhất của mẫu số liệu lợn con giống $B$ là
   \[Q_{1B}=1{,}1+\dfrac{\dfrac{65}{4}-13}{14}  (1{,}2-1{,}1)=1{,}12.\]
  \end{itemchoice}
 } 
\end{ex}

% nguồn: Ok_dungsai_lop11_Lan1
\begin{ex}%[1D5H2-3]%[ma:8cb91a9a]
 Người ta đo đường kính của $61$ cây gỗ được trồng sau $12$ năm (đơn vị: centimét), họ thu được bảng tần số ghép nhóm sau 
 \begin{center}
  \begin{tabular}{|c|c|c|c|c|c|}
   \hline 
   Đường kính & $[20;25)$ & $[25;30)$  & $[30;35)$ & $[35;40)$ & $[40;45)$ \\
    \hline
   Số cây & $4$ & $12$ & $26$ & $13$  & $6$ \\
   \hline 
  \end{tabular}
 \end{center}
 \choiceTF 
  {\True Cỡ mẫu của mẫu số liệu là  $n=61$}
  {\True Tứ phân vị thứ hai của mẫu số liệu  ghép nhóm là $Q_2=32{,}79$}
  {Tứ phân vị thứ  nhất của mẫu số liệu ghép nhóm là $Q_1\approx 19{,}69$}
   {\True Tứ phân vị thứ ba của mẫu số liệu ghép nhóm là $Q_3=36{,}44$}
 \loigiai{ 
  \begin{itemchoice}
   \itemch Đúng.Cỡ mẫu của mẫu số  liệu là $n=61$.
    \itemch Đúng Gọi $x_1{,} x_2{,} \ldots{,} x_{61}$ là mẫu số liệu được sắp xếp theo thứ tự không giảm.\\
   Trung vị của mẫu số liệu này là $x_{31} \in [30;35)$.\\ 
   Ta  có $n_m=26$; $C_1=4+12=16$; $u_m=30$; $u_{m+1}=35$.\\
   Tứ phân vị thứ hai  chính là trung vị của mẫu số liệu ghép nhóm là
   \[Q_2=M_e=30+\dfrac{\dfrac{61}{2}-16}{26}  (35-30)=\dfrac{1705}{52} \approx 32{,}79.
    \]
   \itemch Sai. Xét nửa mẫu số liệu bên trái $x_1{,} x_2{,} \ldots{,}  x_{30}$ có trung vị $\dfrac{x_{15}+x_{16}}{2} \in [25;30)$.\\
   Ta  có $n_i=12$; $C_1=4$; $x_i=25$; $x_{i+1}=30$.\\
   Suy ra tứ phân vị thứ nhất của mẫu số liệu là\\  $$Q_1=25+\dfrac{\dfrac{61}{4}-4}{12} (30-25)=\dfrac{475}{16} \approx 29{,}69.$$
   \itemch Đúng. Xét  nửa mẫu số liệu bên trái $x_{32}{,} x_{33}{,} \ldots{,} x_{61}$ có trung vị $\dfrac{x_{46}+x_{47}}{2} \in [35;40)$.\\
   Ta có $n_j=13$; $C_3=4+12+26=42$; $x_i=35$;  $x_{i+1}=40$.\\
   Suy ra tứ phân vị thứ ba của mẫu số liệu là $$Q_3=35+\dfrac{\dfrac{3\cdot61}{4}-42}{13} (40-35)=\dfrac{1895}{52} \approx  36{,}44.$$
  \end{itemchoice}
  Vậy các tứ phân vị của  mẫu số liệu ghép nhóm là
   \[Q_1\approx 29{,}69;Q_2=32{,}79;Q_3=36{,}44.\]
 } 
\end{ex}

% nguồn: Ok_dungsai_lop11_Lan1
\begin{ex}%[1D5H2-3]%[ma:6ea232ca]
 Dựa vào bảng tần số mẫu số liệu ghép nhóm  sau, hãy tìm tứ phân vị của nó.
 \begin{center}
  \begin{tabular}{|c|c|c|c|c|c|c|}
    \hline
   Nhóm & $[30;40)$ & $[40;50)$ & $[50;60)$ &  $[60;70)$ & $[70;80)$ & $[80;90)$ \\
    \hline
   Tần số  & $2$ & $10$ & $16$ & $8$ & $2$ & $2$ \\
   \hline 
  \end{tabular}
 \end{center}
  \choiceTF
  {\True Cỡ mẫu  của mẫu số liệu là $n=40$}
  {\True Tứ phân vị thứ hai của mẫu số liệu  ghép nhóm là $Q_2=45$}
  {Tứ phân vị thứ nhất của mẫu số liệu ghép nhóm là  $Q_1=48$}
  {Tứ phân vị  thứ ba của mẫu số liệu ghép nhóm là $Q_3=61{,}5$}
  \loigiai{
  \begin{itemchoice}
   \itemch Đúng. Cỡ mẫu của mẫu  số liệu là $n=40$.
   \itemch Đúng. Gọi $x_1, x_2, x_3, \ldots, x_{40}$ là mẫu số  liệu được xếp theo thứ tự không giảm.\\
   Trung vị của  mẫu số liệu $\dfrac{x_{20}+x_{21}}{2} \in [50;60)$.\\
   Ta có $n_m=16$; $C_2=2+10=12$;  $u_m=50$; $u_{n+1}=60$.\\
   Tứ phân vị thứ hai của mẫu số liệu ghép nhóm cũng là trung vị của  mẫu số liệu đó là
    $$Q_2=M_e=50+\dfrac{\dfrac{40}{2}-12}{16} (60-50)=55.$$
   \itemch Sai. Xét nửa mẫu số liệu bên trái $x_1, x_2, x_3, \ldots, x_{20}$ có trung vị  $\dfrac{x_{10}+x_{11}}{2} \in [40;50)$.\\
   Ta  có $n_i=10$; $C_1=2$; $u_i=40$; $u_{i+1}=50$.\\
    Tứ phân vị thứ nhất của mẫu số liệu ghép nhóm là
   $$Q_1=40+\dfrac{\dfrac{40}{4}-2}{10} (50-40)=48.$$ 
   \itemch Sai. Xét  nửa mẫu số liệu bên phải $x_{21}, x_{22}, x_{23}, \ldots, x_{40}$ có trung vị $\dfrac{x_{30}+x_{31}}{2} \in [60;70)$.\\
   Ta có $n_j=8$; $C_3=2+10+16=28$; $u_j=60$;  $u_{j+1}=70$.\\
   Tứ phân vị thứ ba  của mẫu số liệu ghép nhóm là
    $$Q_3=60+\dfrac{\dfrac{3\cdot 40}{4}-28}{8} (70-60)=62{,}5.$$
  \end{itemchoice}
  Vậy tứ phân vị của mẫu số liệu ghép nhóm trên là 
  \[Q_1=48, Q_2=55, Q_3=62{,}5.\] 
 } 
\end{ex}

% nguồn: Ok_dungsai_lop11_Lan1
\begin{ex}%[1D5H2-3]%[ma:0d5e0166]
 Số điểm một cầu thủ ghi được trong $20$ trận đấu được cho  ở bởi bảng sau
 \begin{center}
  \begin{tabular}{|c|c|c|c|c|c|c|c|c|c|}
   \hline 
   25  & 23 & 21 & 13 & 8 & 14 & 15 & 18 & 22 & 11 \\
    \hline
    24 & 12 & 14 & 14 & 18 & 6 & 8 & 25 &10 & 11\\
   \hline 
  \end{tabular}
 \end{center}
  \choiceTF
  {\True Tứ phân  vị thứ hai của dãy số liệu là $Q_2= 14$}
  {Tứ phân vị thứ ba của dãy  số liệu là $Q_3= 11{,}5$}
  {Tổng hợp dãy số liệu trên vào bảng  tần số ghép nhóm theo mẫu sau
  \begin{tabular}{|c|c|c|c|c|}
   \hline 
   Điểm số & $[6;10)$ & $[11; 15)$ & $[16; 20)$ &  $[21;25)$\\
    \hline
   Số trận & 4 & 8  & 2 & 6\\
    \hline
 \end{tabular}}
 {Ước lượng tứ phân vị của số liệu ở bảng tần số ghép nhóm trên ta được tứ phân vị thứ hai của mẫu số liệu là  $Q_2=8{,}25$}
  \loigiai{
   Sắp xếp mẫu số liệu theo thứ tự không giảm, ta có
  \begin{center}
   \begin{tabular}{cccccccccccccccccccc}
    6 & 8 & 8 & 10 & 11 & 11 & 12 & 13 & 14 & 14 & 14 & 15 & 18 &  18 & 21 & 22 & 23 & 24 & 25 & 25 \\
   \end{tabular}
  \end{center}
  \begin{itemchoice}
   \itemch  \textbf{Đúng}. \\
   Ta có cỡ mẫu của dãy số liệu là $n = 20$ nên trung vị của dãy  số liệu là $$Q_2 =\dfrac{x_{10}+x_{11}}{2} = \dfrac{14+14}{2}= 14.$$
   \itemch \textbf{Sai}.  \\
   Tứ phân vị thứ ba  là trung vị của dãy số liệu
   \begin{center}
    \begin{tabular}{cccccccccc}
     14 & 15 & 18 & 18 &  21 & 22 & 23 & 24 & 25 & 25 \\
    \end{tabular}
   \end{center}
   Tứ phân vị thứ ba của dãy số liệu này là  $Q_3 =\dfrac{21+22}{2}=21{,}5$.\\
    Vậy $Q_3= 21{,}5.$
   \itemch  \textbf{Đúng}. \\
   Tổng hợp lại dãy số liệu trên vào bảng tần số ghép nhóm  theo mẫu sau
   \begin{center}
    \begin{tabular}{|c|c|c|c|c|}
      \hline
     \text{Điểm số}  & $[6; 10)$ & $[11; 15)$ & $[16; 20)$ & $[21; 25)$ \\
      \hline
     \text{Số trận} & 4 & 8  & 2 & 6 \\
      \hline
    \end{tabular}
   \end{center}
    \itemch \textbf{Đúng}. \\
   Vì số trận là số nguyên nên ta hiệu chỉnh lại bảng số  liệu sau
   \begin{center}
    \begin{tabular}{|c|c|c|c|c|}
      \hline
     \text{Điểm số}  & $[5,5; 10,5)$ & $[10,5; 15,5)$ & $[15,5; 20,5)$ & $[20,5; 25,5)$ \\
      \hline
     \text{Số trận} & 4 &  8 & 2 & 6 \\
     \hline 
    \end{tabular}
   \end{center}
   Gọi $x_1, x_2, \ldots, x_{20}$ lần lượt là số điểm ghi được ở mỗi trận đấu xếp theo  thứ tự không giảm.
   Do  $x_1, \ldots, x_4 \in [5,5; 10,5); \; x_5, \ldots, x_{12} \in [10,5; 15,5); \; x_{13}, x_{14} \in [15,5; 20,5); \; x_{15}, \ldots, x_{20} \in [20,5; 25,5)$.
   Nên trung vị  của mẫu số liệu $x_1, \ldots, x_{20}$ là
   $$\dfrac{1}{2} \left( x_{10} + x_{11} \right) \in  [10,5; 15,5).$$
   Ta xác định  được $n = 20, \; n_m = 8, \; C = 4, \; u_m = 10,5; \; u_{m+1} = 15,5.$
   Suy ra tứ phân vị thứ hai của mẫu số liệu là 
   $$Q_2 = 10,5 + \dfrac{\dfrac{20}{2} - 4}{8}(15,5  - 10,5) = 14,25.$$
  \end{itemchoice}
  }
\end{ex}

% nguồn: Ok_dungsai_lop11_Lan1
\begin{ex}%[1D5H2-3]%[ma:938aa583]
 Người ta đo đường kính của $61 $cây gỗ được trồng sau $12$ năm (đơn vị: centimét), họ thu  được bảng tần số ghép nhóm sau:
 \begin{center}
  \begin{tabular}{|c|c|c|c|c|c|}
   \hline Đường kính & {$[20; 25)$} &  {$[25; 30)$} & {$[30; 35)$} & {$[35; 40)$} & {$[40; 45)$} \\
   \hline Số cây & $4$  & $12$ & $26$ & $13$ & $6$ \\
    \hline
  \end{tabular}
 \end{center}
 Các khẳng định sau \textbf{đúng}  hay \textbf{sai}?
 \choiceTF 
  {\True Cỡ  mẫu của mẫu số liệu là $n=61$}
  {Độ  dài của nhóm $[30; 35)$ bằng $4$}
  {Tứ phân vị thứ nhất của mẫu số liệu ghép nhóm  là $Q_1 \approx 19{,}69$}
  {\True Tứ phân vị thứ hai của mẫu số liệu ghép nhóm là $Q_2\approx32{,}79$} 
 \loigiai{ 
  Cỡ mẫu  của mẫu số liệu là $n=61$.\\
  Gọi $x_1, x_2, \ldots, x_{61}$ là mẫu số  liệu được sắp xếp theo thứ tự không giảm.\\
  Trung vị của mẫu số liệu này là $x_{31} \in[30; 35)$.\\ 
  Ta có $n_m=26; C_1=4+12=16; u_m=30;  u_{m+1}=35$.\\
  Tứ phân vị thứ hai chính là trung vị của mẫu số  liệu ghép nhóm là
  $$Q_2=M_e=30+\dfrac{\dfrac{61}{2}-16}{26}(35-30)=\dfrac{1705}{52}  \approx 32{,}79\text{.}$$
  Xét nửa mẫu số liệu bên trái $x_1, x_2, \ldots, x_{30}$  có trung vị $\dfrac{x_{15}+x_{16}}{2} \in[25; 30)$.\\
  Ta có $n_i=12; C_1=4; x_i=25; x_{i+1}=30$.\\ 
  Suy ra tứ phân vị thứ  nhất của mẫu số liệu là \\$$Q_1=25+\dfrac{\dfrac{61}{4}-4}{12}(30-25)=\dfrac{475}{16} \approx 29{,}69.$$
  \begin{itemchoice}
   \itemch Đúng. 
   \itemch Sai. Vì độ dài của nhóm $[30;  35)$ bằng $5$.
   \itemch Sai. Vì Tứ phân vị thứ nhất của mẫu số liệu ghép nhóm là $Q_1 \approx 29{,}69$  cm.
   \itemch  Đúng.
  \end{itemchoice}
 } 
\end{ex}

% nguồn: Ok_dungsai_lop11_Lan1
\begin{ex}%[1D5H2-3]%[ma:95120cb3]
 Kiểm tra điện lượng của một số viên pintiểu do một hãng sản xuất  thu được kết quả sau:
 \begin{center}
  \begin{tabular}{|l|c|c|c|c|c|}
   \hline Điện lượng (nghìn mAh) & $[0{,}9; 0{,}95)$ & $[0{,}95; 1{,}0)$ &  $[1{,}0; 1{,}05)$ & $[1{,}05; 1{,}1)$ & $[1{,}1; 1{,}15)$ \\
   \hline Số  viên pin & $10$ & $20$ & $35$ & $15$ & $5$ \\
    \hline
  \end{tabular}
 \end{center}
  \choiceTF
  {\True  $\overline{x}=1{,}016$}
   {$M_o>1{,}5$}
  {$Q_1=0{,}89$} 
  {\True $\Delta Q=0{,}068$} 
  \loigiai{
  Ta có bảng số liệu  sau
  \begin{center}
   \begin{tabular}{|l|c|c|c|c|c|}
    \hline Điện lượng (nghìn mAh)  & {$[0{,}9; 0{,}95)$} & {$[0{,}95; 1{,}0)$} & {$[1{,}0; 1{,}05)$} & {$[1{,}05; 1{,}1)$} & {$[1{,}1; 1{,}15)$} \\
    \hline Giá trị đại diện & $0{,}925$ & $0{,}975$ & $1{,}025$ & $1{,}075$ & $1{,}125$  \\
    \hline Số viên pin & $10$  & $20$ & $35$ & $15$ & $5$ \\
    \hline 
   \end{tabular}
  \end{center}
  Số trung bình của dãy số liệu xấp xỉ  bằng
   $$\overline{x}\approx (0{,}925\cdot 10+0{,}975\cdot 20+1{,}025\cdot 35+1{,}075\cdot 15+1{,}125\cdot 5): 85=1{,}016.$$
  Nhóm  chứa mốt của dãy số liệu là nhóm $[1{,}0; 1{,}05)$. \\
  Mốt của mẫu số liệu trên  là
   $M_o=1{,}0+\dfrac{35-20}{(35-20)+(35-15)} \cdot(1{,}05-1{,}0)=1{,}02$. \\
  Gọi $x_1; x_2; x_3; \ldots; x_{85}$ lần lượt là số viên pintheo thứ tự không giảm.  \\
  Ta có  \\
  $\begin{aligned}
  &x_1, \ldots, x_{10} \in [0{,}9;  0{,}95); x_{11},\ldots, x_{30} \in [0{,}95; 1{,}0); x_{31},\ldots, x_{65} \in [1{,}0; 1{,}05); \\
  &x_{66},\ldots, x_{80} \in  [1{,}05; 1{,}1); x_{81},\ldots, x_{85} \in [1{,}1; 1{,}15).
 \end{aligned}$ \\
 Tứ phân vị  thứ hai của dãy số liệu là $\dfrac{1}{2}\left(x_{42}+x_{43}\right)\in [1{,}0; 1{,}05)$ nên tứ phân vị thứ hai của mẫu số liệu là
  $$Q_2=1{,}0+\dfrac{\dfrac{85}{2}-30}{35}(1{,}05-1{,}0)\approx 1,02.$$
 Tứ phân vị thứ nhất  của dãy số liệu là $\dfrac{1}{2}\left(x_{21}+x_{22}\right)\in [0{,}95; 1{,}0)$ nên tứ phân vị thứ nhất của mẫu số liệu là
 $$Q_1=0{,}95+\dfrac{\dfrac{85}{4}-10}{20}(1{,}0-0{,}95)=0{,}98.$$ 
 Tứ phân vị thứ ba của dãy số liệu  là $\dfrac{1}{2}\left(x_{63}+x_{64}\right)\in [1{,}0; 1{,}05)$ nên tứ phân vị thứ ba của mẫu số liệu là $$Q_3=1{,}0+\dfrac{\dfrac{3\cdot 85}{4}-30}{35}(1{,}05-1{,}0)=1{,}048.$$
 \begin{itemchoice}
   \itemch $\overline{x}\approx 1{,}016$.
  \itemch $M_o\approx 1{,}02$. 
  \itemch  $Q_1\approx 0{,}98$.
  \itemch  $\Delta Q\approx 0{,}068$.
 \end{itemchoice}
 }
\end{ex}

% nguồn: Ok_dungsai_lop11_Lan1
\begin{ex}%[1D5H2-3]%[ma:91b41763]
 Kiểm tra điện  lượng của một số viên pintiểu do một hãng sản xuất thu được kết quả sau
 \begin{tabular}{|c|c|c|c|c|c|}
  \hline 
  Điện lượng (Nghìn mAh) &  $[0{,}9;0{,}95)$ & $[0{,}95;1{,}0)$ & $[1{,}0;1{,}05)$ & $[1{,}05;1{,}1)$ & $[1{,}1;1{,}15)$ \\
   \hline
  Số pin & $10$  & $20$ & $35$ & $15$ & $5$ \\
   \hline
 \end{tabular}
  \choiceTF
  {\True Số  trung bình của dãy số liệu là $1{,}016$}
  {Nhóm chứa mốt của dãy số liệu là  $[1{,}05;1{,}1)$}
  {\True Tứ phân vị thứ  nhất của mẫu số liệu nhóm là $Q_1=0{,}98$}
  {Tứ phân vị thứ ba của mẫu số liệu nhóm là  $Q_3=1{,}248$}
  \loigiai{
  \begin{center}
   \begin{tabular}{|c|c|c|c|c|c|}
    \hline 
    Điện lượng (Nghìn mAh)  & $[0{,}9;0{,}95)$ & $[0{,}95;1{,}0)$ & $[1{,}0;1{,}05)$ & $[1{,}05;1{,}1)$ & $[1{,}1;1{,}15)$ \\
    \hline 
    Giá trị đại diện & $0{,}925$ & $0{,}975$ & $1{,}025$ &  $1{,}075$ & $1{,}125$ \\
    \hline 
    Số pin & $10$ &  $20$ & $35$ & $15$ & $5$ \\
    \hline 
   \end{tabular}
  \end{center}
  \begin{itemchoice}
    \itemch Đúng. Số trung bình của dãy số liệu là
   \[\dfrac{0{,}925\cdot 10+0{,}975\cdot 20+1{,}025\cdot 35+1{,}075\cdot 15+1{,}125\cdot  5}{85}=1{,}016.\]
   \itemch Sai.  Nhóm chứa mốt của dãy số liệu là $[1;1{,}05)$
    \[M_0=1+\dfrac{35-20}{(35-20)+(35-15)} (1{,}05-1)=1{,}02.\]
   Gọi $x_1; x_2; \ldots; x_{85}$ lần lượt  là điện lượng mỗi viên pin xếp theo thứ tự không giảm.\\
   Do $x_1; x_2; \ldots;  x_{10} \in [0{,}9;0{,}95);x_{11};\ldots;x_{30} \in [0{,}95;1{,}0);x_{31};\ldots;x_{65} \in [1{,}0;1{,}05)$; $x_{66};\ldots;x_{80} \in [1{,}05;1{,}1);x_{81};\ldots;x_{85} \in [1{,}1;1{,}15)$ nên trung vị của mẫu số liệu $x_1$; $x_2;\ldots;x_{85}$ là $x_{43} \in [1;1{,}05)$.
   \itemch Đúng. Tứ phân vị thứ nhất của mẫu số liệu là  $\dfrac{1}{2} \left(x_{21}+x_{22} \right)$.
   Do $x_{21}, x_{22} \in  [0{,}95;1)$ nên tứ phân vị thứ nhất của mẫu số liệu nhóm là\\ $Q_1=0{,}95+\dfrac{\dfrac{85}{4}-10}{20} (1-0{,}95)=0{,}98$.
   \itemch Sai. Tứ phân vị thứ ba của mẫu số  liệu là $\dfrac{1}{2} \left(x_{64}+x_{65} \right)$.\\
   Do $x_{64}, x_{65} \in [1;1,05)$ nên tứ phân vị thứ ba của mẫu số liệu nhóm là\\ $Q_3=1+\dfrac{\dfrac{3\cdot 85}{4}-30}{35}  (1{,}05-1)=1{,}048$.
  \end{itemchoice}
  }
\end{ex}

% nguồn: Ok_dungsai_lop11_Lan1
\begin{ex}%[1D5H2-3]%[ma:393b822a]
 Khi đo mắt cho một số học sinh khối  $11$ ở một trường THPT, nhân viên y tế ghi nhận lại ở bảng sau:
 \begin{longtable}{|c|c|c|c|c|c|}
  \hline 
  Độ  cận&$[0{,}25;0{,}75]$&$[0{,}75;1{,}25]$&$[1{,}25;1{,}75]$&$[1{,}75;2{,}25]$&$[2{,}25;2{,}75]$\\
   \hline
  Số học  sinh&$25$&$32$&$14$&$12$&$7$\\
   \hline
 \end{longtable}
  Khi đó
 \choiceTF 
  {\True Độ cận  trung bình của số học sinh trên xấp xỉ $1{,}19$ độ}
  {\True $0{,}89$ độ là độ cận mà số học sinh được đo ở trên chiếm nhiều nhất} 
  {$50\%$ số học sinh được đo có độ cận không vượt quá  $0{,}75$ độ}
  {\True $50\%$ số học sinh được đo có độ cận từ $0{,}7$  độ đến $1{,}625$ độ}
  \loigiai{
   Ta có mẫu số liệu ghép nhóm với giá trị đại diện như sau:
  \begin{longtable}{|c|c|c|c|c|c|}
    \hline
   Độ cận&$[0{,}25;0{,}75]$&$[0{,}75;1{,}25]$&$[1{,}25;1{,}75]$&$[1{,}75;2{,}25]$&$[2{,}25;2{,}75]$\\ 
   \hline 
   Giá trị  đại diện&$0{,}5$&$1{,}0$&$1{,}5$&$2{,}0$&$2{,}5$\\
    \hline
   Số  học sinh&$25$&$32$&$14$&$12$&$7$\\
   \hline 
  \end{longtable}
  Gọi $x_1,x_2,\ldots,x_{90}$ lần lượt là độ cận của $90$ học sinh được sắp xếp  theo thứ tự không giảm.
  \begin{itemchoice}
   \itemch{\bf  Đúng.}\\
   Độ cận trung bình $\overline{x}=\dfrac{0{,}5\cdot  25+1\cdot 32+1{,}5\cdot 14+2\cdot 12+2{,}5\cdot 7}{90}\approx 1{,}19$.
    \itemch{\bf Đúng}\\
   Ta có nhóm  chứa mốt là $[0{,}75;1{,}25)\Rightarrow \heva{&u_m=0{,75}\\&u_{m+1}=1{,}25\\&n_m=32\\&n_{m-1}=25\\&n_{m+1}=14}$.\\
   Do đó $M_0=0{,}75+\dfrac{32-25}{(32-25)+(32-14)}\cdot  (1{,}25-0{,}75)=0{,}89$.
    \itemch{\bf Sai}\\
   Cỡ mẫu $n=90$ nên trung vị là $\dfrac{1}{2}(x_{45}+x_{46})\in [0{,}75;1{,}25)$.\\ 
    Do đó $M_e=0{,}75+\dfrac{\dfrac{90}{2}-25}{32}\cdot 0{,}5=1{,}0625$.\\
   Vậy $50\%$ số học sinh được đo có độ cận không vượt quá $1{,}0625$  độ.
   \itemch{\bf  Đúng}\\
   Cỡ mẫu $n=90$  nên tứ phân vị thứ nhất là $x_{23}\in [0{,}25;0{,}75)$.\\
    Do đó $Q_1=0{,}25+\dfrac{\dfrac{90}{4}-0}{25}\cdot 0{,}5=0{,}7$.\\
   Tứ phân vị  thứ ba là $x_{68}\in [1{,}25;1{,}75)$.\\
   Do đó $Q_3=1{,}25+\dfrac{\dfrac{3\cdot90}{4}-(25+32)}{14}\cdot 0{,}5=1{,}625$.\\ 
   Vậy $50\%$  số học sinh được đo có độ cận từ $0{,}7$ độ đến $1{,}625$ độ.
  \end{itemchoice}
  }
\end{ex}

% nguồn: Ok_dungsai_lop11_Lan1
\begin{ex}%[1D5H2-3]%[ma:a54e97d6]
 Cân nặng của một số lợn con mới sinh thuộc hai giống $A$ và $B$ được cho  ở bảng đây (đơn vị: kg)
 \begin{center}
  \begin{tabular}{|c|c|c|c|c|}
   \hline 
   Cân nặng $(kg)$ & $[1{,}0;1{,}1)$ & $[1{,}1;1{,}2)$ & $[1{,}2;1{,}3)$  & $[1{,}3;1{,}4)$ \\
    \hline
   Số con giống A & $8$ & $28$  & $32$ & $17$ \\
    \hline
   Số con giống B & $13$  & $14$ & $24$ & $14$ \\
    \hline
  \end{tabular}
 \end{center}
 \choiceTF 
   {\True Cân nặng trung bình của giống $A$ là $1{,}22$}
  {\True Cân nặng trung bình của giống $B$ là  $1{,}21$}
  {\True Tứ phân vị thứ nhất của mẫu số liệu lợn  con giống $A$ là $Q_{1A}=1{,}15$}
  {Tứ phân vị thứ nhất của mẫu số liệu  lợn con giống $B$ là $Q_{1B}=1{,}62$}
  \loigiai{
  \begin{itemchoice}
   \itemch Đúng. Cân  nặng trung bình của giống $A$ là
    \[\dfrac{1{,}05\cdot8+1{,}15\cdot28+1{,}25\cdot32+1{,}35\cdot17}{85}=1{,}22.\]
   \itemch Đúng.  Cân nặng trung bình của giống $B$ là
    \[\dfrac{1{,}05\cdot13+1{,}15\cdot14+1{,}25\cdot24+1{,}35\cdot14}{65}=1{,}21.\]
   Vậy  theo số trung bình thì cân nặng của lợn con giống $A$ lớn hơn giống $B$.\\
   Gọi $x_1;x_2;\ldots;x_{85}$ lần  lượt là cân nặng của mỗi con lợn giống $A$ theo thứ tự không giảm.\\
   Do $x_1;\ldots;x_8\in [1{,}0;1{,}1);x_9;x_{10};\ldots;x_{36} \in [1{,}1;1{,}2);x_{37};\ldots;x_{68} \in  [1{,}2;1{,}3);x_{69};\ldots;x_{85} \in [1{,}3;1{,}4)$ nên trung vị của mẫu số liệu lợn con giống $A$ là $x_{43} \in [1{,}2;1{,}3)$.
   \[Q_{2A}=1{,}2+\dfrac{\dfrac{85}{2}-36}{32}  (1{,}3-1{,}2)=1{,}22.\]
   Gọi $y_1;y_2;\ldots;y_{65}$ lần lượt là cân nặng của mỗi con  lợn giống $B$ theo thứ tự không giảm.\\
   Do $y_1;\ldots;y_{13} \in [1{,}0;1{,}1);y_{14};\ldots;y_{27} \in [1{,}1;1{,}2);y_{28};\ldots;y_{51} \in [1{,}2;1{,}3);y_{52};\ldots;y_{65} \in [1{,}3;1{,}4)$ nên trung vị của mẫu số liệu lợn con giống  $B$ là $y_{33} \in [1{,}2;1{,}3)$.
    \[Q_{2A}=1{,}2+\dfrac{\dfrac{65}{2}-27}{24} (1{,}3-1{,}2)=1{,}223.\]
    Vậy theo số trung vị thì cân nặng trung bình của lợn giống $A$ nhỏ hơn lợn giống $B$.
   \itemch Đúng. Tứ phân vị thứ nhất của  dãy số liệu lợn con giống $A$ là $\dfrac{1}{2} \left(x_{21}+x_{22} \right)\in [1{,}1;1{,}2)$ nên tứ phân vị thứ nhất của mẫu số liệu lợn con giống $A$ là
    \[Q_{1A}=1{,}1+\dfrac{\dfrac{85}{4}-8}{28} (1{,}2-1{,}1)=1{,}15.\]
   \itemch Sai. Tứ phân vị thứ nhất của dãy số liệu lợn con giống $B$ là $\dfrac{1}{2} \left(y_{16}+y_{17} \right)\in [1{,}1;1{,}2)$ nên tứ  phân vị thứ nhất của mẫu số liệu lợn con giống $B$ là:
   \[Q_{1B}=1{,}1+\dfrac{\dfrac{65}{4}-13}{14}  (1{,}2-1{,}1)=1{,}12.\]
  \end{itemchoice}
  }
\end{ex}

% nguồn: Ok_dungsai_lop11_Lan1
\begin{ex}%[1D5H2-3]%[ma:c825441a]
 Trong một đề tài nghiên cứu về bệnh A, người ta ghi lại lại tuổi của bệnh  nhân mắc bệnh này, số liệu thông kê được trình bày trong bảng sau
 \begin{center}
  \begin{tabular}{|c|c|c|c|c|c|}
   \hline 
   Độ tuổi & $\left[15;25\right)$  & $\left[25;35\right)$ &$\left[35;45\right)$ &$\left[45;55\right)$ &$\left[55;65\right)$ \\
    \hline
   Số bệnh nhân & $10$ &$12$ &$14$ &$9$ &$5$  \\
    \hline
  \end{tabular}
 \end{center}
 Khi đó 
  \choiceTF
  {\True Cỡ  mẫu là $n=50$}
   {Trung vị của mẫu số liệu thuộc nhóm $\left[25;35\right)$}
  {\True Trung vị của mẫu số liệu gần bằng $37{,}14$} 
  {\True Khoảng tứ phân vị gần  bằng $19$}
 \loigiai{ 
  \begin{itemchoice}
   \itemch{\bf  Đúng.}\\
   Cỡ mẫu  là $n=50$.
    \itemch{\bf Sai}\\
   Gọi $x_1,x_2,\ldots,x_{50}$ là độ tuổi của $50$ bệnh nhân và giả sử  dãy này được sắp xếp theo thứ tự tăng dần.\\
   Khi đó trung vị của mẫu số liệu  là $\dfrac{x_{25}+x_{26}}{2}$. Do $x_{25},x_{26}$ thuộc nhóm $\left[35;45\right)$ nên nhóm này chứa trung vị.
    \itemch{\bf Đúng.}\\
   Gọi $x_1,x_2,\ldots,x_{50}$ là độ tuổi của $50$ bệnh nhân và  giả sử dãy này được sắp xếp theo thứ tự tăng dần.\\
   Khi đó trung vị của mẫu số liệu là $\dfrac{x_{25}+x_{26}}{2}$. Do $x_{25},x_{26}$ thuộc nhóm $\left[35;45\right)$ nên nhóm này chứa trung vị.\\ 
   Ta xác định được $n=50$; $p=3$; $m_3=14$; $m_1+m_2=22$, $a_3=35$, $a_4=45$.\\ 
   Vậy trung vị của mẫu số liệu ghép  nhóm là $M_e=35+\dfrac{\dfrac{50}{2}-22}{14}\cdot (45-35)\approx 37{,}14$.
    \itemch{\bf Đúng}\\
   Cỡ mẫu $n=50$ nên  tứ phân vị thứ nhất là $Q_1=\dfrac{x_{11}+x_{13}}{2}\in [25;35)$.\\
   Do đó  $Q_1=25+\dfrac{\dfrac{50}{4}-10}{12}\cdot 10=27{,}08$.\\
   Tứ phân vị thứ ba  là $Q_1=\dfrac{x_{36}+x_{38}}{2}\in [35;45)$.\\
   Do đó $Q_3=35+\dfrac{\dfrac{3\cdot 50}{4}-22}{14}\cdot  10=46{,}07$.\\
    Vậy $Q_3-Q_1\approx 46{,}07-27{,}08\approx 19$.
  \end{itemchoice}
 } 
\end{ex}

% nguồn: Ok_dungsai_lop11_Lan1
\begin{ex}%[1D5H2-3]%[ma:aa8453af]
 Cho bảng  số liệu ghép nhóm như sau
 \begin{center}
  \begin{tabular}{|c|c|c|c|c|c|c|}
   \hline 
   Nhóm & $[30;40)$ & $[40;50)$ & $[50;60)$ & $[60;70)$ &  $[70;80)$ & $[80;90)$ \\
   \hline 
    Tần số & $2$ & $10$ & $16$ & $8$ & $2$ & $2$ \\
   \hline 
  \end{tabular}
 \end{center}
 Khi  đó
 \choiceTF 
  {\True Cỡ mẫu  của mẫu số liệu là $n=40$}
  {Tứ phân  vị thứ hai của mẫu số liệu ghép nhóm là $Q_2=45$}
  {\True Tứ phân vị thứ nhất của mẫu số liệu ghép nhóm  là $Q_1=48$}
  {Tứ phân vị thứ ba của mẫu  số liệu ghép nhóm là $Q_3=61{,}5$}
 \loigiai{ 
  \begin{itemchoice}
   \itemch \textbf{Đúng.} \\ 
    Cỡ mẫu của mẫu số liệu là $n=40$.
   \itemch \textbf{Sai.} \\ 
    Gọi $x_1, x_2, x_3, \ldots, x_{40}$ là mẫu số liệu được xếp theo thứ tự không giảm.\\
   Trung vị của mẫu số  liệu: $\dfrac{x_{20}+x_{21}}{2} \in [50;60)$.\\
   Ta  có $n_m=16;C_2=2+10=12;u_m=50;u_{n+1}=60$.\\
   Tứ phân  vị thứ hai của mẫu số liệu ghép nhóm cũng là trung vị của mẫu số liệu đó là
   \[Q_2=M_e=50+\dfrac{\dfrac{40}{2}-12}{16} (60-50)=55. 
   \] 
   \itemch  \textbf{Đúng.} \\
   Xét nửa mẫu số liệu bên trái $x_1, x_2, x_3, \ldots, x_{20}$  có trung vị $\dfrac{x_{10}+x_{11}}{2} \in [40;50)$.\\
   Ta có  $n_i=10;C_1=2;u_i=40;u_{i+1}=50$.\\
    Tứ phân vị thứ nhất của mẫu số liệu ghép nhóm là $Q_1=40+\dfrac{\dfrac{40}{4}-2}{10} (50-40)=48$.
   \itemch \textbf{Sai.}  \\
   Xét nửa mẫu số liệu bên phải $x_{21}, x_{22}, x_{23}, \ldots, x_{40}$  có trung vị $\dfrac{x_{30}+x_{31}}{2} \in [60;70)$.\\
   Ta có: $n_j=8;C_3=2+10+16=28;u_j=60;u_{j+1}=70$.\\ 
   Tứ phân vị thứ ba của mẫu số liệu  ghép nhóm là
   \[Q_3=60+\dfrac{\dfrac{3.40}{4}-28}{8} (70-60)=62{,}5. 
    \]
   Vậy tứ  phân vị của mẫu số liệu ghép nhóm trên là $Q_1=48$, $Q_2=55$, $Q_3=62{,}5$.
  \end{itemchoice}
  }
\end{ex}

% nguồn: Ok_dungsai_lop11_Lan1
\begin{ex}%[1D5H2-3]%[ma:dce2f644]
 Cho mẫu số liệu ghép  nhóm về lương của nhân viên trong một công ty như sau:
 \begin{center}
  \begin{tabular}{|l|c|c|c|c|c|}
   \hline Lương (triệu đồng) & {$[9;  12)$} & {$[12; 15)$} & {$[15; 18)$} & {$[18; 21)$} & {$[21; 24)$} \\
   \hline Số nhân viên & $6$ & $12$ & $4$ & $2$  & $1$ \\
   \hline 
  \end{tabular}
 \end{center}
 Xét tính đúng sai của các mệnh  đề sau:
 \choiceTF 
  {\True Giá trị  đại diện của nhóm $[9; 12)$ là $10{,}5$}
  {Trung bình lương  các nhân viên là $16{,}5$ triệu đồng}
  {Nhóm chứa trung vị là $[15;  18)$}
  {\True Tứ phân vị thứ ba là  $15{,}56$}
  \loigiai{
  \begin{itemchoice}
   \itemch Đúng. Vì giá trị đại diện của nhóm $[9;  12)$ là $\dfrac{9+12}{2}=10{,}5$.
   \itemch Sai. Vì trung bình  lương các nhân viên là $$\overline{x}=\dfrac{1}{25}(6 \cdot 10{,}5+12 \cdot 13{,}5+4\cdot 16{,}5+2\cdot 19{,}5+22{,}5)=14{,}1\text{ triệu đồng.}$$
   \itemch Sai. Vì công ty có $25$ nhân sự. Vì $x_{13}\in[12; 15)$ nên nhóm  này chứa trung vị.
   \itemch  Đúng. Vì $x_{19}\in[15; 18)$ nên ta có $Q_3=15+\dfrac{\dfrac{3 \cdot 25}{4}-18}{4} \cdot 3\approx 15{,}56$.
  \end{itemchoice}
 } 
\end{ex}

% nguồn: Ok_dungsai_lop11_Lan1
\begin{ex}%[1D5H2-3]%[ma:8fc73894]
 Thống kê điểm trung bình môn Toán của một số học sinh  lớp $11$ được cho ở bảng sau
 \begin{tabular}{|c|c|c|c|c|c|c|c|}
  \hline 
  Khoảng điểm & $[6{,}5;7)$ & $[7;7{,}5)$ & $[7{,}5;8)$ & $[8;8{,}5)$ & $[8{,}5;9)$ & $[9;9{,}5)$  & $[9{,}5;10)$ \\
  \hline 
  Số học sinh & $8$ & $10$ & $16$ & $24$ & $13$ & $7$ & $4$  \\
   \hline
 \end{tabular}
  \choiceTF
  {Cỡ mẫu của mẫu số liệu là $n=80$} 
  {\True Tứ phân vị thứ nhất của  mẫu số liệu ghép nhóm là $Q_1=7{,}58$}
  {\True Tứ phân vị thứ hai của  mẫu số liệu ghép nhóm là $Q_2=8{,}15$}
  {\True Tứ phân vị thứ ba của mẫu số liệu  ghép nhóm là $Q_3=8{,}63$}
 \loigiai{ 
  \begin{itemchoice}
   \itemch Sai. Cỡ mẫu của  mẫu số liệu là \\
   $n=8+10+16+24+13+7+4=82$. 
   \itemch  Đúng. Gọi $x_1;x_2;\ldots;x_{82}$ lần lượt là điểm trung bình môn Toán của các học sinh sắp xếp theo thứ tự không giảm.\\
   Ta có $x_1;\ldots;x_8\in [6{,}5;7);x_9;\ldots;x_{18} \in [7;7{,}5);x_{19};\ldots;x_{34} \in [8;8{,}5)$; $x_{59};\ldots;x_{71}  \in [8{,}5;9);x_{72};\ldots;x_{78} \in [7;7{,}5);x_{79};\ldots;x_{82} \in [9{,}5;10)$.\\
   Tứ phân  vị thứ nhất của mẫu số liệu là $x_{21} \in [7{,}5;8)$.\\
   Do đó tứ phân vị thứ nhất của  mẫu số liệu là\\
    $Q_1=7{,}5+\dfrac{\dfrac{82}{4}-18}{16} (8-7{,}5)=7{,}58$.
   \itemch Đúng. Tứ phân vị thứ hai của mẫu số liệu là $\dfrac{1}{2} \left(x_{41}+x_{42} \right)$ mà $x_{41};x_{42} \in [8;8{,}5)$ nên tứ phân  vị thứ hai của mẫu số liệu là\\
    $Q_2=8+\dfrac{\dfrac{82}{2}-34}{24} (8{,}5-8)=8{,}15$.
    \itemch Đúng. Tứ phân vị thứ ba của mẫu số liệu là $x_{62} \in [8{,}5;9)$.
   Do đó tứ phân vị thứ ba của mẫu số  liệu là\\
   $Q_3=8{,}5+\dfrac{\dfrac{3.82}{4}-58}{13} (9-8{,}5)=8{,}63$. 
  \end{itemchoice}
  }
\end{ex}

% nguồn: Ok_dungsai_lop11_Lan1
\begin{ex}%[1D5H2-3]%[ma:fe801f48]
 Người ta đo đường kính của $61$ cây gỗ được trồng sau $12$ năm (đơn vị: centimét),  họ thu được bảng tần số ghép nhóm sau
 \begin{center}
  \begin{tabular}{|c|c|c|c|c|c|}
   \hline 
   Đường kính & $[20;25)$ & $[25;30)$ &  $[30;35)$ & $[35;40)$ & $[40;45)$ \\
    \hline
   Số cây & $4$  & $12$ & $26$ & $13$ & $6$ \\
    \hline
  \end{tabular}
 \end{center}
  Khi đó
  \choiceTF
   {\True Cỡ mẫu của mẫu số liệu là $n=61$}
  {\True Tứ phân vị  thứ hai của mẫu số liệu ghép nhóm là $Q_2\approx32{,}79$}
  {Tứ  phân vị thứ nhất của mẫu số liệu ghép nhóm là $Q_1 \approx 19{,}69$}
  {\True Tứ phân vị thứ ba của mẫu số liệu  ghép nhóm là $Q_3\approx36{,}44$}
  \loigiai{
  \begin{itemchoice}
   \itemch  \textbf{Đúng.}\\
   Cỡ mẫu  của mẫu số liệu là $n=61$.
    \itemch \textbf{Đúng.} \\
   Gọi $x_1$, $x_2$, $\ldots$, $x_{61}$ là mẫu số liệu được sắp xếp theo thứ  tự không giảm.\\
   Trung vị của mẫu số liệu  này là $x_{31} \in [30;35)$.\\
   Ta  có $n_m=26$; $C_1=4+12=16$; $u_m=30$; $u_{m+1}=35$.\\
   Tứ phân vị thứ hai chính là trung vị của mẫu số liệu ghép nhóm  là
    \[Q_2=M_e=30+\dfrac{\dfrac{61}{2}-16}{26} (35-30)=\dfrac{1705}{52} \approx 32,79\,(\rm{cm}).\]
   \itemch \textbf{Sai.} \\ 
   Xét nửa mẫu số liệu bên trái $x_1$, $x_2$, $\ldots$, $x_{30}$  có trung vị $\dfrac{x_{15}+x_{16}}{2} \in [25;30)$.\\
   Ta có $n_i=12$; $C_1=4$;  $x_i=25$; $x_{i+1}=30$.\\
   Suy ra tứ phân  vị thứ nhất của mẫu số liệu là
   \[Q_1=25+\dfrac{\dfrac{61}{4}-4}{12} (30-25)=\dfrac{475}{16} \approx 29,69\,(\rm{cm}).\] 
    \itemch \textbf{Đúng.} \\
   Xét nửa mẫu  số liệu bên trái $x_{32}$, $x_{33}$, $\ldots$, $x_{61}$ có trung vị $\dfrac{x_{46}+x_{47}}{2} \in [35;40)$.\\
   Ta có: $n_j=13$; $C_3=4+12+26=42$; $x_i=35$; $x_{i+1}=40$.\\ 
   Suy ra tứ phân vị thứ ba của mẫu số  liệu là
   \[Q_3=35+\dfrac{\dfrac{3\cdot 61}{4}-42}{13} (40-35)=\dfrac{1\,895}{52} \approx 36{,}44\,(\rm{cm}).\] 
   Vậy các tứ phân vị của mẫu  số liệu ghép nhóm là\\
   \[Q_1 \approx 29{,}69; \, Q_2\approx32{,}79;\,  Q_3\approx36{,}44.\]
  \end{itemchoice}
 } 
\end{ex}

% nguồn: Ok_dungsai_lop11_Lan1
\begin{ex}%[1D5H2-3]%[ma:2bbc75e0]
 Kiểm tra điện lượng của một số viên pintiểu do một  hãng sản xuất thu được kết quả sau:
 \begin{center}
  \begin{tabular}{|c|c|c|c|c|c|}
    \hline
   Điện lượng (Nghìn mAh) & $[0{,}9;0{,}95)$ & $[0{,}95;1{,}0)$ & $[1{,}0;1{,}05)$ & $[1{,}05;1{,}1)$ &  $[1{,}1;1{,}15)$ \\
   \hline 
   Số pin & $10$ & $20$ &  $35$ & $15$ & $5$ \\
    \hline
  \end{tabular}
 \end{center}
 Khi đó 
  \choiceTF
  {\True Số trung bình của dãy  số liệu xấp xỉ $1{,}016$}
  {Nhóm  chứa mốt của dãy số liệu là $[1{,}05;1{,}1)$}
  {\True Tứ phân vị thứ nhất  của mẫu số liệu nhóm là $Q_1 \approx 0{,}98$}
  {Tứ phân vị thứ ba của mẫu số liệu nhóm  là $Q_3 \approx 1{,}248$}
  \loigiai{
  \begin{center}
   \begin{tabular}{|c|c|c|c|c|c|}
    \hline 
    Điện lượng  (Nghìn mAh) & $[0{,}9;0{,}95)$ & $[0{,}95;1{,}0)$ & $[1{,}0;1{,}05)$ & $[1{,}05;1{,}1)$ & $[1{,}1;1{,}15)$ \\
    \hline 
    Giá trị đại diện & $0{,}925$ & $0{,}975$ & $1{,}025$  & $1{,}075$ & $1{,}125$ \\
    \hline 
    Số pin & 10 & 20 & 35 & 15  & 5 \\
    \hline 
   \end{tabular}
  \end{center}
  \begin{itemchoice}
   \itemch \textbf{Đúng}.\\ 
   Vì số trung bình của mẫu số liệu ghép nhóm trên là $$\overline{x} =\dfrac{0{,}925 \cdot 10 + 0{,}975 \cdot 20 + 1{,}025 \cdot 35 +  1{,}057 \cdot 15 + 1{,}125 \cdot 5}{85} \approx 1{,}016.$$
    \itemch \textbf{Sai}.\\
   Vì nhóm chứa mốt của  mẫu số liệu ghép nhóm trên là $[1{,}05; 1{,}1)$ $$M_0= 1+ \dfrac{35-20}{(35-20)+(35-15)}(1{,}05-1)\approx 1{,}02.$$
    \itemch \textbf{Đúng}.\\
   Gọi $x_1, x_2, \ldots, x_{85}$  lần lượt là điện lượng mỗi viên pin xếp theo thứ tự không giảm.\\
   Do $x_1, x_2, \ldots, x_{10} \in [0{,}9;0{,}95)$; $x_{11}, \ldots, x_{30}  \in [0{,}95;1{,}0)$; $x_{31}, \ldots, x_{65} \in [1{,}0;1{,}05)$; $x_{66},
   \ldots, x_{80} \in [1{,}05;1{,}1)$; $x_{81}, \ldots, x_{85} \in [1{,}1;1{,}15)$ nên trung vị của mẫu số liệu $x_1, x_2,  \ldots, x_{85}$ là $x_{43} \in [1;1{,}05)$.\\
   Ta xác định được $n =  85, n_m = 35, C = 30, u_m = 1, u_{m+1} = 1{,}05$.\\
   Suy ra tứ phân vị thứ hai của mẫu số liệu là $$Q_2 = 1 + \dfrac{\dfrac{85}{2} - 30}{35} (1{,}05 -  1) \approx 1{,}02.$$
   Tứ phân vị thứ nhất của mẫu số liệu là $\dfrac{1}{2} (x_{21} + x_{22})$.\\ 
   Do $x_{21}, x_{22} \in [0{,}95;1)$ nên tứ phân vị thứ nhất của mẫu số liệu nhóm là $$Q_1 = 0{,}95 + \dfrac{\dfrac{85}{4} - 10}{20} \cdot (1  - 0{,}95) \approx 0{,}98.$$
    \itemch \textbf{Sai}.\\
   Tứ phân  vị thứ ba của mẫu số liệu là $\dfrac{1}{2} (x_{64} + x_{65})$.\\
   Do $x_{64}, x_{65} \in  [1;1{,}05)$ nên tứ phân vị thứ ba của mẫu số liệu nhóm là
   $$Q_3 = 1 + \dfrac{\dfrac{3 \cdot 85}{4} - 30}{35} (1{,}05  - 1) \approx 1{,}048.$$
  \end{itemchoice}
  }
\end{ex}

% nguồn: Ok_dungsai_lop11_Lan2
\begin{ex}%[1D5H2-3]%[ma:1bdde387]
	Người ta đo đường kính của 61 cây gỗ được trồng sau 12 năm  (đơn vị: centimét), họ thu được bảng tần số ghép nhóm sau:
	\begin{center}\begin{tabular}{|c|c|c|c|c|c|}
		 \hline
		 Đường kính&$[20;25)$&$[25;30)$&$[30;35)$&$[35;40)$&$[40;45)$\\
		 \hline
		Số cây&4&12&26&13&6\\ 
		\hline 
	\end{tabular}
\end{center}
 \choiceTF
	{\True Cỡ mẫu của mẫu số liệu  là $n=61$}
	{Tứ phân  vị thứ hai của mẫu số liệu ghép nhóm là$Q_2=32{,}79$}
	{\True Tứ phân vị thứ nhất của  mẫu số liệu ghép nhóm là $Q_1\approx 19{,}69$}
	{\True Tứ phân vị thứ ba của mẫu số liệu ghép  nhóm là $Q_3=36{,}44$}
 \loigiai{
	\begin{itemchoice}
		\itemch Cỡ mẫu của mẫu  số liệu là $n=61$.\\
		Gọi $x_1,x_2,\ldots,x_{61}$ là mẫu số liệu được sắp xếp theo thứ tự  không giảm.\\
		Trung vị của mẫu số liệu này là $x_{31}\in \left[30;35\right)$.\\ 
		 Ta có $n_m=26;C_1=4+12=16;u_m=30;u_{m+1}=35$.
		\itemch Tứ phân vị thứ hai chính là trung vị của mẫu số liệu  ghép nhóm là
		 $$Q_2=M_e=30+\dfrac{\dfrac{61}{2}-16}{26}\left(35-30\right)=\dfrac{1705}{52}\approx 32{,}79 \text{ cm}.$$
		\itemch Xét nửa mẫu số  liệu bên trái $x_1,x_2,\ldots,x_{30}$ có trung vị $\dfrac{x_{15}+x_{16}}{2}\in \left[25;30\right)$.\\
		Ta  có $n_i=12;C_1=4;x_i=25;x_{i+1}=30$.\\
		Suy ra tứ phân vị thứ nhất của mẫu số liệu  là $$Q_1=25+\dfrac{\dfrac{61}{4}-4}{12}\left(30-25\right)=\dfrac{475}{16}\approx 29{,}69 \text{ cm}.$$
		\itemch Xét nửa mẫu số liệu bên trái $x_{32},x_{33},\ldots,x_{61}$ có  trung vị $\dfrac{x_{46}+x_{47}}{2}\in \left[35;40\right)$.\\
		 Ta có $n_j=13;C_3=4+12+26=42;x_i=35;x_{i+1}=40$.\\
		Suy ra tứ phân vị thứ ba của mẫu số liệu  là $$Q_3=35+\dfrac{\dfrac{3\cdot61}{4}-42}{13}\left(40-35\right)=\dfrac{1895}{52}\approx 36{,}44 \text{ cm}.$$
		Vậy các tứ phân  vị của mẫu số liệu ghép nhóm là
		$Q_1\approx 29{,}69;Q_2=32{,}79;Q_3=36{,}44$. 
	\end{itemchoice}
} 
\end{ex}

% nguồn: Ok_dungsai_lop11_Lan2
\begin{ex}%[1D5H2-3]%[ma:44e60d1b]
	Một hãng xe ô tô thống kê lại số  lần gặp sự cố về động cơ của 100 chiếc xe cùng loại sau 2 năm sử dụng đầu tiên ở bảng sau:
	\begin{center}\begin{tabular}{|c|c|c|c|c|c|}
		\hline 
		Số lần  gặp sự cố&$[1;2]$&$[3;4]$&$[5;6]$&$[7;8]$&$[9;10]$\\
		 \hline
		Số  xe&17&33&25&20&5\\
		 \hline
	\end{tabular}
\end{center}
 \choiceTF
	{\True Cỡ mẫu của  mẫu số liệu là $n=100$.}
	{Tứ phân vị  thứ nhất của mẫu số liệu ghép nhóm là $Q_1\approx 1{,}98$}
	{\True Tứ phân vị thứ  hai của mẫu số liệu ghép nhóm là $Q_2=4{,}5$}
	{\True Tứ phân vị thứ ba của mẫu số liệu ghép nhóm là  $Q_3=6{,}5$}
\loigiai{ 
	\begin{itemchoice}
		\itemch Số lần gặp sự cố là số nguyên nên ta có thể sử  dụng bảng tần số ghép nhóm sau:
		\begin{center}\begin{tabular}{|c|c|c|c|c|c|}
			\hline 
			 Số lần gặp sự cố&$[0{,}5;2{,}5)$&$[2{,}5;4{,}5)$&$[4{,}5;6{,}5)$&$[6{,}5;8{,}5)$&$[8{,}5;10{,}5)$\\
			\hline 
			 Số xe&17&33&25&20&5\\
			 \hline
		\end{tabular}
	\end{center}
	Cỡ mẫu của mẫu số liệu là  $n=100$.\\
	Gọi $x_1,x_2,x_3,\ldots,x_{100}$ là mẫu số liệu  được xếp theo thứ tự không giảm.
	\itemch Trung vị của của mẫu số liệu  là $\dfrac{x_{50}+x_{51}}{2}$ với $x_{50}\in \left[2{,}5;4{,}5\right),x_{51}\in \left[4{,}5;6{,}5\right)$.\\
	Vì vậy, tứ phân vị thứ nhất của mẫu số liệu ghép nhóm  là
	 $$Q_1=2{,}5+\dfrac{\dfrac{100}{4}-17}{33}\left(4{,}5-2{,}5\right)=\dfrac{197}{66}\approx 2{,}98.$$
	\itemch Suy ra tứ phân vị thứ hai cũng là trung vị của mẫu số  liệu ghép nhóm là
	$Q_2=4{,}5. $Xét nửa mẫu số liệu bên trái $x_1,x_2,x_3,\ldots,x_{50}$ có trung vị $\dfrac{x_{25}+x_{26}}{2}\in \left[2{,}5;4{,}5\right)$.\\ 
	Ta có  $n_m=33;C=17;u_{m+1}=4{,}5;u_m=2{,}5$.
	\itemch Xét nửa mẫu số liệu bên phải $x_{51},x_{52},\ldots,x_{100}$ có trung vị 
	$\dfrac{x_{75}+x_{76}}{2}\text{;}x_{75}\in \left[4{,}5;6{,}5\right),x_{76}\in \left[6{,}5;8{,}5\right)$ nên tứ phân  vị thứ ba của mẫu số liệu ghép nhóm $Q_3=6{,}5$.\\
	Vậy các tứ phân vị của mẫu số  liệu ghép nhóm là $Q_1\approx 2{,}98;Q_2=4{,}5;Q_3=6{,}5. $.
	\end{itemchoice}
 }
\end{ex}

% nguồn: Ok_dungsai_lop11_Lan2
\begin{ex}%[1D5H2-3]%[ma:b7a561c3]
	Số điểm một  cầu thủ ghi được trong 20 trận đấu được cho ở bảng sau
	\begin{center}\begin{tabular}{|c|c|c|c|c|c|c|c|c|c|}
		 \hline
		$25$&$23$&$21$&$13$&$8$&$14$&$15$&$18$&$22$&$11$\\ 
		 \hline
		$24$&$12$&$14$&$14$&$18$&$6$&$8$&$25$&$10$&$11$\\ 
		 \hline
	\end{tabular}
\end{center}
 \choiceTF
	{\True Tứ phân vị thứ hai của dãy số  liệu là $Q_2=14$}
	{\True Tứ  phân vị thứ ba của dãy số liệu là $Q_3=21{,}5$}
	{\True Tổng hợp lại  dãy số liệu trên vào bảng tần số ghép nhóm theo mẫu sau
	\begin{center}\begin{tabular}{|c|c|c|c|c|}
		 \hline
		Điểm số &$[6;10)$&$[11;15)$&$[16;20)$&$[21;25)$\\ 
		 \hline
		 Số trận&$4$&$8$&$2$&$6$\\
		 \hline
\end{tabular}
\end{center}}
{Ước lượng tứ phân vị của  số liệu ở bảng tần số ghép nhóm trên ta được tứ phân vị thứ hai của mẫu số liệu là $Q_2=8{,}25$}
\loigiai{ 
	Sắp xếp mẫu số liệu theo  thứ tự không giảm ta được
	 $$6;8;8;10;11;11;12;13;14;14;14;15;18;18;21;22;23;24;25;25.$$
	\begin{itemchoice}
		\itemch Tứ phân  vị thứ hai của dãy số liệu là $Q_2=\dfrac{14+14}{2}=14.$
		\itemch Tứ phân vị thứ nhất của dãy  số liệu $6;8;8;10;11;11;12;13;14;14$;
		$Q_1=\dfrac{11+11}{2}=11.$ 
		Tứ phân vị  thứ ba của dãy số liệu $14;15;18;18;21;22;23;24;25;25$
		 $Q_3=\dfrac{21+22}{2}=21{,}5.$
		\itemch Tổng hợp lại dãy số liệu trên vào bảng tần số  ghép nhóm theo mẫu sau
		\begin{center}\begin{tabular}{|c|c|c|c|c|}
			\hline 
			 Điểm số &$[6;10)$&$[11;15)$&$[16;20)$&$[21;25)$\\
			\hline 
			 Số trận&$4$&$8$&$2$&$6$\\
			 \hline
		\end{tabular}
	\end{center}
	\itemch Vì số trận là số nguyên nên ta hiệu chỉnh lại bảng số liệu  sau:
	\begin{center}\begin{tabular}{|c|c|c|c|c|}
		\hline 
		Điểm số  &$[5{,}5;10{,}5)$&$[10{,}5;15{,}5)$&$[15{,}5;20{,}5)$&$[20{,}5;25{,}5)$\\
		 \hline
		Số trận&$4$&$8$&$2$&$6$\\ 
		\hline 
	\end{tabular}
	\end{center}
	Gọi $x_1;x_2;\ldots;x_{20}$ lần lượt là số điểm ghi được ở mỗi trận đấu xếp theo thứ tự  không giảm.\\
	Do $x_1;\ldots;x_4\in [5{,}5;10{,}5);x_5;\ldots;x_{12}\in [10{,}5;15{,}5);x_{13},x_{14}\in [15{,}5;20{,}5);x_{15};\ldots;x_{20}\in  [20{,}5;25{,}5)$ nên trung vị của mẫu số liệu $x_1;\ldots;x_{20}$ là $\dfrac{1}{2}\left(x_{10}+x_{11}\right)\in [10{,}5;15{,}5)$.\\
	Ta xác  định được $n=20,n_m=8,C=4,u_m=10{,}5;u_{m+1}=15{,}5$.\\
	Suy ra tứ phân vị thứ hai của mẫu  số liệu là $Q_2=10{,}5+\dfrac{\dfrac{20}{2}-4}{8}(15{,}5-10{,}5)=14{,}25$.\\
	Tứ phân  vị thứ nhất của mẫu số liệu là $\dfrac{1}{2}\left(x_5+x_6\right)$.\\
	Do $x_5,x_6\in [10{,}5;15{,}5)$  nên tứ phân vị thứ nhất của mẫu số liệu nhóm là
	 $Q_1=10{,}5+\dfrac{\dfrac{20}{4}-4}{8}(15{,}5-10{,}5)=11{,}125$.\\
	Tứ phân vị thứ ba của mẫu số  liệu là $\dfrac{1}{2}\left(x_{15}+x_{16}\right)$.\\
	Do  $x_{15},x_{16}\in [20{,}5;25{,}5)$ nên tứ phân vị thứ ba của mẫu số liệu nhóm là\\
	$Q_3=20{,}5+\dfrac{\dfrac{3\cdot20}{4}-14}{6}\left(25{,}5-20{,}5\right)=21{,}3$. 
	\end{itemchoice}
} 
\end{ex}

% nguồn: Ok_dungsai_lop11_Lan2
\begin{ex}%[1D5H2-3]%[ma:a4d0af62]
	Kiểm tra điện lượng của một số viên pintiểu do một hãng sản xuất  thu được kết quả sau:
	\begin{center}\begin{tabular}{|c|c|c|c|c|c|}
		 \hline
		Điện lượng  (Nghìn mAh)&$[0{,}9;0{,}95)$&$[0{,}95;1{,}0)$&$[1{,}0;1{,}05)$&$[1{,}05;1{,}1)$&$[1{,}1;1{,}15)$\\
		 \hline
		Số pin&10&20&35&15&5\\ 
		 \hline
	\end{tabular}
\end{center}
\choiceTF 
	{\True  Số trung bình của dãy số liệu là $1{,}016$}
	{Nhóm chứa mốt của dãy số  liệu là $[1{,}05;1{,}1)$}
	{\True Tứ phân vị thứ nhất của mẫu số  liệu nhóm là $Q_1=0{,}98$}
	{Tứ phân vị thứ ba của mẫu số liệu nhóm  là $Q_3=1{,}248$}
\loigiai{ 
	\begin{center}\begin{tabular}{|c|c|c|c|c|c|}
		\hline 
		Điện lượng  (Nghìn mAh)&$[0{,}9;0{,}95)$&$[0{,}95;1{,}0)$&$[1{,}0;1{,}05)$&$[1{,}05;1{,}1)$&$[1{,}1;1{,}15)$\\
		\hline 
		Giá trị  đại diện&$0{,}925$&$0{,}975$&$1{,}025$&$1{,}075$&$1{,}125$\\
		 \hline
		Số  pin&10&20&35&15&5\\
		 \hline
	\end{tabular}
	\end{center}
	\begin{itemchoice}
		\itemch Số trung bình của dãy số  liệu là
		$$\dfrac{0{,}925\cdot10+0{,}975\cdot20+1{,}025\cdot35+1{,}075\cdot15+1{,}125\cdot5}{85}=1{,}016.$$ 
		\itemch  Nhóm chứa mốt của dãy số liệu là $\left[1;1{,}05\right)$
		Gọi $x_1;x_2;\ldots;x_{85}$ lần lượt là  điện lượng mỗi viên pin xếp theo thứ tự không giảm.\\
		Do $x_1;x_2;\ldots;x_{10}\in \left[0{,}9;0{,}95\right);x_{11};\ldots;x_{30}\in \left[0{,}95;1{,}0\right);x_{31};\ldots;x_{65}\in  \left[1{,}0;1{,}05\right)$;
		$x_{66};\ldots;x_{80}\in \left[1{,}05;1{,}1\right);x_{81};\ldots;x_{85}\in \left[1{,}1;1{,}15\right)$  nên trung vị của mẫu số liệu $x_1$; $x_2;\ldots;x_{85}$ là $x_{43}\in \left[1;1{,}05\right)$.
		\itemch  Ta xác định được $n=85,n_m=35,C=30{,}u_m=1{,}u_{m+1}=1{,}05$.\\
		Suy ra tứ phân vị  thứ hai của mẫu số liệu là $Q_2=1+\dfrac{\dfrac{85}{2}-30}{35}\left(1{,}05-1\right)=1{,}02$.\\
		 Tứ phân vị thứ nhất của mẫu số liệu là $\dfrac{1}{2}\left(x_{21}+x_{22}\right)$.\\
		Do $x_{21},x_{22}\in \left[0{,}95;1\right)$ nên tứ phân  vị thứ nhất của mẫu số liệu nhóm là
		$$Q_1=0{,}95+\dfrac{\dfrac{85}{4}-10}{20}\left(1-0{,}95\right)=0{,}98.$$ 
		\itemch  Tứ phân vị thứ ba của mẫu số liệu là $\dfrac{1}{2}\left(x_{64}+x_{65}\right)$.\\
		Do $x_{64},x_{65}\in \left[1;1{,}05\right)$ nên  tứ phân vị thứ ba của mẫu số liệu nhóm là
		$$Q_3=1+\dfrac{\dfrac{3\cdot85}{4}-30}{35}\left(1{,}05-1\right)=1{,}048.$$ 
	\end{itemchoice}
} 
\end{ex}

% nguồn: Ok_dungsai_lop11_Lan2
\begin{ex}%[1D5H2-3]%[ma:9b02b4ab]
	Thời gian tự học  mỗi ngày của học sinh khối $11$ được thống kê trong bảng sau
	\begin{center}
		\begin{tabular}{|c|c|c|c|c|c|}
			\hline 
			Thời gian (phút) & $[60; 90)$ & $[90;  120)$ & $[120; 150)$ & $[150; 180)$ & $[180; 210)$ \\
			 \hline
			 Số học sinh & $8$ & $10$ & $13$ & $7$ & $2$ \\
			\hline 
		\end{tabular}
	\end{center}
	\choiceTF 
		{\True Tổng số học sinh được khảo sát  là $40$}
		{\True Giá trị đại diện của nhóm $[90; 120)$ là $105$} 
		{Trung bình thời gian tự học mỗi ngày  của các bạn học sinh được khảo sát là $135$ phút}
		{Có $13$ học sinh tự học ít nhất $2$ giờ  mỗi ngày}
	 \loigiai{
		Bảng tần số ghép nhóm theo  giá trị đại diện là
		\begin{center}
			\begin{tabular}{|c|c|c|c|c|c|}
				 \hline
				Thời gian (phút) & $[60; 90)$ & $[90; 120)$ & $[120; 150)$ & $[150;  180)$ & $[180; 210)$ \\
				 \hline
				Giá trị đại diện  & $75$ & $105$ & $135$ & $165$ & $195$ \\
				\hline 
				Số học sinh & $8$ & $10$ & $13$ & $7$ &  $2$ \\
				 \hline
			\end{tabular}
		\end{center}
		\begin{itemchoice}
			\itemch \textbf{Đúng.}\\ 
			Tổng số học sinh được khảo sát là $8+10+13+7+2=40$  (học sinh).
			\itemch  \textbf{Đúng.}\\
			Giá trị đại diện của nhóm $[90;120)$ là $\dfrac{90+120}{2}=105$. 
			\itemch \textbf{Sai.}\\ 
			Thời gian tự học trung bình mỗi ngày của các bạn học sinh được khảo sát là 
			\[\overline{x}=\dfrac{75\cdot 8+105\cdot 10+135\cdot  13+165\cdot 7+195\cdot 2}{40}=123{,}75.\]
			\itemch  \textbf{Sai.}\\
			Số học sinh tự học ít nhất $2$ giờ  mỗi ngày là $13+7+2=22$ (học sinh).
		\end{itemchoice}
	 }
\end{ex}

% nguồn: Ok_dungsai_lop11_Lan2
\begin{ex}%[1D5H2-3]%[ma:74bf312c]
	Dựa vào bảng tần số mẫu số liệu ghép nhóm sau: 
	\begin{center}\begin{tabular}{|c|c|c|c|c|c|c|}
		\hline 
		Nhóm&$[30;40)$&$[40;50)$&$[50;60)$&$[60;70)$&$[70;80)$&$[80;90)$\\ 
		\hline 
		Tần số&2&10&16&8&2&2\\ 
		\hline 
	\end{tabular}
\end{center}
\choiceTF 
	{\True Cỡ mẫu của mẫu số liệu là  $n=40$}
	{\True Tứ phân vị thứ nhất của mẫu số liệu ghép  nhóm là $Q_1=48$}
	{Tứ phân vị thứ hai của  mẫu số liệu ghép nhóm là $Q_2=45$}
	{Tứ phân vị thứ ba của mẫu số liệu ghép  nhóm là $Q_3=61{,}5$}
\loigiai{ 
	\begin{itemchoice}
		\itemch Cỡ mẫu của mẫu số liệu là  $n=40$.\\
		Gọi $x_1,x_2,x_3,\ldots,x_{40}$ là mẫu số liệu được xếp theo thứ tự  không giảm.\\
		Trung vị của mẫu số liệu:  $\dfrac{x_{20}+x_{21}}{2}\in \left[50;60\right)$.\\
		Ta có $n_m=16;C_2=2+10=12;u_m=50;u_{n+1}=60$.\\ 
		Xét nửa mẫu số liệu  bên trái $x_1,x_2,x_3,\ldots,x_{20}$ có trung vị $\dfrac{x_{10}+x_{11}}{2}\in \left[40;50\right)$.\\
		Ta có $n_i=10;C_1=2;u_i=40;u_{i+1}=50$. 
		\itemch Tứ phân  vị thứ nhất của mẫu số liệu ghép nhóm là $Q_1=40+\dfrac{\dfrac{40}{4}-2}{10}\left(50-40\right)=48$.
		\itemch Tứ phân vị thứ hai của mẫu  số liệu ghép nhóm cũng là trung vị của mẫu số liệu đó là
		$Q_2=M_e=50+\dfrac{\dfrac{40}{2}-12}{16}\left(60-50\right)=55$. 
		\itemch Xét nửa mẫu  số liệu bên phải $x_{21},x_{22},x_{23},\ldots,x_{40}$ có trung vị $\dfrac{x_{30}+x_{31}}{2}\in \left(60;70\right)$.\\
		Ta có $n_j=8$;  $C_3=2+10+16=28$; $u_j=60$; $u_{j+1}=70$.\\
		Tứ phân  vị thứ ba của mẫu số liệu ghép nhóm là $Q_3=60+\dfrac{\dfrac{3\cdot40}{4}-28}{8}\left(70-60\right)=62{,}5$.
		Vậy tứ phân vị của mẫu số liệu ghép nhóm trên là  $Q_1=48$, $Q_2=55$, $Q_3=62{,}5$.
	\end{itemchoice}
} 
\end{ex}

% nguồn: Ok_dungsai_lop11_Lan2
\begin{ex}%[1D5H2-3]%[ma:33fa636c]
	Cho  bảng tần số ghép nhóm sau:
	\begin{center}\begin{tabular}{|c|c|c|c|c|c|}
		\hline 
		 Nhóm&$[0;2)$&$[2;4)$&$[4;6)$&$[6;8)$&$[8;10)$\\
		\hline 
		Tần  số&3&8&12&12&4\\
		 \hline
	\end{tabular}
\end{center}
\choiceTF 
	{Cỡ mẫu của mẫu số liệu là $n=38$} 
	{Tứ phân vị thứ nhất của mẫu số liệu  ghép nhóm là $Q_1\approx 2{,}69$}
	{\True Tứ phân vị thứ hai của mẫu số liệu ghép nhóm là  $Q_2\approx 5{,}42$}
	{\True Tứ phân vị thứ ba của mẫu số liệu ghép nhóm là $Q_3=7{,}04$} 
 \loigiai{
	\begin{itemchoice}
		 \itemch Cỡ mẫu của mẫu số liệu là $n=3+8+12+12+4=39$.
		\itemch Gọi $x_1,x_2,\ldots,x_{39}$  là mẫu số liệu được sắp xếp theo thứ tự không giảm.\\
		Trung vị của mẫu số liệu này là  $x_{20}\in \left[4;6\right)$.\\
		 Ta có $n_m=12;C_2=3+8=11;u_m=4;u_{m+1}=6$.\\
		Xét nửa mẫu số liệu bên trái  $x_1,x_2,\ldots,x_{19}$ có trung vị $x_{10}\in \left[2;4\right)$.\\
		Ta có $n_i=8;C_1=3;x_i=2;x_{i+1}=4$. Suy ra tứ phân vị  thứ nhất của mẫu số liệu là $Q_1=2+\dfrac{\dfrac{39}{4}-3}{8}\left(4-2\right)=\dfrac{59}{16}\approx 3{,}69$.
		\itemch Tứ phân vị thứ hai chính là trung vị của mẫu  số liệu ghép nhóm là
		 $$Q_2=M_e=4+\dfrac{\dfrac{39}{2}-11}{12}\left(6-4\right)=\dfrac{65}{12}\approx 5{,}42.$$
		\itemch 
		Xét nửa mẫu số liệu bên phải $x_{21},x_{22},\ldots,x_{39}$ có trung  vị $x_{30}\in \left[6;8\right)$.\\
		Ta có $n_j=12;C_3=3+8+12=23;x_j=6;x_j=8$.\\ 
		Suy ra  tứ phân vị thứ ba của mẫu số liệu là $Q_3=6+\dfrac{\dfrac{3\cdot39}{4}-23}{12}\left(8-6\right)=\dfrac{169}{24}\approx 7{,}04$.\\
		Vậy các tứ phân vị của mẫu số liệu  ghép nhóm là $Q_1\approx 3{,}69;Q_2=5{,}42;Q_3=7{,}04$.
	\end{itemchoice}
} 
\end{ex}

% nguồn: Ok_dungsai_lop11_Lan2
\begin{ex}%[1D5H2-3]%[ma:cc6360b7]
	Thống kê điểm trung bình môn Toán  của một số học sinh lớp 11 được cho ở bảng sau:
	\begin{center}\begin{tabular}{|c|c|c|c|c|c|c|c|}
		\hline 
		Khoảng  điểm &$[6{,}5;7)$&$[7;7{,}5)$&$[7{,}5;8)$&$[8;8{,}5)$&$[8{,}5;9)$&$[9;9{,}5)$&$[9{,}5;10)$\\
		\hline 
		Số  học sinh&$8$&$10$&$16$&$24$&$13$&$7$&$4$\\
		\hline 
	\end{tabular}
\end{center}
\choiceTF 
	{Cỡ mẫu của mẫu số liệu là  $n=80$}
	{\True Tứ phân vị thứ hai của mẫu số liệu ghép nhóm là  $Q_2=8{,}15$}
	{\True Tứ phân vị thứ nhất  của mẫu số liệu ghép nhóm là $Q_1=7{,}58$}
	{\True Tứ phân vị thứ ba của mẫu số  liệu ghép nhóm là $Q_3=8{,}63$}
 \loigiai{
	\begin{itemchoice}
		\itemch Gọi $x_1;x_2;\ldots;x_{82}$ lần  lượt là điểm trung bình môn Toán của các học sinh sắp xếp theo thứ tự không giảm.\\
		 \itemch Ta có $x_1;\ldots;x_8\in \left[6{,}5;7\right);x_9;\ldots;x_{18}\in \left[7;7{,}5\right);x_{19};\ldots;x_{34}\in \left[8;8{,}5\right)$;
		$x_{59};\ldots;x_{71}\in  \left[8{,}5;9\right);x_{72};\ldots;x_{78}\in \left[7;7{,}5\right);x_{79};\ldots;x_{82}\in \left[9{,}5;10\right)$.\\
		Tứ phân vị thứ hai của mẫu số liệu là $\dfrac{1}{2}\left(x_{41}+x_{42}\right)$  mà $x_{41};x_{42}\in \left[8;8{,}5\right)$ nên tứ phân vị thứ hai của mẫu số liệu là $Q_2=8+\dfrac{\dfrac{82}{2}-34}{24}\left(8{,}5-8\right)=8{,}15$.\\
		\itemch Tứ phân vị thứ  nhất của mẫu số liệu là $x_{21}\in \left[7{,}5;8\right)$.\\
		 Do đó tứ phân vị thứ nhất của mẫu số liệu là $Q_1=7{,}5+\dfrac{\dfrac{82}{4}-18}{16}\left(8-7{,}5\right)=7{,}58$.
		\itemch  Tứ phân vị thứ ba của mẫu số liệu là $x_{62}\in \left[8{,}5;9\right)$.\\
		Do đó tứ phân vị thứ ba  của mẫu số liệu là $Q_3=8{,}5+\dfrac{\dfrac{3\cdot82}{4}-58}{13}\left(9-8{,}5\right)=8{,}63$.
	\end{itemchoice}
} 
\end{ex}

% nguồn: Ok_dungsai_lop11_Lan1
\begin{ex}%[1D5V2-3]%[ma:7de26f7b]
 Điều tra về chiều cao của học sinh khối lớp  10, ta có kết quả sau:
 \begin{center}
  \begin{tabular}{|c|c|c|}
   \hline Nhóm & Chiều cao $(\mathrm{cm})$ & Số  học sinh \\
   \hline 1 & {$[150; 152)$}  & $5$ \\
   \hline 2 & {$[152; 154)$}  & $18$ \\
    \hline 3 & {$[154; 156)$} & $40$ \\
   \hline 4 & {$[156;  158)$} & $26$ \\
   \hline 5 & {$[158;  160)$} & $8$ \\
   \hline 6 & {$[160; 162)$} &  $3$ \\
    \hline
  \end{tabular}
 \end{center}
  \choiceTF
  {\True Tổng số học sinh  điều tra bằng $100$}
  {\True Chiều cao trung  bình của các em học sinh trên là $155{,}46$}
  {Số trung vị của mẫu số liệu trên  bằng $156{,}35$}
  {\True Tứ phân vị thứ nhất  của mẫu số liệu trên bằng $154{,}1$}
 \loigiai{ 
  \begin{itemchoice}
   \itemch Đúng. Vì ta có tổng số học sinh là  $n=100$.
   \itemch Đúng. Vì chiều cao trung bình bằng  $$\overline{x}=\dfrac{5\cdot 151+18\cdot 153+40\cdot 155+26\cdot 157+8.159+3\cdot 161}{100}=155{,}46.$$
   \itemch Sai. Vì $n=100$ nên trung vị là $\dfrac{x_{50}+x_{51}}{2}$, do hai giá  trị $x_{50}, x_{51}$ thuộc nhóm 3.\\
   Áp dụng công thức  $M_e=154+\dfrac{50-(18+5)}{40} \cdot 2 \approx 155{,}35$.
   \itemch Đúng. Vì tứ phân vị thứ  nhất $Q_3$ là $\dfrac{x_{25}+x_{26}}{2}$ mà $x_{25}, x_{26}$ thuộc nhóm $[154; 156)$.\\
   Do đó, $p=3$, $a_p=154$, $m_p=40$, $m_1+m_2=5+18=23$,  $a_{p+1}-a_p=2$ và ta có
   $$Q_1=154+\dfrac{\dfrac{100}{4}-23}{40} \times 2=154{,}1.$$ 
  \end{itemchoice}
 } 
\end{ex}

% nguồn: Ok_dungsai_lop11_Lan1
\begin{ex}%[1D5V2-3]%[ma:e93a377e]
 Kết quả điều tra về số giờ làm thêm trong một tuần của sinh viên một trường đại học $\mathrm{X}$ được cho bởi bảng  sau:
 \begin{center}
  \begin{tabular}{|l|c|c|c|c|c|}
   \hline Số giờ làm thêm & {$[2; 4)$} & {$[4; 6)$} & {$[6; 8)$} &  {$[8; 10)$} & {$[10; 12)$}\\
   \hline Số sinh viên &  $12$ & $20$ & $37$ & $21$ & $10$ \\
   \hline 
  \end{tabular}
 \end{center}
 \choiceTF 
  {\True Số sinh viên được điều tra là  $100$}
  {\True Số giờ làm thêm trung bình của mỗi sinh viên  trường đại học $\mathrm{X}$ không ít hơn $6$}
  {Mốt của mẫu số liệu trên là  $7{,}5$}
  {Tứ phân vị thứ hai  của dãy số liệu lớn hơn $6{,}5$}
 \loigiai{ 
  \begin{itemchoice}
   \itemch  Đúng. Vì cỡ mẫu $n=100$.
   \itemch Đúng. Vì từ  bảng thống kê ta có
   \begin{center}
    \begin{tabular}{|l|c|c|c|c|c|}
     \hline Số giờ làm thêm & {$[2; 4)$} & {$[4; 6)$} & {$[6; 8)$}  & {$[8; 10)$} & {$[10; 12)$} \\
     \hline Số giờ làm thêm đai  diện & $3$ & $5$ & $7$ & $9$ & $11$ \\
     \hline Số sinh viên  & $12$ & $20$ & $37$ & $21$ & $10$ \\
     \hline 
    \end{tabular}
   \end{center}
   Số trung bình của  mẫu số liệu là
   $$\overline{x}=\dfrac{3\cdot 12+5\cdot 20+7\cdot 37+9\cdot  21+11\cdot 10}{100}=6\cdot 94.$$
   \itemch Sai. Vì nhóm  chứa mốt số liệu trên là nhóm $[6; 8)$.\\
   Do đó $u_m=6$; $n_m=37$, $n_{m-1}=20$; $n_{m+1}=21$;  $u_{m+1}=8$.\\
   Vậy mốt của mẫu số liệu ghép  nhóm là
   $$M_0=6+\dfrac{37-20}{(37-20)+(37-21)}\cdot(8-6)\approx  7{,}03.$$
   \itemch Sai. Vì $x_1; x_2; \ldots; x_{100}$ là  mẫu số liệu theo thứ tự không giảm.\\
   Tứ phân vị thứ hai của dãy số  liệu $x_1; x_2; \ldots; x_{100}$ là $\dfrac{1}{2}\left(x_{50}+x_{51}\right)$.\\
    Do $x_{50}$ và $x_{51}$ thuộc nhóm $[6; 8)$ nên $Q_2=6+\dfrac{\dfrac{2 \cdot 100}{4}-32}{37} \cdot(8-6) \approx 6{,}97$.
  \end{itemchoice}
  }
\end{ex}

% nguồn: Ok_dungsai_lop11_Lan1
\begin{ex}%[1D5V2-3]%[ma:ba05f152]
 Cho bảng biểu diễn mẫu số liệu ghép nhóm về chiều cao của $50$ mẫu cây ở  một vườn thực vật (đơn vị: centimet). Xét tính đúng sai của các mệnh đề sau.
 \begin{center}
  \begin{tabular}{|c|c|c|}
   \hline Nhóm  & Tần số & Tần số tích luỹ \\
   \hline $[40; 45)$ & $7$ &  $7$ \\
   $[45; 50)$& $12$  & $19$ \\
   $[50; 55)$&  $9$ & $28$ \\
   $[55; 60)$& $10$ &  $38$ \\
   $[60; 65)$& $8$  & $46$ \\
   $[65; 70)$& $4$ & $50$  \\
   \hline & $n=50$  & \\
    \hline
  \end{tabular}
 \end{center}
  \choiceTF
  {\True Cỡ mẫu  $n=50$}
   {$Q_1=45$}
   {$Q_3=67{,}5$}
   {$\Delta_Q\approx 54{,}075$}
  \loigiai{
  \begin{itemchoice}
    \itemch
   Số phần tử của mẫu  là $n=50$.
   \itemch Ta có $\dfrac{n}{4}=\dfrac{42}{4}=10{,}5$ mà $5<10{,}5<15$. 
    $$Q_1=45+\left(\dfrac{10{,}5-5}{10}\right) \cdot 5=47{,}75~(\mathrm{cm}).$$
   \itemch Ta có $\dfrac{3  n}{4}=\dfrac{3 \cdot 42}{4}=31{,}5$ mà $31<31{,}5<38$.\\
   Áp dụng công thức, ta có tứ phân vị thứ ba là 
   $$Q_3=60+\left(\dfrac{31{,}5-31}{7}\right) \cdot 5 \approx 60{,}4~(\mathrm{cm}).$$ 
   \itemch Từ các kết quả trên  ta suy ra khoảng tứ phân vị của mẫu số liệu ghép nhóm đã cho là\\
   $$ \Delta_Q=Q_3-Q_1 \approx 60{,}4-47{,}75  \approx 12{,}65~(\mathrm{cm}).
   $$ 
 \end{itemchoice}}
\end{ex}

% nguồn: Ok_dungsai_lop11_Lan1
\begin{ex}%[1D5V2-3]%[ma:c90cb281]
 Một hãng xe ô tô thống kê lại số lần gặp sự cố về động cơ của $100$  chiếc xe cùng loại sau $2$ năm sử dụng đầu tiên ở bảng số liệu sau:
 \begin{center}
  \begin{tabular}{|c|c|c|c|c|c|}
   \hline 
   Số lần gặp  sự cố & $[1;2]$ & $[3;4]$ & $[5;6]$ & $[7;8]$ & $[9;10]$ \\
    \hline
   Số xe & $17$ & $33$ &  $25$ & $20$ & $5$ \\
    \hline
  \end{tabular}
 \end{center}
 \choiceTF 
   {\True Cỡ mẫu của số liệu là $n=100$}
  {Tứ phân vị thứ nhất  của mẫu số liệu ghép nhóm là: $Q_1 \approx 1{,}98$}
   {\True Tứ phân vị thứ hai của mẫu số liệu ghép nhóm là: $Q_2=4{,}5$}
  {\True Tứ  phân vị thứ ba của mẫu số liệu ghép nhóm là: $Q_3=6{,}5$}
 \loigiai{ 
  Gọi $x_1, x_2, x_3, \ldots, x_{100}$ là mẫu số liệu được xếp theo thứ tự  không giảm.\\
  Trung vị của của mẫu số liệu là $\dfrac{x_{50}+x_{51}}{2}$ với $x_{50} \in[2{,}5; 4{,}5), x_{51}  \in[4{,}5; 6{,}5)$.
  \begin{itemchoice}
   \itemch Đúng. Vì cỡ mẫu của mẫu  số liệu là $n=100$.
   \itemch Sai. Vì xét nửa mẫu số liệu bên trái $x_1, x_2, x_3, \ldots,  x_{50}$ có trung vị\\ $\dfrac{x_{25}+x_{26}}{2} \in[2{,}5; 4{,}5)$.\\
    Ta có: $n_m=33; C=17; u_{m+1}=4{,}5; u_m=2{,}5$.\\
   Vì vậy, tứ phân vị thứ  nhất của mẫu số liệu ghép nhóm là:
   $$ 
   Q_1=2{,}5+\dfrac{\dfrac{100}{4}-17}{33}(4{,}5-2{,}5)=\dfrac{197}{66} \approx  2{,}98
   $$ 
   \itemch Đúng. Vì tứ phân vị thứ hai cũng là trung vị của  mẫu số liệu ghép nhóm là
    $Q_2=4{,}5$.
   \itemch Đúng. Vì Xét nửa mẫu số liệu bên phải $x_{51}, x_{52}, \ldots, x_{100}$ có trung vị\\ 
   $\dfrac{x_{75}+x_{76}}{2}$ với $x_{75} \in[4{,}5; 6{,}5), x_{76} \in[6{,}5; 8{,}5)$ nên tứ phân vị thứ ba của  mẫu số liệu ghép nhóm $Q_3=6{,}5$.
  \end{itemchoice}
  }
\end{ex}

% nguồn: Ok_dungsai_lop11_Lan1
\begin{ex}%[1D5V2-3]%[ma:ccf0e3d7]
 Thống kê điểm  trung bình môn Toán của một số học sinh lớp 11 được cho ở bảng sau
 \begin{center}
  \begin{tabular}{|c|c|c|c|c|c|c|c|}
   \hline Khoảng điểm & {$[6,5; 7)$} & {$[7; 7,5)$} & {$[7,5; 8)$} & {$[8;  8,5)$} & {$[8,5,9)$} & {$[9; 9,5)$} & {$[9,5; 10)$} \\
    \hline Tần số & 8 & 10 & 16 & 24 & 13 & 7 & 4 \\
    \hline
  \end{tabular}
 \end{center}
  \choiceTF
  {\True Số  trung bình của mẫu số liệu là $\overline{x}=8{,}122$}
  {Tứ phân vị thứ nhất bằng  $Q_1=6{,}75$}
  {\True Tứ  phân vị thứ ba bằng $Q_1=8{,}75$}
  {Mốt của mẫu số  liệu bằng $M_o=8{,}12$}
 \loigiai{ 
  \begin{itemize}
   \item  Số trung bình\\ $\overline{x}=\dfrac{6{,}75\cdot 8+7{,}25\cdot 10+7{,}75\cdot 16+8{,}25\cdot 24+8{,}75\cdot 13+9{,}25\cdot 7+9{,}75\cdot}{82}=8{,}12195122$.
   \item Tứ  phân vị thứ hai chính là trung vị\\
   $Q_2= M_e=u_m+\dfrac{\dfrac{n}{2}-C}{n_m} \cdot\left(u_{m+1}-u_m\right)=8+\dfrac{\dfrac{82}{2}-34}{24}\cdot  \left(8{,}5-8\right)=8{,}25$.
    \item Tứ phân vị thứ nhất \\
    $
    Q_1=u_m+\dfrac{\dfrac{n}{4}-C}{n_m} \cdot\left(u_{m+1}-u_m\right)=7{,}5+\dfrac{\dfrac{82}{4}-18}{16}\left(8-7{,}5\right)=7{,}75
    $.
   \item Tứ phân vị thứ ba  \\ $
   Q_3=u_j+\dfrac{\dfrac{3n}{4}-C}{n_j}  \cdot\left(u_{j+1}-u_j\right)=8{,}5+\dfrac{\dfrac{3\cdot 82}{4}-58}{13} \cdot\left(9-8{,}5\right)=8{,}75
   $. 
   \item  Mốt\\
   $ M_o=u_m+\dfrac{n_m-n_{m-1}}{\left(n_m-n_{m-1}\right)+\left(n_m-n_{m+1}\right)}  .\left(u_{m+1}-u_m\right)\\
   M_o=8+\dfrac{24-16}{(24-16)+(24-13)}\left(8{,}5-8\right)=\dfrac{156}{19}$. 
  \end{itemize}
  }
\end{ex}

% nguồn: Ok_dungsai_lop11_Lan2
\begin{ex}%[1D5V2-3]%[ma:0c8b3f3e]
	Khảo sát thời gian tập thể dục của một số học sinh khối  $11$ thu được mẫu số liệu ghép nhóm sau
	\begin{center}
		\begin{tabular}{|c|c|c|c|c|c|}
			 \hline
			Thời gian (phút) & $[0;  20)$ & $[20; 40)$ & $[40; 60)$ & $[60; 80)$ & $[80; 100)$ \\ \hline
			Số học sinh & $5$ & $9$ & $12$  & $10$ & $6$ \\ \hline
		\end{tabular}
	\end{center}
	 \choiceTF
		{\True Nhóm chứa trung vị của mẫu số liệu trên là $[40;60)$} 
		{Tứ phân vị thứ ba của mẫu số liệu  bằng $70$}
		{\True Thời gian trung bình học sinh tập  thể dục thuộc khoảng $(50;53)$}
		{\True Mốt của mẫu số liệu  trên là $52$}
	\loigiai{ 
		 Bảng tần số ghép nhóm theo giá trị đại diện là
		\begin{center}
			\begin{tabular}{|c|c|c|c|c|c|}
				\hline 
				Thời gian (phút) & $[0;20)$ & $[20;40)$ & $[40;60)$ & $[60;80)$ & $[80;100)$  \\ \hline
				Giá trị đại diện & $10$ & $30$ &  $50$ & $70$ & $90$ \\ \hline
				Số học sinh & $5$ & $9$ &  $12$ & $10$ & $6$ \\ \hline
			\end{tabular}
		\end{center}
		\begin{itemchoice}
			\itemch  \textbf{Đúng.}\\
			Cỡ mẫu  là $n=5+9+12+10+6=42$.
			Gọi $x_{1}$, $x_{2}$, $\ldots$, $x_{42}$ là thời  gian tập của $42$ học sinh, giả sử dãy này đã được sắp xếp theo thứ tự tăng dần, khi đó trung vị là $\dfrac{x_{21}+x_{22}}{2}$. Do hai giá trị $x_{21}$, $x_{22}$ thuộc nhóm $[40;60)$ nên nhóm này chứa trung vị.
			\itemch \textbf{Sai.}\\ 
			Tứ  phân vị thứ ba $Q_{3}$ là $x_{33}$, do nhóm chứa $x_{33}$ là $[60;80)$, do đó $p=4$; $a_{4}=60;m_{1}+m_{2}+m_{3}=5+9+12=26$; $a_{5}-a_{4}=20$. Ta có
			 \[Q_{3}=a_{4}+\dfrac{\dfrac{3n}{4}-(m_{1}+m_{2}+m_{3})}{m_{4}} (a_{5}-a_{4})=60+\dfrac{\dfrac{3\cdot42}{4}-26}{10}\cdot20=71.\]
			\itemch  \textbf{Đúng.}\\
			Thời gian trung bình  học sinh tập thể dục là
			 \[\overline{x}=\dfrac{5\cdot10+9\cdot30+12\cdot50+10\cdot70+6\cdot90}{42}=\dfrac{360}{7}.\]
			\itemch  \textbf{Đúng.}\\
			Tần số lớn nhất là $12$ nên nhóm chứa mốt là $[40;60)$. Ta có $j=3$, $a_{3}=40$,  $m_{3}=12$, $m_{2}=9$, $m_{4}=10$, $h=20$. Do đó
			\[M_{o}=a_{3}+\dfrac{m_{3}-m_{2}}{(m_{3}-m_{2})+(m_{3}-m_{4})} \cdot  h=40+\dfrac{12-9}{(12-9)+(12-10)} \cdot 20=52.\]
		\end{itemchoice}
	} 
\end{ex}

% nguồn: Ok_dungsai_lop11_Lan2
\begin{ex}%[1D5V2-3]%[ma:ebdba05a]
	Dưới đây là hai bảng thống kê doanh số bán hàng của $20$ nhân viên tại một cửa hàng  điện thoại trong tháng $9$ đối với hai nhãn hàng Oppo và Samsung.\\
	Bảng  thống kê doanh số điện thoại bán được của Oppo trong tháng $9$.
	\begin{center}
		\begin{tabular}{|c|c|c|c|c|c|}
			 \hline
			Doanh số &  $[18;20]$ & $[21;23]$ & $[24;26]$ & $[27;29]$ & $[30;32]$ \\ \hline
			Số nhân viên  & $2$ & $5$ & $8$ & $3$ & $2$ \\ \hline
		\end{tabular}
	\end{center}
	Bảng thống kê doanh số điện thoại bán được của Samsung  trong tháng $9$.
	\begin{center}
		\begin{tabular}{|c|c|c|c|c|c|}
			 \hline
			Doanh số & $[15;19]$ & $[20;24]$  & $[25;29]$ & $[30;34]$ & $[35;39]$ \\ \hline
			Số nhân viên & $5$ & $8$ & $5$ & $1$ & $1$ \\ \hline 
		\end{tabular}
	\end{center}
	Chủ cửa hàng điện thoại muốn dành phần thưởng khích lệ cho  các nhân viên bán được doanh số cao. Điều kiện được nhận quà là phải nằm trong Top $5$ nhân viên đạt doanh số điện thoại Samsung cao nhất và đồng thời phải nằm trong Top $10$ nhân viên đạt doanh số điện thoại Oppo cao nhất
	 \choiceTF
		{\True Đối với $20$ nhân viên bán hàng được khảo sát thì bảng thống  kê cho thấy điện thoại của nhãn hàng Oppo dễ bán hơn so với điện thoại của hãng Samsung. (so sánh dựa trên giá trị trung bình của $2$ bảng thống kê)}
		{Đối với điện thoại của hãng Samsung, khả năng một nhân  viên bán được $26$ chiếc là cao nhất}
		{Anh An nghĩ mình sẽ nhận được thưởng vì anh An bán được $25$ chiếc điện thoại Oppo và  $26$ chiếc điện thoại Samsung}
		{Chị Bình nghĩ mình cũng sẽ nhận được thưởng dù chị chỉ bán được  $24$ chiếc điện thoại Oppo nhưng chị bán được tới $27$ chiếc điện thoại Samsung}
	\loigiai{ 
		Ta có bảng số  liệu thống kê sau chuẩn hóa.\\
		Bảng thống kê doanh số điện thoại bán được  của Oppo trong tháng $9$.
		\begin{center}
			\begin{tabular}{|c|c|c|c|c|c|}
				 \hline
				Doanh số & $[17{,}5; 20{,}5)$ & $[20{,}5; 23{,}5)$ & $[23{,}5; 26{,}5)$ & $[26{,}5; 29{,}5)$ & $[29{,}5; 32{,}5)$  \\ \hline
				Giá trị đại diện & $19$ & $22$ & $25$ & $28$  & $31$ \\ \hline
				Số nhân  viên & $2$ & $5$ & $8$ & $3$ & $2$ \\
				\hline 
			\end{tabular}
		\end{center}
		Bảng thống  kê doanh số điện thoại bán được của Samsung trong tháng $9$.
		\begin{center}
			\begin{tabular}{|c|c|c|c|c|c|}
				\hline 
				Doanh số & $[14{,}5; 19{,}5)$  & $[19{,}5; 24{,}5)$ & $[24{,}5; 29{,}5)$ & $[29{,}5; 34{,}5)$ & $[34{,}5; 39{,}5)$ \\ \hline
				Giá  trị đại diện & $17$ & $22$ & $27$ & $32$ & $37$ \\ \hline
				Số nhân viên & $5$ & $8$ & $5$ & $1$ & $1$  \\
				 \hline
			\end{tabular}
		\end{center}
		Gọi $x_{1}$; $x_{2}$; $x_{3}$; $\ldots$; $x_{20}$ là số lượng điện thoại Oppo mà các nhân viên bán  được trong tháng $9$ (được xếp theo thứ tự không giảm).\\
		Ta có trung vị của mẫu số  liệu là $\dfrac{1}{2} (x_{10}+x_{11})\in [23{,}5;26{,}5)$.\\
		Trung vị của  bảng thống kê doanh số điện thoại bán được của hãng Oppo trong tháng $9$
		\[M_{e}=u_{m}+\dfrac{\dfrac{n}{2}-C}{n_{m}} \cdot  (u_{m+1}-u_{m})=23{,}5+\dfrac{10-(2+5)}{8} \cdot 3=24{,}625.\]
		Gọi $x'_{1}$; $x'_{2}$; $x'_{3}$; $\ldots$; $x'_{20}$ là số lượng  điện thoại Samsung mà các nhân viên bán được trong tháng $9$ (được xếp theo thứ tự không giảm).\\
		Ta có tứ phân vị  thứ ba của mẫu số liệu là $\dfrac{1}{2} (x'_{15}+x'_{16})\in [24{,}5;29{,}5)$.\\
		Tứ phân vị thứ ba của  bảng thống kê doanh số điện thoại bán được của hãng Samsung trong tháng $9$.
		\[Q_{3}=u_{m}+\dfrac{\dfrac{3n}{4}-C}{n_{m}} \cdot  (u_{m+1}-u_{m})=24{,}5+\dfrac{15-(5+8)}{5} \cdot 5=26{,}5.\]
		\begin{itemchoice}
			\itemch  \textbf{Đúng.}\\
			Số điện thoại Oppo trung bình mà một nhân viên cửa hàng  bán được trong tháng $9$
			\[\overline{x}=\dfrac{n_{1}  c_{1}+n_{2} c_{2}+n_{3} c_{3}+n_{4} c_{4}+n_{5} c_{5}}{20}=\dfrac{247}{10}=24{,}7.\]
			Số điện  thoại Samsung trung bình mà một nhân viên cửa hàng bán được trong tháng $9$
			\[\overline{x}=\dfrac{n_{1} c_{1}+n_{2} c_{2}+n_{3} c_{3}+n_{4} c_{4}+n_{5} c_{5}}{20}=\dfrac{465}{20}=23{,}25.\] 
			Dựa vào giá trị trung bình điện thoại bán được, ta thấy điện thoại Oppo dễ bán hơn so với điện thoại  Samsung.
			\itemch \textbf{Sai.}\\ 
			Ta  có
			 \[M_{o}=u_{m}+\dfrac{n_{m}-n_{m-1}}{(n_{m}-n_{m-1})+(n_{m}-n_{m+1})}\cdot(u_{m+1}-u_{m})=19,5+\dfrac{3}{3+3}\cdot5=22.\]
			Vậy khả năng  một nhân viên cửa hàng bán được $22$ chiếc điện thoại Samsung là cao nhất.
			\itemch  \textbf{Sai.}\\
			Vì anh An chỉ bán được $26$ chiếc điện thoại Samsung, không nằm trong Top $5$ nhân viên đạt doanh số  điện thoại Samsung cao nhất
			\itemch  \textbf{Sai.}\\
			Vì chị Bình chỉ bán được $24$ chiếc điện  thoại Oppo, không nằm trong Top 10 nhân viên đạt doanh số điện thoại Oppo cao nhất.
		\end{itemchoice}
	} 
\end{ex}

% nguồn: Ok_dungsai_lop11_Lan2
\begin{ex}%[1D5V2-3]%[ma:6a504e3a]
	Thời gian luyện tập trong một ngày (tính theo  giờ) của một số vận động viên được ghi lại ở bảng sau
	\begin{center}
		\begin{tabular}{|c|c|c|c|c|c|}
			 \hline
			Thời gian (giờ) & $[0; 2)$ & $[2; 4)$ & $[4; 6)$ & $[6; 8)$ &  $[8; 10)$ \\
			\hline 
			Số học sinh & $3$ & $8$ & $12$  & $12$ & $4$ \\
			 \hline
		\end{tabular}
	\end{center}
	 \choiceTF
		{\True Số vận đuộng viên được khảo sát  là $39$}
		{Tứ phân vị thứ nhất là $x_{10}$ thuộc  nhóm $[0;2)$}
		{\True Tứ phân vị thứ hai là $x_{20}$ thuộc  nhóm $[4;6)$}
		{Tứ phân vị thứ  ba là $x_{30}$ thuộc nhóm $[8;10)$}
	\loigiai{ 
		Gọi $x_{1}$, $x_{2}$, $\ldots$, $x_{39}$  là thời gian luyện tập của $39$ vận động viên được xếp theo thứ tự không giảm.
		\begin{itemchoice}
			\itemch \textbf{Đúng.}\\ 
			Số vận động viên  được khảo sát là $n=3+8+12+12+4=39$.
			 \itemch \textbf{Sai.}\\\
			Tứ phân vị thứ nhất là $x_{10}$$ $thuộc nhóm$[2;4)$. 
			\itemch \textbf{Đúng.}\\ 
			Tứ phân vị thứ hai là $x_{20}$ thuộc  nhóm $[4;6)$.
			\itemch  \textbf{Sai.}\\
			Tứ phân vị thứ ba là $x_{30}$  thuộc nhóm $[6;8)$.
		\end{itemchoice}
	} 
\end{ex}

% nguồn: Ok_dungsai_lop11_Lan2
\begin{ex}%[1D5V2-3]%[ma:1f75a010]
	Cân nặng của một số lợn con mới sinh thuộc  hai giống $A$ và $B$ được cho ở bảng đây (đơn vị: kg)
	\begin{center}\begin{tabular}{|c|c|c|c|c|}
		\hline 
		Cân nặng  (kg) &$[1{,}0;1{,}1)$&$[1{,}1;1{,}2)$&$[1{,}2;1{,}3)$&$[1{,}3;1{,}4)$\\
		\hline 
		Số con  giống A&$8$&$28$&$32$&$17$\\
		\hline 
		Số  con giống B&$13$&$14$&$24$&$14$\\
		\hline 
	\end{tabular}
\end{center}
\choiceTF 
	{\True Cân nặng trung bình  của giống $A$ là $1{,}22$}
	{\True Cân nặng trung  bình của giống $B$ là $1{,}21$}
	{\True Tứ phân vị thứ nhất của mẫu số  liệu lợn con giống $A$ là $Q_{1A}=1{,}15$}
	{Tứ phân vị thứ nhất của mẫu  số liệu lợn con giống $B$ là $Q_{1B}=1{,}62$}
\loigiai{ 
	\begin{itemchoice}
		\itemch Cân nặng trung bình của giống  $A$ là
		$\dfrac{1{,}05\cdot8+1{,}15\cdot28+1{,}25\cdot32+1{,}35\cdot17}{85}=1{,}22. $ 
		 \itemch Cân nặng trung bình của giống $B$ là
		$\dfrac{1{,}05\cdot13+1{,}15\cdot14+1{,}25\cdot24+1{,}35\cdot14}{65}=1{,}21. $\\ 
		\itemch Tứ phân vị thứ nhất của dãy  số liệu lợn con giống $A$ là $\dfrac{1}{2}\left(x_{21}+x_{22}\right)\in \left[1{,}1;1{,}2\right)$ nên tứ phân vị thứ nhất của mẫu số liệu lợn con giống $A$ là $Q_{1A}=1{,}1+\dfrac{\dfrac{85}{4}-8}{28}\left(1{,}2-1{,}1\right)=1{,}15. $
		 Tứ phân vị thứ ba của dãy số liệu lợn con giống $A$ là $\dfrac{1}{2}\left(x_{64}+x_{65}\right)\in \left[1{,}2;1{,}3\right)$ nên tứ phân vị thứ ba của mẫu số liệu lợn con giống $A$ là $Q_{3A}=1{,}2+\dfrac{\dfrac{3\cdot85}{4}-36}{32}\left(1{,}3-1{,}2\right)=1{,}29. $
		\itemch Tứ phân vị thứ nhất của dãy số liệu lợn con giống $B$ là $\dfrac{1}{2}\left(y_{16}+y_{17}\right)\in \left[1{,}1;1{,}2\right)$ nên tứ phân vị thứ nhất của mẫu số liệu lợn con giống $B$ là $Q_{1B}=1{,}1+\dfrac{\dfrac{65}{4}-13}{14}\left(1{,}2-1{,}1\right)=1{,}12$.\\ 
		Tứ phân vị thứ ba của dãy số liệu lợn con giống $B$ là $\dfrac{1}{2}\left(y_{49}+y_{50}\right)\in \left[1{,}2;1{,}3\right)$ nên tứ phân vị thứ  ba của mẫu số liệu lợn con giống $B$ là $Q_{3B}=1{,}2+\dfrac{\dfrac{3\cdot65}{4}-27}{24}\left(1{,}3-1{,}2\right)=1{,}29.$
	\end{itemchoice}
} 
\end{ex}

% nguồn: Ok_dungsai_lop11_Lan2
\begin{ex}%[1D5V2-3]%[ma:3d48ebf3]
	Điều  tra về chiều cao của $100$ học sinh khối lớp $11$, ta có kết quả sau
	\begin{center}
		\begin{tabular}{|c|c|c|c|c|c|c|}
			\hline 
			Chiều cao (cm) & $[150; 152)$ & $[152; 154)$ & $[154; 156)$ &  $[156; 158)$ & $[158; 160)$ & $[160; 162)$ \\ \hline
			Số học sinh & $5$  & $18$ & $40$ & $26$ & $8$ & $3$ \\ \hline
		\end{tabular}
	\end{center}
	 \choiceTF
		{\True Giá trị đại diện  của nhóm thứ tư là $157$}
		 {Chiều cao trung bình của mẫu số liệu là $155$}
		{Tứ phân vị thứ ba của mẫu số liệu là $156$} 
		 {Mốt của mẫu số liệu là $154{,}61$}
	\loigiai{ 
		Ta có bảng giá  trị đại diện theo tần số ghép nhóm như sau
		\begin{center}
			\begin{tabular}{|c|c|c|c|c|c|c|}
				 \hline
				Chiều cao (cm) & $[150; 152)$ & $[152; 154)$  & $[154; 156)$ & $[156; 158)$ & $[158; 160)$ & $[160; 162)$ \\ \hline
				Số học sinh & $5$ & $18$  & $40$ & $26$ & $8$ & $3$ \\ \hline
				Giá trị đại diện  & $151$ & $153$ & $155$ & $157$ & $159$ & $161$ \\
				 \hline
			\end{tabular}
		\end{center}
		\begin{itemchoice}
			 \itemch \textbf{Đúng.}\\
			Giá trị đại diện của nhóm  thứ tư là $157$.
			\itemch \textbf{Sai.}\\ 
			 Chiều cao trung bình của mẫu số liệu là
			\[\overline{x}=\dfrac{5 \cdot 151+18 \cdot 153+40 \cdot  155+26 \cdot 157+8 \cdot 159+3 \cdot 161}{100}=155{,}46.\]
			\itemch  \textbf{Sai.}\\
			Cỡ mẫu là $n=100$. Gọi $x_{1}$, $x_{2}$, $\ldots$, $x_{100}$ là chiều cao của $100$ học sinh,  giả sử dãy này đã được sắp xếp theo thứ tự tăng dần. Tứ phân vị thứ ba $Q_{3}$ là $\dfrac{x_{75}+x_{76}}{2}$. Do $x_{75} ,x_{76}$ đều thuộc nhóm $[156;158)$ nên nhóm này chứa $Q_{3}$.\\
			Do  đó $p=4$; $a_{4}=156$; $m_{1}+m_{2}+m_{3}=63$; $a_{5}-a_{4}=2$.\\
			Ta có 
			\[Q_{3}=a_{4}+\dfrac{\dfrac{3n}{4}-(m_{1}+m_{2}+m_{3})}{m_{4}} (a_{5}-a_{4})=156+\dfrac{\dfrac{3 \cdot 100}{4}-63}{26} \cdot 2 \approx 156{,}9.\] 
			 \itemch \textbf{Sai}.\\
			Mốt của mẫu số liệu. 
			Nhóm $[154;156)$  là nhóm có tần số lớn nhất, nên nhóm chứa mốt là nhóm $[154;156)$.\\
			Do đó $a_{j}=154$,  $m_{j}=40$, $m_{j-1}=18$, $m_{j+1}=26$, $h=2$. Ta có
			\[M_{o}=a_{j}+\dfrac{m_{j}-m_{j-1}}{m_{j}-m_{j-1}+m_{j}+m_{j+1}}  \cdot h=154+\dfrac{40-18}{40-18+40-26} \cdot 2\approx 155{,}22.\]
		\end{itemchoice}
	 }
\end{ex}

% nguồn: Ok_dungsai_lop11_Lan2
\begin{ex}%[1D5V2-3]%[ma:eca571b8]
	 Điều tra về chiều cao của học sinh khối lớp $10$, ta có kết quả sau
	\begin{center}
		\begin{tabular}{|c|c|}
			\hline 
			Nhóm chiều cao (cm) & Số  học sinh \\ \hline
			 $[150;152)$ & $5$ \\ \hline
			$[152;154)$ & $18$ \\  \hline
			$[154;156)$ &  $40$ \\ \hline
			 $[156;158)$ & $26$ \\ \hline
			$[158;160)$ &  $8$ \\ \hline
			$[160;162)$ &  $3$ \\ \hline
		\end{tabular}
	\end{center}
	\choiceTF 
		{\True Tổng  số học sinh điều tra bằng $100$}
		{\True Chiều  cao trung bình của các em học sinh trên là $155{,}46$}
		{Số trung vị của  mẫu số liệu trên bằng $156{,}35$}
		 {\True Tứ phân vị thứ nhất của mẫu số liệu trên bằng $154{,}1$}
	 \loigiai{
		Bảng  tần số ghép nhóm theo giá trị đại diện là
		\begin{center}
			\begin{tabular}{|c|c|c|}
				 \hline
				Nhóm chiều cao (cm) & Giá trị đại diện & Số học sinh  \\ \hline
				$[150;152)$  & $151$ & $5$ \\ \hline
				$[152;154)$ & $153$ & $18$ \\ \hline 
				$[154;156)$ & $155$  & $40$ \\ \hline
				 $[156;158)$ & $157$ & $26$ \\ \hline
				$[158;160)$ & $159$ & $8$ \\  \hline
				$[160;162)$ &  $161$ & $3$ \\ \hline
			\end{tabular}
		\end{center}
		\begin{itemchoice}
			\itemch \textbf{Đúng.}\\ 
			Ta có tổng số học sinh là $n=100$. 
			\itemch \textbf{Đúng.}\\ 
			Chiều cao  trung bình bằng
			\[\overline{x}=\dfrac{5\cdot 151+18\cdot 153+40\cdot  155+26\cdot 157+8\cdot 159+3\cdot 161}{100}=155{,}46.\]
			 \itemch \textbf{Sai.}\\
			Do $n=100$ nên trung vị là $\dfrac{x_{50}+x_{51}}{2}$,  do hai giá trị $x_{50}$, $x_{51}$ thuộc nhóm $3$. Ta có
			\[M_e=154+\dfrac{50-(18+5)}{40}\cdot 2=155{,}35.\] 
			 \itemch \textbf{Đúng.}\\
			Tứ phân vị thứ nhất $Q_3$ là $\dfrac{x_{25}+x_{26}}{2}$ mà $x_{25}$, $x_{26}$  thuộc nhóm $[154;156)$.\\
			Do đó, $p=3$, $a_p=154$,  $m_p=40$, $m_1+m_2=5+18=23$, $a_{p+1}-a_p=2$ và ta có
			\[Q_1=154+\dfrac{\dfrac{100}{4}-23}{40}\cdot 2=154{,}1.\] 
		\end{itemchoice}
	 }
\end{ex}

% nguồn: Ok_dungsai_lop11_Lan2
\begin{ex}%[1D5V2-3]%[ma:f218d1cf]
	Thống kê số trẻ em dưới $10$  tuổi trong một khu dân cư ta được bẳng số liệu như sau
	\begin{center}
		\begin{tabular}{|c|c|c|c|c|c|}
			\hline 
			Tuổi & $[0; 2)$ & $[2; 4)$ & $[4; 6)$ & $[6; 8)$ & $[8; 10)$ \\ 
			\hline 
			 Số trẻ & $3$ & $8$ & $12$ & $12$ & $4$ \\
			 \hline
		\end{tabular}
	\end{center}
	 \choiceTF
		{\True Cỡ  mẫu của bảng số liệu trên là $39$}
		{$[2;4)$ là nhóm  chứa mốt}
		 {\True Tuổi trung bình trẻ em được thống kê là $\overline{x}\approx 5{,}31$}
		{\True Tứ phân vị  thứ nhất của mẫu số liệu ghép nhóm trên là $Q_1=3{,}6875$}
	 \loigiai{
		 Bảng tần số ghép nhóm theo giá trị đại diện là
		\begin{center}
			\begin{tabular}{|c|c|c|c|c|c|}
				 \hline
				 Nhóm & $[0;2)$ & $[2;4)$ & $[4;6)$ & $[6;8)$ & $[8;10)$ \\ \hline
				Giá  trị đại diện & $1$ & $3$ & $5$ & $7$ & $9$ \\ \hline
				Tần số & $3$ & $8$ &  $12$ & $12$ & $4$ \\ \hline
			\end{tabular}
		\end{center}
		\begin{itemchoice}
			\itemch  \textbf{Đúng.}\\
			Cỡ mẫu của bảng số  liệu trên là $n=3+8+12+12+4=39$.
			 \itemch \textbf{Sai.}\\
			Nhóm chứa mốt là nhóm có tần số lớn  nhất.
			Nhóm $[2;4)$ có tần số bằng $8$; không phải là  tần số lớn nhất.
			\itemch  \textbf{Đúng.}\\
			Chiều cao trung  bình học sinh nam trong lớp
			 \[\overline{x}=\dfrac{1\cdot 3+3\cdot 8+5\cdot 12+7\cdot 12+9\cdot 4}{39} \approx 5{,}31\]
			 \itemch \textbf{Đúng.}\\
			Gọi $x_{1}$; $x_{2}$; $\ldots$; $x_{39}$ là  số tuổi của $39$ trẻ em sắp xếp theo thứ tự không giảm.\\
			Tứ phân vị thứ nhất của  dãy mẫu số liệu là $x_{10} \in [2;4)$. Xác định được $n=39$; $n_{m}=12$; $C=3$; $u_{m}=2$; $u_{m+1}=4$.\\
			Tứ phân vị  thứ nhất của mẫu số liệu ghép nhóm trên là
			\[Q_{1}=u_{m}+\dfrac{\dfrac{n}{4}-C}{n_{m}}\cdot(u_{m+1}-u_{m})=2+\dfrac{\dfrac{1\cdot39}{4}-3}{8}  \cdot(4-2)=3{,}6875.\]
		\end{itemchoice}
	} 
\end{ex}

% nguồn: Ok_dungsai_lop11_Lan2
\begin{ex}%[1D5V2-3]%[ma:86734ab4]
	Doanh thu bán hàng trong $20$ ngày được lựa chọn ngẫu nhiên của một của hàng được ghi lại ở bảng sau (đơn vị: triệu  đồng)
	\begin{center}
		\begin{tabular}{|c|c|c|c|c|c|}
			 \hline
			Doanh thu & $[5; 7)$ & $[7;  9)$ & $[9; 11)$ & $[11; 13)$ & $[13; 15)$ \\
			 \hline
			Số ngày & $2$ & $7$ & $7$ & $3$ & $1$ \\ 
			\hline 
		\end{tabular}
	\end{center}
	 \choiceTF
		{Số ngày doanh thu bán hàng ít hơn $9$ triệu là  $7$ ngày}
		{\True Số ngày doanh thu vượt quá $11$ triệu là  $4$ ngày}
		{\True Doanh thu trung bình trong  $20$ ngày bán hàng là $9{,}4$ triệu}
		{Tứ phân vị thứ  nhất của mẫu số liệu (làm tròn đến hàng phần trăm) là $7{,}87$}
	 \loigiai{
		 Bảng tần số ghép nhóm theo giá trị đại diện là
		\begin{center}
			\begin{tabular}{|c|c|c|c|c|c|}
				\hline 
				Doanh thu & $[5; 7)$ & $[7; 9)$ & $[9; 11)$ & $[11;  13)$ & $[13; 15)$ \\
				 \hline
				Giá trị đại diện & $6$ & $8$ & $10$ & $12$ &  $14$ \\
				 \hline
				Số ngày & $2$ & $7$ & $7$ & $3$ & $1$  \\
				 \hline
			\end{tabular}
		\end{center}
		\begin{itemchoice}
			\itemch \textbf{Sai.}\\ 
			Số ngày doanh  thu bán hàng ít hơn $9$ triệu là $9$ ngày.
			\itemch  \textbf{Đúng.}\\
			Số ngày  doanh thu vượt quá $11$ triệu là $4$ ngày.
			 \itemch \textbf{Đúng.}\\
			Số  trung bình
			\[\overline{x}=\dfrac{2\cdot 6+7\cdot 8+7\cdot 10+3\cdot  12+1\cdot 14}{20}=9{,}4.\]
			\itemch  \textbf{Sai.}\\
			Gọi $x_{1}$; $x_{2}$; $\ldots$; $x_{20}$ là doanh thu bán hàng trong $20$ ngày xếp theo thứ  tự không giảm.\\
			Tứ phân vị thứ nhất của mẫu số liệu  thuộc nhóm $[7;9)$. Ta có $n=20$, $n_{m}=7$, $C=2$, $u_{m}=7$, $u_{m+1}=9$, khi đó
			 \[Q_{1}=7+\dfrac{\dfrac{1\cdot20}{4}-2}{7}\cdot(9-7)\approx 7{,}86\]
		\end{itemchoice}
	} 
\end{ex}

% nguồn: Ok_dungsai_lop11_Lan2
\begin{ex}%[1D5V2-3]%[ma:cd8c91c8]
	Khảo sát thời gian tập thể dục của một số học sinh khối $11$  thu được mẫu số liệu ghép nhóm sau
	\begin{center}
		\begin{tabular}{|c|c|c|c|c|c|}
			\hline 
			Thời gian (phút) & $[0; 20)$ & $[20; 40)$ &  $[40; 60)$ & $[60; 80)$ & $[80; 100)$ \\
			\hline 
			Số học sinh  & $5$ & $9$ & $12$ & $10$ & $6$ \\
			 \hline
		\end{tabular}
	\end{center}
	 \choiceTF
		{\True Nhóm chứa mốt của mẫu số  liệu trên là $[40;60)$}
		{Giá trị đại diện của nhóm $[80;100)$  là $80$}
		{\True Mốt của mẫu số  liệu trên là $52$}
		{Nhóm chứa tứ phân  vị thứ nhất của mẫu số liệu trên là $[40;60)$}
	 \loigiai{
		\begin{itemchoice}
			\itemch  \textbf{Đúng.}\\
			Vì nhóm $[40;60)$ có số học sinh bằng $12$ có tần số  lớn nhất trong các nhóm.
			 \itemch \textbf{Sai.}\\
			Giá  trị đại diện của nhóm $[80;100)$là $\dfrac{80+100}{2}=90$.
			\itemch  \textbf{Đúng.}\\
			Mốt $M_{o}$ chứa trong nhóm $[40;60)$. Do đó $u_{m}=40$; $u_{m+1}=60$ suy ra $u_{m+1}-u_{m}=60-40=20$, $n_{m-1}=9$; $n_{m}=12$; $n_{m+1}=10$.  Ta có
			 \[M_{o}=40+\dfrac{12-9}{(12-9)+(12-10)}\cdot(60-40)=52.\]
			 \itemch \textbf{Sai.}\\
			Ta có $n=42$. Nên tứ phân vị thứ nhất của mẫu số liệu trên là $Q_{1}=x_{11}$. Mà $x_{11} \in  [20;40)$.
			Vậy nhóm chứa tứ phân vị thứ nhất của mẫu số  liệu trên là nhóm $[20;40)$.
		\end{itemchoice}
	} 
\end{ex}

% nguồn: Ok_dungsai_lop11_Lan2
\begin{ex}%[1D5V2-3]%[ma:88c558e5]
	Số liệu  thống kê chiều cao của $40$ mẫu cây ở một vườn thực vật (đơn vị: cm) được cho trong bảng sau
	\begin{center}
		\begin{tabular}{|c|c|c|c|c|c|c|}
			 \hline
			Chiều cao (cm) &  $[30;40)$ & $[40;50)$ & $[50;60)$ & $[60;70)$ & $[70;80)$ & $[80;90)$ \\
			 \hline
			Tần số tích lũy & $4$ & $14$ & $28$ & $34$  & $38$ & $40$ \\
			\hline 
		\end{tabular}
	\end{center}
	\choiceTF 
		{\True  Tần số của nhóm chiều cao $[50;60)$ là $14$}
		 {\True Chiều cao trung bình của mẫu số liệu là $55{,}5$}
		{\True Trung  vị của mẫu số liệu ghép nhóm là $54{,}286$}
		{Tứ phân vị $Q_2$ của mẫu số liệu  trên là $66{,}4$}
	 \loigiai{
		Bảng thống kê tần số tích lũy  suy ra cỡ mẫu là $n=40$.
		\begin{center}
			\begin{tabular}{|c|c|c|c|c|c|c|} \hline
				Chiều cao (cm) & $[30;40)$ & $[40;50)$ & $[50;60)$ & $[60;70)$  & $[70;80)$ & $[80;90)$ \\ \hline
				Giá trị đại diện & $35$ & $45$ & $55$ & $65$ & $75$ & $85$ \\  \hline
				 Tần số tích lũy & $4$ & $14$ & $28$ & $34$ & $38$ & $40$ \\ \hline
				Tần số  & $4$ & $10$ & $14$ & $6$ & $4$ & $2$ \\ \hline
			\end{tabular}
		\end{center}
		\begin{itemchoice}
			\itemch  \textbf{Đúng.}\\
			Tần số của nhóm chiều  cao $[50;60)$ là $14$.
			\itemch \textbf{Đúng.}\\ 
			 Số trung bình của bảng là
			\[\overline{x}=\dfrac{35\cdot 4+45\cdot 10+55\cdot 14+65\cdot  6+75\cdot 4+85\cdot 2}{40}=\dfrac{2220}{40}=55{,}5.\]
			\itemch \textbf{Đúng.}\\ 
			Gọi $x_{1}$; $x_{2}$; $\ldots$; $x_{40}$ là chiều cao của $40$ mẫu cây và giả sử dãy  này đã được sắp xếp theo thứ tự tăng dần. Khi đó, trung vị là $\dfrac{x_{20}+x_{21}}{2}$. Do hai giá trị $x_{20} ,x_{21}$ thuộc nhóm $[50;60)$ nên nhóm này chứa trung vị. Do đó, $p=3$, $a_{3}=50$, $m_{3}=14$, $m_{1}+m_{2}=4+10=14$; $a_{4}-a_{3}=60-50=10$ và ta có:
			\[M_{e}=a_{3}+\dfrac{\dfrac{n}{2}-(m_{1}+m_{2})}{m_{3}}\cdot(a_{4}-a_{3})=50+\dfrac{20-14}{14}\cdot10=54{,}286.\] 
			 \itemch \textbf{Sai.}\\
			Do  $Q_2=M_e=54{,}285=6$.
		\end{itemchoice}
	} 
\end{ex}

% nguồn: Ok_dungsai_lop11_Lan2
\begin{ex}%[1D5V2-3]%[ma:12653167]
	Phỏng vấn một số học sinh khối $11$ về  thời gian (giờ) ngủ của một buổi tối, thu được bảng số liệu sau
	\begin{center}
		\begin{tabular}{|c|c|c|}
			\hline 
			Thời gian & Số học sinh nam & Số học  sinh nữ \\
			 \hline
			$[4;5)$ & $8$ & $4$ \\ \hline 
			$[5;6)$ & $10$ &  $6$ \\ \hline
			$[6;7)$ & $13$ &  $12$ \\ \hline
			 $[7;8)$ & $15$ & $18$ \\ \hline
			$[8;9)$ & $7$ & $8$ \\  \hline
		\end{tabular}
	\end{center}
	\choiceTF 
		 {\True Độ dài của nhóm bằng $1$}
		{Thời gian ngủ trung bình của các bạn học sinh nam nhiều hơn các bạn học sinh  nữ}
		{\True Phần lớn học sinh được khảo sát trong khối $11$ ngủ nhiều hơn $6{,}5$  giờ}
		{$75\%$ học sinh được khảo sát trong khối $11$ ngủ  ít nhất $5{,}5$ giờ}
	\loigiai{ 
		\begin{itemchoice}
			\itemch  \textbf{Đúng.}\\
			Độ dài của nhóm là $5-4=1$. 
			\itemch  \textbf{Sai.}\\
			Từ  bảng thống kê ta có
			\begin{center}
				\begin{tabular}{|c|c|c|c|}
					\hline 
					Thời gian & Giá  trị đại diện & Số học sinh nam & Số học sinh nữ \\
					 \hline
					$[4;5)$ & $4{,}5$ & $8$ &  $4$ \\ \hline
					$[5;6)$ & $5{,}5$ & $10$ &  $6$ \\ \hline
					$[6;7)$ & $6{,}5$  & $13$ & $12$ \\ \hline
					$[7;8)$ & $7{,}5$ & $15$ & $18$  \\ \hline
					$[8;9)$ & $8{,}5$ & $7$  & $8$ \\ \hline
				\end{tabular}
			\end{center}
			Thời gian ngủ trung bình của các bạn nam  là
			\[\overline{x}_\text{nam}=\dfrac{4{,}5\cdot  8+5{,}5\cdot 10+6{,}5\cdot 13+7{,}5\cdot 15+8{,}5\cdot 7}{8+10+13+15+7} \approx 6{,}56.\]
			Thời gian ngủ trung bình của các bạn  nữ là
			\[\overline{x}_\text{nữ}=\dfrac{4{,}5\cdot 4+5{,}5\cdot  6+6{,}5\cdot 12+7{,}5\cdot 18+8{,}5\cdot 8}{4+6+12+18+8} \approx 6{,}92\]
			Do $6{,}56<6{,}92$ như vậy thời gian  ngủ trung bình của các bạn nữ nhiều hơn các bạn nam.
			 \itemch \textbf{Đúng.}\\
			Cỡ mẫu $n=101$. Trung vị của mẫu số liệu  là $x_{51}$ thuộc nhóm $[6;7)$.\\
			Do đó, $p=3$, $a_{3}=6$, $a_{2}=5$, $m_{2}=16$, $m_{1}=12$, $a_{3}-a_{2}=1$.  Ta có trung vị của mẫu số liệu là
			\[M_{e}=6+\dfrac{\dfrac{101}{2}-(12+16)}{25} \cdot  (6-5)=6{,}9.\]
			Vậy $50\%$ học sinh có thời gian ngủ nhiều hơn $6{,}9$  giờ.
			\itemch  \textbf{Sai.}\\
			Cỡ mẫu $n=101$. Tứ phân vị thứ nhất $Q_{1}$ là $\dfrac{x_{25}+x_{26}}{2}$. Do $x_{25}$, $x_{26}$ đều thuộc nhóm $[5;6)$ nên nhóm này chứa $Q_{1}$. Do đó, $p=2$, $a_{2}=5$, $m_{2}=16$, $m_{1}=12$, $a_{3}-a_{2}=1$. Ta có tứ  phân vị thứ nhất của mẫu số liệu là
			\[Q_{1}=5+\dfrac{\dfrac{101}{4}-12}{16} \cdot 1\approx  5{,}83.\]
			Vậy $75\%$ học sinh khối  $11$ ngủ ít nhất $5{,}83$ giờ.
		\end{itemchoice}
	 }
\end{ex}

% nguồn: Ok_dungsai_lop11_Lan2
\begin{ex}%[1D5V2-3]%[ma:a077ffd9]
	Kết quả điều tra về số giờ làm thêm trong một tuần của sinh  viên một trường đại học X được cho bởi bảng sau
	\begin{center}
		\begin{tabular}{|c|c|c|c|c|c|}
			 \hline
			Số giờ làm thêm & $[2; 4)$ &  $[4; 6)$ & $[6; 8)$ & $[8; 10)$ & $[10; 12)$ \\ \hline
			Số sinh viên & $12$ & $20$ & $37$ & $21$ &  $10$ \\ \hline
		\end{tabular}
	\end{center}
	\choiceTF 
		{\True Số  sinh viên được điều tra là $100$}
		{\True Số giờ làm thêm trung bình của mỗi sinh viên trường  đại học X không ít hơn $6$}
		{Mốt của mẫu số  liệu trên là $7{,}5$}
		{\True Tứ phân vị thứ hai của dãy số liệu  lớn hơn $6{,}5$}
	\loigiai{ 
		\begin{itemchoice}
			\itemch \textbf{Đúng.}\\ 
			 Cỡ mẫu $n=100$.
			 \itemch \textbf{Đúng.}\\
			Từ bảng thống kê ta  có
			\begin{center}
				\begin{tabular}{|c|c|c|c|c|c|}
					 \hline
					Số giờ làm thêm & $[2;4)$ & $[4;6)$ & $[6;8)$  & $[8;10)$ & $[10;12)$ \\ \hline
					Số giờ làm thêm  đại diện & $3$ & $5$ & $7$ & $9$ & $11$ \\ \hline
					Số sinh viên & $12$ & $20$ & $37$ & $21$ & $10$ \\  \hline
				\end{tabular}
			\end{center}
			Số trung bình  của mẫu số liệu là
			\[\overline{x}=\dfrac{3 \cdot 12+5 \cdot 20+7 \cdot 37+9 \cdot  21+11 \cdot 10}{100}=6{,}94.\]
			\itemch  \textbf{Sai.}\\
			Nhóm  chứa mốt số liệu trên là nhóm $[6;8)$. Do đó $u_{m}=6$; $n_{m}=37$; $n_{m-1}=20$; $n_{m+1}=21$; $u_{m+1}=8$.\\
			Vậy mốt của mẫu  số liệu ghép nhóm là
			\[M_{o}=6+\dfrac{37-20}{(37-20)+(37-21)} \cdot(8-6)\approx 7{,}03.\] 
			 \itemch \textbf{Đúng.}\\
			Gọi $x_{1}$; $x_{2}$; $\ldots$; $x_{100}$ là mẫu số liệu theo  thứ tự không giảm.
			Tứ phân vị thứ hai của  dãy số liệu $x_{1}$; $x_{2}$; $\ldots$;$x_{100}$ là $\dfrac{1}{2} (x_{50}+x_{51}).$\\
			Do $x_{50}$  và $x_{51}$ thuộc nhóm $[6;8)$ nên $Q_{2}=6+\dfrac{\dfrac{2.100}{4}-32}{37} \cdot(8-6)\approx 6{,}97.$
		\end{itemchoice}
	} 
\end{ex}

% nguồn: Ok_dungsai_lop11_Lan2
\begin{ex}%[1D5V2-3]%[ma:fa456cff]
	Thống kê cân nặng của $40$ học sinh lớp 11A trong  một trường trung học phổ thông (đơn vị: kilôgam) được cho ở bảng sau
	\begin{center}
		\begin{tabular}{|c|c|c|c|c|c|}
			\hline 
			Nhóm & $[30; 40)$ & $[40; 50)$ & $[50; 60)$ &  $[60; 70)$ & $[70; 80)$ \\
			 \hline
			Tần số  & $4$ & $10$ & $18$ & $10$ & $2$ \\
			 \hline
		\end{tabular}
	\end{center}
	\choiceTF 
		{Số học sinh có số cân nặng  từ $40$ kg đến dưới $50$ kg là $18$}
		{\True Nhóm $[50;60)$ là nhóm chứa  mốt}
		{\True Cân nặng  trung bình của các bạn học sinh trong lớp 11A là $54$ kg}
		{Tứ phân vị thứ nhất của mẫu số liệu ghép nhóm trên là  $50$ kg}
	 \loigiai{
		Bảng tần số ghép nhóm theo giá trị đại  diện là
		\begin{center}
			\begin{tabular}{|c|c|c|c|c|c|}
				\hline 
				Nhóm & $[30;  40)$ & $[40; 50)$ & $[50; 60)$ & $[60; 70)$ & $[70; 80)$ \\
				\hline 
				Giá trị đại diện & $35$ &  $45$ & $55$ & $65$ & $75$ \\
				 \hline
				Tần số & $4$ & $10$ & $18$ & $10$ &  $2$ \\
				\hline 
			\end{tabular}
		\end{center}
		\begin{itemchoice}
			\itemch  \textbf{Sai.}\\
			Số học sinh có số  cân nặng từ 40 kg đến dưới 50 kg là 10.
			\itemch \textbf{Đúng.}\\ 
			Nhóm chứa mốt là nhóm có tần số lớn nhất nên $[50;60)$là nhóm chứa  mốt.
			\itemch \textbf{Đúng.}\\ 
			Cân nặng trung bình  của các bạn học sinh trong lớp 11A là
			\[\dfrac{4\cdot 35+10\cdot 45+18\cdot 55+10\cdot 65+2\cdot 75}{44} \approx  54.\]
			 \itemch \textbf{Sai.}\\
			Số phần tử của mẫu là $n=44$. Ta có $\dfrac{n}{4}=\dfrac{44}{4}=11$ mà $14>11$. Suy ra nhóm $2$ là nhóm đầu tiên có tần số tích lũy lớn hơn hoặc bằng  $11$.\\
			Xét nhóm $2$ là nhóm $[40;50)$ có $s=40$; $h=10$; $n_{2}=10$ và nhóm $1$ là  nhóm có $cf_{1}=4$.\\
			 Áp dụng công thức, ta có tứ phân vị thứ nhất là
			 \[Q_{1}=40+\dfrac{11-4}{10}\cdot10=47.\]
		\end{itemchoice}
	 }
\end{ex}

% nguồn: Ok_dungsai_lop11_Lan2
\begin{ex}%[1D5V2-3]%[ma:4f6a1491]
	Thống kê điểm trung  bình môn Toán của một số hoc sinh lớp $11$ được cho ở bảng sau
	\begin{center}
		\begin{tabular}{|c|c|c|c|c|c|c|c|}
			 \hline
			Khoảng điểm & $[6{,}5; 7)$ & $[7; 7{,}5)$ & $[7{,}5; 8)$ & $[8; 8{,}5)$ & $[8{,}5; 9)$ & $[9; 9{,}5)$ & $[9{,}5; 10)$  \\
			 \hline
			Tần số & $8$ &  $10$ & $16$ & $24$ & $13$ & $7$ & $4$ \\
			 \hline
		\end{tabular}
	\end{center}
	 \choiceTF
		{\True Điểm trung bình  môn Toán là $8{,}12$}
		{Tứ phân vị thứ  nhất bằng $8{,}15$}
		{Tứ phân vị  thứ hai bằng $7{,}58$}
		 {\True Mốt của mẫu số liệu bằng $8{,}21$}
	 \loigiai{
		Bảng tần số ghép nhóm theo  giá trị đại diện là
		\begin{center}
			\begin{tabular}{|c|c|c|c|c|c|c|c|}
				 \hline
				Khoảng điểm & $[6{,}5; 7)$ & $[7; 7{,}5)$ & $[7{,}5; 8)$ & $[8; 8{,}5)$ & $[8{,}5; 9)$ &  $[9; 9{,}5)$ & $[9{,}5; 10)$ \\
				\hline 
				Giá trị đại diện & $6{,}75$ & $7{,}25$ &  $7{,}75$ & $8{,}25$ & $8{,}75$ & $9{,}25$ & $9{,}75$ \\
				 \hline
				Tần số & $8$ &  $10$ & $16$ & $24$ & $13$ & $7$ & $4$ \\
				\hline 
			\end{tabular}
		\end{center}
		\begin{itemchoice}
			\itemch  \textbf{Đúng.}\\
			Điểm trung bình môn Toán của một số  học sinh lớp $11$ là
			\[\overline{x}=\dfrac{8\cdot 6{,}75+10\cdot 7{,}25+16\cdot 7{,}75+24\cdot 8{,}25+13\cdot 8{,}75+7\cdot 9{,}25+4\cdot 9{,}75}{82} \approx  8{,}12.\]
			 \itemch \textbf{Sai.}\\
			Tứ phân vị thứ nhất thuộc  nhóm $[7{,}5;8)$, ta có
			\[Q_{1}=7{,}5+\dfrac{\dfrac{2\cdot 82}{4}-(8+10)}{16}\cdot(8-7;5)\approx 7{,}58.\] 
			\itemch \textbf{Sai.}\\ 
			Tứ  phân vị thứ hai thuộc nhóm $[8;8{,}5)$, ta có
			\[Q_{2}=8+\dfrac{\dfrac{2\cdot82}{4}-(8+10+16)}{24}\cdot(8{,}5-8)\approx  8{,}15.\]
			\itemch  \textbf{Đúng.}\\
			Mốt $M_{o}$ chứa  trong nhóm $[8;8,5)$ và $j=4$. Do đó $u_{4}=8$; $u_{5}=8{,}5$ suy ra $h=u_{5}-u_{4}=8{,}5-8=0{,}5$; $n_{3}=16$;$n_{4}=24$; $n_{5}=13$, ta có
			 \[M_{o}=u_{4}+\dfrac{n_{4}-n_{3}}{(n_{4}-n_{3})+(n_{4}-n_{5})}\cdot h=8+\dfrac{24-16}{(24-16)+(24-13)}\cdot(8{,}5-8)=8{,}21.\]
		\end{itemchoice}
	} 
\end{ex}

% nguồn: Ok_dungsai_lop11_Lan2
\begin{ex}%[1D5V2-3]%[ma:eb7d3877]
	Khảo sát cân nặng trung bình của học sinh lớp 11A cho trong bảng  dưới đây
	\begin{center}
		\begin{tabular}{|c|c|c|c|c|c|c|}
			 \hline
			Cân nặng (kg) & $[40; 45{,}5)$ & $[45{,}5; 50{,}5)$ & $[50{,}5; 55{,}5)$ & $[55{,}5; 60{,}5)$ & $[60{,}5;  65{,}5)$ & $[65{,}5; 70{,}5)$ \\
			 \hline
			Số học sinh & $9$ &  $7$ & $16$ & $5$ & $2$ & $3$ \\
			\hline 
		\end{tabular}
	\end{center}
	\choiceTF 
		{\True Số học sinh được  khảo sát là $42$ học sinh}
		{\True Cân nặng trung bình của học  sinh lớp 11A không ít hơn $52$ kg}
		{Mốt của mẫu số liệu  trên là $51{,}75$}
		{Tứ phân vị  thứ nhất lớn hơn $47$}
	\loigiai{ 
		\begin{itemchoice}
			\itemch \textbf{Đúng.}\\ 
			Cỡ  mẫu là $n=9+7+16+5+2+3=42$.
			 \itemch \textbf{Đúng.}\\
			Cân nặng trung bình (kg) của học sinh lớp  11A là
			\[\overline{x}=\dfrac{9\cdot43+7\cdot48+16\cdot53+5\cdot58+2\cdot63+3\cdot68}{42} \approx 52{,}17.\] 
			\itemch \textbf{Sai.}\\ 
			Tần số lớn nhất là  $16$ nên nhóm chứa mốt là $[50{,}5;55{,}5)$. Ta có
			\[M_{o}=50{,}5+\dfrac{16-7}{(16-7)+(16-5)}\cdot(55{,}5-50{,}5)=52{,}75.\] 
			\itemch \textbf{Sai.}\\ 
			 Tứ phân vị thứ nhất $Q_{1}$ là $x_{11}$ thuộc nhóm $[45{,}5;50{,}5)$. Ta có
			\[Q_{1}=45{,}5+\dfrac{\dfrac{42}{4}-9}{7}\cdot(50{,}5-45{,}5)\approx  46{,}57.\]
		\end{itemchoice}
	} 
\end{ex}

% nguồn: Ok_dungsai_lop11_Lan2
\begin{ex}%[1D5V2-3]%[ma:edda2bfa]
	Một bảng xếp hạng đã tính điểm chuần hoá cho chỉ số nghiên cứu của một số trường đại học ở Việt Nam và thu được  kết quả sau
	\begin{center}
		\begin{tabular}{|c|c|c|c|c|c|c|}
			\hline 
			Điểm & Dưới $20$  & $[20; 30)$ & $[30; 40)$ & $[40; 60)$ & $[60; 80)$ & $[80; 100)$ \\
			 \hline
			Số trường & $4$ &  $19$ & $6$ & $2$ & $3$ & $1$ \\
			 \hline
		\end{tabular}
	\end{center}
	 \choiceTF
		{\True  Cỡ của mẫu là $35$}
		 {Nhóm chứa tứ phân vị thứ $3$ là $[40;60)$}
		{Mốt của mẫu số liệu trên là $54{,}286$} 
		{\True Tứ phân vị $Q_3$ của mẫu số liệu trên là $35{,}42$} 
	\loigiai{ 
		\begin{itemchoice}
			\itemch \textbf{Đúng.}\\ 
			Cỡ mẫu $n=35$. 
			 \itemch \textbf{Sai.}\\
			Tứ phân vị thứ ba $Q_{3}$ là $x_{27}$ thuộc  nhóm $[30;40)$.
			 \itemch \textbf{Sai.}\\
			Tần số lớn nhất là $19$, nên nhóm chứa mốt là nhóm $[20;30)$, ta có $j=2$, $a_{2}=20$,  $m_{2}=19$, $m_{1}=4$, $m_{3}=6$, $h=10$. Do đó
			\[M_{o}=a_{2}+\dfrac{m_{2}-m_{1}}{(m_{2}-m_{1})+(m_{2}-m_{3})}\cdot h=20+\dfrac{19-4}{(19-4)+(19-6)}\cdot  10=25{,}36.\]
			\itemch  \textbf{Đúng.}\\
			Ta có $p=3$, $a_{3}=30$, $m_{3}=6$, $m_{1}+m_{2}=4+19=23$,  $a_{4}-a_{3}=10$
			\[Q_{3}=30+\dfrac{\dfrac{3\cdot35}{4}-23}{6}  \cdot10=35{,}42.\]
		\end{itemchoice}
	} 
\end{ex}

% nguồn: Ok_dungsai_lop11_Lan2
\begin{ex}%[1D5V2-3]%[ma:081a57ab]
	Điều tra  về chiều cao của học sinh khối lớp $10$, ta có kết quả sau
	\begin{center}
		\begin{tabular}{|c|c|c|}
			\hline 
			Nhóm &  Chiều cao (cm) & Số học sinh \\
			 \hline
			 $1$ & $[150;152)$ & $5$ \\
			\hline 
			$2$ & $[152;154)$  & $18$ \\
			\hline 
			$3$ & $[154;156)$  & $40$ \\
			 \hline
			$4$ & $[156;158)$ & $26$ \\ 
			\hline 
			$5$ & $[158;160)$ & $8$ \\ 
			\hline 
			$6$ & $[160;162)$  & $3$ \\
			\hline 
		\end{tabular}
	\end{center}
	\choiceTF 
		{Nhóm chứa $Q_1$ là nhóm $[152;154)$} 
		{Cỡ của mẫu  bằng $40$}
		{\True Nhóm chứa $Q_1$ là nhóm  thứ $4$}
		{\True Tứ phân vị $Q_1$ bằng $63{,}16$} 
	\loigiai{ 
		\begin{itemchoice}
			 \itemch \textbf{Sai.}\\
			Cỡ mẫu $n=100$, $\dfrac{3n}{4}=75$, nhóm $[156;158)$ là nhóm  đầu tiên có tần số tích lũy lớn hơn $75$ nên nhóm này chứa $Q_{1}$.
			\itemch \textbf{Sai.}\\ 
			 Cỡ mẫu $n=100$.
			 \itemch \textbf{Đúng.}\\
			Nhóm thứ $4$ là nhóm  $[156;158)$ nên nhóm này chứa $Q_{1}$.
			\itemch \textbf{Đúng.}\\ 
			Ta có $p=4$, $a_{4}=156$,  $m_{4}=40$, $d=a_{5}-a_{4}=2$, khi đó
			 \[Q_{1}=a_{4}+\dfrac{\dfrac{n}{4}-(m_{1}+m_{2}+m_{3})}{m_{4}}\cdot(a_{5}-a_{4})=60+\dfrac{15-9}{19}\cdot10=63{,}16.\]
		\end{itemchoice}
	 }
\end{ex}


%%%%%%%%%%%%%%%% TRẢ LỜI NGẮN %%%%%%%%%%%%%%%%

% nguồn: Ok_TraLoiNgan_Lop11_Lan1
\begin{ex}%[1D5H2-3]%[ma:329fa40d]
	Tổng hợp tiền lương tháng của một số nhân viên văn phòng được  ghi lại như sau (đơn vị: triệu đồng):
	\begin{center}
		\begin{tabular}{|c|c|c|c|c|}
			\hline 
			Lương tháng  (triệu đồng) & $[6;8)$ & $[8;10)$ & $[10;12)$ & $[12;14)$ \\
			 \hline
			Số nhân viên & $3$ & $6$ & $8$ & $7$ \\ 
			\hline 
		\end{tabular}
	\end{center}
	Tứ phân vị  thứ nhất của mẫu số liệu ở bảng tần số ghép nhóm trên bằng bao nhiêu?
	\shortans{9} 
	\loigiai{ 
		Gọi $x_1;x_2;\ldots;x_{24}$  lần lượt là lương tháng của mỗi nhân viên được xếp theo thứ tự không giảm.\\
		Do $x_1,\ldots,x_3\in [6;8); x_4,\ldots,x_9\in  [8;10); x_{10},\ldots,x_{17}\in[10;12);x_{18},\ldots,x_{24}\in[12;14)$.\\
		Nên trung vị của mẫu  số liệu $x_1;x_2;\ldots;x_{24}$ là $\dfrac{1}{2}(x_{12}+x_{13})\in[10;12)$.\\
		Ta xác định được  $n=24$; $n_m=8$; $C=3+6=9$; $u_m=10$; $u_{m+1}=12$.\\
		Suy ra tứ phân vị thứ hai của mẫu  số liệu ghép nhóm là
		\[Q_2=u_m+\dfrac{\dfrac{n}{2}-C}{n_m}\cdot(u_{m+1}-u_{m})=10+\dfrac{\dfrac{24}{2}-9}{8}\cdot(12-10)=10{,}75.\] 
		Tứ phân vị thứ nhất của mẫu số  liệu là $\dfrac{1}{2}(x_6+x_7)\in[8;10)$.\\
		 Ta xác định được $n=24$; $n_m=6$; $C=3$; $u_m=8$; $u_{m+1}=10$.\\
		 \[Q_1=u_m+\dfrac{\dfrac{n}{4}-C}{n_m}\cdot(u_{m+1}-u_{m})=8+\dfrac{\dfrac{24}{4}-3}{6}\cdot(10-8)=9.\]
	 }
\end{ex}

% nguồn: Ok_TraLoiNgan_Lop11_Lan1
\begin{ex}%[1D5H2-3]%[ma:40e00cbc]
	Kết quả thi đấu thể thao của  $20$ đội được cho bởi bảng sau:
	\begin{center}
		\begin{tabular}{|c|c|c|c|c|}
			 \hline
			Điểm số &  $[5{,}5;10{,}5)$ & $[10{,}5;15{,}5)$ & $[15{,}5;20{,}5)$ & $[20{,}5;25{,}5)$ \\
			 \hline
			Số đội & $4$ & $8$ &  $2$ & $6$ \\
			 \hline
		\end{tabular}
	\end{center}
	Tứ phân vị thứ hai của số liệu ở bảng tần số ghép nhóm trên bằng bao  nhiêu? (\textit{Kết quả được làm tròn đến hàng đơn vị}).
	 \shortans{14}
	\loigiai{ 
		Gọi $x_1;x_2;\ldots;x_{20}$ lần lượt là số  điểm được xếp theo thứ tự không giảm.\\
		 Do $x_1,\ldots,x_4\in [5{,}5;10{,}5); x_5,\ldots,x_{12}\in [10{,}5;15{,}5); x_{13},x_{14}\in[15{,}5;20{,}5); \break x_{15},\ldots,x_{20}\in [20{,}5;25{,}5)$.\\
		Nên trung vị của mẫu số liệu  $x_1;x_2;\ldots;x_{20}$ là $\dfrac{1}{2}(x_{10}+x_{11})\in[10{,}5;15{,}5)$.\\
		Ta xác định được $n=20$; $n_m=8$; $C=4$; $u_m=10{,}5$;  $u_{m+1}=15{,}5$.\\
		Vậy tứ phân vị thứ hai của mẫu số liệu  ghép nhóm là
		 \[
		 Q_2=u_m+\dfrac{\dfrac{n}{2}-C}{n_m}\cdot(u_{m+1}-u_{m})=10{,}5+\dfrac{\dfrac{20}{2}-4}{8}\cdot(15{,}5-10{,}5)\approx14.
		 \]
	 }
\end{ex}

% nguồn: Ok_TraLoiNgan_Lop11_Lan2
\begin{ex}%[1D5H2-3]%[ma:f7fccd6c]
	Số điểm của một đội bóng chuyền ghi được  trong 23 séc đấu được cho ở bảng sau:
	\begin{center}\begin{tabular}{|c|c|c|c|c|c|c|c|c|c|c|c|}
		 \hline
		$8$&$12$&$25$&$23$&$22$&$18$&$20$&$21$&$16$&$15$&$26$&$24$\\ 
		\hline 
		$16$&$23$&$20$&$19$&$21$&$10$&$24$&$15$&$19$&$22$&$27$&\\ 
		 \hline
	\end{tabular}
\end{center}
Tìm trung vị của dãy số  trên.
\shortans{20} 
 \loigiai{
	Sắp xếp mẫu số liệu theo thứ tự không giảm, ta được: 
	$$8;10;12;15;15;16;16;18;19;19;20;20;21;21;22;22;23;23;24;24;25;26;27.$$ 
	Vì  cỡ mẫu $n=23$ là số lẻ nên trung vị của dãy số liệu trên là $M_e=20$.
 }
\end{ex}

% nguồn: Ok_TraLoiNgan_Lop11_Lan2
\begin{ex}%[1D5H2-3]%[ma:9e1c7622]
	Lượng nước tiêu thụ của một số  hộ gia đình được ghi lại kết quả như sau:
	\begin{center}\begin{tabular}{|c|c|c|c|c|c|c|c|c|}
		\hline 
		 Lượng nước (m$^3$) &$[0;10)$&$[10;20)$&$[20;30)$&$[30;40)$&$[40;50)$&$[50;60)$&$[60;70)$&$[70;80)$\\
		 \hline
		Số hộ gia đình&$2$&$4$&$6$&$15$&$10$&$8$&$4$&$1$\\ 
		 \hline
	\end{tabular}
\end{center}
Tìm trung vị của mẫu số liệu  ghép nhóm trên.
 \shortans{38,7}
\loigiai{ 
	Gọi $x_1\le x_2\le x_3\le\ldots x_{50}$ là mẫu số liệu được  xếp theo thứ tự không giảm.\\
	Vì cỡ mẫu $n=50$ là số chẵn  nên trung vị của mẫu số liệu số $x_1\le x_2\le x_3\le\ldots x_{50}$ là $\dfrac{1}{2}\left(x_{25}+x_{26}\right)\in [30;40)$.\\
	Vậy trung vị của mẫu số  liệu ghép nhóm trên là $M_e=30+\dfrac{\dfrac{50}{2}-12}{15}\left(40-30\right)=\dfrac{116}{3}\approx 38{,}7$.
} 
\end{ex}

% nguồn: Ok_TraLoiNgan_Lop11_Lan2
\begin{ex}%[1D5H2-3]%[ma:959cc923]
	Sản lượng sữa mỗi ngày của một  số con bò thuộc hai giống $A$ và $B$ được cho bởi bảng sau:
	\begin{center}\begin{tabular}{|c|c|c|c|c|c|}
		\hline 
		Sản lượng sữa (lít)  &$[18;20)$&$[20;22)$&$[22;24)$&$[24;26)$&$[26;28)$\\
		 \hline
		Giống A&$2$&$8$&$25$&$15$&$4$\\ 
		\hline 
		Giống  B&$8$&$28$&$10$&$6$&$4$\\
		 \hline
	\end{tabular}
\end{center}
Hãy ước lượng tứ  phân vị thứ ba của lượng sữa mỗi ngày của bò giống $B$.
 \shortans{23,2}
 \loigiai{
	Ta có bảng  thống kê lượng sữa theo giá trị đại diện:
	\begin{center}\begin{tabular}{|c|c|c|c|c|c|}
		 \hline
		Sản lượng sữa (lít)  &$[18;20)$&$[20;22)$&$[22;24)$&$[24;26)$&$[26;28)$\\
		 \hline
		Sản lượng sữa đại diện (lít) &$19$&$21$&$23$&$25$&$27$\\ 
		 \hline
		Giống A&$2$&$8$&$25$&$15$&$4$\\ 
		 \hline
		Giống  B&$8$&$28$&$10$&$6$&$4$\\
		 \hline
	\end{tabular}
	\end{center}
	Gọi $y_1\le y_2\le y_3\le\ldots y_{56}$ là mẫu số  liệu được xếp theo thứ tự không giảm của giống bò $B$.\\
	Vì cỡ mẫu $n=56$ là số chẵn  nên trung vị của mẫu số liệu số $y_1\le y_2\le y_3\le\ldots y_{56}$ là $\dfrac{1}{2}\left(x_{28}+x_{29}\right)\in [20;22)$.\\
	Vậy trung vị của mẫu số liệu ghép nhóm giống  bò $B$ là $M_e=20+\dfrac{\dfrac{56}{2}-8}{28}\left(22-20\right)=\dfrac{150}{7}$.\\
	 Giống bò $B$.\\
	Tứ phân vị thứ ba của dãy số liệu $y_1\le y_2\le y_3\le\ldots  y_{56}$ là $\dfrac{1}{2}\left(y_{42}+y_{43}\right)\in [22;24)$.\\
	Do đó tứ phân vị thứ ba  của mẫu số liệu ghép nhóm là $Q_3=22+\dfrac{\dfrac{3\cdot56}{4}-36}{10}.\left(24-22\right)=23{,}2$.
 }
\end{ex}

% nguồn: Ok_TraLoiNgan_Lop11_Lan2
\begin{ex}%[1D5H2-3]%[ma:edeed2d6]
	Điều tra $42$ học sinh của một lớp $11$ về số giờ tự học ở nhà, người ta có bảng sau  đây
	\begin{center}\begin{tabular}{|c|c|c|c|c|c|}
		\hline 
		Số  giờ tự học &$[1;2)$&$[2;3)$&$[3;4)$&$[4;5)$&$[5;6)$\\
		 \hline
		 Tần số &$8$&$10$&$12$&$9$&$3$\\
		\hline 
		Tần số tích  lũy&$8$&$18$&$30$&$39$&$42$\\
		 \hline
	\end{tabular}
\end{center}
Số trung vị  của mẫu số liệu là bao nhiêu?
 \shortans{3,25}
\loigiai{ 
	Ta có số phần tử của mẫu là  $n=42\Rightarrow \dfrac{n}{2}=21$.\\
	Mà $cf_2=18<21<cf_3=30$ suy ra nhóm $3$ là nhóm đầu tiên có tần số tích lũy lớn hơn hoặc  bằng $21$.\\
	Xét nhóm $3$ là nhóm $\left[3;4\right)$  có $r=3;d=1;n_3=12$ và nhóm $2$ là nhóm $\left[2;3\right)$ có $cf_2=18$.\\
	Áp dụng công thức  ta có trung vị của mẫu số liệu là
	$M_e=3+\left(\dfrac{21-18}{12}\right).1=3{,}25$ (giờ). 
 }
\end{ex}

% nguồn: Ok_TraLoiNgan_Lop11_Lan2
\begin{ex}%[1D5H2-3]%[ma:aef3ba54]
	Lượng mưa trung bình trong một tháng của một số tỉnh được ghi lại ở bảng sau  (đơn vị: mm)
	\begin{center}\begin{tabular}{|c|c|c|c|c|c|c|c|c|c|c|c|}
		 \hline
		$5{,}3$&$4{,}5$&$10{,}0$&$32{,}4$&$36{,}5$&$45{,}8$&$40{,}7$&$33{,}2$&$5{,}7$&$12{,}5$&$16{,}8$&$9{,}9$\\ 
		 \hline
		$6{,}5$&$7{,}8$&$9{,}1$&$15{,}7$&$30{,}5$&$49{,}7$&$28{,}4$&$48{,}3$&$18{,}8$&$20{,}7$&$16{,}2$&$9{,}7$\\ 
		 \hline
	\end{tabular}
\end{center}
Tìm tứ  phân vị thứ ba của mẫu số liệu trên
 \shortans{32,8}
 \loigiai{
	Sắp xếp mẫu số liệu theo thứ tự không giảm, ta được 
	$$\begin{aligned}
	& 4{,}5;5{,}3;5{,}7;6{,}5;7{,}8;9{,}1;10{,}0;12{,}5;15{,}7;16{,}2;16{,}8;18{,}8;19{,}7;19{,}9;20{,}7; \\ 
	& 28{,}4;30{,}5;32{,}4;33{,}2;36{,}5;40{,}7;45{,}8;48{,}3;49{,}7 \\ 
	\end{aligned}$$
	Vì cỡ mẫu $n=24$  là số chẵn nên tứ phân vị thứ hai là $Q_2=\dfrac{1}{2}\left(18{,}8+19{,}7\right)=19{,}25$.\\
	Tứ phân vị thứ nhất là trung vị của  mẫu:$4{,}5;5{,}3;5{,}7;6{,}5;7{,}8;9{,}1;10{,}0;12{,}5;15{,}7;16{,}2;16{,}8;18{,}8$.\\
	Do  đó $Q_1=\dfrac{1}{2}\left(9{,}1+10\right)=9{,}55$.\\
	Tứ phân vị thứ ba  là trung vị của mẫu:$19{,}7;19{,}9;20{,}7;28{,}4;30{,}5;32{,}4;33{,}2;36{,}5;40{,}7;45{,}8;48{,}3;49{,}7$.\\
	Do  đó $Q_3=\dfrac{1}{2}\left(32{,}4+33{,}2\right)=32{,}8$.
 }
\end{ex}

% nguồn: Ok_TraLoiNgan_Lop11_Lan2
\begin{ex}%[1D5H2-3]%[ma:ad32b4fd]
	Điều tra $42$ học sinh của một lớp $11$ về số giờ tự học ở nhà, người ta  có bảng sau đây:
	\begin{center}\begin{tabular}{|c|c|c|c|c|c|}
		 \hline
		Số giờ tự học &$[1;2)$&$[2;3)$&$[3;4)$&$[4;5)$&$[5;6)$\\ 
		\hline 
		Tần  số &$8$&$10$&$12$&$9$&$3$\\
		\hline 
		Tần  số tích lũy&$8$&$18$&$30$&$39$&$42$\\
		\hline 
	\end{tabular}
\end{center}
Tứ phân vị thứ nhất của mẫu số liệu trên là bao  nhiêu?
\shortans{2,25} 
 \loigiai{
	Ta có số  phần tử của mẫu là $n=42\Rightarrow \dfrac{n}{4}=10{,}5$.\\
	Suy ra nhóm $2$ là nhóm đầu  tiên có tần số tích lũy lớn hơn hoặc bằng $10{,}5$.\\
	Xét nhóm $2$ là nhóm $\left[2;3\right)$  có $s=2;h=1;n_2=10$ và nhóm $1$ là nhóm $\left[1;2\right)$ có $cf_1=8$.\\
	 Áp dụng công thức ta có trung vị của mẫu số liệu là
	$Q_1=2+\left(\dfrac{10{,}5-8}{10}\right).1=2{,}25$  (giờ).
 }
\end{ex}

% nguồn: Ok_TraLoiNgan_Lop11_Lan1
\begin{ex}%[1D5V2-3]%[ma:af7913bf]
	Phỏng vấn một  số học sinh khối $11$ về thời gian (giờ) ngủ của một buổi tối, thu được bảng số
	liệu  sau:
	\begin{center}
		\begin{tabular}{|c|c|c|c|c|c|}
			 \hline
			 Thời gian (Giờ) & $\left[4;5\right)$ & $\left[5;6\right)$ &$\left[6;7\right)$ &$\left[7;8\right)$ &$\left[8;9\right)$ \\
			\hline 
			Số  lượng & $6$ &$12$ &$13$ &$10$ &$3$ \\
			\hline 
		\end{tabular}
	\end{center}
	 Hãy cho biết $75\%$ học sinh khối $11$ ngủ nhiều nhất bao nhiêu giờ?
	 \shortans{7,2}
	 \loigiai{
		Cỡ  mẫu là $n=44$.\\
		Tứ phân vị thứ hai $Q_2$ là giá  trị của $\dfrac{x_{22}+x_{23}}{2}$.\\
		Tứ phân vị thứ nhất  $Q_1$ là giá trị của $\dfrac{x_{11}+x_{12}}{2}$.\\
		Tứ phân vị thứ ba $Q_2$ là giá trị của  $\dfrac{x_{33}+x_{34}}{2}$.\\
		Do $x_{33};x_{34}$ thuộc nhóm $\left[7;8\right)$  nên nhóm này chứa $Q_3$.\\
		Ta xác định được $n=44$; $p=4$;  $m_4=10$; $m_1+m_2+m_3=31$, $a_5=8$, $a_4=7$.\\
		Vậy  $Q_3=7+\dfrac{33-31}{10}\cdot (8-7)=7{,}2$.\\
		Vậy tứ  phân vị thứ ba nên $75\%$ học sinh khối $11$ ngủ nhiều nhất là $7{,}2$ giờ.
	 }
\end{ex}

% nguồn: Ok_TraLoiNgan_Lop11_Lan1
\begin{ex}%[1D5V2-3]%[ma:21e10761]
	Phỏng vấn một số học  sinh khối 11 vể thời gian (giờ) ngủ của một buổi tối, thu được số liệu như bảng bên dưới.
	\begin{center}
		\begin{tabular}{|c|c|c|c|c|c|}
			 \hline
			Thời gian (giờ)  & $\left[4;5\right)$ & $\left[5;6\right)$ & $\left[6;7\right)$ & $\left[7;8\right)$ & $\left[8;9\right)$ \\
			 \hline
			Số học sinh & 6 & 10 & 13 & 9 & 7  \\
			\hline 
		\end{tabular}
	\end{center}
	Hãy cho biết $75 \%$ học sinh khối 11 ngủ ít nhất bao nhiêu giờ? (kết quả làm tròn  đến hàng phần trăm)
	 \shortans{7,53}
	 \loigiai{
		 Cỡ mẫu là $n=45$.\\
		Gọi $x_1, \ldots, x_{45}$  là mẫu số liệu được sắp xếp theo thứ tự không giảm. \\
		Khi đó, trung vị  là $x_{23}$.\\
		Do đó,  tứ phân vị thứ ba $Q_3$ là $x_{34}$. \\
		Do $x_{34}$ đều thuộc nhóm $[7;8)$ nên  nhóm này chứa $Q_3$.\\
		Do đó,  $p=4; \;a_4=7;\; m_4=9; \;m_1+m_2+m_3=29; \;a_5-a_4=1$ và ta có
		$$ 
		Q_3=7+\dfrac{\dfrac{3  \cdot 45}{4}-29}{9}\cdot 1\approx7{,}53.
		 $$
		Vậy $75\%$ học sinh khối  11 ngủ ít nhất $7{,}53$ giờ.
	 }
\end{ex}

% nguồn: Ok_TraLoiNgan_Lop11_Lan1
\begin{ex}%[1D5V2-3]%[ma:3ce0c797]
	Thời gian luyện tập trong một ngày (tính theo giờ) của một số vận động viên được ghi lại ở  bảng sau:
	\begin{center}
		\begin{tabular}{|c|c|c|c|c|c|}
			\hline 
			\textbf{Thời gian luyện tập  (giờ)} & $ \left[0; 2\right) $ & $ \left[2; 4\right) $ & $ \left[4; 6\right) $ & $ \left[6; 8\right) $ & $ \left[8; 10 \right) $ \\
			 \hline
			\textbf{Số vận động viên} & $ 3 $ & $ 8 $ & $ 12 $ & $ 12  $ & $ 4 $ \\
			 \hline
		\end{tabular}
	\end{center}
	Huấn luyện viên muốn xác định nhóm gồm $ 25\% $ các vận  động viên có số giờ luyện tập cao nhất. Hỏi huấn luyện viên nên chọn các vận động viên có thời gian luyện tập từ bao nhiêu giờ trở lên vào nhóm này?
	\shortans{7,04} 
	 \loigiai{
		Số vận động viên được khảo sát là $ n=3+8+12+12+4=39$.\\ 
		Gọi $ x_1 $; $ x_2 $; \ldots;$ x_{39} $ là thời gian luyện tập của $ 39 $ vận động viên được xếp theo thứ  tự không giảm.\\
		Ta có $ x_1, x_2, x_3 \in  \left[0; 2\right) $; $ x_4, \ldots, x_{11}\in \left[2; 4 \right) $; $ x_{12}, \ldots, x_{23}\in \left[4; 6\right) $; $ x_{24}, \ldots, x_{35}\in \left[6;8\right) $; $ x_{36},\ldots , x_{39}\in \left[8; 10\right) $. Do đó đối với dãy số liệu $ x_1 $; $ x_2 $; \ldots; $ x_{39} $ thì
		\begin{itemize}
			\item Tứ phân vị thứ nhất là $ x_{10} $ thuộc  nhóm $\left[2; 4\right)$;
			 \item Tứ phân vị thứ hai là $ x_{20} $ thuộc nhóm $ \left[4; 6\right)$;
			\item  Tứ phân vị thứ ba là $ x_{30} $ thuộc nhóm $ \left[6; 8\right) $.
		\end{itemize}
		Ta nói nhóm $  \left[2; 4 \right)$ là nhóm chứa tứ phân vị thứ nhất; nhóm $ \left[4; 6\right) $ là nhóm tứ phân vị thứ hai; nhóm $ \left[6; 8\right) $ là nhóm chứa tứ phân vị thứ ba.\\
		Tứ phân vị thứ hai của dãy số liệu $ x_1; x_2; \ldots; x_{39}$ là $ x_{20}\in \left[4; 6\right) $. Do đó tứ phân vị thứ hai của mẫu số  liệu ghép nhóm là $$Q_2=4+\dfrac{\dfrac{2\cdot 39}{4}-(3+8)}{12}\cdot(6-4)=\dfrac{65}{12}\approx 5{,}42.$$
		Tứ phân vị thứ nhất của dãy số liệu $ x_1; x_2; \ldots; x_{39}$ là $ x_{10}\in \left[2; 4\right) $. Do đó tứ phân vị thứ nhất của mẫu số liệu ghép nhóm là  $$Q_1=2+\dfrac{\dfrac{1\cdot 39}{4}-3}{8}\cdot(4-2)=\dfrac{59}{16}\approx 3{,}69.$$
		Tứ phân vị thứ ba của dãy số liệu $ x_1; x_2; \ldots; x_{39}$ là $ x_{30}\in \left[6; 8\right) $. Do đó tứ phân vị thứ  ba của mẫu số liệu ghép nhóm là $$Q_3=6+\dfrac{\dfrac{3\cdot 39}{4}-(3+8+12)}{12}\cdot(8-6)=\dfrac{169}{24}\approx 7{,}04.$$
		Huấn luyện viên muốn xác  định nhóm gồm $ 25\% $ các vận động viên có số giờ luyện tập cao nhất thì huấn luyện viên nên chọn các vận động viên có thời gian tập luyện từ $ Q_3 \approx 7{,}04$ (giờ).
	} 
\end{ex}

% nguồn: Ok_TraLoiNgan_Lop11_Lan1
\begin{ex}%[1D5V2-3]%[ma:d0a51b97]
	Cân nặng (đơn vị: kg) của $29$ học sinh nam ở một lớp  11 được cho như sau
	\setlength\extrarowheight{2pt} 
	\begin{longtable}{|c|c|c|c|c|c|c|c|c|c|}
		\hline 
		52{,}3 & 61{,}7 &  51{,}2 & 55{,}9 & 56{,}5 & 60{,}4 & 60{,}7 & 50{,}2 & 58{,}5 & 64{,}6 \\
		\hline 
		61{,}8 & 52{,}3 & 63{,}4 & 57{,}9 & 56{,}2 & 50{,}3 & 45{,}8 &  53{,}3 & 64{,}1 & 47{,}7 \\
		\hline 
		49{,}6 & 59{,}1 &  55{,}3 & 55{,}8 & 45{,}5 & 46{,}8 & 54 & 49{,}2 & 52{,}6 & \\
		\hline 
	\end{longtable}
	\noindent Chia mẫu dữ liệu thành 5 nhóm với [45;49) là nhóm chứa cân nặng thấp nhất của học sinh, từ  đó ước lượng tứ phân vị thứ ba của mẫu dữ liệu đã ghép nhóm ta được kết quả sau khi làm tròn đến hàng phần chục là
	 \shortans{58,4}
	\loigiai{Chia mẫu dữ liệu thành 5 nhóm với [45;49) là nhóm chứa cân nặng thấp nhất của học sinh, ta được bảng tần số ghép nhóm theo cân nặng của học sinh như sau 
		\begin{longtable}{|c|c|c|c|c|c|}
			 \hline
			Cân nặng (đơn vị: kg) &  [45;49) & [49;53) & [53;57) & [57;61) & [61;65) \\
			 \hline
			Số  học sinh & 4 & 9 & 7 & 5 & 4 \\
			\hline 
		\end{longtable}
		\noindent Gọi $x_1$; $x_2$; $\ldots$; $x_{29}$ là mẫu số liệu được sắp xếp theo thứ tự  không giảm.\\
		Ta có $x_1; x_2; x_3; x_4 \in [45;49)$, $x_4; x_5; \ldots; x_{12} \in [49;53)$, $x_{13}; x_{14}; \ldots; x_{20} \in [53;57)$, $x_{21}; x_{22}; \ldots; x_{25} \in  [57;61)$, $x_{26}; x_{27}; x_{28}; x_{29} \in [61;65)$.\\
		Tứ phân vị thứ ba của dãy số liệu $x_1;  x_2; \ldots; x_{29}$ là $x_{21}$ nên nó thuộc vào nhóm $[57;61)$.\\
		Vậy tứ phân vị thứ ba của  mẫu số liệu ghép nhóm là
		$$Q_3 = 57 + \dfrac{\dfrac{3\cdot 29}{4}-(4+9+7)}{5} \cdot (61 -  57) = 58{,}4.$$
	} 
\end{ex}

% nguồn: Ok_TraLoiNgan_Lop11_Lan1
\begin{ex}%[1D5V2-3]%[ma:1c3fa5a8]
	Khảo sát về cân nặng của các học sinh lớp $11$D$3$  người ta được một mẫu dữ liệu ghép nhóm như sau
	\begin{center}
		\begin{tabular}{|c|c|c|c|c|c|c|}
			\hline Cân nặng & {$[30;  40)$} & {$[40; 50)$} & {$[50; 60)$} & {$[60; 70)$} & {$[70; 80)$} & {$[80; 90)$} \\
			\hline Số học sinh & $2$ & $10$ & $16$ & $8$  & $2$ & $2$ \\
			 \hline
		\end{tabular}
	\end{center}
	 Tìm tứ phân vị thứ nhất của mẫu số liệu trên.
	 \shortans{48}
	\loigiai{ 
		Ta có số  phần tử của mẫu là $n=40\Rightarrow\dfrac{n}{4}=10$.\\
		Suy ra nhóm $\left[40\,;\,50\right)$ chứa tứ phân  vị thứ nhất.\\
		Do đó tứ phân  vị thứ nhất của mẫu số liệu trên là\\
		 $Q_1=40+\dfrac{10-2}{10}\cdot10=48$.
	} 
\end{ex}

% nguồn: Ok_TraLoiNgan_Lop11_Lan1
\begin{ex}%[1D5V2-3]%[ma:990937c8]
	Một hãng xe ô tô thống kê lại số lần gặp sự cố về động cơ về động cơ của $100$ chiếc xe cùng loại sau $2$ năm sử dụng đầu tiên ở  bảng sau
	\begin{center}
		\begin{tabular}{|l|c|c|c|c|c|}
			\hline Số lần gặp  sự cố &{$[0{,}5;2{,}5)$}&{$[2{,}5;4{,}5)$}&{$[4{,}5;6{,}5)$}&{$[6{,}5; 8{,}5)$}&{$[8{,}5;10[,]5)$}\\
			\hline Số  xe & $17$ & $33$ & $25$ & $20$ & $5$ \\
			\hline 
		\end{tabular}
	\end{center}
	 Tìm tứ phân vị thứ hai của mẫu số liệu.
	 \shortans{4,5}
	\loigiai{Cỡ mẫu là $n=100$.\\ 
		Gọi $x_1, \ldots, x_{100}$ là mẫu số  liệu được sắp xếp theo thứ tự không giảm.\\
		Khi đó, trung vị  là $\dfrac{x_{50}+x_{51}}{2}$.\\
		Do $x_{50} \in [2{,}5;4{,}5)$, $x_{51} \in [4,5;6,5)$ nên tứ phân vị thứ hai của mẫu số liệu ghép nhóm  là $Q_2=4{,}5$.
	 }
\end{ex}

% nguồn: Ok_TraLoiNgan_Lop11_Lan1
\begin{ex}%[1D5V2-3]%[ma:f4fda15d]
	Một người thống kê lại thời gian thực hiện  các cuộc gọi điện thoại của người đó trong một tuần ở bảng sau:
	\begin{center}
		\begin{tabular}{|c|c|c|c|c|c|c|}
			\hline 
			Thời gian (giây)& $ [0;\,60) $ & $ [60;\,120) $ & $ [120;\,180) $ & $ [180;\,240)  $ & $ [240;\,300) $ & $ [300;\,360) $ \\
			 \hline
			Số cuộc gọi & $ 8 $ & $ 10 $ & $ 7  $ & $ 5 $ & $ 2 $ & $ 1 $ \\
			 \hline
		\end{tabular}
	\end{center}
	 Tứ phân vị thứ 3 của mẫu số liệu bằng
	\shortans{180} 
	\loigiai{ 
		Gọi $ x_1,\,\ldots\,,  x_{33} $ là thời gian thực hiện các cuộc gọi trong một tuần theo thứ tự không giảm.\\
		Tứ phân vị thứ ba của mẫu số liệu là  $ \dfrac{1}{2}(x_{25}+x_{26}) $ nên $ Q_3=180 $.
	 }
\end{ex}

% nguồn: Ok_TraLoiNgan_Lop11_Lan1
\begin{ex}%[1D5V2-3]%[ma:f755826e]
	Khảo sát thời gian (phút) tập thể dục trong ngày của 1 số học  sinh lớp 11 thu được mẫu số liệu ghép nhóm sau
	\begin{center}
		\begin{tabular}{|l|c|c|c|c|c|}
			\hline Thời gian (phút)  & { $[0; 10)$ } & { $[10; 20)$ } & { $[20; 30)$ } & { $[30; 40)$ } & { $[40; 50)$ } \\
			 \hline Số học sinh & 5 & 10 & 13 & 9 & 7 \\
			 \hline
		\end{tabular}
	\end{center}
	Tìm  số $a$ sao cho có $25 \%$ số học sinh tập thể dục với thời gian ít nhất là $a$ (phút).
	\shortans{35,6} 
	 \loigiai{
		Trong mỗi khoảng thời gian, giá trị đại điện là trung bình  cộng của giá trị hai đầu mút nên ta có bảng sau
		\begin{center}
			\begin{tabular}{|l|c|c|c|c|c|}
				\hline Thời gian (phút) & 5 & 15 & 25 & 35 & 45  \\
				\hline Số học sinh & 5 & 10 & 13  & 9 & 7 \\
				 \hline
			\end{tabular}
		\end{center}
		Số $a$  chính là tứ phân vị thứ $3$.\\
		Gọi $x_1, \ldots, x_{44}$ là thời gian tập thể dục (phút) trong ngày của $44$ học sinh lớp $11$ và giả sử  dãy này được xếp theo thứ tự tăng dần.
		 \\
		Khi đó,  tứ phân vị thứ $3$ là $\dfrac{x_{33}+x_{34}}{2}$. \\
		Do $x_{33}$; $x_{34}$ đều thuộc  nhóm $[30; 40)$ nên nhóm này chứa $Q_3$. \\
		Do đó $p=4$; $a_4=30$; $m_4=9$; $m_1+m_2+m_3=5+10+13=28$; $a_5-a_4=10$. 
		$$ 
		Q_3=30+\dfrac{\frac{3\cdot  44}{4}-28}{9} \cdot 10=\dfrac{320}{9} \approx 35{,}6.
		$$ 
		Vậy có  $25\%$ số học sinh tập thể dục với thời gian ít nhất là $35{,}6$ (phút).
	} 
\end{ex}

% nguồn: Ok_TraLoiNgan_Lop11_Lan1
\begin{ex}%[1D5V2-3]%[ma:d874890b]
	Thời gian (đơn vị: phút) di chuyển đến trường của nhóm học  sinh trường THPT A được tổng hợp trong bảng sau:
	\begin{longtable}{|c|c|c|c|c|c|c|c|}
		 \hline
		 Thời gian (phút)&$[15;20)$&$[20;25)$&$[25;30)$&$[30;35)$&$[35;40)$&$[40;45)$&$[45;50)$\\
		 \hline
		 Số học sinh&$6$&$14$&$25$&$37$&$13$&$9$&$21$\\
		 \hline
	\end{longtable}
	Tìm khoảng tứ phân vị của mẫu số liệu ghép nhóm trên (kết  quả làm tròn đến hàng phần mười).
	 \shortans{12,3}
	 \loigiai{
		 Cỡ mẫu $n=125$.\\
		Gọi $x_1,x_2,\ldots,x_{125}$ lần lượt là thời gian di chuyển đến trường  của $125$ học sinh xếp theo thứ tự không giảm.\\
		Tứ phân vị thứ nhất của mẫu số liệu $x_1,x_2,\ldots,x_{125}$  là $Q_1=\dfrac{1}{2}(x_{31}+x_{32})\in [25;30)$.\\
		Do đó giá trị của tứ  phân vị thứ nhất là $$Q_1=25+\dfrac{\dfrac{125}{4}-(6+14)}{25}\cdot(30-25)=27{,}25\,\text{(phút)}.$$
		Tứ phân vị thứ ba của mẫu số liệu $x_1,x_2,\ldots,x_{125}$ là $Q_3=\dfrac{1}{2}(x_{94}+x_{95})\in [35;40)$.\\ 
		Do đó giá trị của tứ phân  vị thứ ba là $$Q_3=35+\dfrac{\dfrac{3\cdot 125}{4}-(6+14+25+37)}{13}\cdot(40-35)\approx 39{,}52\,\text{(phút)}.$$
		Khoảng tứ phân vị của  mẫu số liệu $\Delta Q=Q_3-Q_1=39{,}52-27{,}25\approx 12{,}3$ (phút).
	 }
\end{ex}

% nguồn: Ok_TraLoiNgan_Lop11_Lan1
\begin{ex}%[1D5V2-3]%[ma:309ad374]
	 Số điểm một cầu thủ bóng rổ ghi được trong 20 trận đấu được cho ở bảng sau:
	\begin{center}
		\begin{tabular}{|l|c|c|c|c|c|}
			\hline Điểm  số &{$[5,5;10,5)$}&{$[10,5;15,5)$}&{$[15,5;20,5)$}&{$[20,5;25,5)$}\\
			\hline Số trận & $3$ & $9$ &  $2$ & $6$ \\
			 \hline
		\end{tabular}
	\end{center}
	Ước lượng số điểm cao nhất mà cầu thủ đạt được trong $5$ trận đấu đầu tiên. \textit{(làm tròn  kết quả đến hàng phần chục)}
	\shortans{11,6} 
	\loigiai{ 
		Dễ thấy số điểm cao nhất mà cầu thủ đạt được trong $5$ trận đấu  đầu tiên chính là giá trị của tứ phân vị thứ nhất $Q_1$ trong mẫu số liệu. Do đó:\\
		Với  cỡ mẫu là $n=20$.\\
		Gọi $x_1, \ldots, x_{20}$ là mẫu số  liệu được sắp xếp theo thứ tự không giảm. Khi đó, trung vị là $\dfrac{x_{10}+x_{11}}{2}$. \\
		Tứ phân  vị thứ nhất $Q_1$ là $\dfrac{x_{5}+x_{6}}{2}$. Do $x_{5}$, $x_{6}$ đều thuộc nhóm $[10,5;15,5)$ nên nhóm này chứa $Q_1$. \\ Do đó, $p=2; \;a_2=10{,}5;\; m_2=9; \;m_1=3, \;a_3-a_2=5$ và ta có
		 $$
		Q_1=10{,}5+\dfrac{\frac{20}{4}-3}{9}\cdot 5\approx 11{,}6. 
		 $$
		Vậy số điểm cao nhất mà cầu thủ đạt được trong $5$ trận đấu  đầu tiên là $11{,}6$.
	} 
\end{ex}

% nguồn: Ok_TraLoiNgan_Lop11_Lan1
\begin{ex}%[1D5V2-3]%[ma:7ce92027]
	Khảo sát thời gian tập thể dục của một số học sinh khối 11 thu được mẫu số liệu ghép  nhóm sau:
	\begin{center}
		\begin{tabular}{|l|c|c|c|c|c|}
			\hline Thời gian (phút) & {$[0; 20)$} & {$[20; 40)$} & {$[40; 60)$} & {$[60;  80)$} & {$[80; 100)$} \\
			\hline  Số học sinh & $5$ & $9$ & $12$ & $10$ & $6$ \\
			 \hline
		\end{tabular}
	\end{center}
	Tứ phân vị thứ ba của mẫu số  liệu là bao nhiêu?
	 \shortans{71}
	\loigiai{ 
		Cỡ mẫu là  $n=5+9+12+10+6=42$.\\
		Tứ phân vị thứ ba $Q_3$ là $x_{33}$, do nhóm chứa $x_{33}$ là  $[60; 80)$, do đó $p=4$; $a_4=60$; $m_1+m_2+m_3=5+9+12=26$; $a_5-a_4=20$.\\
		Ta có  $Q_3=a_4+\dfrac{\dfrac{3 n}{4}-\left(m_1+m_2+m_3\right)}{m_4}\left(a_5-a_4\right)=60+\dfrac{\dfrac{3.42}{4}-26}{10} \cdot 20=71$.
	 }
\end{ex}

% nguồn: Ok_TraLoiNgan_Lop11_Lan1
\begin{ex}%[1D5C2-3]%[ma:26d6ebb8]
	Phỏng vấn một số học sinh khối $11$ về thời gian (giờ) ngủ của một buổi tối, thu được bảng số liệu ở  bên dưới.
	\begin{center}
		\begin{tabular}{|c|c|c|}
			\hline 
			\textbf{Thời gian}&\textbf{Số học sinh nam}&\textbf{Số học sinh nữ}\\ 
			 \hline
			$[4;5)$&$6$&$4$\\ 
			 \hline
			$[5;6)$&$10$&$8$\\ 
			\hline 
			 $[6;7)$&$13$&$10$\\
			 \hline
			$[7;8)$&$9$&$11$\\ 
			 \hline
			 $[8;9)$&$7$&$8$\\
			\hline 
		\end{tabular}
	\end{center}
	Hãy  cho biết $75\%$ học sinh khối $11$ ngủ ít nhất bao nhiêu giờ?
	 \shortans{5,64}
	\loigiai{ 
		Ta có 
		\begin{center}
			\begin{tabular}{|c|c|c|c|}
				 \hline
				\textbf{Thời  gian}&\textbf{Số học sinh nam}&\textbf{Số học sinh nữ}&\textbf{Số học sinh khối $11$}\\
				\hline 
				$[4;5)$&$6$&$4$&$10$\\ 
				 \hline
				 $[5;6)$&$10$&$8$&$18$\\
				\hline 
				$[6;7)$&$13$&$10$&$23$\\ 
				\hline 
				 $[7;8)$&$9$&$11$&$20$\\
				 \hline
				$[8;9)$&$7$&$8$&$15$\\ 
				\hline 
			\end{tabular}
		\end{center}
		Tổng số học sinh khối $11$ được khảo  sát là $n=45+41=86$.\\
		Gọi $x_1$, $x_2$, $x_3$, $\ldots$, $x_{86}$ là thời gian ngủ của các học sinh khối $11$ được khảo sát và giả sử dãy  này đã sắp xếp theo thứ tự tăng dần.\\
		Khi đó trung vị của mẫu số liệu  là $\dfrac{x_{43}+x_{44}}{2}$.\\
		Do đó, tứ phân vị thứ nhất Q1 là $x_{22}$.\\ 
		Vì $x_{22}$ thuộc nhóm $[5;6)$ nên nhóm này chứa  $Q_1$.\\
		Do đó, $p=2$; $a_2=5$; $m_2=18$; $m_1=10$; $a_3–a_2=6–5=1$ và  ta có\\
		\[Q_1=5+\dfrac{\dfrac{86}{4}-10}{18}\cdot  1 \approx 5{,}64.\]
		Tứ phân vị thứ nhất $Q_1$ chia mẫu số liệu thành $2$ phần, phần dưới chiếm $25\%$ số liệu của mẫu và  phần trên chiếm $75\%$ số liệu của mẫu.\\
		Vậy $75\%$ học sinh khối $11$ ngủ ít nhất  $5{,}64$\,giờ.
	} 
\end{ex}


%%%%%%%%%%%%%%%% TỰ LUẬN %%%%%%%%%%%%%%%%

% nguồn: Ok_TuanLuan_Toan11_Lan1
\begin{bt}%[1D5H2-3]%[ma:bcc0e2d4]
	Thống kê điểm trung bình  môn Toán của một số hoc sinh lớp $11$ được cho ở bảng sau
	\begin{center}
		\begin{tabular}{|c|c|c|c|c|c|c|c|}
			\hline 
			Khoảng  điểm
			 &{$[6{,}5;7)$}&{$[7;7{,}5)$}&{$[7{,}5;8)$}&{$[8;8{,}5)$}&{$[8{,}5;9)$}&{$[9;9{,}5)$}&{$[9{,}5;10)$}\\
			\hline Tần số&$8$&$10$&$16$&$24$&$13$&$7$&$4$\\ 
			\hline 
		\end{tabular}
	\end{center}
	Hãy  ước lượng số tứ phân vị của mẫu số liệu ở bảng tần số ghép nhóm trên.
	\loigiai{ 
		Tứ phân vị thứ hai. Nhóm  $[8;8{{,}}5)$
		 $${Q_2}=8+\dfrac{\dfrac{2\cdot 82}{4}-(8+10+16)}{24}(8{,}5-8)\approx 8{,}15.$$
		Tứ phân vị thứ nhất.  Nhóm $[7{{,}}5;8)$
		$${Q_1}=7{{,}}5+\dfrac{\dfrac{2\cdot 82}{4}-(8+10)}{16}(8-7{,}5)\approx  7{,}58.$$
		Tứ phân vị thứ  ba. Nhóm $[8{{,}}5;9)$
		$${Q_3}=8{{,}}5+\dfrac{\dfrac{3\cdot  82}{4}-(8+10+16+24)}{16}(9-8{,}5)\approx 8{,}63.$$
	} 
\end{bt}

% nguồn: Ok_TuanLuan_Toan11_Lan1
\begin{bt}%[1D5H2-3]%[ma:ef20d892]
	Ghi lại  tốc độ bóng trong $200$ lần giao bóng của một vận động viên môn quần vợt cho kết quả như bảng sau
	\begin{center}
		\begin{tabular}{|c|c|}
			\hline  Tốc độ $v$ (km/h)&Số lần\\
			\hline  $150\le v<155$&$18$\\
			\hline  $155\le v<160$&$28$\\
			\hline $160\le v<165$&$35$\\ 
			\hline $165\le v<170$&$43$\\ 
			\hline  $170\le v<175$&$41$\\
			\hline $175\le  v<180$&$35$\\
			 \hline
		\end{tabular}
	\end{center}
	Tìm tứ phân vị  thứ nhất $Q_1$ và tứ phân vị thứ ba $Q_3$ của mẫu số liệu ghép nhóm trên.
	 \loigiai{
		\begin{center}
			\begin{tabular}{|c|c|c|}
				\hline  Tốc độ $v$ (km/h)&Số lần&TSTL\\
				\hline $150\le v<155$&$18$&$18$\\ 
				 \hline $155\le v<160$&$28$&$46$\\
				\hline $160\le  v<165$&$35$&$81$\\
				\hline $165\le  v<170$&$43$&$124$\\
				\hline  $170\le v<175$&$41$&$165$\\
				 \hline $175\le v<180$&$35$&$200$\\
				 \hline
			\end{tabular}
		\end{center}
		Ta có $\dfrac{n}{4}=\dfrac{200}{4}=50$ nên $Q_1\in [160;165)$. Khi đó 
		$$Q_1=160+\dfrac{50-46}{35}\cdot 5=\dfrac{1124}{7}.$$ 
		 Ta có $\dfrac{3n}{4}=\dfrac{3\cdot 200}{4}=150$ nên $Q_3\in [170;175)$. Khi đó
		 $$Q_3=170+\dfrac{150-124}{41}\cdot 5=\dfrac{7100}{41}.$$
	 }
\end{bt}

% nguồn: Ok_TuanLuan_Toan11_Lan2
\begin{ex}%[1D5H2-3]%[ma:3890ca60]
	Thống  kê điểm trung bình môn Toán của một số hoc sinh lớp $11$ được cho ở bảng sau
	\begin{center}
		\begin{tabular}{|c|c|c|c|c|c|c|c|}
			\hline 
			Khoảng  điểm
			&{$[6{,}5;7)$}&{$[7;7{,}5)$}&{$[7{,}5;8)$}&{$[8;8{,}5)$}&{$[8{,}5;9)$}&{$[9;9{,}5)$}&{$[9{,}5;10)$}\\ 
			\hline  Tần số&$8$&$10$&$16$&$24$&$13$&$7$&$4$\\
			\hline 
		\end{tabular}
	\end{center}
	Hãy ước lượng số tứ phân vị  của mẫu số liệu ở bảng tần số ghép nhóm trên.
	\loigiai{ 
		Tứ phân vị thứ hai. Nhóm  $[8;8{{,}}5)$
		 $${Q_2}=8+\dfrac{\dfrac{2\cdot 82}{4}-(8+10+16)}{24}(8{,}5-8)\approx 8{,}15.$$
		Tứ phân vị thứ  nhất. Nhóm $[7{{,}}5;8)$
		 $${Q_1}=7{{,}}5+\dfrac{\dfrac{2\cdot 82}{4}-(8+10)}{16}(8-7{,}5)\approx 7{,}58.$$
		Tứ phân vị thứ  ba. Nhóm $[8{{,}}5;9)$
		$${Q_3}=8{{,}}5+\dfrac{\dfrac{3\cdot 82}{4}-(8+10+16+24)}{16}(9-8{,}5)\approx 8{,}63.$$ 
	 }
\end{ex}

% nguồn: Ok_TuanLuan_Toan11_Lan2
\begin{ex}%[1D5H2-3]%[ma:094e5ba2]
	Ghi lại tốc độ bóng trong $200$ lần giao bóng của một vận động viên môn quần vợt cho kết  quả như bảng sau
	\begin{center}
		\begin{tabular}{|c|c|}
			\hline Tốc độ $v$ (km/h)&Số  lần\\
			 \hline $150\le v<155$&$18$\\
			\hline  $155\le v<160$&$28$\\
			\hline  $160\le v<165$&$35$\\
			 \hline $165\le v<170$&$43$\\
			 \hline $170\le v<175$&$41$\\
			\hline  $175\le v<180$&$35$\\
			\hline 
		\end{tabular}
	\end{center}
	Tìm tứ phân vị thứ nhất $Q_1$ và tứ phân vị thứ ba  $Q_3$ của mẫu số liệu ghép nhóm trên.
	\loigiai{ 
		\begin{center}
			\begin{tabular}{|c|c|c|}
				\hline Tốc  độ $v$ (km/h)&Số lần&TSTL\\
				\hline $150\le v<155$&$18$&$18$\\ 
				\hline  $155\le v<160$&$28$&$46$\\
				\hline $160\le  v<165$&$35$&$81$\\
				\hline $165\le v<170$&$43$&$124$\\ 
				\hline  $170\le v<175$&$41$&$165$\\
				 \hline $175\le v<180$&$35$&$200$\\
				\hline 
			\end{tabular}
		\end{center}
		Ta có  $\dfrac{n}{4}=\dfrac{200}{4}=50$ nên $Q_1\in [160;165)$. Khi đó
		 $$Q_1=160+\dfrac{50-46}{35}\cdot 5=\dfrac{1124}{7}.$$
		Ta  có $\dfrac{3n}{4}=\dfrac{3\cdot 200}{4}=150$ nên $Q_3\in [170;175)$. Khi đó
		$$Q_3=170+\dfrac{150-124}{41}\cdot 5=\dfrac{7100}{41}.$$ 
	 }
\end{ex}

% nguồn: Ok_TuanLuan_Toan11_Lan2
\begin{ex}%[1D5H2-3]%[ma:0301c36d]
	Trong tuần lễ bảo vệ môi trường, các học sinh khối 12 tiến hành thu  nhặt vỏ lon nước ngọt để tái chế. Nhà trường thống kê kết quả thu nhặt vỏ lon nước ngọt của học sinh khối 12 ở bảng sau:\[
	\begin{array}{|c|c|c|c|c|c|}
		\hline \text { Số vỏ lon } & {[11; 15]} & {[16; 20]} & {[21;  25]} & {[26; 30]} & {[31; 35]} \\
		\hline \text { Số  học sinh } & 58 & 87 & 54 & 44 & 23 \\
		 \hline
	\end{array}
	\] 
	Hãy tìm trung vị của mẫu số liệu ghép  nhóm trên
	 \loigiai{
		 Do số vỏ chai là số nguyên nên ta hiệu chỉnh lại như sau:\[
		\begin{array}{|c|c|c|c|c|c|}
			\hline \text { Số vỏ lon } & {[10,5; 15,5)} & {[15,5; 20,5)} & {[20,5; 25,5)} & {[25,5; 30,5)} & {[30,5;  35,5)} \\
			\hline  \text { Số học sinh } & 58 & 87 & 54 & 44 & 23 \\
			\hline 
		\end{array}
		\] 
		Số học sinh tham gia thu nhặt  vỏ lon nước ngọt là $n=58+87+54+44+23=266$.\\
		Gọi $x_1$; $x_2$; $\ldots$; ${x_{266}}$ lần lượt là số vỏ chai $266$ học  sinh khối 12 thu nhặt được xếp theo thứ tự không giảm.\\
		Do $x_1$; $x_2$; $\ldots$; $x_{58}\in \left[10,5;15,5 \right)$; ${x_{59}}$; ${x_{55}}$; $\ldots$; ${x_{145}}\in  \left[15,5;20,5 \right)$ nên trung vị của mẫu số liệu $x_1$; $x_2$; $\ldots$; ${x_{240}}$ là $\dfrac{1}{2}\left( {x_{133}}+{x_{134}} \right)\in \left[15,5;20,5 \right)$.\\
		Ta xác định được $n=266$, $n_m=87$, $C=58$, $u_m=15,5$ và ${u_{m+1}}=20,5$.\\ 
		Trung vị  của mẫu số liệu ghép nhóm là\\
	$M_e=15,5+\dfrac{\dfrac{266}{2}-58}{87}\cdot \left( 20,5-15,5 \right)\approx  19,81$.}
\end{ex}

% nguồn: Ok_TuanLuan_Toan11_Lan2
\begin{ex}%[1D5H2-3]%[ma:500202e6]
	Lương  tháng của một số nhân viên văn phòng được ghi lại như sau (đơn vị: triệu đồng)
	\begin{longtable}{|*{13}{c|}}
		 \hline
		$12{,}5$ & $9{,}6$ & $11{,}7$  & $12{,}7$ & $10{,}0$ & $10{,}0$ & $12{,}2$ & $9{,}8$ & $10{,}9$ & $6{,}7$ & $13{,}6$ & $9{,}2$ \\
		\hline 
		$13{,}1$ & $6{,}5$ & $10{,}7$ & $8{,}9$ & $11{,}2$ & $13{,}2$ & $8{,}3$ & $11{,}1$ & $11{,}9$ & $8{,}4$ & $6{,}7$ & $13{,}8$ \\ 
		\hline 
	\end{longtable}
	\begin{enumerate}
		 \item Tìm tứ phân vị của dãy số liệu trên.
		\item Tổng hợp lại dãy số liệu trên vào bảng tần số ghép nhóm theo  mẫu sau
		\begin{longtable}{|*{5}{c|}}
			 \hline
			Lương tháng (triệu  đồng)
			& [$6$; $8$)  & [$8$; $10$) & [$10$; $12$) & [$12$; $14$) \\
			\hline 
			 Số nhân viên & ? & ? & ? & ? \\
			\hline 
		\end{longtable}
		\item Hãy ước lượng tứ phân vị của số liệu ở bảng tần số ghép  nhóm trên.
	\end{enumerate}
	\loigiai{ 
		\begin{enumerate}
			\item Sắp xếp lại dãy số liệu theo thứ tự không giảm ta được dãy $x_1, x_2,  \ldots, x_{24}$.\\
			Giá trị tứ phân vị thứ hai là $Q_2=\dfrac{1}{2} \left(x_{12}+x_{13} \right)=\dfrac{1}{2} \left(10{,}7+10{,}9\right)=10{,}8$.\\ 
			Giá trị tứ  phân vị thứ nhất là $Q_1=\dfrac{1}{2} \left(x_6+x_7 \right)=\dfrac{1}{2} \left(8{,}9+9{,}2\right)=9{,}05$.\\
			Giá trị tứ phân  vị thứ ba là $Q_3=\dfrac{1}{2} \left(x_{18}+x_{19} \right)=\dfrac{1}{2} \left(12{,}2+12{,}5\right)=12{,}35$.
			\item  Hoàn thành bảng sau
			\begin{longtable}{|*{5}{c|}}
				\hline 
				Lương  tháng (triệu đồng)
				& [$6$; $8$)  & [$8$; $10$) & [$10$; $12$) & [$12$; $14$) \\
				 \hline
				Số nhân viên & $8$ & $3$ & $6$ &  $7$\\
				 \hline
			\end{longtable}
			 \item Ta có $x_1, x_2, \ldots, x_3 \in \left[6; 8\right)$; $x_4, x_5, \ldots, x_9 \in \left[8; 10\right)$; $x_{10}, x_{12}, \ldots, x_{17} \in \left[10; 12\right)$; $x_{18}, x_{19}, \ldots, x_{24} \in \left[12; 14\right)$.\\
			Tứ phân vị thứ hai của dãy số liệu $x_1, x_2, \ldots, x_{24}$ là $x_{12}$ và  $x_{13}$.\\
			Vì $x_{12}, x_{13} \in \left[10; 12\right)$  nên tứ phân vị thứ hai của mẫu số liệu là
			\begin{align*}
				 Q_2=10+\dfrac{\dfrac{2\cdot24}{4}-\left(3+6\right)}{8} \left(12-10\right)=10{,}75.
			\end{align*}
			Tứ phân vị thứ nhất của dãy số liệu $x_1, x_2, \ldots, x_{24}$  là $x_6$ và $x_7$. \\
			Vì $x_6, x_7 \in \left[8; 10\right)$ nên tứ phân vị thứ  nhất của mẫu số liệu là
			\begin{align*}
				Q_1=8+\dfrac{\dfrac{1\cdot24}{4}-3}{6}  \left(10-8\right)=9.
			\end{align*}
			Tứ phân vị thứ ba của dãy số liệu  $x_1, x_2, \ldots, x_{24}$ là $x_{18}$ và $x_{19}$.\\
			Vì $x_{18}, x_{19} \in \left[12; 14\right)$ nên tứ phân  vị thứ ba của mẫu số liệu là
			\begin{align*}
				 Q_3=12+\dfrac{\dfrac{3\cdot24}{4}-\left(3+6+8\right)}{7} \left(14-12\right)\approx 12{,}28.
			\end{align*}
		\end{enumerate}
	} 
\end{ex}

% nguồn: Ok_TuanLuan_Toan11_Lan2
\begin{ex}%[1D5H2-3]%[ma:25cfb6ec]
	Ghi lại tốc độ bóng trong $200$ lần  giao bóng của một vận động viên môn quần vợt cho kết quả như bảng bên
	\begin{center}
		\begin{tabular}{|c|c|}
			 \hline
			\multicolumn{1}{|c|}{\textbf{Tốc độ $v$ (km/h)}}  & \multicolumn{1}{c|}{\textbf{Số lần}} \\
			\hline 
			\multicolumn{1}{|c|}{$150\leq v<155$} &  \multicolumn{1}{c|}{$18$} \\
			 \hline
			\multicolumn{1}{|c|}{$155\leq  v<160$} & \multicolumn{1}{c|}{$28$} \\
			\hline 
			\multicolumn{1}{|c|}{$160\leq v<165$} & \multicolumn{1}{c|}{$35$}  \\
			 \hline
			\multicolumn{1}{|c|}{$165\leq v<170$} & \multicolumn{1}{c|}{$43$} \\ 
			\hline 
			\multicolumn{1}{|c|}{$170\leq v<175$} & \multicolumn{1}{c|}{$41$} \\ 
			 \hline
			\multicolumn{1}{|c|}{$175\leq v<180$}  & \multicolumn{1}{c|}{$35$} \\
			 \hline
		\end{tabular}
	\end{center}
	Tìm tứ phân vị thứ nhất và tứ phân vị thứ ba cho mẫu số  liệu ghép nhóm
	 \loigiai{
		Cỡ mẫu là $n=200$.\\ 
		Tứ phân vị thứ nhất $Q_1$ là $\dfrac{x_{50}+x_{51}}{2}$. Do $x_{50}$, $x_{51}$ đều thuộc nhóm $\left[160;165\right.)$  nên nhóm này chứa $Q_1$.\\
		Do đó $p=3$; $a_3=160$; $m_3=35$;  $m_1+m_2=18+28=46$; $a_4-a_3=165-160=5$ và ta có \[Q_1=160+\dfrac{\dfrac{200}{4}-46}{35}\cdot5\approx 160{,}57.\]
		Tứ  phân vị thứ ba $Q_3$ là $\dfrac{x_{150}+x_{151}}{2}$. Do $x_{150}$, $x_{151}$ đều thuộc nhóm $[170;175)$ nên nhóm này chứa $Q_3$. Do đó $p=5$; $a_5=170$; $m_5=41$; $m_1+m_2+m_3+m_4=18+28+35+43=124$; $a_6-a_5=175-170=5$ và ta có \[Q_3=170+\dfrac{\dfrac{3.200}{4}-124}{41}\cdot5\approx 173{,}17.\]
	} 
\end{ex}

% nguồn: Ok_TuanLuan_Toan11_Lan2
\begin{ex}%[1D5H2-3]%[ma:88cc9855]
	Một công  ty bảo hiểm thống kê lại độ tuổi các khách hàng mua bảo hiểm xe ô tô ở bảng sau:\[
	\begin{array}{|c|c|c|c|c|c|c|}
		\hline \text { Độ tuổi } & {[25; 30)} & {[30; 35)} & {[35; 40)} & {[40; 45)} & {[45; 50)} & {[50; 55)}  \\
		\hline \text { Số khách hàng } & 25 & 38 & 62 & 42 &  37 & 29 \\
		\hline 
	\end{array}
	\] 
	Hãy ước lượng số trung bình, mốt và các tứ phân vị của mẫu số liệu ghép nhóm  trên
	 \loigiai{
	$\bar{x} \approx 39,97$;  $M_o=\dfrac{415}{11}$; $Q_1=\dfrac{275}{8}$; $Q_2=M_e=\dfrac{4875}{124}$, $Q_3=\dfrac{6815}{148}$.}
\end{ex}

% nguồn: Ok_TuanLuan_Toan11_Lan2
\begin{ex}%[1D5H2-3]%[ma:8099c2a1]
	Bảng sau thống kê doanh số bán hàng của các nhân viên  một trung tâm thương mại trong một ngày.\[
	\begin{array}{|c|c|c|c|c|c|}
		\hline \text { Doanh số (triệu đồng) }  & {[20; 30)} & {[30; 40)} & {[40; 50)} & {[50; 60)} & {[60; 70)} \\
		\hline \text { Sô nhân viên } & 4 & 8 &  12 & 7 & 5 \\
		 \hline
	\end{array}
	\] 
	a) Hãy ước lượng số trung bình, mốt và trung vị của mẫu số liệu ghép  nhóm trên.\\
	b) Trung tâm thương mại dự định sẽ thưởng cho ${25 \%}$ số nhân viên có doanh số bán hàng cao nhất. Theo mẫu số liệu trên, trung tâm thương  mại nên khen thưởng các nhân viên có doanh số bán hàng ít nhất là bao nhiêu?
	\loigiai{ 
		a) $\bar{x}=\dfrac{815}{18}$; $M_o=\dfrac{400}{9}$; $M_e=45$.\\ 
		b) Tứ phân vị thứ ba của mẫu số liệu  ghép nhóm là ${Q_3=\dfrac{380}{7} \approx 54,29}$.\\
	Do đó, trung tâm thương mại nên khen thưởng các nhân viên có doanh số bán hàng một ngày ít nhất là $54,29$ triệu đồng.} 
\end{ex}

% nguồn: Ok_TuanLuan_Toan11_Lan2
\begin{ex}%[1D5H2-3]%[ma:f9d628a3]
	Kết quả điều tra về số giờ làm thêm trong một tuần của 100 sinh viên  được cho ở biểu đồ bên.\\
	Hãy ước  lượng số trung bình, mốt và các tứ phân vị của số liệu đó.\\
	\begin{tikzpicture}
			\draw[->]  (0,0) -- (11,0) node[right] {(số giờ)};
			\draw[->] (0,0) --  (0,5.5) node[above] {(số sinh viên)};
			\node at (1.5,-0.3)  {[2; 4)};
			\node at (3.5,-0.3)  {[4; 6)};
			\node  at (5.5,-0.3) {[6; 8)};
			\node at (7.5,-0.3) {[8;  10)};
			\node  at (9.5,-0.3) {[10; 12)};
			\foreach \y/\label in  {0/0,1/10,2/20,3/30,4/40,5/50}{
				\draw  (-0.2,\y) -- (0,\y) node[left] {\label};
			 }
			\fill[orange]  (1,0) rectangle (2,1.2);
			\node at  (1.5,1.4) {12};
			\fill[orange] (3,0) rectangle (4,2); 
			 \node at (3.5,2.2) {20};
			\fill[orange]  (5,0) rectangle (6,3.7);
			\node at  (5.5,3.9) {37};
			\fill[orange] (7,0) rectangle  (8,2.1);
			\node at (7.5,2.3)  {21};
			\fill[orange] (9,0)  rectangle (10,1);
			\node at (9.5,1.2)  {10};
			\node at  (6,5.5) {Số giờ làm thêm trong một tuần của sinh viên};
		\end{tikzpicture}
	 \loigiai{
		Từ mẫu số liệu ghép nhóm, ta có bảng thống kê số giờ làm thêm trong một tuần của 100 sinh viên như sau:\[ 
		\begin{array}{|c|c|c|c|c|c|}
			\hline \text { Số giờ  làm thêm } & {[2; 4)} & {[4; 6)} & {[6; 8)} & {[8; 10)} & {[10; 12)} \\
			\hline \text  { Số giờ làm thêm đại diện } & 3 & 5 & 7 & 9 & 11 \\
			\hline \text { Số sinh viên } & 12 & 20 & 37 & 21  & 10 \\
			 \hline
		\end{array}
		 \]
		Cỡ mẫu  ${n=100}$.\\
		Số trung bình của mẫu số liệu ghép nhóm  trên là\\
		${\bar{x}=\dfrac{3 \cdot 12  + 5 \cdot 20 + 7 \cdot 37 + 9 \cdot 21 + 11 \cdot 10}{100}=6,94 .
		 }$\\
		Nhóm chứa mốt của mẫu số liệu trên là  nhóm ${[6; 8)}$.\\
		Do đó: $u_m=6$; $n_m=37$;  $n_{m-1}=20$; $n_{m+1}=21$; $u_{m+1}=8$.\\
		Vậy mốt của mẫu số liệu ghép  nhóm là\\
		${M_o=6+\dfrac{37-20}{(37-20)+(37-21)} \cdot(8-6) \approx  7,03}$.\\
		Gọi $x_1$; $x_2$; $\ldots$; $x_{100}$ là mẫu số liệu được xếp theo thứ tự không giảm.\\ 
		Tứ phân vị thứ hai của dãy  số liệu $x_1$; $x_2$; $\ldots$; $x_{100}$ là ${\dfrac{1}{2}\left(x_{50}+x_{51}\right)}$.\\
		 Do ${x_{50}}$ và ${x_{51}}$ thuộc nhóm ${[6; 8)}$ nên ${Q_2=6+\dfrac{\dfrac{2 \cdot 100}{4}-32}{37} \cdot(8-6) \approx 6,97}$.\\
		Tứ phân vị thứ  nhất của dãy số liệu $x_1$; $x_2$; $\ldots$; $x_{100}$ là ${\dfrac{1}{2}\left(x_{25}+x_{26}\right)}$. Do ${x_{25}}$ và ${x_{26}}$ thuộc nhóm ${[4; 6)}$ nên ${Q_1=4+\dfrac{\dfrac{1.100}{4}-12}{20} \cdot(6-4)=5,3}$.\\
	Tứ phân vị thứ ba của dãy số liệu $x_1$; $x_2$; $\ldots$; $x_{100}$ là  ${\dfrac{1}{2}\left(x_{75}+x_{76}\right)}$. Do ${x_{75}}$ và ${x_{76}}$ thuộc nhóm ${[8; 10)}$ nên ${Q_3=8+\dfrac{\dfrac{3 \cdot 100}{4}-69}{21} \cdot(10-8) \approx 8,57}$.}
\end{ex}

% nguồn: Ok_TuanLuan_Toan11_Lan2
\begin{ex}%[1D5H2-3]%[ma:be31f011]
	Một phòng khám thống kê số bệnh nhân đến khám bệnh mỗi ngày trong tháng 4 năm 2022 ở  bảng sau
	\begin{longtable}{|c|c|c|c|c|c|}
		\hline 
		Số bệnh nhân & [$1$; $10$] & [$11$; $20$] &  [$21$; $30$] & [$31$; $40$] & [$41$; $50$] \\
		\hline 
		Số ngày & $7$ & $8$  & $7$ & $6$ & $2$ \\
		\hline 
	\end{longtable}
	\begin{enumerate}
		\item  Hãy ước lượng trung vị và tứ phân vị của mẫu số liệu ghép nhóm trên.
		\item Quản lí phòng khám cho rằng có khoảng 25\% số ngày khám có nhiều hơn 35 bệnh nhân đến khám. Nhận định trên  có hợp lí không?
	\end{enumerate}
	\loigiai{ 
		\begin{enumerate}
			\item Do  số bệnh nhân đến khám là số nguyên nên ta hiệu chỉnh lại như sau
			\begin{longtable}{|c|c|c|c|c|c|}
				 \hline
				 \textbf{Số bệnh nhân} & [$0{,}5$; $10{,}5$) & [$10{,}5$; $20{,}5$) & [$20{,}5$; $30{,}5$) & [$30{,}5$; $40{,}5$) & [$40{,}5$; $50{,}5$) \\
				\hline 
				\textbf{Số ngày} & $7$ & $8$  & $7$ & $6$ & $2$ \\
				\hline 
			\end{longtable}
			Tổng  số ngày khám là $7+8+7+6+2=30$ (cuộc gọi).\\
			Gọi $x_1, x_2, \ldots, x_{30}$ là số  bệnh nhân đến khám mỗi ngày xếp theo thứ tự không giảm.\\
			Ta có $x_1, x_2, \ldots, x_7 \in \left[0{,}5; 10{,}5\right)$; $x_8, x_9, \ldots, x_{15} \in \left[10{,}5; 20{,}5\right)$; $x_{16}, x_{17}, \ldots, x_{22} \in  \left[20{,}5; 30{,}5\right)$; \\
			$x_{23}, x_{24}, \ldots, x_{28} \in \left[30{,}5; 40{,}5\right)$, \ldots\\ 
			Tứ phân vị thứ hai của dãy số liệu $x_1, x_2,  \ldots, x_{30}$ là $x_{15}$ và $x_{16}$.\\
			Vì $x_{15} \in \left[10{,}5; 20{,}5\right)$ và $x_{16} \in \left[20{,}5; 30{,}5\right)$ nên tứ phân vị thứ hai của mẫu số liệu là $Q_2=20{,}5$.\\ 
			 Tứ phân vị thứ nhất của dãy số liệu $x_1, x_2, \ldots, x_{30}$ là $x_8$.\\
			Vì $x_8 \in  \left[10{,}5; 20{,}5\right)$ nên tứ phân vị thứ nhất của mẫu số liệu là
			\begin{align*}
				 Q_1=10{,}5+\dfrac{\dfrac{1\cdot30}{4}-7}{8} \left(20{,}5-10{,}5\right)=11{,}125.
			\end{align*}
			Tứ phân vị thứ ba của dãy số liệu $x_1, x_2, \ldots, x_{30}$  là $x_{22}$.\\
			Vì $x_{22} \in \left[20{,}5; 30{,}5\right)$ nên tứ  phân vị thứ ba của mẫu số liệu là
			\begin{align*}
				Q_3=20{,}5+\dfrac{\dfrac{3\cdot 30}{4}-\left(8+10\right)}{7}  \left(30{,}5-20{,}5\right)\approx 26{,}9.
			\end{align*}
			\item Vì $Q_1$, $Q_2$, $Q_3$ đều nhỏ hơn $35$ nên nhận định của đề bài không hợp  lí.
		\end{enumerate}
	 }
\end{ex}

% nguồn: Ok_TuanLuan_Toan11_Lan2
\begin{ex}%[1D5H2-3]%[ma:08a4531e]
	Các  bạn học sinh một lớp thống kê số túi nhựa mà gia đình bạn đó sử dụng trong một tuần. Kết quả được tổng hợp lại ở bảng sau:\[
	\begin{array}{|c|c|c|c|c|c|}
		\hline \text { Số túi  } & {[5; 9]} & {[10; 14]} & {[15; 19]} & {[20; 24]} & {[25; 29]} \\
		\hline \text { Số gia đình } & 8 & 15 &  12 & 7 & 2 \\
		 \hline
	\end{array}
	\] 
	a) Hãy ước lượng số trung  bình và mốt của mẫu số liệu trên.\\
	b) Cô giáo dự định trao danh hiệu "Gia đình xanh" cho ${25 \%}$ gia  đình các bạn sử dụng ít túi nhựa nhất. Cô giáo nên trao danh hiệu cho các gia đình dùng không quá bao nhiêu túi nhựa?
	\loigiai{ 
		a) Số trung bình của mẫu số liệu ghép nhóm là  khoảng $14,73$.\\
		Mốt của mẫu  số liệu là $13$.\\
		b) Tứ phân  vị thứ nhất của mẫu số liệu ghép nhóm là ${Q_1=\dfrac{21}{2}=10,5}$.\\
	Do đó, cô giáo nên trao danh hiệu cho gia đình các bạn dùng không  quá $10$ túi nhựa.}
\end{ex}

% nguồn: Ok_TuanLuan_Toan11_Lan2
\begin{ex}%[1D5H2-3]%[ma:afdf4789]
	Cho mẫu số liệu ghép nhóm  về chiều cao của $21$ cây na giống.
	\begin{center}
		\begin{tabular}{|c|c|c|c|c|}
			\hline 
			Chiều cao (cm) & $[0;5)$ & $[5;10)$  & $[10;15)$ & $[15;20)$ \\
			 \hline
			Số cây  & $3$ & $8$ & $7$ & $3$ \\
			\hline 
		\end{tabular}
	\end{center}
	Hãy cho biết tứ phân vị thứ nhất  $Q_1$ và tứ phân vị thứ ba $Q_3$ thuộc nhóm nào.
	 \loigiai{
		Vì $n=21$ nên tứ phân vị thứ nhất là trung vị của dãy gồm $10$ số liệu đầu tiên và chính  là trung bình cộng của giá trị ở vị trí thứ $5$ và thứ $6$, do đó $Q_1=\dfrac{x_5+x_6}{2}$,\\
		Tứ phân vị thứ ba là trung vị  của dãy gồm $10$ số liệu nằm bên phải trung vị là dãy $x_{12}$, $x_{13}$, $\ldots$, $x_{21}$ nên $Q_3=\dfrac{x_{16}+x_{17}}{2}$. \\
		Ta có $3+8+7=18$, do đó  $x_{16}$, $x_{17}$ thuộc nhóm $[10;15)$ nên tứ phân vị thứ ba $Q_3$ thuộc nhóm $[10;15)$.
	 }
\end{ex}

% nguồn: Ok_TuanLuan_Toan11_Lan2
\begin{ex}%[1D5H2-3]%[ma:2cd59929]
	Một cửa hàng đã ghi lại số tiền bán xăng cho $35$ khách hàng đi xe máy. Mẫu số liệu gốc có dạng $x_1, x_2, \ldots, x_{35}$ trong đó $x_i$ là số tiền bán xăng  cho khách hàng thứ $i$. Vì một lí do nào đó, cửa hàng chỉ có mẫu số liệu ghép nhóm dạng sau:
	\begin{center}
		\begin{tabular}{|c|c|c|c|c|}
			\hline 
			Số tiền (nghìn đồng) & $[0;30)$ & $[30;60)$  & $[60;90)$ & $[90;120)$ \\
			\hline 
			Số khách hàng & $3$ & $15$ & $10$ & $7$ \\ 
			\hline 
		\end{tabular}\\
		Bảng $3.1$. Số tiền khách hàng  mua xăng
	\end{center}
	Hãy tính các số đặc trưng cho mẫu số liệu trong Bảng $3.1$ và giải  thích ý nghĩa của các giá trị thu được.
	\loigiai{ 
		\begin{enumerate}
			 \item Số trung bình\\
			Trong mỗi khoảng số tiền, giá trị đại diện là trung bình cộng của giá trị  hai đầu mút nên ta có bảng sau:
			\begin{center}
				\begin{tabular}{|c|c|c|c|c|}
					 \hline
					Số tiền (nghìn đồng) & $15$ & $45$ & $75$ & $105$ \\ 
					\hline 
					Số khách hàng &  $3$ & $15$ & $10$ & $7$ \\
					 \hline
				\end{tabular}
			\end{center}
			 Tổng số khách hàng là $n=35$.\\
			Số trung bình  là $\overline{x}=\dfrac{3\cdot15+15\cdot45+10\cdot75+7\cdot105}{35}=63$.\\
			Số trung bình của mẫu số liệu ghép nhóm xấp xỉ cho số trung bình của mẫu số liệu gốc. Từ  đó, ta thấy số tiền bán xăng trung bình của $35$ khách hàng xấp xỉ $63$ nghìn đồng và có thể dùng làm đại diện cho mẫu số liệu.
			 \item Số trung vị, tứ phân vị\\
			Cỡ mẫu là  $n=35$.\\
			Gọi $x_1, x_2, \ldots, x_{35}$ là số  tiền xăng của $35$ khách hàng và giả sử dãy này đã được sắp xếp theo thứ tự tăng dần. Khi đó, trung vị là $x_{18}$. Do $x_{18}$ thuộc nhóm $[30;60)$ nên nhóm này chứa trung vị. Do đó $p=2$; $a_2=30$; $m_2=15$; $m_1=3$; $a_3-a_2=60-30=30$ và ta có
			 \[M_e=30+\dfrac{\dfrac{35}{2}-3}{15}\cdot30=59.
			\] 
			Trung vị của mẫu số liệu ghép nhóm xấp xỉ cho trung vị của mẫu số liệu gốc, nó chia mẫu số liệu gốc thành  hai phần, mỗi phần chứa $50\%$ giá trị. Từ đó ta thấy trung vị của mẫu số liệu gốc xấp xỉ bằng $59$, giá trị này là ngưỡng để chia mẫu số liệu gốc thành $2$ phần.\\
			Tứ phân vị thứ nhất $Q_1$ là $x_9$. Do $x_9$ thuộc nhóm $[30;60)$ nên nhóm này chứa $Q_1$. Do  đó, $p=2$; $a_2=30$; $m_2=15$; $m_1=3$; $a_3-a_2=60-30=30$ và ta có \[Q_1=30+\dfrac{\dfrac{35}{4}-3}{15}\cdot30=41{,}5.\]\\
			Tứ phân vị thứ ba $Q_3$ là $x_{27}$. Do $x_{27}$ thuộc nhóm $[60;90)$ nên nhóm này chứa $Q_3$. Do đó  $p=3$; $a_3=60$; $m_3=10$; $m_1+m_2=3+15=18$; $a_4-a_3=90-60=30$ và ta có
			\[Q_3=60+\dfrac{\dfrac{3\cdot 35}{4}-18}{10}\cdot30=84{,}75. 
			\] 
			Tứ phân vị  thứ hai $Q_2=M_e=59$.
			Do đó, các tứ phân vị của mẫu số liệu gốc xấp xỉ là $Q_1=41{,}5$; $Q_2=59$ và $Q_3=84{,}75$, mỗi giá trị này là các ngưỡng  để chia mẫu số liệu thành 4 phần, mỗi phần chứa $25\%$ giá trị.
			 \item Mốt\\
			Tần số lớn nhất  là $15$ nên nhóm chứa mốt là nhóm $[30;60)$.\\
			 Ta có $j=2$, $a_2=30$, $m_2=15$, $m_1=3$, $m_3=10$, $h=30$.\\
			Do  đó $M_0=30+\dfrac{15-3}{(15-3)+(15-10)}\cdot30\approx 51{,}18$.\\
			Mốt của mẫu số liệu gốc xấp xỉ bằng $51{,}18$. Vậy số khách hàng mua xăng với giá  tiền khoảng $51{,}18$ nghìn đồng là nhiều nhất.
		\end{enumerate}
	 }
\end{ex}

% nguồn: Ok_TuanLuan_Toan11_Lan2
\begin{ex}%[1D5H2-3]%[ma:d16dc243]
	Cho mẫu số liệu ghép nhóm thống kê thời gian sử dụng điện thoại trước  khi ngủ (đơn vị: phút) của một người trong $120$ ngày như ở Bảng dưới đây. Xác định các số đặc trưng đo xu thế trung tâm cho mẫu số liệu đó (làm tròn các kết quả đến hàng phần mười).\[
	\begin{array}{|c|c|}
		\hline \text { Nhóm }  & \text { Tần số } \\
		 \hline[0; 4) & 13 \\
		 \hline[4; 8) & 29 \\
		\hline[8;  12) & 48 \\
		\hline[12; 16) & 22  \\
		\hline[16; 20) & 8  \\
		\hline & n=120  \\
		 \hline
	\end{array}
	\] 
	\loigiai{ 
		Số  trung bình cộng là:\\
		 ${\bar{x}=\dfrac{13 \cdot 2+29 \cdot 6+48 \cdot 10+22 \cdot 14+8 \cdot 18}{120} \approx 9,4 .
		 }$\\
		Bảng tần số ghép nhóm bao gồm cả  tần số tích luỹ được cho như ở Bảng dưới đây:\[
		\begin{array}{|c|c|c|}
			\hline \text { Nhóm } & \text { Tần số } & \text  { Tần số tích luỹ } \\
			\hline[0; 4) & 13 &  13 \\
			\hline[4; 8) & 29 &  42 \\
			\hline[8; 12) &  48 & 90 \\
			\hline[12; 16) &  22 & 112 \\
			\hline[16; 20) & 8 & 120 \\ 
			\hline & n=120 &  \\
			 \hline
		\end{array}
		\] 
		Ta có: $\dfrac{n}{2}=60$, $\dfrac{n}{4}=30$,  $\dfrac{3 n}{4}=90$.\\
		Vì ${42<60<90}$ nên nhóm $3$ là nhóm  đầu tiên có tần số tích luỹ lớn hơn hoặc bằng $60$.\\
		Suy ra trung vị là: ${M_e=8+\left(\dfrac{60-42}{48}\right) \cdot 4=9,5  }$.\\
		Tứ phân vị  thứ hai là: ${Q_2=M_e=9,5}$.\\
		Vì ${13<30<42}$ nên nhóm $2$ là nhóm đầu tiên có tần số tích luỹ lớn hơn hoặc bằng $30$. Suy ra tứ  phân vị thứ nhất là: ${Q_1=4+\left(\dfrac{30-13}{29}\right) \cdot 4 \approx 6,3}$.\\
		Vì ${42<90 \leq 90}$ nên nhóm $3$ là nhóm đầu tiên có tần số tích  luỹ lớn hơn hoặc bằng $90$. Suy ra tứ phân vị thứ ba là: ${Q_3=8+\left(\dfrac{90-42}{48}\right) \cdot 4=12
		 }$.\\
	Trong các nhóm, nhóm $3$ có tần số lớn nhất. Suy ra mốt là: ${M_0=8+\left(\dfrac{48-29}{2.48-29-22}\right) \cdot 4 \approx 9,7}$.} 
\end{ex}

% nguồn: Ok_TuanLuan_Toan11_Lan2
\begin{ex}%[1D5H2-3]%[ma:2eda7711]
	Một trang báo điện tử thống kê thời  gian người sử dụng đọc thông tintrên trang trong mỗi lần truy cập ở bảng sau:\[
	\begin{array}{|c|c|c|c|c|c|}
		\hline \text { Thời gian đọc (phút) } & {[0; 2)} & {[2; 4)} &  {[4; 6)} & {[6; 8)} & {[8; 10)} \\
		\hline \text { Số luợt truy cập } & 45 & 34 & 23 & 18 & 5 \\ 
		 \hline
	\end{array}
	\] 
	Hãy ước lượng các tứ  phân vị của mẫu số liệu ghép nhóm trên
	\loigiai{ 
	$Q_1=\dfrac{25}{18}$; $Q_2=\dfrac{103}{34}$; $Q_3=\dfrac{243}{46}$.} 
\end{ex}

% nguồn: Ok_TuanLuan_Toan11_Lan2
\begin{ex}%[1D5H2-3]%[ma:0bcd5c07]
	Kết quả khảo sát cân nặng  của 1 thùng táo ở một lô hàng cho trong bảng sau:\[
	\begin{array}{|c|c|c|c|c|c|}
		\hline \text { Cân nặng (g) } & {[150; 155)} & {[155; 160)} & {[160;  165)} & {[165; 170)} & {[170; 175)} \\
		\hline \text { Số quả táo } &  4 & 7 & 12 & 6 & 2 \\
		\hline 
	\end{array}
	 \]
	 Hãy tìm trung vị của mẫu số liệu ghép nhóm trên
	 \loigiai{
		Gọi $x_1$; $x_2$; $\ldots$; ${x_{31}}$ là cân nặng của $31$  quả bơ xếp theo thứ tự không giảm.\\
		Do $x_1$; $x_2$;  $\ldots$; $x_4\in \left[150;155 \right)$; $x_5$; $x_6$; $\ldots$; ${x_{11}}\in \left[155;160 \right)$ nên trung vị của mẫu số liệu $x_1$; $x_2$; $\ldots$; ${x_{31}}$ là ${x_{16}}\in \left[160;165 \right)$.\\
		Ta  xác định được $n=31$, $n_m=12$, $C=4+7=11$, $u_m=160$, ${u_{m+1}}=165$.\\
		 Vậy trung vị của mẫu số liệu ghép nhóm là\\
	$M_e=160+\dfrac{\dfrac{31}{2}-11}{12}\cdot  \left( 165-160 \right)=161,875$.}
\end{ex}

% nguồn: Ok_TuanLuan_Toan11_Lan2
\begin{ex}%[1D5H2-3]%[ma:6c4fa847]
	Người ta thống kê tốc độ của một số  xe ô tô di chuyển qua một trạm kiểm soát trên đường cao tốc trong một khoảng thời gian ở bảng sau:\[
	\begin{array}{|c|c|c|c|c|c|}
		\hline \text { Tốc độ }(k m / h) & {[75;  80)} & {[80; 85)} & {[85; 90)} & {[90; 95)} & {[95; 100)} \\
		\hline \text { Số xe } & 5 & 12 & 18 & 24 & 19 \\ 
		\hline 
	\end{array}
	 \]
	Hãy ước lượng các tứ phân vị của mẫu số liệu ghép nhóm trên 
	\loigiai{ 
	$Q_1=\dfrac{3085}{36}$;  $Q_2=\dfrac{545}{6}$; $Q_3=\dfrac{4555}{48}$.}
\end{ex}

% nguồn: Ok_TuanLuan_Toan11_Lan2
\begin{ex}%[1D5H2-3]%[ma:03e46173]
	Điểm thi môn Toán (thang điểm $100$, điểm được làm tròn đến $1$) của $60$ thí  sinh được cho trong bảng sau:
	\begin{center}
		\begin{tabular}{|c|c|c|c|c|c|}
			 \hline
			Điểm  & $0-9$ & $10-19$ & $20-29$ & $30-39$ & $40-49$ \\
			 \hline
			Số thí sinh & $1$ & $2$ & $4$  & $6$ & $15$ \\
			 \hline
			Điểm & $50-59$ & $60-69$ & $70-79$ & $80-89$ & $90-99$  \\
			\hline 
			 Số thí sinh & $12$ & $10$ & $6$ & $3$ & $1$ \\
			\hline 
		\end{tabular}
	\end{center}
	\begin{enumerate}[a)]
		\item Hiệu chỉnh để  thu được mẫu số liệu ghép nhóm.
		\item Tìm các tứ phân vị và giải thích  ý nghĩa của chúng.
	\end{enumerate}
	\loigiai{ 
		\begin{enumerate}[a)]
			\item \,\\ 
			\begin{center}
				\begin{tabular}{|c|c|c|c|c|c|}
					\hline 
					 Điểm & $[0;9{,}5)$ & $[9{,}5;19{,}5)$ & $[19{,}5;29{,}5)$ & $[29{,}5;39{,}5)$ & $[39{,}5;49{,}5)$ \\
					\hline 
					Số thí sinh  & $1$ & $2$ & $4$ & $6$ & $15$ \\
					 \hline
					Điểm & $[49{,}5;59{,}5)$ &  $[59{,}5;69{,}5)$ & $[69{,}5;79{,}5)$ & $[79{,}5;89{,}5)$ & $[89{,}5;99{,}5)$ \\
					\hline 
					Số thí sinh & $12$  & $10$ & $6$ & $3$ & $1$ \\
					 \hline
				\end{tabular}
			\end{center}
			\item Cỡ mẫu là  $n=60$.
		\end{enumerate}
		Gọi  $x_1, x_2, \ldots, x_{60}$ là điểm thi môn Toán của 60 thí sinh và giả sử dãy này đã được sắp xếp theo thứ tự tăng dần. Khi đó, trung vị là $\dfrac{x_{30}+x_{31}}{2}$. Do hai giá trị $x_{30}$, $x_{31}$ thuộc nhóm $[49,5;59,5)$ nên nhóm này chứa trung vị.\\
		Do đó $p=6$; $a_6=49,5$; $m_6=12$; $m_1+m_2+m_3+m_4+m_5=1+2+4+6+15=28$; $a_7-a_6=59$, $5-49,5=10$ và ta có  \[M_e=49{,}5+\dfrac{\dfrac{60}{2}-28}{12}\cdot10\approx 51{,}17.\]
		Tứ phân vị thứ nhất $Q_1$ là $\dfrac{x_{15}+x_{16}}{2}$. Do $x_{15}$ và $x_{16}$ đều thuộc nhóm $[39,5;49,5)$ nên nhóm này chứa $Q_1$. Do đó $p=5$; $a_5=39,5$; $m_5=15$; $m_1+m_2+m_3+m_4=13$; $a_6-a_5=10$ và  ta có \[Q_1=39,5+\dfrac{\dfrac{60}{4}-13}{15}\cdot10\approx 40{,}83.\]
		Tứ phân vị thứ ba $Q_3$ là $\dfrac{x_{45}+x_{46}}{2}$. Do $x_{45}$ và $x_{46}$ đều thuộc nhóm $[59,5;69,5)$ nên nhóm này chứa $Q_3$. Do  đó $p=7$; $a_7=59,5$; $m_7=10$; $m_1+m_2+m_3+m_4+m_5+m_6=40$; $a_6-a_5=10$ và ta có \[Q_3=59,5+\dfrac{\dfrac{3\cdot60}{4}-40}{10}\cdot10=64{,}5.\]
		Tứ phân  vị thứ hai $Q_2=M_e \approx 51{,}17.$
		Vậy các tứ phân vị của mẫu số liệu là $Q_1 \approx 40{,}83$; $Q_2 \approx 51{,}17$ và $Q_3=64{,}5$.  Các giá trị này các là ngưỡng để phân điểm của $60$ học sinh thành $4$ phần để xếp loại học sinh.
	} 
\end{ex}

% nguồn: Ok_TuanLuan_Toan11_Lan2
\begin{ex}%[1D5H2-3]%[ma:7d28d1eb]
	Xác định trung vị và tứ phân vị của mẫu số liệu ghép nhóm được cho ở Bảng dưới  đây (làm tròn các kết quả đến hàng phần mười).\[
	\begin{array}{|c|c|c|}
		\hline \text { Nhóm } & \text {  Tần số } & \text { Tần số tích luỹ } \\
		\hline[40; 45)  & 5 & 5 \\
		\hline[45; 50) &  10 & 15 \\
		\hline[50; 55) & 7 & 22  \\
		\hline[55; 60) & 9 &  31 \\
		\hline[60; 65) & 7 & 38  \\
		 \hline[65; 70) & 4 & 42 \\
		\hline  & n=42 & \\
		\hline 
	\end{array}
	 \]
	\loigiai{ 
		 Số phần tử của mẫu là ${n=42}$.\\
		- Ta có: ${\dfrac{n}{2}=\dfrac{42}{2}=21}$ mà ${15<21<22}$. Suy ra nhóm 3 là nhóm đầu tiên có tần số tích luỹ lớn hơn hoặc bằng  $21$.\\
		Xét nhóm 3 là nhóm ${[50; 55)}$ có  ${r=50, d=5, n_3=7}$ và nhóm 2 là nhóm ${[45; 50)}$ có ${c f_2=15}$.\\
		Áp dụng công thức, ta có trung vị của mẫu số liệu là: $M_e=50+\left(\dfrac{21-15}{7}\right)  \cdot 5 \approx 54,3$ (km/h).\\
		Tứ phân vị thứ hai của mẫu  số liệu là: $Q_2=M_e \approx 54,3$ (km/h).\\
		- Ta có: ${\dfrac{n}{4}=\dfrac{42}{4}=10,5}$ mà ${5<10,5<15}$ nên nhóm 2 là nhóm đầu tiên có tần số tích luỹ lớn  hơn hoặc bằng 10,5.\\
		Xét nhóm 2 là nhóm ${[45; 50)}$ có ${s=45, h=5,  n_2=10}$ và nhóm 1 là nhóm ${[40; 45)}$ có ${c f_1=5}$.\\
		Áp dụng công thức, ta có  tứ phân vị thứ nhất của mẫu số liệu là:\\
		$Q_1=45+\left(\dfrac{10,5-5}{10}\right) \cdot 5 \approx 47,8$  (km/h).\\
		- Ta có: ${\dfrac{3 n}{4}=\dfrac{3 \cdot 42}{4}=31,5}$ mà ${31<31,5<38}$ nên nhóm 5 là  nhóm đầu tiên có tần số tích luỹ lớn hơn hoặc bằng $31,5$.\\
		Xét nhóm 5 là nhóm ${[60; 65)}$ có ${t=60, l=5, n_5=7}$ và nhóm 4 là nhóm ${[55; 60)}$ có ${c  f_4=31}$.\\
		Áp dụng công thức, ta  có tứ phân vị thứ ba của mẫu số liệu là:\\
		$Q_3=60+\left(\dfrac{31,5-31}{7}\right)  \cdot 5 \approx 60,4$ (km/h).\\
		Vậy  tứ phân vị của mẫu số liệu trên là: $Q_1 \approx 47,8({~km} / {h}); Q_2 \approx 54,3({~km} / {h}); Q_3 \approx 60,4$ (km/h).
	 }
\end{ex}

% nguồn: Ok_TuanLuan_Toan11_Lan2
\begin{ex}%[1D5H2-3]%[ma:66956915]
	Thầy giáo thống kê lại số lần kéo xà đơn của các học sinh  nam khối 11 ở bảng sau:\[
	\begin{array}{|c|c|c|c|c|c|}
		\hline \text { Số lần } &  {[6; 10]} & {[11; 15]} & {[16; 20]} & {[21; 25]} & {[26; 30]} \\
		\hline  \text { Số học sinh } & 35 & 54 & 32 & 17 & 5 \\
		\hline 
	\end{array}
	 \]
	a) Hãy ước lượng  số trung bình, mốt và trung vị của mẫu số liệu ghép nhóm trên.\\
	b) Thầy giáo dự định chọn ${25 \%}$ học sinh có  số lần kéo thấp nhất để bồi dưỡng thể lực thêm. Thầy giáo nên chọn học sinh có thành tích kéo xà đơn dưới bao nhiêu lần để bồi dưỡng thể lực?
	 \loigiai{
		a) $\bar{x}=\dfrac{2089}{143}$; $M_o=\dfrac{1051}{82}$; $M_e=\dfrac{1499}{108}$.\\ 
	b) ${Q_1=\dfrac{797}{72} \approx 11,07}$. Thầy giáo nên chọn các bạn có thành tích kéo xà dưới $12$ lần để  bồi dưỡng thể lực thêm.}
\end{ex}

% nguồn: Ok_TuanLuan_Toan11_Lan2
\begin{ex}%[1D5H2-3]%[ma:79f14a46]
	Thâm niên công tác  của các công nhân hai nhà máy ${A}$ và ${B}$.\[
	\begin{array}{|l|c|c|c|c|c|}
		\hline \text { Thâm niên công tác (năm) } & {[0; 5)} &  {[5; 10)} & {[10; 15)} & {[15; 20)} & {[20; 25)} \\
		\hline \text { Số công nhân nhà  máy } A & 35 & 13 & 12 & 12 & 8 \\
		\hline \text  { Số công nhân nhà máy } B & 14 & 26 & 24 & 11 & 5 \\
		 \hline
	\end{array}
	\] 
	a) Hãy so  sánh thâm niên công tác của nhân viên hai nhà máy theo số trung bình và trung vị.\\
	 b) Hãy ước lượng tứ phân vị thứ nhất và thứ ba của hai mẫu số liệu ghép nhóm trên
	 \loigiai{
		a) So sánh theo số trung bình: $\bar x_A=9,0625$; $\bar x_B=10,4375$, suy ra ${\bar x_A<\bar  x_B}$.\\
		So sánh theo trung vị:  $M_e(A)=\dfrac{90}{13}$; $M_e(B)=10$; $M_e(A)<M_e(B)$.\\
	b) $Q_1(A)=\dfrac{20}{7}$; $Q_3(A)=15$; $Q_1(B)=\dfrac{80}{13}$; $Q_3(B)=\dfrac{85}{6}$.} 
\end{ex}

% nguồn: Ok_TuanLuan_Toan11_Lan2
\begin{ex}%[1D5H2-3]%[ma:0f2104b8]
	Kết quả kiểm tra cân nặng của một số quả trứng chim cút được lựa chọn ngẫu nhiên ở hai trang trại chăn  nuôi ${A}$ và ${B}$ được biểu diễn ở biểu đồ sau (đơn vị: g).\\
	\begin{tikzpicture}
			\draw[->] (0,0) --  (11,0) node[right] {(g)};
			\draw[->] (0,0) --  (0,5.5) node[above] {Số quả};
			\node  at (1.5,-0.3) {[8,2; 8,4]};
			 \node at (3.5,-0.3) {[8,4; 8,6]};
			\node at (5.5,-0.3) {[8,6; 8,8]}; 
			\node at  (7.5,-0.3) {[8,8; 9,0]};
			\node at (9.5,-0.3) {[9,0;  9,2]};
			\foreach \y/\label  in {0/0,1/10,2/20,3/30,4/40}{
				\draw (-0.2,\y) -- (0,\y) node[left] {\label}; 
			} 
			\fill[gray] (0.5,0) rectangle  (1.5,0.7);
			\node at (1,0.9) {7}; 
			\fill[gray]  (2.5,0) rectangle (3.5,1.8);
			\node  at (3,2) {18};
			 \fill[gray] (4.5,0) rectangle (5.5,3.4);
			\node at (5,3.6) {34}; 
			\fill[gray] (6.5,0)  rectangle (7.5,2.1);
			 \node at (7,2.3) {21};
			\fill[gray] (8.5,0) rectangle  (9.5,0.9);
			 \node at (9,1.1) {9};
			 \fill[orange] (1.5,0) rectangle (2.5,1.5);
			\node at (2,1.7) {15}; 
			\fill[orange] (3.5,0) rectangle  (4.5,3.7);
			 \node at (4,3.9) {37};
			\fill[orange] (5.5,0) rectangle  (6.5,1.2);
			\node at  (6,1.5) {12};
			 \fill[orange] (7.5,0) rectangle (8.5,0.7);
			\node  at (8,0.9) {7};
			\fill[orange] (9.5,0) rectangle  (10.5,0.2);
			\node at  (10,0.4) {2};
			 \draw[gray,fill=gray] (0.5,5) rectangle (1,4.7);
			\node at (3.2,4.85) {Trứng  của trang trại A};
			\draw[orange,fill=orange]  (5.9,5) rectangle (6.4,4.7);
			 \node at (8.6,4.85) {Trứng của trang trại B};
		\end{tikzpicture}
	\\ 
	a) Hãy so sánh cân nặng của trứng chim cút của hai trang trại ${A}$ và ${B}$ theo  số trung bình và trung vị.\\
	b) Hãy ước lượng tứ phân vị thứ nhất và tứ phân vị thứ ba của cân nặng trứng chim cút của trang  trại ${A}$.
	 \loigiai{
		a) Từ biểu  đồ đã cho, ta lập được bảng số liệu ghép nhóm sau:\[
		\begin{array}{|c|c|c|c|c|c|}
			\hline \text { Cân nặng (gam) } & {[8,2;  8,4)} & {[8,4; 8,6)} & {[8,6; 8,8)} & {[8,8; 9,0)} & {[9,0; 9,2)} \\
			\hline \text { Số quả trứng của trại } A & 7 & 18 & 34 & 21 & 9 \\ 
			\hline \text { Số quả trứng của trại } B &  15 & 37 & 12 & 7 & 2 \\
			\hline 
		\end{array}
		 \]
		Từ đó, ta có bảng thống kê số quả trứng chim cút của hai trang trại theo giá trị đại diện như sau:\[\begin{array}{|c|c|c|c|c|c|}
		\hline \text { Cân nặng đại diện (gam)  } & 8,3 & 8,5 & 8,7 & 8,9 & 9,1 \\
		\hline \text { Số quả trứng của trang trai } A &  7 & 18 & 34 & 21 & 9 \\
		\hline \text { Số  quả trứng của trang trại } B & 15 & 37 & 12 & 7 & 2 \\
		\hline 
	\end{array}\]
	- Đối với trang trại ${A}$:  Cỡ mẫu ${n_A=89}$.\\
	Cân nặng trung bình của mỗi quả trứng của mẫu số liệu ghép nhóm  là\\
	$\bar x_A=\dfrac{8,3 \cdot 7+8,5 \cdot18+8,7 \cdot34+8,9 \cdot21+9,1 \cdot9}{89} \approx 8,72$.\\ 
	Trung vị của dãy số liệu là số liệu thứ $45$ theo thứ tự  không giảm của dãy và thuộc nhóm ${[8,6; 8,8)}$.\\
	Trung vị của mẫu số liệu  ghép nhóm là ${M_e(A)=8,6+\dfrac{\dfrac{89}{2}-25}{34} \cdot(8,8-8,6) \approx 8,71}$.\\
	-  Đối với trang trại ${B}$:\\
	Cỡ mẫu ${n_B=73}$.\\ 
	Cân nặng trung bình của mỗi quả  trứng của mẫu số liệu ghép nhóm là\\
	${\bar x_B=\dfrac{8,3 \cdot15+8,5  \cdot37+8,7 \cdot12+8,9 \cdot7+9,1 \cdot2}{73} \approx 8,55}$.\\
	Trung vị của dãy số liệu là  số liệu thứ 37 theo thứ tự không giảm của dãy và thuộc nhóm ${[8,4; 8,6)}$.\\
	Trung vị của mẫu số liệu ghép nhóm là  ${M_e(B)=8,4+\dfrac{\dfrac{73}{2}-15}{37} \cdot(8,6-8,4) \approx 8,52}$.\\
	Ta thấy ${\bar x_A>\bar x_B}$ và ${M_e(A)>M_e(B)}$. Vậy khi so  sánh theo số trung bình hay theo trung vị, cân nặng của trứng chim cút của trang trại ${A}$ đều lớn hơn cân nặng của trứng chim cút của trang trại ${B}$.\\
	b) Đối với  trang trại ${A}$:\\
	Tứ phân vị thứ nhất của mẫu số liệu là trung bình cộng của số liệu thứ 22 và 23 theo thứ tự không giảm của dãy và thuộc nhóm ${[8,4; 8,6)}$, nên tứ phân  vị thứ nhất của mẫu số liệu ghép nhóm là\\
	${Q_1(A)=8,4+\dfrac{\dfrac{1.89}{4}-7}{18}  \cdot(8,6-8,4) \approx 8,57}$.\\
Tứ phân vị thứ ba của mẫu số liệu ghép nhóm là trung bình cộng của số liệu thứ $67$ và $68$ theo thứ tự không giảm của dãy và thuộc nhóm ${[8,8; 9,0)}$, nên tứ phân vị thứ ba của mẫu số liệu ghép nhóm là ${Q_3(A)=8,8+\dfrac{\dfrac{3.89}{4}-59}{21} \cdot(9,0-8,8)  \approx 8,87}$.}
\end{ex}

% nguồn: Ok_TuanLuan_Toan11_Lan2
\begin{ex}%[1D5H2-3]%[ma:5a416fe7]
	 Trong một hội thao, thời gian chạy $200$ m của một nhóm vận động viên được ghi lại ở bảng sau
	\begin{longtable}{|*{6}{c|}}
		\hline 
		Thời gian (giây) & [$21$; $21{,}5$) &  [$21{,}5$; $22$) & [$22$; $22{,}5$) & [$22{,}5$; $23$) & [$23$; $23{,}5$) \\
		\hline 
		Số vận động  viên & $5$ & $12$ & $32$ & $45$ & $30$ \\
		 \hline
	\end{longtable}
	Dựa vào bảng số liệu trên, ban tổ chứ muốn chọn ra khoảng $50\%$ số vận động viên chạy nhanh nhất để tiếp tục thi vòng $2$. Ban tổ chức nên chọn các vận  động viên có thời gian chạy không quá bao nhiêu giây?
	 \loigiai{
		Gọi $x_1, x_2, \ldots, x_{124}$ là thời gian chạy của $124$  vận động viên xếp theo thứ tự không giảm.\\
		Số trung  vị của mẫu số liệu trên là $\dfrac{1}{2} \left(x_{62}+x_{63} \right)$\\
		Ta có $x_1, x_2, \ldots, x_5 \in \left[21; 21{,}5\right)$; $x_6,  x_7, \ldots, x_{17} \in \left[21{,}5; 22\right)$; $x_{18}, x_{19}, \ldots, x_{49} \in \left[22; 22{,}5\right)$; \\
		$x_{50}, x_{51},  \ldots, x_{94} \in \left[22{,}5; 23\right)$; $x_{95}, x_{96}, \ldots, x_{124} \in \left[23; 23{,}5\right)$.\\
		Nhận thấy $x_{62}$, $x_{63}$ thuộc $\left[22{,}5; 23\right)$.\\ 
		Do đó số trung vị của mẫu số liệu trên là $M_e=22{,}5+\dfrac{\dfrac{124}{2}-\left(5+12+32\right)}{45} \left(23-22{,}5\right)\approx  22{,}64$.\\
		Vậy ban tổ chức nên chọn  vận động viên có thời gian chạy không quá $22{,}64$ (giây).
	} 
\end{ex}

% nguồn: Ok_TuanLuan_Toan11_Lan2
\begin{ex}%[1D5H2-3]%[ma:96da8261]
	Số điểm một cầu thủ bóng rổ ghi được trong $20$ trận đấu được cho ở bảng sau 
	\begin{longtable}{|*{10}{c|}}
		\hline 
		$25$ & $24$  & $21$ & $13$ & $8$ & $14$ & $15$ & $18$ & $22$ & $11$ \\
		 \hline
		$24$ & $12$ & $14$ & $14$ & $18$ & $6$ & $8$ & $25$ &  $10$ & $11$ \\
		\hline 
	\end{longtable}
	\begin{enumerate}
		\item Tìm tứ phân vị của dãy số liệu  trên.
		\item  Tổng hợp lại dãy số liệu trên vào bảng tần số ghép nhóm theo mẫu sau
		\begin{longtable}{|*{5}{c|}}
			\hline 
			Điểm số & [$6$;  $10$] & [$11$; $15$] & [$16$; $20$] & [$21$; $25$] \\
			 \hline
			Số trận & ? &  ? & ? & ? \\
			\hline 
		\end{longtable}
		\item Hãy ước lượng  tứ phân vị của số liệu ở bảng tần số ghép nhóm trên.
	\end{enumerate}
	 \loigiai{
		\begin{enumerate}
			\item Sắp xếp lại dãy số liệu theo thứ  tự không giảm ta được dãy $x_1, x_2, \ldots, x_{20}$.\\
			Giá trị  tứ phân vị thứ hai là $Q_2=\dfrac{1}{2} \left(x_{10}+x_{11} \right)=\dfrac{1}{2} \left(14+14\right)=14$.\\
			Giá  trị tứ phân vị thứ nhất là $Q_1=\dfrac{1}{2} \left(x_5+x_6 \right)=\dfrac{1}{2} \left(11+11\right)=11$.\\
			Giá trị tứ phân vị  thứ ba là $Q_3=\dfrac{1}{2} \left(x_{15}+x_{16} \right)=\dfrac{1}{2} \left(21+22\right)=21{,}5$.
			\item Hoàn thành bảng  sau
			\begin{longtable}{|*{5}{c|}}
				 \hline
				Điểm  số & [$6$; $10$] & [$11$; $15$] & [$16$; $20$] & [$21$; $25$] \\
				 \hline
				Số trận &  $8$ & $4$ & $2$ & $6$ \\
				 \hline
			\end{longtable}
			\item Vì điểm số là số nguyên ta hiệu chỉnh bảng tần  số ghép nhóm như sau\\
			Ta có $x_1, x_2, \ldots, x_4 \in \left[5{,}5; 10{,}5\right)$; $x_5, x_6, \ldots, x_{12}  \in \left[10{,}5; 15{,}5\right)$; $x_{13}, x_{14} \in \left[15{,}5; 20{,}5\right)$;\\
			$x_{15}, x_{16}, \ldots, x_{20}  \in \left[20{,}5; 25{,}5\right)$\\
			Tứ phân vị thứ  hai của dãy số liệu $x_1, x_2, \ldots, x_{20}$ là $x_{10}$ và $x_{11}$.\\
			 Vì $x_{10}, x_{11} \in \left[10{,}5; 15{,}5\right)$ nên tứ phân vị thứ hai của mẫu số liệu là
			\begin{align*}
				Q_2=10{,}5+\dfrac{\dfrac{2\cdot20}{4}-4}{8} \left(15{,}5-10{,}5\right)=14{,}25. 
			\end{align*}
			Tứ phân vị thứ nhất của dãy số liệu $x_1, x_2, \ldots, x_{20}$  là $x_5$ và $x_6$.\\
			Vì $x_5, x_6 \in  \left[10{,}5; 15{,}5\right)$ nên tứ phân vị thứ nhất của mẫu số liệu là
			\begin{align*}
				 Q_1=10{,}5+\dfrac{\dfrac{1\cdot20}{4}-4}{8} \left(15{,}5-10{,}5\right)=11{,}125.
			\end{align*}
			Tứ phân vị thứ ba của dãy số liệu  $x_1, x_2, \ldots, x_{20}$ là $x_{15}$ và $x_{16}$.\\
			Vì  $x_{16}, x_{16} \in \left[20{,}5; 25{,}5\right)$ nên tứ phân vị thứ ba của mẫu số liệu là
			\begin{align*}
				Q_3=20{,}5+\dfrac{\dfrac{3\cdot20}{4}-\left(4+8+2\right)}{6}  \left(25{,}5-20{,}5\right)\approx 21{,}33.
			\end{align*}
		\end{enumerate}
	} 
\end{ex}

% nguồn: Ok_TuanLuan_Toan11_Lan2
\begin{ex}%[1D5H2-3]%[ma:4a851a58]
	 Một người thống kê lại thời gian thực hiện các cuộc gọi điện thoại của người đó trong một tuần ở bảng sau
	\begin{longtable}{|*{7}{c|}}
		\hline 
		 Thời gian (giây)
		& [$0$;  $60$) & [$60$; $120$) & [$120$; $180$) & [$180$; $240$) & [$240$; $300$) & [$300$; $360$) \\
		 \hline
		Số cuộc  gọi & $8$ & $10$ & $7$ & $5$ & $2$ & $1$ \\
		 \hline
	\end{longtable}
	Hãy ước lượng các tứ phân vị của  mẫu số liệu ghép nhóm trên.
	 \loigiai{
		Tổng số cuộc gọi là $8+10+7+5+2+1=33$  (cuộc gọi).\\
		Gọi $x_1,  x_2, \ldots, x_{33}$ là thời gian thực hiện cuộc gọi xếp theo thứ tự không giảm.\\
		Ta có $x_1, x_2, \ldots, x_8 \in \left[0; 60\right)$; $x_9, x_{10}, \ldots, x_{18} \in \left[60; 120\right)$; $x_{19}, x_{20}, \ldots, x_{25} \in \left[120;  180\right)$; \\
		 $x_{26}, x_{27}, \ldots, x_{30} \in \left[180; 240\right)$.\\
		Tứ phân vị thứ hai của dãy số liệu $x_1,  x_2, \ldots, x_{33}$ là $x_{17}$.\\
		Vì $x_{17} \in \left[60; 120\right)$ nên tứ phân vị thứ hai của  mẫu số liệu là $Q_2=60+\dfrac{\dfrac{2\cdot 33}{4}-8}{10} \left(120-60\right)=111$.\\
		Tứ phân vị thứ nhất của dãy số liệu $x_1, x_2, \ldots, x_{33}$ là  $x_9$.\\
		Vì $x_9 \in \left[60; 120\right)$ nên  tứ phân vị thứ nhất của mẫu số liệu là $Q_1=60+\dfrac{\dfrac{1\cdot 33}{4}-8}{10} \left(120-60\right)=61{,}5$.\\
		Tứ phân vị thứ ba của dãy số liệu $x_1, x_2, \ldots, x_{33}$  là $x_{26}$.\\
		Vì $x_{26} \in  \left[180; 240\right)$ nên tứ phân vị thứ ba của mẫu số liệu là
		\begin{align*}
			Q_3=180+\dfrac{\dfrac{3\cdot 33}{4}-\left(8+10+7\right)}{5} \left(240-180\right)=177. 
		\end{align*}
	 }
\end{ex}

% nguồn: Ok_TuanLuan_Toan11_Lan2
\begin{ex}%[1D5H2-3]%[ma:5beea523]
	Một cửa hàng sách thống kê  số truyện thiếu nhi bán được trong hai tháng ở bảng sau:\[
	\begin{array}{|l|c|c|c|c|c|}
		\hline \text { Số sách } & {[14; 20]} & {[21; 27]} & {[28; 34]} & {[35; 41]}  & {[42; 48]} \\
		\hline \text { Số ngày } & 5 & 7 &  25 & 15 & 9 \\
		\hline 
	\end{array}
	\] 
	Hãy ước lượng số trung bình, mốt và các  tứ phân vị của mẫu số liệu ghép nhóm trên
	 \loigiai{
	$\bar{x}=\dfrac{2003}{61}$; $M_o=32$; $Q_1=28,41$; $Q_2=M_e=32,68$; $Q_3=\dfrac{4967}{90}$.} 
\end{ex}

% nguồn: Ok_TuanLuan_Toan11_Lan2
\begin{ex}%[1D5H2-3]%[ma:7402a0aa]
	Tiền lương nhận được trong 1 giờ làm việc của nhân viên  công ty A được thống kê theo mẫu số liệu ghép nhóm sau (đơn vị: ngàn đồng):
	 \[
	\begin{array}{|c|c|}
		\hline \text { Các lớp tiền lương } & \text { Số nhân viên  } \\
		 \hline[50; 60) & 8 \\
		 \hline[60; 70) & 10 \\
		 \hline[70; 80) & 16 \\
		\hline[80;  90) & 14 \\
		\hline[90; 100) & 10  \\
		\hline[100;  110) & 5 \\
		\hline[110; 120) & 2 \\ 
		 \hline
	\end{array}
	\] 
	Hãy xác định các tứ phân vị của mẫu số liệu trên. 
	 \loigiai{
		Gọi  $x_1\le {x_2}\le \ldots \le {x_{65}}$ là tiền lương của 65 nhân viên nhận được trong $1$ giờ.\\
		Tứ phân vị thứ hai của dãy số liệu $x_1$; $x_2$; $\ldots$; $x_{65}$  là ${x_{33}}\in \left[70;80 \right)$. Do đó tứ phân vị thứ hai của mẫu số liệu ghép nhóm là\\
		$Q_2=70+\dfrac{\dfrac{2.65}{4}-\left(  8+10 \right)}{16}\cdot \left( 80-70 \right)=79,0625$.\\
		Tứ phân vị thứ nhất của dãy số liệu $x_1$; $x_2$; $\ldots$; $x_{65}$ là ${x_{17}}\in \left[60;70  \right)$. Do đó tứ phân vị thứ nhất của mẫu số liệu ghép nhóm là\\
		$Q_1=60+\dfrac{\dfrac{1\cdot 65}{4}-8}{10}\cdot \left( 70-60  \right)=68,25$.\\
		Tứ phân vị thứ ba của dãy số liệu $x_1$;  $x_2$; $\ldots$; $x_{65}$ là ${x_{49}}\in \left[90;100 \right)$. Do đó tứ phân vị thứ ba của mẫu số liệu ghép nhóm là\\
		$Q_3=90+\dfrac{\dfrac{3\cdot  65}{4}-\left( 8+10+16+14 \right)}{10}\cdot \left( 100-90 \right)=90,75$.\\
		Chú ý: Nếu tứ phân vị thứ $k$ là $\dfrac{1}{2}\left( {x_m}+{x_{m+1}} \right)$, trong đó $x_m\in \left[{u_{j-1}};u_j \right)$ và  ${x_{m+1}}\in \left[{u_j};{u_{j+1}} \right)$ thuộc hai nhóm liên tiếp thì ta lấy $Q_k=u_j$.\\
	} 
\end{ex}

% nguồn: Ok_TuanLuan_Toan11_Lan2
\begin{ex}%[1D5H2-3]%[ma:0f20bbb9]
	Một nhóm gồm 45  học sinh làm một bài kiểm tra trắc nghiệm gồm 40 câu hỏi. Số câu trả lời đúng của mỗi bạn được ghi lại ở bảng sau:\[
	\begin{array}{|l|l|l|l|l|l|l|l|l|l|l|l|l|l|l|}
		\hline 24 & 35 & 37 & 24 &  30 & 23 & 21 & 39 & 28 & 20 & 32 & 37 & 17 & 40 & 34 \\
		\hline 27  & 34 & 30 & 21 & 26 & 26 & 38 & 37 & 16 & 35 & 19 & 20 & 22 & 25 & 38 \\
		\hline 34 & 29 & 39 & 40 &  36 & 18 & 31 & 24 & 36 & 33 & 24 & 24 & 36 & 26 & 37 \\
		\hline 
	\end{array}
	\] 
	\begin{enumerate}[a)]
		\item Tìm  các tứ phân vị của dãy số liệu trên.
		\item Tổng hợp lại dãy số liệu trên  vào bảng tần số ghép nhóm theo mẫu sau:
	\end{enumerate}
	 \[
	\begin{array}{|c|c|c|c|c|c|}
		\hline \text { Số câu trả lời đúng } & {[16; 20]} & {[21; 25]} &  {[26; 30]} & {[31; 35]} & {[36; 40]} \\
		\hline \text { Số học sinh } & ? &  ? & ? & ? & ? \\
		\hline 
	\end{array}
	\] 
	c) Hãy ước lượng các tứ phân vị của mẫu số liệu ghép nhóm  trên
	\loigiai{ 
		a) Mẫu số liệu đã cho được sắp xếp theo thứ tự không giảm như sau:\[{\begin{aligned}
		& 16;  17; 18; 19; 20; 20; 21; 21; 22; 23; 24; 24; 24; 24; 24;\\
		& 25; 26; 26; 26;  27; 28; 29; 30; 30; 31; 32; 33; 34; 34; 34;\\
		& 35; 35;  36; 36; 36; 37; 37; 37; 37; 38; 38; 39; 39; 40; 40 .
	\end{aligned}
	 }\]\\
	Cỡ mẫu là ${n=45}$ là số lẻ nên trung vị của mẫu số liệu là ${Q_2=x_{23}=30}$.\\ 
	Tứ phân vị thứ nhất là  ${Q_1=\dfrac{1}{2}\left(x_{11}+x_{12}\right)=\dfrac{1}{2}(24+24)=24}$.\\
	Tứ phân vị thứ ba là  ${Q_3=\dfrac{1}{2}\left(x_{34}+x_{35}\right)=\dfrac{1}{2}(36+36)=36}$.\\
	b) Bảng tần  số ghép nhóm của mẫu số liệu trên như sau:\[
	\begin{array}{|c|c|c|c|c|c|}
		\hline \text { Số câu trả  lời đúng } & {[16; 20]} & {[21; 25]} & {[26; 30]} & {[31; 35]} & {[36; 40]} \\
		\hline \text { Số học sinh  } & 6 & 10 & 8 & 8 & 13 \\
		\hline 
	\end{array}
	 \]
	c) Do số câu trả lời đúng của học sinh là số nguyên nên ta hiệu chỉnh lại bảng  số liệu như sau:\[
	\begin{array}{|c|c|c|c|c|c|}
		\hline \text { Số câu trả lời đúng } & {[15,5; 20,5)} & {[20,5; 25,5)} & {[25,5; 30,5)} & {[30,5; 35,5)} & {[35,5; 40,5)}  \\
		\hline \text { Số học sinh } & 6 & 10 & 8 &  8 & 13 \\
		\hline 
	\end{array}
	\] 
	Gọi $x_1$; $x_2$; $\ldots$; $x_{45}$ là số  câu trả lời đúng của 45 học sinh được xếp theo thứ tự không giảm.\\
	Ta có: $x_1$, $\ldots$, $x_6  \in[15,5; 20,5)$; $x_7$, $\ldots$, $x_{16} \in[20,5; 25,5)$; $x_{17}$, $\ldots$, $x_{24} \in[25,5; 30,5)$; $x_{25}$, $\ldots$, $x_{32} \in[30,5; 35,5)$; $x_{33}$, $\ldots$, $x_{45} \in[35,5; 40,5)$.\\
	Tứ phân vị thứ hai của mẫu số liệu $x_1$; $x_2$; $\ldots$;  $x_{45}$ là ${x_{23} \in[25,5; 30,5)}$.\\
	Do đó, tứ phân vị thứ hai của  mẫu số liệu ghép nhóm là\\
	 ${Q_2=25,5+\dfrac{\dfrac{45}{2}-(6+10)}{8} \cdot(30,5-25,5)=29,5625 .
	}$\\ 
	Tứ phân vị thứ nhất  của mẫu số liệu $x_1$; $x_2$; $\ldots$; $x_{45}$ là ${\dfrac{1}{2}\left(x_{11}+x_{12}\right)}$.\\
	Do ${x_{11}}$ và ${x_{12}}$ thuộc nhóm ${[20,5; 25,5)}$ nên tứ phân vị thứ nhất  của mẫu số liệu ghép nhóm là\\
	${Q_1=20,5+\dfrac{\dfrac{45}{4}-6}{10} \cdot(25,5-20,5)=23,125 \text  {. }
	 }$\\
	Tứ phân vị  thứ ba của mẫu số liệu $x_1$; $x_2$; $\ldots$; $x_{45}$ là ${\dfrac{1}{2}\left(x_{34}+x_{35}\right)}$. Do ${x_{34}}$ và ${x_{35}}$ thuộc nhóm ${[35,5; 40,5)}$ nên tứ phân vị thứ ba của mẫu số liệu ghép nhóm là\\
	${Q_3=35,5+\dfrac{\dfrac{3.45}{4}-(6+10+8+8)}{13} \cdot(40,5-35,5) \approx  36,173 .
 }$}
\end{ex}

% nguồn: Ok_TuanLuan_Toan11_Lan2
\begin{ex}%[1D5V2-3]%[ma:50e3e8ca]
	Để chuẩn bị cho đồ án tốt nghiệp, một sinh viên y khoa đã khảo sát huyết áp tối đa của một số bệnh nhân và  lập được bảng tần số ghép nhóm sau:
	\begin{center}
		\begin{tabular}{|c|c|}
			 \hline Huyết áp & Tần số \\
			\hline$[90; 110)$ &  6 \\
			\hline$[110; 130)$  & 20 \\
			\hline$[130; 150)$ &  35 \\
			\hline$[150; 170)$ &  45 \\
			\hline$[170; 190)$  & 30 \\
			\hline$[190; 210)$ & 16 \\ 
			\hline 
		\end{tabular}
	\end{center}
	\begin{enumerate}
		\item Xác định trung bình, trung vị, tứ phân vị  và mốt của mẫu số liệu.
		\item Hãy giải thích vì sao trong truờng hợp này, cả ba giá trị trung bình, trung vị và mốt  tìm được đều đại diện tốt cho huyết áp của nhưng bệnh nhân được khảo sát.
	\end{enumerate}
	 \loigiai{
		\begin{enumerate}
			\item Bảng tần  số ghép nhóm
			\begin{center}
				\begin{tabular}{|c|c|c|}
					\hline Huyết áp  & Giá trị đại diện & Tần số \\
					\hline$[90; 110)$  & 100 & 6 \\
					\hline$[110; 130)$ & 120 &  20 \\
					\hline$[130; 150)$ & 140  & 35 \\
					\hline$[150; 170)$ & 160 & 45  \\
					 \hline$[170; 190)$ & 180 & 30 \\
					 \hline$[190; 210)$ & 200 & 16 \\
					 \hline
				\end{tabular}
			\end{center}
			Giá  trị trung bình
			 \[
			\overline{x}=\dfrac{100 \cdot 6+120 \cdot 20+140  \cdot 35+160 \cdot 45+180 \cdot 30+200 \cdot 16}{152}=155{,}9.
			\] 
			Trung vị $M_{\mathrm{e}}=150+\dfrac{76-61}{45} \cdot 20=156{,}7$. 
			\\ 
			Tứ phân vị thứ nhất và thứ ba 
			\begin{align*}
				 & Q_1=130+\dfrac{38-26}{35} \cdot 20=136{,}9; \\
				&  Q_3=170+\dfrac{114-106}{30} \cdot 20=175{,}3.
			\end{align*}
			Tứ phân vị của mẫu số liệu  $Q_1=136,9$; $Q_2=156{,}7$; $Q_3=175{,}3$.
			 \\
			Mốt $M_0=150+\dfrac{10}{10+15} \cdot 20=158$. 
			\item Vì ba giá trị này cùng thuộc vào một nhóm ghép có tần số cao nhất và chúng gần như không có  sự chênh lệch nhiều.
		\end{enumerate}
	} 
\end{ex}

% nguồn: Ok_TuanLuan_Toan11_Lan2
\begin{ex}%[1D5V2-3]%[ma:f1b1a717]
	Mẫu số liệu sau  ghi lại cân nặng của 30 bạn học sinh (đơn vị: kilogam)
	\begin{center}
		\begin{tabular}{|c|c|c|c|c|c|c|c|c|c|}
			\hline 17 & 40 & 39 & 40{,}5 & 42 & 51 & 41{,}5 & 39  & 41 & 30 \\
			 \hline 40 & 42 & 40{,}5 & 39{,}5 & 41 & 40{,}5 & 37 & 39{,}5 & 40 & 41 \\
			\hline 38{,}5 & 39{,}5 &  40 & 41 & 39 & 40{,}5 & 40 & 38{,}5 & 39{,}5 & 41{,}5 \\
			 \hline
		\end{tabular}
	\end{center}
	\begin{enumerate}
		\item Lập bảng tần số ghép nhóm cho mẫu số liệu trên có tám nhóm ứng với tám nửa khoảng $[15; 20)$, $[20; 25)$, $[25; 30)$, $[30; 35)$, $[35; 40)$, $[40; 45)$,  $[45; 50)$, $[50; 55)$.
		\item Xác định số trung bình cộng, trung vị, tứ phân vị của mẫu  số liệu ghép nhóm trên.
		\item Mốt của mẫu số liệu  ghép nhóm trên là bao nhiêu?
	\end{enumerate}
	 \loigiai{
		\begin{enumerate}
			\item  Bảng tần số ghép nhóm cho mẫu số liệu
			\begin{center}
				\begin{tabular}{|c|c|c|}
					\hline  Nhóm & Giá trị đại diện & Tần số \\
					\hline$[15; 20)$ & 17,5 & 1 \\ 
					\hline$[20; 25)$ & 22,5 & 0 \\ 
					\hline$[25; 30)$ & 27,5 & 0  \\
					 \hline$[30; 35)$ & 32,5 & 1 \\
					\hline$[35; 40)$ & 37,5 & 10 \\ 
					\hline$[40; 45)$ & 42,5  & 17 \\
					\hline$[45; 50)$  & 47,5 & 0 \\
					\hline$[50; 55)$ & 52,5  & 1 \\
					\hline & & $n=30$ \\ 
					\hline 
				\end{tabular}
			\end{center}
			\item Số  trung bình cộng
			\[ 
			\overline{x}=\dfrac{17{,}5+32{,}5+37{,}5 \cdot 10+42{,}5 \cdot 17+52{,}5}{30}=40. 
			 \]
			 Trung vị $Q_2=40+\dfrac{15-12}{17} \cdot 5 \approx 40{,}9$.
			 \\
			Trung vị $Q_1$ và $Q_3$ 
			\begin{align*}
				& Q_1=35+\dfrac{7{,}5-2}{10} \cdot 5=37{,}75; \\ 
				& Q_3=40+\dfrac{22{,}5-12}{17}  \cdot 5 \approx 43{,}1.
			\end{align*}
			Tứ phân vị $Q_1=37{,}75$, $Q_2=40{,}9$, $Q_3=43{,}1$. 
			\item Mốt của mẫu  số liệu
			\[ 
			M_0=40+\dfrac{7}{7+17} \cdot 5 \approx  41{,}5.
			\] 
		\end{enumerate}
	 }
\end{ex}

% nguồn: Ok_TuanLuan_Toan11_Lan2
\begin{ex}%[1D5V2-3]%[ma:83631651]
	Anh Văn ghi lại cự li $30$  lần ném lao của mình ở bảng sau (đơn vị: mét):
	\begin{center}
		\begin{tabular}{|c|c|c|c|c|c|c|c|c|c|}
			\hline $72{,}1$ & $72{,}9$ & $70{,}2$ & $70{,}9$ & $72{,}2$ & $71{,}5$ & $72{,}5$ & $69{,}3$ & $72{,}3$ & $69{,}7$ \\ 
			\hline $72{,}3$ & $71{,}5$ & $71{,}2$ & $69{,}8$ & $72{,}3$ & $71{,}1$ & $69{,}5$ &  $72{,}2$ & $71{,}9$ & $73{,}1$ \\
			\hline $71{,}6$ & $71{,}3$ & $72{,}2$ & $71{,}8$ & $70{,}8$ & $72{,}2$ & $72{,}2$ & $72{,}9$ & $72{,}7$ & $70{,}7$  \\
			 \hline
		\end{tabular}
	\end{center}
	\begin{enumerate}
		\item Tính cự li  trung bình của mỗi lần ném.
		\item Tổng hợp lại  kết quả ném của anh Văn vào bảng tần số ghép nhóm theo mẫu sau:
		\begin{center}
			\begin{tabular}{|c|c|c|c|c|c|}
				\hline Cư li $( m)$ & $[69,2; 70)$ & $[70; 70,8)$ & $[70,8; 71,6)$ & $[71,6; 72,4)$ &  $[72,4; 73,2)$ \\
				\hline Số lần & $?$ & $?$ & $?$ & $?$  & $?$ \\
				\hline 
			\end{tabular}
		\end{center}
		\item Hãy ước lượng cự li trung bình mỗi lần ném từ bảng tần  số ghép nhóm trên.
		\item  Khả năng anh Văn ném được bao nhiêu mét là cao nhất?
	\end{enumerate}
	 \loigiai{
		\begin{enumerate}
			\item Cự li  trung bình mỗi lần ném là
			\[ 
			\overline{x}=\dfrac{72{,}1+72{,}9+70{,}2+\ldots+70{,}7}{30} \approx 71{,}56\, \text{(mét)}. 
			\] 
			\item Tổng hợp lại ta có bảng  tần số ghép nhóm là
			\begin{center}
				\begin{tabular}{|c|c|c|c|c|c|}
					\hline Cự li $( m)$ &  $[69,2; 70)$ & $[70; 70,8)$ & $[70,8; 71,6)$ & $[71,6; 72,4)$ & $[72,4; 73,2)$ \\
					\hline Số lần &  4 & 2 & 7 & 12 & 5 \\
					 \hline
				\end{tabular}
			\end{center}
			 \item Ta có
			\begin{center}
				\begin{tabular}{|c|c|c|c|c|c|}
					\hline Cự li (m) & $[69,2; 70)$ & $[70; 70,8)$  & $[70,8; 71,6)$ & $[71,6; 72,4)$ & $[72,4; 73,2)$ \\
					\hline  Giá trị đại diện & 69{,}6 & 70{,}4 & 71{,}2 & 72 & 72{,}8 \\
					\hline Sồ lần & 4 & 2 & 7 & 12 &  5 \\
					 \hline
				\end{tabular}
			\end{center}
			Cự li trung bình  ước lượng là
			 \[
			\overline{x}=\dfrac{69{,}6 \cdot  4+70{,}4 \cdot 2+\ldots+72{,}8 \cdot 5}{30} \approx 71{,}52 \, \text{(mét)}.
			 \]
			\item Anh Văn ném được khoảng $73$ mét là  cao nhất.
		\end{enumerate}
	 }
\end{ex}

% nguồn: Ok_TuanLuan_Toan11_Lan2
\begin{ex}%[1D5V2-3]%[ma:bc58f88e]
	Một thư viện thống kê số lượng sách được mượn mỗi ngày trong ba tháng ở bảng sau. 
	\begin{center}
		\begin{tabular}{|c|c|c|c|c|c|c|c|}
			\hline Số  sách & $[16; 20]$ & $[21; 25]$ & $[26; 30]$ & $[31; 35]$ & $[36; 40]$ & $[41; 45]$ & $[46; 50]$ \\
			\hline Số ngày & 3 & 6 & 15  & 27 & 22 & 14 & 5 \\
			 \hline
		\end{tabular}
	\end{center}
	 Hãy ước lượng số trung bình và mốt của mẫu số liệu nghép nhóm trên.
	 \loigiai{
		 Ta có bảng giá trị đại diện sau.
		\begin{center}
			\begin{tabular}{|c|c|c|c|c|c|c|c|}
				\hline Số sách & $[16; 20]$ & $[21; 25]$ & $[26; 30]$ &  $[31; 35]$ & $[36; 40]$ & $[41; 45]$ & $[46; 50]$ \\
				\hline Giá trị đại diện & 18  & 23 & 28 & 33 & 38 & 43 & 48 \\
				\hline Sốngày & 3 & 6 & 15 & 27 & 22  & 14 & 5 \\
				 \hline
			\end{tabular}
		\end{center}
		Ước lượng số  trung bình của mẫu số liệu là
		\[ 
		\overline{x}=\dfrac{18 \cdot 3+23 \cdot 6+28 \cdot 15+33 \cdot 27+38 \cdot  22+43 \cdot 14+48 \cdot 5}{92} \approx 35 \text {~(quyển).}
		 \]
		Mốt của  mẫu số liệu trên là
		 \[
		 M_0=31+\dfrac{27-15}{27-15+27-22} \cdot(36-31) \approx 35.
		 \]
	 }
\end{ex}

% nguồn: Ok_TuanLuan_Toan11_Lan2
\begin{ex}%[1D5V2-3]%[ma:e32f089b]
	An tìm hiểu hàm lượng chất béo (đơn vị: g) có trong $100$ g mỗi loại thực phẩm.  Sau khi thu thập dữ liệu về $60$ loại thực phẩm, An lập được bảng thống kê
	\begin{center}
		\begin{tabular}{|c|c|c|c|c|c|c|}
			\hline Hàm lượng chất béo (g) & $[2; 6)$ & $[6; 10)$ & $[10; 14)$ & $[14; 18)$ & $[18; 22)$ &  $[22; 26)$ \\
			\hline Tần số & 2 & 6 & 10 &  13 & 16 & 13 \\
			 \hline
		\end{tabular}
	\end{center}
	\begin{enumerate}
		\item Xác định giá trị trung bình, trung vị, tứ phân vị của mẫu số liệu. 
		\item  Tìm mốt của mẫu số liệu.
	\end{enumerate}
	 \loigiai{
		\begin{enumerate}
			\item Bảng tần số ghép  nhóm
			\begin{center}
				\begin{tabular}{|c|c|c|c|c|c|c|}
					\hline Hàm lượng chất béo $( g)$ & $[2;  6)$ & $[6; 10)$ & $[10; 14)$ & $[14; 18)$ & $[18; 22)$ & $[22; 26)$ \\
					\hline Giá trị đại diện  & 4 & 8 & 12 & 16 & 20 & 24 \\
					\hline Tần số & 2 & 6 & 10 & 13 & 16 & 13 \\ 
					 \hline
				\end{tabular}
			\end{center}
			Giá trị trung bình 
			\[ 
			\overline{x}=\dfrac{4 \cdot 2+8 \cdot 6+12 \cdot 10+16 \cdot 13+20 \cdot 16+24 \cdot 13}{60} \approx  16{,}9.
			 \]
			Trung vị $M_{\mathrm{e}}=14+\dfrac{30-18}{13}  \cdot 4=17{,}7$.
			\\ 
			Tứ phân  vị thứ nhất và thứ ba
			\begin{align*}
				& Q_1=10+\dfrac{15-8}{10}  \cdot 4=12{,}8; \\
				& Q_3=18+\dfrac{45-31}{16} \cdot 4=21{,}5. 
			\end{align*}
			Tứ phân vị của mẫu số liệu là $Q_1=12,8$;  $Q_2=17{,}7$; $Q_3=21{,}5$.
			 \item Mốt của mẫu số liệu $M_0=18+\dfrac{3}{3+3} \cdot 4=20$.
		\end{enumerate}
	 }
\end{ex}

% nguồn: Ok_TuanLuan_Toan11_Lan2
\begin{ex}%[1D5V2-3]%[ma:5fb07e84]
	Lương tháng của một số giáo viên THPT được ghi lại  như sau (đơn vị: triệu đồng):
	\begin{center}
		\begin{tabular}{|c|c|c|c|c|c|c|c|c|c|c|c|}
			\hline 
			$12{,}2$ & $9{,}8$ & $10{,}9$ & $6{,}7$ & $13{,}6$ & $9{,}2$ & $12{,}5$ & $9{,}6$ & $11{,}7$ & $12{,}7$ & $10{,}0$ &  $10{,}0$
			 \\
			\hline 
			$8{,}3$ &  $11{,}1$ & $11{,}9$ & $8{,}4$ & $6{,}7$ & $13{,}8$ & $13{,}1$ & $6{,}5$ & $10{,}7$ & $8{,}9$ & $11{,}2$ & $13{,}2$
			\\ 
			\hline 
		\end{tabular}
	\end{center}
	\begin{enumerate}
		 \item Tìm tứ phân vị của dãy số liệu trên.
		\item Tổng hợp lại  dãy số liệu trên vào bảng tần số ghép nhóm theo mẫu sau.
		\begin{center}
			\begin{tabular}{|c|c|c|c|c|}
				 \hline
				Lương tháng (triệu đồng)  & $[6; 8)$ &$[8; 10)$ &$[10; 12)$ &$[12; 14)$\\
				\hline 
				Số nhân viên & $?$ & $?$ &  $?$ & $?$ \\
				\hline 
			\end{tabular}
		\end{center}
		\item Hãy  ước lượng tứ phân vị của số liệu ở bảng tần số ghép nhóm trên.
	\end{enumerate}
	 \loigiai{
		\begin{enumerate}
			\item Các tứ phân vị  là $Q_2=10{,}8; Q_1=9{,}05; Q_3=12{,}35$.
			\item  Bảng tần số ghép nhóm theo mẫu.
			\begin{center}
				\begin{tabular}{|c|c|c|c|c|}
					 \hline
					Lương tháng  &$[6; 8)$ &$[8; 10)$ &$[10; 12)$ & $[12; 14)$\\
					\hline 
					Số  giáo viên & $3$ & $6$ & $8$ & $7$ \\
					\hline 
				\end{tabular}
			\end{center}
			\item Tứ  phân vị thứ nhất là $Q_1=8+\dfrac{\dfrac{24}{4}-3}{6} \cdot (10-8)=9$.
			\\ 
			Tứ phân vị thứ hai là $Q_2=10+\dfrac{\dfrac{24}{2}-9}{8}  \cdot (12-10)=10{,}75$.
			\\ 
			Tứ phân vị thứ ba  là $Q_3=12+\dfrac{\dfrac{3 .24}{4}-17}{7} \cdot (14-12) \approx 12{,}29$.
		\end{enumerate}
	} 
\end{ex}

% nguồn: Ok_TuanLuan_Toan11_Lan2
\begin{ex}%[1D5V2-3]%[ma:8e2b2408]
	Kết quả đo chiều cao của 200 cây keo $3$ năm tuổi ở một nông trường được biểu diễn ở  biểu đồ dưới đây.
	\begin{center}
		\begin{tikzpicture}[scale=.8]
				\draw[thick,->]  (0,0) -- (11,0) node[below] {(m)};
				\draw[thick,->] (0,0) -- (0,8)  node[left] {(Số cây)};
				\foreach \y in {0,10,...,70}{ 
					\draw (0.1,\y/10)  -- (-0.1,\y/10) node[left] {\y};
				 }
				\foreach \x/\label in {1/[8{,}5;  8.8],3/[8{,}8; 9.1],5/[9{,}1; 9.4],7/[9{,}4; 9.7],9/[9{,}7; 10]}{
					\draw (\x,0) -- (\x,-0.2) node[below] {\label}; 
				} 
				\fill[cyan!40] (0.5,0) rectangle (1.5,2); 
				\draw(0.5,0) rectangle  (1.5,2);
				\fill[cyan!40] (2.5,0) rectangle  (3.5,3.5);
				\draw(2.5,0) rectangle  (3.5,3.5);
				\fill[cyan!40] (4.5,0)  rectangle (5.5,6);
				\draw(4.5,0)  rectangle (5.5,6);
				 \fill[cyan!40] (6.5,0) rectangle (7.5,5.5);
				 \draw(6.5,0) rectangle (7.5,5.5);
				\fill[cyan!40] (8.5,0) rectangle (9.5,3); 
				\draw(8.5,0) rectangle (9.5,3); 
				\node[above] at (1,2)  {20};
				\node[above] at (3,3.5)  {35};
				\node[above] at (5,6) {60}; 
				\node[above] at (7,5.5)  {55};
				\node[above]  at (9,3) {30};
				\node at (5,7.5) {\textbf{Chiều cao 200 cây keo 3 năm  tuổi}};
			\end{tikzpicture}
	\end{center}
	Hãy ước lượng số trung bình và  mốt của mẫu số liệu nghép nhóm trên.
	 \loigiai{
		Ta có bảng  giá trị đại diện sau
		\begin{center}
			\begin{tabular}{|c|c|c|c|c|c|}
				\hline Chiều cao của cây & $[8{,}5; 8{,}8)$ & $[8{,}8; 9{,}1)$ & $[9{,}1; 9{,}4)$ & $[9{,}4; 9{,}7)$  & $[9{,}7; 10)$ \\
				\hline Giá trị đại  diện & 8,65 & 8,95 & 9,25 & 9,55 & 9,85 \\
				\hline Số cây & 20 & 35 & 60 & 55 &  30 \\
				 \hline
			\end{tabular}
		\end{center}
		 Ước lượng trung bình của mẫu số liệu là
		\[ 
		\overline{x}=\dfrac{8,65 \cdot 20+8,95  \cdot 35+9,25 \cdot 60+9,55 \cdot 55+9,85 \cdot 30}{200} \approx 9{,}31\, \text{(m)}.
		\] 
		Mốt của mẫu số liệu ghép nhóm  là
		\[ 
		M_0=9{,}1+\dfrac{60-35}{60-35+60-55} \cdot(9{,}4-9{,}1) \approx  9{,}35 \, \text{(m)}.
		 \]
	 }
\end{ex}

% nguồn: Ok_TuanLuan_Toan11_Lan2
\begin{ex}%[1D5V2-3]%[ma:743f2d06]
	Người ta đếm số xe ô tô đi qua một trạm thu phí mỗi phút trong khoảng thời gian từ $9$ giờ đến $9$ giờ $30$ phút sáng. Kết quả được ghi lại ở  bảng sau.
	\begin{center}
		\begin{tabular}{|c|c|c|c|c|c|c|c|c|c|c|c|c|c|c|}
			\hline 15 & 16 & 13 & 21 & 17 & 23 & 15 & 21 & 6 & 11 &  12 & 23 & 19 & 25 & 11 \\
			 \hline 25 & 7 & 29 & 10 & 28 & 29 & 24 & 6 & 11 & 23 & 11 & 21 & 9 & 27 & 15 \\
			\hline 
		\end{tabular}
	\end{center}
	\begin{enumerate}
		\item Tính số xe trung bình đi qua trạm thu phí  trong mỗi phút.
		\item Tổng hợp lại số liệu trên vào bảng tần số ghép nhóm  theo mẫu sau.
		\begin{center}
			\begin{tabular}{|c|c|c|c|c|c|}
				\hline Số xe & $[6; 10]$ & $[11; 15]$ & $[16; 20]$ & $[21; 25]$ & $[26; 30]$ \\ 
				\hline Số lần & $?$ & $?$ & $?$ & $?$ &  $?$ \\
				\hline 
			\end{tabular}
		\end{center}
		\item Hãy ước lượng trung  bình số xe đi qua trạm thu phí trong mỗi phút từ bảng tần số ghép nhóm trên.
	\end{enumerate}
	 \loigiai{
		\begin{enumerate}
			\item Số xe trung  bình đi qua trạm thu phí trong mỗi phút là
			\[ 
			\overline{x}=\dfrac{15+16+13+21+\ldots+27+15}{30} \approx  17 \, \text{(xe)}.
			\] 
			\item Ta có  bảng sau
			\begin{center}
				\begin{tabular}{|c|c|c|c|c|c|}
					\hline Số xe & $[6; 10]$ & $[11; 15]$ & $[16;  20]$ & $[21; 25]$ & $[26; 30]$ \\
					\hline Số lần & 5  & 9 & 3 & 9 & 4 \\
					\hline 
				\end{tabular}
			\end{center}
			\item Ta có bảng giá trị  đại diện sau.
			\begin{center}
				\begin{tabular}{|c|c|c|c|c|c|}
					\hline Số xe & $[6; 10]$ & $[11; 15]$ & $[16; 20]$ &  $[21; 25]$ & $[26; 30]$ \\
					\hline Giá trị đại diện & 8 &  13 & 18 & 23 & 28 \\
					 \hline Sồ lần & 5 & 9 & 3 & 9 & 4 \\
					\hline 
				\end{tabular}
			\end{center}
			 Số xe trung bình ước lượng đi qua trạm thu phí từ bảng tần số ghép nhóm trên là
			\[ 
			\overline{x}=\dfrac{8  \cdot 5+13 \cdot 9+18 \cdot 3+23 \cdot 9+28 \cdot 4}{30} \approx 18\, \text{(xe)}.
			\] 
		\end{enumerate}
	} 
\end{ex}

% nguồn: Ok_TuanLuan_Toan11_Lan2
\begin{ex}%[1D5V2-3]%[ma:70097acd]
	Một cửa hàng đã thống kê số ba lô bán được  mỗi ngày trong tháng $9$ với kết quả cho như sau.
	\begin{center}
		\begin{tabular}{|c|c|c|c|c|c|c|c|c|c|c|c|c|c|c|}
			\hline 12 & 29 & 12 & 19 & 15 & 21 & 19 & 29 & 28  & 12 & 15 & 25 & 16 & 20 & 29 \\
			\hline 21  & 12 & 24 & 14 & 10 & 12 & 10 & 23 & 27 & 28 & 18 & 16 & 10 & 20 & 21 \\
			\hline 
		\end{tabular}
	\end{center}
	Hãy chia mẫu số liệu trên thành $5$ nhóm, lập bảng tần số ghép nhóm, hiệu chỉnh bảng tần số ghép nhóm và xác định giá  trị đại điện cho mỗi nhóm.
	\loigiai{ 
		Khoảng biến  thiên của mẫu số liệu trên là: $R=29-10=19$.
		\\ 
		Độ dài mỗi  nhóm $L>\dfrac{R}{k}=\dfrac{19}{5}=3{,}8$.
		\\ 
		Chọn $L=4$ và chia  các dữ liệu thành $5$ nhóm
		$[10; 14)$, $[14; 18)$, $[18; 22)$, $[22; 26)$, $[26; 30)$. 
		\\ 
		Khi đó ta có bảng tần số ghép nhóm  sau
		\begin{center}
			\begin{tabular}{|c|c|c|c|c|c|}
				\hline Số  ba lô đã bán & $[10; 14)$ & $[14; 18)$ & $[18; 22)$ & $[22; 26)$ & $[26; 30)$ \\
				\hline Giá trị đại diện & 12 & 16 & 20  & 24 & 28 \\
				\hline Số ngày & 8 & 5 & 8  & 3 & 6 \\
				 \hline
			\end{tabular}
		\end{center}
	} 
\end{ex}

% nguồn: Ok_TuanLuan_Toan11_Lan2
\begin{ex}%[1D5V2-3]%[ma:03d2a32f]
	Một câu lạc bộ thể dục thể thao đã ghi lại số giờ các thành viên của mình sử dụng cơ sở vật chất của câu lạc bộ để tập luyện trong một tháng.  Họ tổ chức dữ liệu thu được vào bảng sau
	\begin{center}
		\begin{tabular}{|c|c|c|c|c|c|c|}
			\hline  Thời gian (giờ) & $[1; 5)$ & $[5; 9)$ & $[9; 13)$ & $[13; 17)$ & $[17; 21)$ & $[21; 25)$ \\
			\hline Tần số (Số người) & 10 & 14 & 31  & 2 & 5 & 23 \\
			 \hline
		\end{tabular}
	\end{center}
	Hãy tính (làm tròn kết quả đến hàng  phần mười).
	\begin{enumerate}
		\item Trung vị của  mẫu số liệu.
		\item Trung bình của mẫu số liệu. Trong truờng hợp này thì trung bình hay trung vị đại diện  tốt hơn cho mẫu số liệu?
	\end{enumerate}
	 \loigiai{
		\begin{enumerate}
			\item Ta  lập được bảng sau
			\begin{center}
				\begin{tabular}{|c|c|c|c|c|c|c|}
					\hline Thời gian (giờ) & $[1; 5)$ & $[5; 9)$ & $[9; 13)$ & $[13; 17)$ & $[17; 21)$ & $[21; 25)$  \\
					\hline Giá trị đại diện &  3 & 7 & 11 & 15 & 19 & 23 \\
					\hline Tần số (Số người) & 10 & 14 &  31 & 2 & 5 & 23 \\
					\hline 
				\end{tabular}
			\end{center}
			Trung vị của mẫu số  liệu
			\[ 
			M_{\mathrm{e}}=9+\dfrac{42,5-24}{31} \cdot  4 \approx 11{,}4.
			 \]
			\item Trung bình của  mẫu số liệu là $\overline{x}=13{,}2$.
			 \\
			Trung bình thuộc nhóm $[13; 17$)cho thấy trong $85$ số liệu đã có ít nhất $55$ số liệu nhỏ hơn số trung  bình. Suy ra, trong trường hợp này thì trung vị là số đại diện tốt hơn cho mẫu số liệu.
		\end{enumerate}
	} 
\end{ex}

% nguồn: Ok_TuanLuan_Toan11_Lan2
\begin{ex}%[1D5V2-3]%[ma:5f441cfb]
	Trong một cuộc đua Marathon được  tổ chức ở thành phố $A$ người ta thống kê lại được như sau
	\begin{center}
		\begin{tabular}{|c|c|c|c|c|c|}
			\hline Thời gian &$[120; 140)$  &$[140; 160)$ &$[160; 180)$ &$[180; 200)$ &$[200; 220)$ \\
			\hline Số người & $4$ & $6$ &  $10$ & $15$ & $25$ \\
			 \hline
		\end{tabular}
	\end{center}
	Hãy ước lượng số trung bình và tứ phân vị của mẫu số liệu ghép nhóm  trên.
	 \loigiai{
		Ta có $n=4+6+10+15+25=60$. 
		 \\
		Trung  vị là $M_{\mathrm{e}}=180+\dfrac{\dfrac{60}{2}-20}{15} \cdot (200-180)=\dfrac{580}{3} \approx 193{,}33$.
		\\ 
		Tứ  phân vị thứ nhất là $Q_1=160+\dfrac{\dfrac{60}{4}-10}{10} \cdot (180-160)=170$.
		 \\
		Tứ phân vị thứ ba  là $Q_3=200+\dfrac{3 \cdot \dfrac{60}{4}-35}{25} \cdot (220-200)=208$.
	} 
\end{ex}

% nguồn: Ok_TuanLuan_Toan11_Lan2
\begin{ex}%[1D5V2-3]%[ma:2d4e686c]
	Mức lương hàng tháng ở 1 công ty được  Công đoàn thu thập theo bảng sau (đơn vị triệu đồng).
	\begin{center}
		\begin{tabular}{|c|c|c|c|c|c|}
			\hline Mức lương & $[6; 10)$ & $[10; 15)$ & $[15; 20)$ &  $[20; 25)$ & $[25; 30)$ \\
			\hline Nhân viên & 17 &  38 & 27 & 21 & 7 \\
			\hline 
		\end{tabular}
	\end{center}
	\begin{enumerate}
		 \item Hãy ước lượng các tứ phân vị của mẫu số liệu ghép nhóm trên.
		\item Chủ tịch Công đoàn muốn đề nghị hỗ trợ cho nhóm $25 \%$ số nhân viên có mức lương thấp  nhất và ước lượng rằng số nhân viên này không ít hơn $10$. Nhận định của chủ tịch có hợp lí hay không?
	\end{enumerate}
	\loigiai{ 
		\begin{enumerate}
			\item Gọi $x_1 \leq x_2 \leq \ldots \leq x_{110}$ là mức lương của $110$ nhân viên được  công ty trả trong 1 tháng.
			\\ 
			Tứ phân vị thứ hai của dãy số liệu $x_1 \leq x_2 \leq  \ldots \leq x_{110}$ là $\dfrac{1}{2}\left(x_{55}+x_{56}\right)$.
			 \\
			Do  $x_{55} \in[10; 15)$ và $x_{56} \in[15; 20)$.
			\\ 
			Nên đó tứ phân vị thứ  hai của mẫu số liệu ghép nhóm là $Q_2=15$.
			 \\
			Tứ phân vị thứ nhất của dãy số liệu $x_1; x_2;  \ldots; x_{110}$ là $\dfrac{1}{2}\left(x_{27}+x_{28}\right)$.
			 \\
			Do $x_{27} \in[10; 15)$ và  $x_{28} \in[10; 15)$.
			 \\
			Nên  tứ phân vị thứ nhất của mẫu số liệu ghép nhóm là
			 \[
			Q_1=10+\dfrac{\dfrac{1  \cdot 110}{4}-17}{38} \cdot(15-10) \approx 11{,}38.
			 \]
			Tứ phân vị thứ ba của dãy số liệu $x_1;  x_2; \ldots; x_{110}$ là $\dfrac{1}{2}\left(x_{82}+x_{83}\right)$.
			 \\
			Do  $x_{82} \in[15; 20)$ và $x_{83} \in[20; 25)$.
			 \\
			Nên tứ phân vị thứ ba của mẫu số liệu ghép nhóm là $Q_3=20$. 
			\item Do tứ  phân vị thứ nhất $Q_1 \approx 11,38$ nên nhận định trên là không hợp lí.
		\end{enumerate}
	} 
\end{ex}

% nguồn: Ok_TuanLuan_Toan11_Lan2
\begin{ex}%[1D5V2-3]%[ma:1a70d594]
	Thống kê cân nặng của 30 học sinh trong một lớp học thu được kết quả như sau 
	\begin{center}
		\begin{tabular}{|c|c|c|c|c|c|c|c|c|c|}
			 \hline \multicolumn{10}{|c|}{Cân nặng (kg)} \\
			\hline 54 & 42 & 53 & 63 & 58 & 54 & 63 &  70 & 59 & 49 \\
			\hline 73 & 49 & 66 & 60 & 55 & 48 & 65 & 63  & 56 & 51 \\
			\hline 61  & 57 & 69 & 67 & 63 & 52 & 60 & 56 & 49 & 44 \\
			\hline 
		\end{tabular}
	\end{center}
	\begin{enumerate}
		\item Hãy tính cân nặng  trung bình của học sinh lớp học trên.
		\item Hãy thu gọn dữ liệu bằng cách lập bảng tần số ghép nhóm với các nhóm ghép  là $[40; 50)$; $[50; 60)$; $[60; 70)$; $[70; 80)$.
		\item Hãy ước lượng cân nặng trung  bình của học sinh trong lớp từ bảng tần số ghép nhóm trên.
	\end{enumerate}
	 \loigiai{
		\begin{enumerate}
			\item  Cân nặng trung bình của học sinh lớp học trên
			 \[
			\overline{x}=(54+42+53+  \hdots +49+44) \colon 30 \approx 57{,}63.
			\] 
			Bảng tần số ghép nhóm. 
			\begin{center}
				\begin{tabular}{|c|c|c|c|c|}
					\hline Cân nặng (kg) & $[40; 50)$ & $[50; 60)$ & $[60; 70)$ & $[70; 80)$ \\ 
					\hline Số học sinh  & 6 & 12 & 10 & 2 \\
					\hline 
				\end{tabular}
			\end{center}
			\item Ta có bảng thống kê cân nặng học  sinh theo giá trị đại diện
			\begin{center}
				\begin{tabular}{|c|c|c|c|c|}
					\hline Cân nặng (kg) & $[40; 50)$ & $[50; 60)$  & $[60; 70)$ & $[70; 80)$ \\
					\hline Cân nặng đại diện  & 45 & 55 & 65 & 75 \\
					\hline Số học sinh & 6  & 12 & 10 & 2 \\
					 \hline
				\end{tabular}
			\end{center}
			 Cân nặng trung bình của học sinh lớp học trên
			 \[
			(45 \cdot 6+55 \cdot 12+65 \cdot 10+75 \cdot 2) \colon 30 \approx 57{,}67. 
			 \]
		\end{enumerate}
	} 
\end{ex}

% nguồn: Ok_TuanLuan_Toan11_Lan2
\begin{ex}%[1D5V2-3]%[ma:8090f0da]
	Một học sinh làm bài kiểm tra môn Toán gồm $30$ câu trong $1$ tháng được thống kê lại như sau: 
	\begin{center}
		\begin{tabular}{|c|c|c|c|c|c|c|c|c|c|}
			\hline $8$ & $18$ & $22$ & $11$ &  $25$ & $23$ & $21$ & $13$ & $14$ & $15$ \\
			\hline $18$ & $25$ & $10$ & $11$ & $24$ & $12$ & $14$ & $14$ & $6$ &  $8$ \\
			 \hline
		\end{tabular}
	\end{center}
	\begin{enumerate}
		 \item Tìm tứ phân vị của dãy số liệu trên.
		\item Tổng hợp lại dãy số liệu trên vào bảng tần số ghép nhóm theo  mẫu sau
		\begin{center}
			\begin{tabular}{|c|c|c|c|c|}
				\hline Điểm số &$[6; 10]$ &$[11; 15]$ &$[16;  20]$ &$[21; 25]$ \\
				\hline Số trận & $?$ & $?$  & $?$ & $?$ \\
				 \hline
			\end{tabular}
		\end{center}
		\item Hãy ước  lượng tứ phân vị của số liệu ở bảng tần số ghép nhóm trên.
	\end{enumerate}
	 \loigiai{
		\begin{enumerate}
			\item Các tứ phân vị là  $Q_2=14$, $Q_1=11$, $Q_3=21{,}5$.
			\item Bảng tần số ghép nhóm  theo mẫu.
			\begin{center}
				\begin{tabular}{|c|c|c|c|c|}
					\hline Điểm &$[6; 10]$ &$[11;  15]$ &$[16; 20]$ &$[21; 25]$ \\
					\hline Tần số &  $4$ & $8$ & $2$ & $6$ \\
					\hline 
				\end{tabular}
			\end{center}
			\item Tứ phân vị thứ nhất  là
			$Q_1=11+\dfrac{\dfrac{1.20}{4}-4}{8} \cdot (15-11)=11{,}5$. 
			 \\
			Tứ phân vị thứ hai là $Q_2=15{,}5$. 
			\\ 
			Tứ  phân vị thứ ba là $Q_3=21+ \dfrac{\dfrac{3 \cdot 20}{4}-14}{6} \cdot (25-21) \approx 21{,}67$.
		\end{enumerate}
	 }
\end{ex}

% nguồn: Ok_TuanLuan_Toan11_Lan2
\begin{ex}%[1D5V2-3]%[ma:4ac91f98]
	Cho bảng tần số ghép nhóm số liệu thống kê chiều cao  của 40 mẫu cây ở một vườn thực vật (đơn vị: centimét).
	\begin{center}
		\begin{tabular}{|c|c|}
			\hline Nhóm  & Tần số \\
			\hline$[30; 40)$ &  4 \\
			 \hline$[40; 50)$ & 10 \\
			\hline$[50; 60)$  & 14 \\
			\hline$[60; 70)$ & 6  \\
			\hline$[70; 80)$ & 4 \\ 
			\hline$[80; 90)$  & 2 \\
			\hline & $n=40$  \\
			 \hline
		\end{tabular}
	\end{center}
	\begin{enumerate}
		 \item Xác định số trung bình cộng, trung vị, tứ phân vị của mẫu số liệu ghép nhóm trên.
		\item  Mốt của mẫu số liệu trên là bao nhiêu?
	\end{enumerate}
	\loigiai{ 
		\begin{enumerate}
			\item Ta có 
			\begin{center}
				\begin{tabular}{|c|c|c|}
					 \hline Nhóm & Giá trị đại diện & Tần số \\
					\hline$[30;  40)$ & 35 & 4 \\
					\hline$[40; 50)$ & 45 & 10 \\ 
					\hline$[50; 60)$ & 55  & 14 \\
					\hline$[60; 70)$  & 65 & 6 \\
					\hline$[70; 80)$ & 75 &  4 \\
					\hline$[80; 90)$ & 85 &  2 \\
					\hline  & & $n=40$ \\
					 \hline
				\end{tabular}
			\end{center}
			Giá trị trung bình của mẫu số liệu là 
			 \[
			\overline{x} = \dfrac{35 \cdot 4+45 \cdot 10+55 \cdot 14+65 \cdot 6+75 \cdot 4+85  \cdot 2}{40}=55{,}5.
			\] 
			Khoảng chứa  trung vị $[50; 60)$. Trung vị của mẫu số liệu là
			 \[
			M_{\mathrm{e}}=50+\dfrac{20-14}{14} \cdot 10 \approx  54{,}29.
			 \]
			Nhóm  chứa tứ phân vị thứ nhất $Q_1$ và nhóm chứa tứ phân vị thứ hai $Q_3$ lần lượt là $[40; 50)$, $[60; 70)$.
			\\ 
			 Khi đó
			\begin{align*}
				& Q_1=40+\dfrac{10-4}{10}  \cdot 10=46 \\
				& Q_3=60+\dfrac{30-28}{6} \cdot  10=63{,}3.
			\end{align*}
			Vậy tứ phân vị của mẫu số liệu  là $Q_1=46$; $Q_2=54{,}29$; $Q_3=63{,}3$.
			\item 
			Nhóm  chứa mốt $[50; 60)$. Mốt của mẫu số liệu
			 \[
			M_0=50+\dfrac{4}{4+8} \cdot  10 \approx 53{,}3.
			\] 
		\end{enumerate}
	} 
\end{ex}

% nguồn: Ok_TuanLuan_Toan11_Lan2
\begin{ex}%[1D5V2-3]%[ma:aebd73e4]
	Chiều cao của hai loài hoa được một người thống kê theo  biểu đồ sau
	\begin{center}
		\begin{tikzpicture}[scale=.8,font=\footnotesize]
				\foreach \y  in {0,10,...,60}{
					\draw (0.1,\y/10) -- (-0.1,\y/10) node[left] {\y}; 
				 }
				\foreach \x/\label in {1/[100; 199],3/[200; 299],5/[300;  399],7/[400; 499]}{
					\draw (\x,0) -- (\x,-0.2) node[below] {\label}; 
				 }
				\foreach  \x/\y/\z in {1/20/35,3/18/30,5/14/20,7/10/15}{
					\fill[pattern={north east lines},draw]  ({\x-1},0) rectangle ++(1,\y/10);
					\node[above] at  ({\x-0.5},{\y/10}) {$\y$};
					\fill[pattern={horizontal lines},draw] ({\x},0) rectangle ++(1,\z/10); 
					 \node[above] at ({\x+0.5},{\z/10}) {$\z$};
					\draw[thick,->] (0,0) -- (12,0) node[below] {Chiều cao (mm)}; 
					\draw[thick,->]  (0,0) -- (0,7) node[left] {Số cây};
				} 
				\fill[pattern={north  east lines},draw] (9.5,6) rectangle ++(1,.5);
				\node[right] at (10.5,6.25) {Loại A}; 
				\fill[pattern={horizontal lines},draw]  (9.5,5) rectangle ++(1,.5);
				\node[right]  at (10.5,5.25) {Loại B};
			\end{tikzpicture}
	\end{center}
	\begin{enumerate}
		\item Hãy so sánh chiều cao  của loài $A$ và loài $B$ theo số trung bình và trung vị.
		\item Hãy ước lượng tứ phân vị thứ nhất và thứ ba của  chiều cao của loài $A$ và loài $B$.
	\end{enumerate}
	 \loigiai{
		\begin{enumerate}
			\item Bảng giá trị đại  diện
			\begin{center}
				\begin{tabular}{|c|c|c|c|c|}
					\hline Chiều cao đại diện (mm) & $149{,}5$ & $249{,}5$ &  $349{,}5$ & $449{,}5$ \\
					\hline Loài  $A$ & $20$ & $18$ & $14$ & $10$ \\
					 \hline Loài $B$ & $35$ & $30$ & $20$ & $15$ \\
					\hline 
				\end{tabular}
			\end{center}
			Chiều trung bình của loài $A$ xấp xỉ  bằng
			 \[
			\overline x_A 
			= \dfrac{149{,}5 \cdot 20 + 249{,}5 \cdot 18 + 349{,}5 \cdot 14 +  449{,}5 \cdot 10}{62}
			\approx  272{,}08.
			 \]
			Chiều trung bình của loài $B$ xấp  xỉ bằng
			\[ 
			\overline  x_B
			=\dfrac{149{,}5 \cdot 35 +  249{,}5 \cdot 30 + 349{,}5 \cdot 20 + 449{,}5 \cdot 15}{100}
			 = 264{,}5.
			\] 
			Chiều cao trung bình của loài $A$ lớn hơn loài $B$. 
			 \\
			Gọi $x_1 \leq x_2 \leq \ldots \leq x_{62}$ là chiều  cao của $62$ cây loài $A$.
			 \\
			Nên trung  vị của mẫu số liệu là $\dfrac{1}{2}\left(x_{31}+x_{32}\right)\in[200; 299)$.
			\\ 
			Suy ra trung vị là $M_{\mathrm{e}}=200+\dfrac{\dfrac{62}{2}-20}{18} \cdot  (299-200)=260{,}5$.
			 \\
			 Gọi $y_1 \leq y_2 \leq \ldots \leq y_{100}$ là chiều cao của $100$ loài $B$.
			\\ 
			Nên trung vị của mẫu số liệu là $\dfrac{1}{2}\left(x_{50}+x_{51}\right) \in[200;  299)$.
			 \\
			Suy ra trung vị là $M_{\mathrm{e}} =  200+\dfrac{\dfrac{100}{2}-35}{30} \cdot (299-200)=249{,}5$.
			\\ 
			Chiều cao trung bình của  loài $A$ lớn hơn $B$.
			\item Tứ phân vị thứ nhất của dãy số liệu  $x_1 \leq x_2 \leq \ldots \leq x_{62}$ là $x_{16}\in[100; 199)$.
			\\ 
			Do đó tứ phân vị thứ nhất  của loài $A$ là
			\[ 
			Q_1= 100  +\dfrac{\dfrac{1 \cdot 62}{4} -0}{20 } \cdot (199-100)= 176{,}725.
			\] 
			Tứ phân vị thứ ba của dãy số liệu $x_1 \leq x_2 \leq \ldots \leq x_{62}$ là  $x_{47}\in[300; 399)$.
			\\ 
			Do  đó tứ phân vị thứ ba của loài $A$ là
			\[ 
			Q_3 = 300+\dfrac{\dfrac{3 \cdot 62}{4}-38}{14} \cdot (399-300)  \approx 360{,}11.
			\] 
			Tứ phân vị thứ nhất  của dãy số liệu $y_1 \leq y_2 \leq \ldots \leq y_{100}$ là $y_{25}\in[100; 199)$.
			\\ 
			Do đó tứ phân vị thứ nhất của  loài $B$ là
			\[ 
			Q_1=100+\dfrac{\dfrac{1 \cdot 62}{4}-0}{35} \cdot (199-100)  \approx 143{,}84.
			\] 
			Tứ phân vị thứ ba của dãy số liệu $y_1 \leq y_2 \leq \ldots \leq y_{100}$ là $y_{75}\in[300; 399)$. 
			 \\
			 Do đó tứ phân vị thứ ba của loài $B$ là
			 \[
			Q_3=300+\dfrac{\dfrac{3 \cdot 62}{4}-65}{20} \cdot (399-300)=208{,}425. 
			\] 
		\end{enumerate}
	 }
\end{ex}

% nguồn: Ok_TuanLuan_Toan11_Lan2
\begin{ex}%[1D5V2-3]%[ma:c242f26a]
	Các bạn học sinh lớp 11A1 trả lời $40$ câu hỏi  trong một bài kiểm tra. Kết quả được thống kê ở bảng sau.
	\begin{center}
		\begin{tabular}{|c|c|c|c|c|c|}
			\hline Số câu trả lời đúng & $[16; 21)$ & $[21; 26)$ & $[26; 31)$ & $[31; 36)$  & $[36; 41)$ \\
			\hline Số học sinh &  4 & 6 & 8 & 18 & 4 \\
			\hline 
		\end{tabular}
	\end{center}
	Hãy ước lượng trung bình số câu trả lời đúng của  các học sinh lớp 11A1?
	 \loigiai{
		Ta có bảng  tần số ghép nhóm là
		\begin{center}
			\begin{tabular}{|c|c|c|c|c|c|}
				\hline Số câu trả lời đúng & $[16; 21)$ & $[21; 26)$ & $[26;  31)$ & $[31; 36)$ & $[36; 41)$ \\
				\hline Giá trị đại diện & 19  & 24 & 29 & 34 & 39 \\
				\hline Sô học  sinh & 4 & 6 & 8 & 18 & 4 \\
				 \hline
			\end{tabular}
		\end{center}
		Ước lượng trung bình số  câu trả lời đúng của các học sinh lớp 11A1 là
		 \[
		\overline{x}=\dfrac{n_1 c_1+n_2 c_2+n_3 c_4+n_4  c_4+n_5 c_4}{40} \approx 30 \,\text{câu hỏi}.
		\] 
	} 
\end{ex}

% nguồn: Ok_TuanLuan_Toan11_Lan2
\begin{ex}%[1D5V2-3]%[ma:b53c42a2]
	Một  đại lí bảo hiểm đã thống kê số lượng khách mua bảo hiểm nhân thọ trong một ngày ở biểu đồ sau.
	\begin{center}
		\begin{tikzpicture}
				\foreach \x/\y  in {0/4,2/6,4/10,6/7,8/3}{
					 \fill[pattern={north east lines},draw] (\x,0) rectangle ++(1,\y/2);
				} 
				\foreach \x/\y  in {1/3,3/9,5/6,7/4,9/2}{
					\fill[pattern={horizontal lines},draw] (\x,0) rectangle  ++(1,\y/2);
				 }
				\fill[pattern={north  east lines},draw] (9,5) rectangle (10,5.5);
				\node at (10.5,5.25)  {Nam};
				\fill[pattern={horizontal  lines},draw] (9,4) rectangle (10,4.5);
				\node  at (10.5,4.25) {Nữ};
				\node at (5,6.2)  {\textbf{Biểu đồ số lượng khách hàng}};
				\node at (5,5.8) {\textbf{theo giới tính và độ tuổi}}; 
				\draw[thick,->] (0,0) --  (11,0) node[below] {Tuổi};
				\draw[thick,->] (0,0)  -- (0,7) node[left] {Người};
				\foreach \y  in {0,2,...,12}{
					\draw (0.1,\y/2) --  (-0.1,\y/2) node[left] {\y};
				} 
				\foreach \x/\label in  {1/[20; 30],3/[30; 40],5/[40; 50],7/[50; 60],9/[60; 70]}{
					\draw (\x,0.1)  -- (\x,-0.1) node[below] {\label};
				} 
			\end{tikzpicture}
	\end{center}
	Hãy sử dụng dữ liệu trên để tư vấn cho đại lí bảo hiểm xác định khách hàng nam và nữ ở độ tuổi nào hay  mua bảo hiểm nhất.
	\loigiai{ 
		Dựa vào bảng dữ liệu, ta thấy mốt của mẫu số liệu  đối với khách hàng nam là
		 \[
		M_o=40+\dfrac{10-6}{(10-6)+(10-7)}  \cdot 10 \approx 46.
		 \]
		Nên với nam ở độ tuổi từ $40$ đến $50$ có nhu cầu mua bảo hiểm lớn nhất,  đặc biệt là độ tuổi $46$.
		 \\
		Mốt của mẫu số  liệu đối với khách hàng nữ là
		\[ 
		M_o=30+\dfrac{9-3}{(9-3)+(9-6)} \cdot  10 \approx 37.
		 \]
		Nên với nử ở độ tuổi từ $30$ đến $40$ có nhu cầu mua bảo hiểm lớn nhất, đặc biệt là  độ tuổi $37$.
	 }
\end{ex}

% nguồn: Ok_TuanLuan_Toan11_Lan2
\begin{ex}%[1D5V2-3]%[ma:6c7451a6]
	Kết quả điều tra thời gian chờ khám (đơn vị: phút) của bệnh nhân tại phòng khám của một  bệnh viện ở Thành phố Hồ Chí Minh được cho trong bảng sau
	\begin{center}
		\begin{tabular}{|c|c|c|c|c|}
			\hline Thời  gian (phút) & $[20; 25)$ & $[25; 30)$ & $[30; 35)$ & $[35; 40)$ \\
			\hline Tần số & &  & & \\
			 \hline
		\end{tabular}
	\end{center}
	Hãy ước lượng thời gian trung bình chờ khám bệnh  và mốt của mẫu số liệu ghép nhóm trên.
	\loigiai{ 
		Ta có bảng thống kê thời gian chờ khám bệnh theo  giá trị đại diện.
		\begin{center}
			\begin{tabular}{|c|c|c|c|c|}
				\hline Thời gian (phút) &  $[20; 25)$ & $[25; 30)$ & $[30; 35)$ & $[35; 40)$ \\
				\hline Thời gian đại diện & 22{,}5  & 27{,}5 & 32{,}5 & 37{,}5 \\
				\hline  Tần số & 4 & 5 & 8 & 3 \\
				 \hline
			\end{tabular}
		\end{center}
		Thời gian trung bình  chờ khám bệnh
		\[ 
		(22{,}5 \cdot 4+27{,}5  \cdot 5+32{,}5 \cdot 8+37{,}5 \cdot 3) \colon 20 = 30.
		\] 
	} 
\end{ex}

% nguồn: Ok_TuanLuan_Toan11_Lan2
\begin{ex}%[1D5V2-3]%[ma:16dedd16]
	Cân nặng của $28$  học sinh nam lớp 11 được cho như sau.
	\begin{center}
		\begin{tabular}{|c|c|c|c|c|c|c|}
			\hline 55{,}4 & 62{,}6 & 54{,}2 & 56{,}8 & 58{,}8 & 59{,}4 &  60{,}7 \\
			\hline 58 & 59{,}5 &  63{,}6 & 61{,}8 & 52{,}3 & 63{,}4 & 57{,}9 \\
			\hline 49{,}7 & 45{,}1 & 56{,}2 &  63{,}2 & 46{,}1 & 49{,}6 & 59{,}1 \\
			\hline 55{,}3 & 55{,}8  & 45{,}5 & 46{,}8 & 54 & 49{,}2 & 52{,}6 \\
			\hline 
		\end{tabular}
	\end{center}
	Hãy ước lượng cân nặng trung bình của học sinh trong bảng trên sau khi ghép nhóm và so sánh kết quả tìm được với cân nặng trung bình  của mẫu số liệu gốc.
	 \loigiai{
		 Ta có bảng ghép nhóm sau.
		\begin{center}
			\begin{tabular}{|c|c|c|c|c|c|}
				\hline  Cân nặng & $[45; 49)$ & $[49; 53)$ & $[53; 57)$ & $[57; 61)$ & $[61; 65)$ \\
				\hline Giá trị đại diện &  47 & 51 & 55 & 59 & 63 \\
				\hline Số học sinh & 4 & 5 & 7  & 7 & 5 \\
				\hline 
			\end{tabular}
		\end{center}
		 Cân nặng trung bình của học sinh xếp xỉ là
		\[ 
		\overline{x}=\dfrac{47 \cdot 4+51 \cdot 5+55 \cdot 7+59 \cdot 7+63 \cdot 5}{28}  \approx 55{,}6 \, (\text{kg}).
		\] 
		Cân nặng  trung bình của mẫu số liệu gốc là
		 \[
		 \overline{x}=\dfrac{54{,}2+56{,}8+58{,}8+59{,}4+\ldots+52{,}6}{28} \approx 53{,}4 \, (\text{kg}).
		\] 
		Ta thấy giá trị cân nặng trung bình ước lượng gần bằng giá  trị cân nặng trung bình của mẫu số liệu gốc.
	} 
\end{ex}
